% ====================================================================
% graphite_ica_full_rebuilt.tex — Ch1~5 통합 full 문건 (RB Phase 7.x, standalone)
%   흑연 음극 ICA(dQ/dV) 이론서 RB-series Chapter 1~5 재구성본을 한 문서로 병합.
%   ★ Ch6 해체(2026-06-01): 별도 Chapter 6 없음 — 구 Ch6 (실데이터 피팅 워크플로)
%     내용은 Ch1 부록 B 에 이미 흡수됨. 따라서 본 통합본은 Ch1~5 만 편입한다.
%     · preamble union(documentclass·kotex·amsmath·\AL/\brk 매크로·theorem 환경
%       groundingbox/boundbox/flagbox/keybox/linkbox/verifybox/practicebox/loopbox
%       — 5개 챕터 preamble 합집합, 중복 정의 1회).
%     · 각 챕터 \documentclass/\begin{document}/\end{document}/개별 \tableofcontents 제거.
%     · 챕터 = \section{Chapter N ...}, 챕터 내부 \section→\subsection, \subsection→\subsubsection
%       (article 유지, \part/\chapter 미사용 — \section 레벨 정합).
%     · 통합 bibliography(refs 문건과 동일 33 키 union) 1회 포함.
%     · title page + table of contents.
%   ★ 원본 5개 챕터 _rebuilt.tex 는 수정 0 — body 는 byte-verbatim 복사(요약 0).
%     label 충돌 0(Ch namespaced; Ch1 의 sec:ch6_* 부록 라벨도 타 챕터와 미충돌),
%     cross-chapter 평문 상호참조는 현행 유지(수정 0).
%   ★ theorem 환경 caption 통일 주의: verifybox/linkbox 는 챕터별 caption 문구가
%     달랐으나 union 1회 정의로 대표 caption 을 쓴다(물리 내용·식 변경 0, heading
%     label 만 통일). 상세는 본 문서 말미 통합 노트 참조.
% Date: 2026-06-02  Author: Project_Anode_Fit (Claude 측, RB-series 통합 재생성)
% ====================================================================
\documentclass[11pt,a4paper]{article}
\usepackage[margin=22mm]{geometry}
\usepackage{kotex}
\usepackage{amsmath,amssymb,mathtools,bm}
\usepackage{booktabs,longtable,array}
\usepackage{enumitem}
\usepackage{amsthm}
\usepackage{hyperref}
\usepackage{fancyhdr}
\usepackage{lastpage}
\usepackage{setspace}
\setstretch{1.13}
\setlength{\parskip}{0.45em}\setlength{\parindent}{0pt}\setlength{\headheight}{14pt}
% 통합본 전역 줄바꿈 여유(문단 overfull/underfull 흡수; 내용 변경 0):
\setlength{\emergencystretch}{3em}\tolerance=1200\hbadness=4000
\hypersetup{colorlinks=true, linkcolor=blue, urlcolor=blue, citecolor=blue,
  pdftitle={흑연 음극 ICA tail — Chapter 1~5 통합본 (RB 재구성본)}}
\pagestyle{fancy}\fancyhf{}
\lhead{흑연 음극 ICA(dQ/dV) tail 이론 — Ch.1~5 통합본 (RB)}\rhead{\thepage/\pageref{LastPage}}
\renewcommand{\headrulewidth}{0.4pt}

% --- theorem 환경 union (Ch1~5 preamble 합집합, 중복 정의 1회) ---
% (verifybox/linkbox 는 챕터별 caption 문구가 달랐으나, union 정의이므로 대표
%  caption 으로 통일 — 물리 내용·식 변경 0, box heading label 만 통일.)
\newtheorem*{groundingbox}{Grounding}
\newtheorem*{boundbox}{유효범위(Validity bound)}
\newtheorem*{flagbox}{FLAGGED (미확정)}
\newtheorem*{keybox}{Keystone}
\newtheorem*{linkbox}{전달식 (앞 장에서 상속)}
\newtheorem*{verifybox}{정합/유도 검증}
\newtheorem*{practicebox}{실무 팁 (본문 물리식 아님)}
\newtheorem*{loopbox}{도입 관찰로의 폐합}

% --- 매크로 union (Ch1~5 preamble 의 \newcommand 합집합, 중복 1회) ---
\newcommand{\dd}{\mathrm{d}}
\newcommand{\eff}{\mathrm{eff}}
\newcommand{\eq}{\mathrm{eq}}
\newcommand{\app}{\mathrm{app}}
\newcommand{\drive}{\mathrm{drive}}
\newcommand{\bg}{\mathrm{bg}}
\newcommand{\cell}{\mathrm{cell}}
\newcommand{\ext}{\mathrm{ext}}
\newcommand{\OCV}{\mathrm{OCV}}
\newcommand{\tot}{\mathrm{tot}}
\newcommand{\irr}{\mathrm{irr}}
\newcommand{\rev}{\mathrm{rev}}
\newcommand{\rxn}{\mathrm{rxn}}
\newcommand{\kin}{\mathrm{kin}}
\newcommand{\ohm}{\mathrm{ohm}}
\newcommand{\conc}{\mathrm{conc}}
\newcommand{\src}{\mathrm{src}}
\newcommand{\loss}{\mathrm{loss}}
\newcommand{\amb}{\mathrm{amb}}
\newcommand{\rlx}{\mathrm{relax}}
\newcommand{\cand}{\mathrm{cand}}
\newcommand{\conv}{\mathrm{conv}}
\newcommand{\prog}{\mathrm{prog}}
\newcommand{\coord}{\mathrm{coord}}
\newcommand{\ct}{\mathrm{ct}}
\newcommand{\sstat}{\mathrm{ss}}
\newcommand{\resid}{\mathrm{resid}}
\newcommand{\film}{\mathrm{film}}
\newcommand{\obs}{\mathrm{obs}}
\newcommand{\trans}{\mathrm{tr}}
\newcommand{\site}{\mathrm{site}}
\newcommand{\curr}{\mathrm{curr}}
\newcommand{\transp}{\mathrm{transport}}
\newcommand{\tr}{\mathrm{tr}}
\newcommand{\chem}{\mathrm{chem}}
\newcommand{\dis}{\mathrm{dis}}
\newcommand{\chg}{\mathrm{ch}}
\newcommand{\rst}{\mathrm{rest}}
\newcommand{\hys}{\mathrm{hys}}
\newcommand{\pol}{\mathrm{pol}}
\newcommand{\tar}{\mathrm{tar}}
\newcommand{\stt}{\mathrm{state}}
\newcommand{\cb}{\mathrm{cb}}
\newcommand{\dyn}{\mathrm{dyn}}
\newcommand{\tol}{\mathrm{tol}}
\newcommand{\noise}{\mathrm{noise}}
\newcommand{\AL}[1]{\textsf{[AL-#1]}}
\newcommand{\brk}{\discretionary{}{}{}}% 표/참고문헌 긴 DOI·ISBN 줄바꿈 허용(zero-width)

\title{\textbf{리튬이온전지 흑연 음극 ICA($dQ/dV$) tail 이론 — Chapter 1$\sim$5 통합본}\\[0.5em]
\large 단일 인과 사슬 $\to$ 가역 반응열 $\to$ 반응속도론 $\to$ entropy production $\to$ 히스테리시스\\[0.3em]
\normalsize 하나의 관찰(저온$\cdot$고율 ICA 꼬리 겹침)에서 출발하는 단일 인과 사슬의 5-Chapter 전개\\
(RB 재구성본 통합 — 구 Ch6 는 Ch1 부록 B 로 흡수)}
\author{Project\_Anode\_Fit (RB 재구성본)}
\date{2026-06-02}

\begin{document}
\sloppy
\maketitle

\begin{abstract}
\noindent 본 문서는 흑연 음극 ICA($dQ/dV$) tail 이론 RB-series 재구성본 Chapter 1$\sim$5 를 한 문서로
통합한 것이다. \textbf{Chapter 1} 은 하나의 관찰(``저온$\cdot$고율에서 ICA peak 꼬리가 길어져 다음 peak 와
겹친다'')에서 출발해 열역학이 무대($V_n$, 평형 target $\xi_{\eq}$)를 깔고 동역학(진행률 완화 $+$
relaxation-length spectrum)이 꼬리의 주역임을 보이고, closure(Plan A/B)$\cdot$심플 근사식$\cdot$
\textbf{실데이터 피팅 워크플로(부록 B — 구 Chapter 6 흡수)}까지 본 장 안에서 닫는다. \textbf{Chapter 2}
는 그 위에서 전지 가역 반응열(평형 전위 온도계수 $=$ reaction entropy)과 비가역 소산열(과전압 flux$\times$force)을
Bernardi 에너지 balance 로 잇는다. \textbf{Chapter 3} 은 mobility $k_j$ 를 미시 forward/backward rate
$r_j^\pm$(Level B)로 확대하고 detailed balance 로 $n_j^{\eff}=RT/(Fw_j)$ 를 고정한다. \textbf{Chapter 4} 는
transition-level entropy production 이 거시 비가역 발열 일반형 $\sum_j n_j^{\eff}I_j\eta_j\ge0$ 으로 정확히
환원됨을 보인다. \textbf{Chapter 5} 는 충전 branch 부호 $s_{\phi,j}^b=-1$ 을 유도하고 히스테리시스를
target 의 이력 의존(metastable branch)으로 설명한다. 통합 참고문헌(33 키)과 통합 Assumption
Ledger(AL-1$\sim$69)는 별도 문건 \texttt{graphite\_ica\_refs\_rebuilt.tex} 에 정리되어 있으며, 본
통합본에도 동일 33 키 bibliography 가 1회 포함된다.
\end{abstract}

\tableofcontents
\newpage


% ====================================================================
% ===== Chapter 1 body (원본 graphite_ica_ch1_rebuilt.tex, byte-verbatim; heading 1단계 demote) =====
% ====================================================================
\clearpage
\section{Chapter 1: 단일 인과 사슬 — 전하보존$\cdot$평형$\cdot$속도론$\cdot$spectrum$\cdot$closure$\cdot$실데이터 피팅(부록 B)}\label{chap:ch1}

% ====================================================================
\subsection*{서 — 하나의 관찰에서 출발하는 단일 인과 사슬}
\addcontentsline{toc}{subsection}{서 — 단일 인과 사슬}
% ====================================================================
본 장은 다음 \textbf{하나의 관찰}에서 출발하여 끝까지 그 끈을 놓지 않는다.
\begin{quote}\itshape
흑연 음극의 $dQ/dV$(ICA, incremental capacity) 곡선에서 상변이 peak 의 꼬리(tail)가 \textbf{온도가
낮을수록 길게 늘어져 다음 peak 와 겹치고}, 높을수록 짧게 끝난다. 또 전류(C-rate)가 커지면 peak 겹침이
심해진다.
\end{quote}
이 관찰을 \textbf{다섯 단계}로 푼다. 각 단계는 학부 수준(화학열역학·전기화학·미적분)으로 따라갈 수 있게
중간을 생략하지 않는다.
\begin{enumerate}[label=\textbf{(\arabic*)},leftmargin=2.2em]
\item \textbf{평형 열역학만으로는 설명이 안 된다.} 뒤에서 보이듯 평형 전이폭은 $w=RT/F$ 라 \emph{저온서
  오히려 좁아진다}(\S\ref{sec:eq}). 평형 이론은 ``저온 $\to$ 짧은 꼬리''를 예측한다 — 관찰과 \emph{반대}.
\item \textbf{따라서 꼬리의 주역은 동역학이다.} 진행률 $\xi_j$ 가 평형 목표 $\xi_{\eq,j}$ 를 \emph{유한
  속도}로 좇으며 생기는 지연(lag)이 꼬리를 만든다. 속도상수 $k\propto e^{-\Delta G_a/RT}$ 가 저온서
  작아져 지연이 커지고 꼬리가 길어진다(\S\ref{sec:dyn}) — 관찰과 \emph{일치}.
\item \textbf{그 배리어는 극판 전위가 낮춘다.} 전류가 흐를 때 내부 전위는 평형(OCV)과 다르며, 그 전위
  구동력 $\mathcal A_j$ 가 유효 배리어를 낮춰 진행 속도를 바꾼다(\S\ref{sec:rate}). 그래서 C-rate 가
  peak 겹침을 키운다.
\item \textbf{무대는 전하 보존이 깐다.} 내부 음극 전위 $V_n$ 은 외부에서 읽어오는 값이 아니라 \emph{전하
  보존식의 해}로 결정된다(\S\ref{sec:charge}). 꼬리의 \emph{깊이}는 배리어 분포가 만드는
  \emph{relaxation-length spectrum} 의 중첩(kernel integral)으로 설명된다(\S\ref{sec:spectrum}--\ref{sec:kernel}).
\item \textbf{결과 = 피팅 가능한 닫힌 논리식}(\S\ref{sec:fiteq}). 이후 발열(Ch.2,4)$\cdot$반응속도론
  (Ch.3)$\cdot$히스테리시스(Ch.5)로 확장하고, 수치$\cdot$피팅 실무는 \emph{본 장 안}(\S\ref{sec:ch1_simplefit} 심플
근사식, \S\ref{sec:planB}--\ref{sec:planA} closure, 부록 \S\ref{sec:ch1_numeric})에서 닫는다.
\end{enumerate}
\noindent\textbf{한 문장 요지: 열역학이 \emph{무대}($V_n$ 과 평형 target)를 깔고, 관찰된 꼬리 거동의
\emph{주역}은 동역학이다.} 본 Chapter 1 은 \textbf{방전(탈리튬화) 방향}만 다룬다. 충전 branch,
히스테리시스, full-cell 전압 분해는 Ch.5 에서 정의한다.

\begin{groundingbox}
\textbf{지배 표준(GS).} \textbf{(GS-1)} 모든 식은 물리적 의미를 먼저 말로 제시한 뒤 수식화하고 중간
단계를 생략하지 않는다(학부 무비약). \textbf{(GS-2)} 결론은 출판 수준 엄밀성. \textbf{(GS-3)} 모든 가정은
통합 Assumption Ledger(\AL{\#}: 논문/교과서 근거 + 검증 DOI)를 인용하며, 근거가 없으면 FLAGGED 로
표기하고 established 로 쓰지 않는다. \textbf{(GS-4)} 임의 clip/cap/softplus/step-function 정의식과 근거
없는 수치 임계값을 물리 모델로 쓰지 않는다(속도 제한은 물리적 병목으로 도입). 수치 해법과 피팅 실무는
본 장 부록 \S\ref{sec:ch1_numeric} 로 분리한다.
\end{groundingbox}

% ====================================================================
\subsection{기호와 단위}\label{sec:notation}
% ====================================================================
진행률 기호는 $\xi$ 로 통일한다(일부 문헌의 $\theta$ 와 동일물). 방전 누적 전하 좌표 $q$, 외부 전류 $I$,
방향 부호 $s_I,s_{\phi,j}\in\{+1,-1\}$(방전 기본 $+1$).
\begin{longtable}{p{0.16\textwidth} p{0.10\textwidth} p{0.66\textwidth}}
\toprule 기호 & 단위 & 의미 \\ \midrule \endhead
$q$ & --- & 방전 진행 좌표, $q=Q_\ext/Q_\cell$ \\
$Q_\ext,\ Q_\cell$ & C & 외부 누적 방전 전하, 기준 셀 용량(쿨롱) \\
$V_n$ & V & 내부 음극 전위 — 전하 보존식의 해(\AL{9}) \\
$V_{n,\app}$ & V & 측정되는 apparent 전위 $=V_n+s_I|I|R_n$ (\AL{6}) \\
$V_{n,\drive}$ & V & 속도식 구동 전위; 기본 $=V_n$ \\
$V_{n,\OCV}$ & V & 평형($\xi_j=\xi_{\eq,j}$)에서의 전하보존식 해 \\
$R_n$ & $\Omega$ & 음극 유효 분극 계수(전압축 이동) \\
$\xi_j,\ \xi_{\eq,j}$ & --- & $j$번째 effective transition 진행률, 평형 진행률 ($0\!\le\!\xi\!\le\!1$) \\
$U_j(T),\ w_j(T)$ & V & 전이 중심 전위(\S\ref{sec:eq} 에서 $U_j\equiv U_j^0-\mu^0/(s_{\phi,j}F)$ 로 $\mu^0$ 흡수), 평형 전이 폭 \\
$Q_{j,\tot}$ & C & $j$번째 transition 의 방전 방향 용량 기여 ($>0$) \\
$Q_\bg(V_n,T)$ & C & staging 으로 분리 안 되는 배경 누적 용량; $\partial Q_\bg/\partial V_n$=잔류 chemical capacitance \\
$\mathcal A_j$ & J/mol & 상변이 구동력 $=s_{\phi,j}F(V_{n,\drive}-U_j)$ \\
$\chi_j$ & --- & 배리어 낮춤 민감도 ($0\le\chi_j\le1$;\brk Ch.3 transfer coefficient $\beta_j$ 와는 \emph{대칭 Marcus 1차$\cdot$정상 target 극한 한정}으로만 동일물 — 일반은 별개) \\
$\Delta G_{a,j},\Delta H_{a,j},\Delta S_{a,j}$ & J/mol & 활성화 자유에너지/엔탈피/엔트로피 \\
$G$ & J/mol & local activation barrier (mode 별 값; \S\ref{sec:spectrum} 에서 분포) \\
$k_j,\ k_0(T)$ & 1/s & 진행률 완화 속도상수(mobility),\brk Eyring prefactor \\
$L$ & --- & local relaxation length $=|I|/(Q_\cell k_j)$ \\
$\rho_G(G;T)$ & mol/J & activation barrier 확률밀도 ($\int_0^\infty\rho_G\,\dd G=1$; $G$ 가 J/mol 이므로 단위는 그 역수 mol/J) \\
$A_L^{\mathrm{prob}},A_L^{\mathrm{amp}}$ & --- & relaxation-length 확률밀도(정규화) / 진폭 스펙트럼 $=A_0 A_L^{\mathrm{prob}}$;\brk kernel 의 $A_L\equiv A_L^{\mathrm{amp}}$ (\S\ref{sec:spectrum}) \\
$A_0(L)$ & --- & mode 별 물리적 진폭 가중(입자크기분포$\cdot$활성 site 밀도 등 독립 물리량) \\
$L_0,\ L_{\min}$ & --- & $L_0=|I|/(Q_\cell k_0)$ (prefactor 기준 길이), $L_{\min}=L_0 e^{-\chi_j\mathcal A_j/RT}$ (support 하한, $G\!\ge\!0$) \\
$q_a$ & --- & 꼬리 시작 좌표(post-peak tail onset);\brk single-mode 지수감쇠 $e^{-(q-q_a)/L}$ 의 기준점, ICA peak 위치 또는 $\dd\xi_\eq/\dd q$ 가 임계 이하로 떨어지는 첫 $q$ 로 고정 \\
$L_\varphi$ & V & 전압축 특성 꼬리 길이 $=|\dd V_n/\dd q|_{q_a}L$ (\S\ref{sec:fiteq}, 식~\eqref{eq:Lphi}) \\
$\Theta,\ Q_p$ & ---,\ C & 전체 effective phase progress, 그 용량 scale \\
$R$ & J/(mol\,K) & 기체상수(molar gas constant), $8.314$; $R=k_B N_A$, 평형 폭 $w_j=RT/F$ \\
$F$ & C/mol & Faraday 상수, $96485$;\brk 전위$\to$몰당 에너지 환산($\mathcal A_j=s_{\phi,j}F(V_{n,\drive}-U_j)$, $w_j=RT/F$) \\
$T$ & K & 절대온도(absolute temperature);\brk $w_j=RT/F$, $k_0=(k_BT/h)\kappa$, $\Delta G_{a,j}=\Delta H_{a,j}-T\Delta S_{a,j}$ 의 독립변수 \\
$k_B$ & J/K & Boltzmann 상수, $1.381\times10^{-23}$;\brk $S=k_B\ln W$, $R=k_B N_A$, Eyring prefactor $k_0=(k_BT/h)\kappa$ \\
$h$ & J\,s & Planck 상수, $6.626\times10^{-34}$;\brk Eyring prefactor $k_0(T)=(k_BT/h)\kappa(T)$ 의 고정 인자 \\
\bottomrule
\end{longtable}
\textbf{단위 주의.} 미분방정식의 전류 $I$ 가 A $=\mathrm{C/s}$ 이므로 $Q_\cell$ 은 쿨롱이어야 한다
($Q_\cell[\mathrm C]=3600\,Q_\cell[\mathrm{Ah}]$). 외부 전하·진행 좌표는 $Q_\ext(t)=\int_0^t|I|\,\dd t'$,
$q=Q_\ext/Q_\cell$, 정전류에서 $\dd q/\dd t=|I|/Q_\cell$.

% ====================================================================
\subsection{흑연 staging 과 effective transition}\label{sec:staging}
% ====================================================================
흑연의 lithiation 은 stage $4\to3\to2\mathrm L\to2\to1$ 의 순서로 진행하며, 각 전이가 $dQ/dV_n$ 에 peak 를
남긴다 \cite{dahn1991,ohzuku1993,funabiki1999ea,funabiki1999jes}(\AL{3}). 인접 전이의 중심 전위 차가
전이폭 이하이면 peak 가 융합되어 관측 peak 수가 $N_p\approx3\sim4$ 가 된다. 본 장은 분리 가능한 유효
transition 을 $j=1,\dots,N_p$ 로 둔다. 각 유효 transition 내부에도 입자 크기$\cdot$국소 환경$\cdot$결정립계
차이에 의한 \emph{local mode 분포}가 존재하며, 이것이 \S\ref{sec:spectrum} spectrum 의 물리적 근원이다.
진행률 방향은 방전 좌표 $q$ 와 같게 잡아 $\xi_j=0$(미진행)$\to\xi_j=1$(완료), $Q_{j,\tot}>0$.

% ====================================================================
\subsection{무대 (I): 전하 보존식과 내부 전위 $V_n$}\label{sec:charge}
% ====================================================================
\textbf{물리.} 외부 회로로 흘린 방전 전하와, 음극 내부의 방전 방향 누적 전하(배경 저장분 + staging 진행분)는
같아야 한다 — 단순한 전하 보존이다. 이 등식이 내부 전위 $V_n$ 을 결정한다(외부에서 주어지는 함수가
아니다; \AL{1},\AL{9}, \cite{doyle1993,newman}):
\begin{equation}
\boxed{\;Q_\cell\,q \;=\; Q_\bg(V_n,T) \;+\; \sum_{j=1}^{N_p} Q_{j,\tot}\,\xi_j\;.}
\label{eq:charge_balance}
\end{equation}
좌변은 외부 입력(아는 값), 우변은 내부 상태($V_n$ 과 $\xi_j$ 의 함수)다. 주어진 $(q,T,\{\xi_j\})$ 에서 이
식을 만족하는 $V_n$ 이 내부 전위다. 이는 다공성 전극 이론의 implicit 화학퍼텐셜 관계와 같은 구조다
\cite{doyle1993,newman}.

\textbf{배경 용량의 의미.} $Q_\bg(V_n,T)$ 는 단순 그래프 기저선이 아니라 뚜렷한 staging 으로 분리하지 않은
방전 방향 잔류 누적 용량이며, $\xi_j$ 가 평형 진행률을 즉시 따라가지 못할 때 전하 보존을 맞추기 위해 $V_n$
이 얼마나 이동해야 하는지를 정한다. 그 전위 기울기 $\partial Q_\bg/\partial V_n$ 이 잔류 chemical
capacitance 다. $V_n$ 을 안정적으로 풀려면 이 기울기가 양이어야 한다(\AL{5}; bazant2013 의 비평형
열역학적 미분용량):
\begin{equation}
\frac{\partial Q_\bg}{\partial V_n}\ \ge\ \epsilon_Q\ >\ 0 .
\label{eq:monotone}
\end{equation}
$\epsilon_Q$ 는 임의의 정규화 상수가 아니라 양의 미분용량이라는 물리 조건이며, 수치적으로는 해의 특이점을
막는 하한이다(\AL{5} BOUNDED).

\textbf{평형 극한 = OCV.} 평형(또는 충분히 느린 조건)에서는 $\xi_j=\xi_{\eq,j}(V_n,T)$ 이고
식~\eqref{eq:charge_balance} 가 $Q_\cell q=Q_\bg+\sum_j Q_{j,\tot}\xi_{\eq,j}$ 가 되어 그 해가 평형 음극
전위 $V_{n,\OCV}(q,T)$ 다. \textbf{즉 OCV 곡선은 임의로 주어진 lookup 이 아니라 전하 보존식의 평형
특수해다} — 중심식이 ``OCV 를 읽는'' 흐름으로 회귀하지 않는다.

\textbf{해 존재 조건(well-posedness).} 식~\eqref{eq:charge_balance} 가 $V_n$ 해를 가지려면 우변 목표가 배경
용량의 치역 안에 있어야 한다(\AL{7}): $Q_\cell q-\sum_j Q_{j,\tot}\xi_j\in\mathrm{Range}[Q_\bg(\cdot,T)]$.
단조성(식~\eqref{eq:monotone})과 이 치역 조건을 함께 만족하면 중간값 정리로 해가 유일하게 존재한다(부분
구간 피팅용 상대 기준형$\cdot$총용량 정합 등 실무는 \S\ref{sec:ch1_simplefit}).

% ====================================================================
\subsection{무대 (II): 평형 진행률 $\xi_\eq$ (logistic)}\label{sec:eq}
% ====================================================================
\textbf{왜 logistic 인가 (무비약 유도).} 흑연 intercalation 의 한 전이를 격자기체(lattice-gas)/정규용액
(regular-solution) 으로 보면, 점유율 $\xi$ 인 상태의 화학퍼텐셜은 \cite{mckinnon1983,bazant2013}(\AL{2})
\begin{equation}
\mu(\xi)=\mu^0+RT\ln\frac{\xi}{1-\xi}+\Omega_j(1-2\xi),
\label{eq:mu_latticegas}
\end{equation}
두 항의 출처를 학부 통계역학으로 한 줄씩 확인한다(식~\eqref{eq:mu_latticegas} 를 인용으로만 받지 않고 본문에서 유도한다).

\textbf{(i) 혼합 엔트로피 항.} $N$ 개의 동등한 자리 중 $n$ 개가 점유된 배열의 가짓수는
$W=\dfrac{N!}{n!\,(N-n)!}$ 이며, 점유율 $\xi\equiv n/N$ 으로 둔다. Boltzmann 관계 $S=k_B\ln W$ 에 Stirling
근사 $\ln m!\simeq m\ln m-m$ 를 적용하면(세 계승의 $-m$ 선형항이 $N=n+(N-n)$ 으로 상쇄되어)
\begin{equation}
S_{\mathrm{mix}}=k_B\big[N\ln N-n\ln n-(N-n)\ln(N-n)\big]
=-k_B N\big[\xi\ln\xi+(1-\xi)\ln(1-\xi)\big],
\label{eq:smix}
\end{equation}
즉 몰 기준($N\!\to\!N_A$, $R=k_B N_A$)으로 $S_{\mathrm{mix}}=-R[\xi\ln\xi+(1-\xi)\ln(1-\xi)]$ 다. 자유
에너지의 혼합 기여 $G_{\mathrm{mix}}=-TS_{\mathrm{mix}}=RT[\xi\ln\xi+(1-\xi)\ln(1-\xi)]$ 를 점유율로 미분하면
\begin{equation}
\frac{\partial G_{\mathrm{mix}}}{\partial\xi}=RT\big[(\ln\xi+1)-(\ln(1-\xi)+1)\big]=RT\ln\frac{\xi}{1-\xi}
\label{eq:mu_mix}
\end{equation}
($+1$ 과 $-1$ 이 상쇄). 여기서 $G_{\mathrm{mix}}$ 는 점유 자리 1몰당(intensive) 자유에너지이고, 화학퍼텐셜의 정의는
$\mu_{\mathrm{config}}=\partial G^{\mathrm{tot}}/\partial n$ ($G^{\mathrm{tot}}=N G_{\mathrm{mix}}$, 자리수 $N$, $\xi=n/N$)
인데 $\partial/\partial n=(1/N)\,\partial/\partial\xi$ 라 $N$ 인자가 상쇄되어 $\mu_{\mathrm{config}}=\partial G_{\mathrm{mix}}/\partial\xi$
다. 따라서 위 ξ-미분이 곧 화학퍼텐셜의 혼합 기여이며, 식~\eqref{eq:mu_latticegas} 의 첫 로그 항이 정확히 재현된다.

\textbf{(ii) 상호작용 항.} 정규용액 근사에서 점유 입자 간 상호작용(혼합 엔탈피)은
$G_{\mathrm{int}}=\Omega_j\,\xi(1-\xi)$ 이고, 그 미분 $\partial G_{\mathrm{int}}/\partial\xi
=\Omega_j(1-2\xi)$ (곱미분 $\dd[\xi-\xi^2]/\dd\xi=1-2\xi$)가 둘째 항이다. 표준 상태 항 $\mu^0$ 까지 더하면
$\mu=\mu^0+RT\ln[\xi/(1-\xi)]+\Omega_j(1-2\xi)$ 로 식~\eqref{eq:mu_latticegas} 가 완성된다(\AL{2};
$\mu^0,\Omega_j$ 의 미시값은 \cite{mckinnon1983,bazant2013} 에 위임).

\textbf{(iii) 전기화학 평형.} 평형은 이 화학퍼텐셜이 전극 전위와 균형을 이루는 조건
$\mu(\xi_\eq)=s_{\phi,j}F(V_n-U_j^0)$ 이다($s_{\phi,j}$=방전 방향 부호 인자, \S\ref{sec:notation};
eq:Aj 와 동일 convention). ideal 극한($\Omega_j=0$)에서 $\mu^0+RT\ln[\xi_\eq/(1-\xi_\eq)]
=s_{\phi,j}F(V_n-U_j^0)$ 이고, 상수 $\mu^0$ 를 기준 전위에 흡수해 $U_j\equiv U_j^0-\mu^0/(s_{\phi,j}F)$ 로
재정의하면
\begin{equation}
RT\ln\frac{\xi_\eq}{1-\xi_\eq}=s_{\phi,j}F(V_n-U_j).
\label{eq:eqbal_ideal}
\end{equation}

\textbf{(iv) $\xi_\eq$ 에 대해 풀기(한 줄씩).} 식~\eqref{eq:eqbal_ideal} 양변을 $RT$ 로 나누고 폭
$w_j\equiv RT/F$ 를 정의하면 $\ln[\xi_\eq/(1-\xi_\eq)]=s_{\phi,j}(V_n-U_j)/w_j$. 양변에 지수를 취해
$\dfrac{\xi_\eq}{1-\xi_\eq}=E$, $E\equiv\exp[s_{\phi,j}(V_n-U_j)/w_j]$. 이를 풀면 $\xi_\eq=E(1-\xi_\eq)$
$\Rightarrow\xi_\eq(1+E)=E\Rightarrow\xi_\eq=E/(1+E)$, 분자·분모를 $E$ 로 나눠 $1/E=\exp[-s_{\phi,j}(V_n-U_j)/w_j]$
이므로
\begin{equation}
\boxed{\;\xi_{\eq,j}(V_n,T)=\frac{1}{1+\exp[-s_{\phi,j}(V_n-U_j(T))/w_j(T)]}\;,\qquad w_j=\frac{RT}{F}\ (\text{ideal})\;.}
\label{eq:logistic}
\end{equation}
모든 중간 단계가 위에 명시됐다(방전 기준 $s_{\phi,j}=+1$ 이 기본이며, 이때 $V_n$ 증가가 $\xi_\eq$ 를 증가시킨다).
$\Omega_j\ne0$ 이면 식~\eqref{eq:eqbal_ideal} 우변에 $-\Omega_j(1-2\xi_\eq)$ 가 더해져 \emph{음함수}(implicit
isotherm)가 되며 단순 logistic 이 아니다 — 이때 $w_j$ 를 effective width 로 두는 것은 닫힌 유도가 아니라
coarse-grained ansatz 임을 명시한다(\AL{2} BOUNDED; $\Omega_j$ 가 크면 다중해·상분리 가능). 이는 매끄러운 logistic 곡선이며 $0\!\to\!1$ 급점프가
아니다. $\Omega_j\ne0$ 이면 폭·형태가 매끄럽게 보정되며, 실제로는 nonideality$\cdot$도메인 불균일$\cdot$유한
전이폭을 흡수하는 effective width $w_j$ 로 두어 $RT/F$ 에 강제로 묶지 않는다(\AL{4}). 이때
$n_j^{\eff}\equiv RT/(Fw_j)$ 가 effective 전자수 역할을 하며 Ch.3$\cdot$Ch.4 로 전달된다.

\textbf{저온 거동(서 단계 (1) 의 근거; 무비약).} logistic 의 peak 형상은 그 도함수로 결정된다. 식~
\eqref{eq:logistic} 에서 $z\equiv s_{\phi,j}(V_n-U_j)/w_j$ 로 두면 $\xi_\eq=(1+e^{-z})^{-1}$ 이고
$\dd z/\dd V_n=s_{\phi,j}/w_j$. 연쇄법칙으로 $\partial\xi_\eq/\partial V_n=(\dd\xi_\eq/\dd z)(\dd z/\dd V_n)$
이며, $\dd\xi_\eq/\dd z=-(1+e^{-z})^{-2}\cdot(-e^{-z})=e^{-z}/(1+e^{-z})^2$ 다. 여기서
$1-\xi_\eq=e^{-z}/(1+e^{-z})$ 이므로 $e^{-z}/(1+e^{-z})^2=[1/(1+e^{-z})]\cdot[e^{-z}/(1+e^{-z})]
=\xi_\eq(1-\xi_\eq)$ (표준 logistic 항등식)이다. 따라서
\begin{equation}
\frac{\partial\xi_{\eq,j}}{\partial V_n}=\frac{s_{\phi,j}\,\xi_{\eq,j}(1-\xi_{\eq,j})}{w_j(T)}
\label{eq:dxidV}
\end{equation}
로 결정되어 peak 높이 $\propto 1/w_j$, 폭 $\propto w_j$ 다($s_{\phi,j}=+1$ 방전 기준). 저온이면 \emph{ideal}
$w_j=RT/F$ 가 \emph{감소}하므로
peak 가 더 좁고 가팔라진다 — 즉 꼬리가 더 빨리 끝난다. 수치로 $25^\circ$C 에서 $w=RT/F\approx25.7$ mV,
$-15^\circ$C 에서 $\approx22.2$ mV($8.314\times298/96485=25.68$ mV, $\times258/96485=22.23$ mV).
\textbf{따라서 평형 열역학의 예측은 ``저온 $\to$ 짧은 꼬리''} 이고, 관찰(저온 긴 꼬리)은 평형 broadening 으로
설명할 수 없다. 꼬리의 주역은 평형이 아니라 동역학이어야 한다(\S\ref{sec:rate}--\ref{sec:dyn}).
\begin{flagbox}
\textbf{평형 형태는 데이터가 정한다.} logistic 대신 erf/Gaussian-cumulative isotherm 을 쓰는 것은 흑연에
대한 열역학적 유도가 문헌에 없고(무질서 반도체 Gaussian-disorder 모형 \cite{baessler1993} 과의 유추 +
경험적 peak-fit 수준) FLAGGED ansatz 다. 실제 형태는 저속 OCV 의 near-peak 감쇠를 $\log(dQ/dV)$ 대
$(V-U)$(지수=logistic) 또는 $(V-U)^2$(가우시안)로 보고 판별한다(\S\ref{sec:falsify}). \textbf{본 장의
꼬리 결론은 평형 형태에 무관}하다 — post-peak 에서 $\dd\xi_\eq/\dd q\simeq0$ 만 쓰기 때문이다.
\end{flagbox}

% ====================================================================
\subsection{주역 (I): 극판 전위에 의한 배리어 낮춤과 속도상수}\label{sec:rate}
% ====================================================================
이 절이 본 장의 핵심이다 — \emph{극판 전위가 어떻게 상변이 속도를 바꾸는가}.

\textbf{세 전위의 구분.} 측정되는 apparent 전위는 내부 전위에 분극이 더해진 값이다(\AL{6}):
\begin{equation}
V_{n,\app}=V_n+s_I|I|R_n(q,T,|I|).
\label{eq:Vapp}
\end{equation}
상변이 속도를 정하는 구동 전위는 $V_{n,\drive}$ 이며 가장 보수적인 기본형은 $V_{n,\drive}=V_n$ 이다(국소
계면 과전압을 분리하는 일반형 $V_{n,\drive}=V_n+\eta_{n,\drive}$ 는 Ch.3). \textbf{$R_n$(전압축 이동)과
아래의 속도 보조항을 같은 데이터로 동시에 자유 피팅하지 않는다} — 역할이 다르다($R_n$=전압 강하,
$\chi_j$=상변이 가속). 이 구분이 무너지면 둘이 서로를 흉내 내어 식별이 불가능해진다(\S\ref{sec:falsify}).

\textbf{구동력.} 한 전이의 구동력은 구동 전위가 전이 중심 전위에서 벗어난 정도다(\AL{11}):
\begin{equation}
\mathcal A_j=s_{\phi,j}F\,[V_{n,\drive}-U_j(T)].
\label{eq:Aj}
\end{equation}
$F[\mathrm V]$ 가 $\mathrm{J/mol}$ 차원이므로 $\mathcal A_j$ 는 몰당 자유에너지 구동력이다. 부호 인자
$s_{\phi,j}$ 는 선택한 방전 방향 전이의 구동력 부호를 정한다: $\chi_j\ge0$ 이고 $s_{\phi,j}[V_{n,\drive}-U_j]>0$
일 때 방전 방향 전이가 촉진된다. (본 장은 ideal 폭 $n_j^{\eff}=RT/(Fw_j)=1$ 을 기본으로 두므로 위 형태를 쓴다;
일반형 $\mathcal A_j=s_{\phi,j}\,n_j^{\eff}F(V_{n,\drive}-U_j)$ 는 effective width 와 함께 Ch.3 에서 쓴다.)

\textbf{유효 배리어(무비약).} 고유 활성화 자유에너지는 전이상태 이론(\AL{10}, \cite{eyring1935})으로
$\Delta G_{a,j}(T)=\Delta H_{a,j}-T\Delta S_{a,j}$ 다. 전위 구동력이 이 배리어를 \emph{선형으로} 낮춘다 —
이것이 Brønsted--Evans--Polanyi/Butler--Volmer 류의 표준 자유에너지 관계다(\AL{11},
\cite{evanspolanyi1938,bardfaulkner,bazant2013}).

\emph{선형의 출처(소구동력 1차 Taylor; 무생략).} Marcus 이론에서 구동력 $\mathcal A_j$ 가 주어졌을 때
활성화 장벽은 포물면 교차로부터 $\Delta G^\ddagger(\mathcal A_j)=\dfrac{(\lambda-\mathcal A_j)^2}{4\lambda}
=\dfrac{\lambda}{4}-\dfrac{\mathcal A_j}{2}+\dfrac{\mathcal A_j^2}{4\lambda}$ ($\lambda$=재구성 에너지)다.
$\mathcal A_j=0$ 둘레 1차 Taylor 전개를 하면 상수항 $\Delta G^\ddagger(0)=\lambda/4$ 와 1차 항
$\big[\dd\Delta G^\ddagger/\dd\mathcal A_j\big]_{0}\mathcal A_j=-\tfrac12\mathcal A_j$ 가 남고, 2차 항
$\mathcal A_j^2/(4\lambda)$ 은 $|\mathcal A_j|\ll\lambda$ 에서 무시된다. 따라서 1차 계수의 크기가 전달 계수
$\chi_j$ 다. \emph{대칭 Marcus 포물면(forward/backward 동일 곡률 $\lambda$)의 1차 전개는 정확히 $\chi_j=1/2$
단일값만 준다}; $\chi_j\ne1/2$(일반 $0\le\chi_j\le1$)은 1차 Taylor 자체에서 나오는 것이 아니라 \emph{곡률 비대칭}
($\lambda_f\ne\lambda_b$) 또는 BEP/Butler--Volmer 선형 자유에너지 관계의 경험적 기울기에서 오며
(\AL{11}, \cite{evanspolanyi1938,bardfaulkner}), 그 일반 범위로 둔다. 이 1차 항이 장벽을 낮추므로($\mathcal A_j>0$ 일 때 $-\chi_j\mathcal A_j<0$)
\begin{equation}
\boxed{\;\Delta G_{\eff,j}=\Delta G_{a,j}(T)-\chi_j\,\mathcal A_j\;=\;G-\chi_j\mathcal A_j\;}
\label{eq:Geff}
\end{equation}
($G\equiv\Delta G_{a,j}$ 를 mode 의 intrinsic barrier 로 표기, 위 $\lambda/4$ 등 상수항 흡수;
\S\ref{sec:spectrum} 에서 분포). 곡률(2차 항)이 중요해지는 $|\mathcal A_j|\sim\lambda$ 영역은 아래 Marcus
bound 로 분리한다. \emph{주의}: 이 $\chi_j\mathcal A_j$ 는 forward/backward 를 같은 배율로 키우는
\emph{방향 없는 common-mode mobility 가속}이며(keystone 참조), forward 장벽만 낮추는 표준 Butler--Volmer
와는 다르다 — 방향성 분해는 Ch.3 Level B 에서 다룬다. 구동력이
클수록($\mathcal A_j>0$) 유효 배리어가 낮아진다. 속도상수는 Eyring rate(\AL{10})
\begin{equation}
\boxed{\;k_j=k_0(T)\exp\!\Big[-\frac{\Delta G_{\eff,j}}{RT}\Big]\;,\qquad k_0(T)=\frac{k_BT}{h}\kappa(T)\;.}
\label{eq:kact}
\end{equation}
학부 수준에서는 Boltzmann 인자 $k_j\propto e^{-\Delta G_{\eff,j}/RT}$ 만으로 이후 논증에 충분하다 — 본 장의
모든 결론이 속도의 \emph{비}로 정의되는 길이 $L=|I|/(Q_\cell k_j)$(\S\ref{sec:dyn})로만 쓰여 prefactor 절대값
$k_0$ 가 상쇄되기 때문이다. prefactor 가 $k_BT/h$ 로 \emph{고정}(임의 상수 아님)이라는 절대값 자체는 전이상태
이론(\cite{eyring1935}) 결과로 \emph{인용}한다. $\Delta G_{\eff,j}\le0$
(강한 전위 구동)은 선형 근사가 강구동 영역에 들어간 것으로 표시하되 $\max(\cdot,0)$ 같은 절단으로 자르지
\emph{않는다}(GS-4; 절단이 필요해 보이면 그것은 아래 물리적 병목으로 다룬다).
\begin{keybox}
\textbf{$\chi_j\mathcal A_j$ 는 ``방향 있는 배리어 낮춤''이 아니라 ``방향 없는 완화 속도(mobility) 가속''이다
(Level A).} 무비약 근거: 2-상태 반응속도론에서 net rate $\dot\xi_j=r_j^+(1-\xi_j)-r_j^-\xi_j$. 정상점은
$\dot\xi_j=0$ 에서 $r_j^+(1-\xi_{\eq,j})=r_j^-\xi_{\eq,j}\Rightarrow r_j^+=(r_j^++r_j^-)\xi_{\eq,j}
\Rightarrow\xi_{\eq,j}=r_j^+/(r_j^++r_j^-)$ 다. $r_j^\pm$ 를 $\xi_j$ 의 작은 변화에 대해 국소 상수로 두면
net rate 를 직접 재배열할 수 있다: $\dot\xi_j=r_j^+-r_j^+\xi_j-r_j^-\xi_j=r_j^+-(r_j^++r_j^-)\xi_j
=(r_j^++r_j^-)\big[r_j^+/(r_j^++r_j^-)-\xi_j\big]=(r_j^++r_j^-)(\xi_{\eq,j}-\xi_j)$. 따라서
$\dot\xi_j=k_j(\xi_{\eq,j}-\xi_j)$, $k_j=r_j^++r_j^-$ 가 인수분해로 동시에 나온다(\AL{12}). \emph{강조}: 비약은
``$r_j^\pm$ 를 ($\xi_j$ 에 무관한) 국소 상수로 본다''는 \emph{가정}에만 있고, 그 가정 위에서는 위 재배열이 Taylor
전개가 아니라 대수적 항등이다(반대로 $r_j^\pm(\xi_j)$ 면 같은 식이 국소 선형화가 된다). 즉 $k_j$=완화 속도(목표에 접근하는 빠르기),
$\xi_{\eq,j}$=목표(방향). 식~\eqref{eq:kact} 처럼 $\chi_j\mathcal A_j$ 가 \emph{$k_j$
전체}에 들어가면 $r_j^+,r_j^-$ 가 같은 배율로 커져 비 $r_j^+/r_j^-=\xi_\eq/(1-\xi_\eq)$ 가 \emph{불변}
— 즉 방향(목표)은 열역학(\S\ref{sec:eq})이 전담하고, $\chi_j\mathcal A_j$ 는 방향 없는 속도 scale 일 뿐이다.
\emph{방향을 가진} 배리어 낮춤(forward 만 $-\beta\mathcal A$,
backward 는 $+(1-\beta)\mathcal A$ 라 비가 바뀜)은 \textbf{Ch.3 의 Level B} 에서 도입하며, 그때 본 장의
$\chi_j$ 가 transfer coefficient $\beta_j$ 와 같아진다(\AL{11}). \textbf{한정}: 본 장 Level A 의 ``공통 배율''
(forward/backward 동일 가속)은 Marcus 의 \emph{비대칭} 장벽 이동($\lambda\mp\mathcal A_j$)을 대칭으로 평균낸
\emph{모형 선택}이며, 이 대칭 극한에서만 forward/backward 가속이 세부 균형(detailed balance) 비를 보존한다. Level B(Ch.3)와
일치하는 것은 \emph{정상 target 극한} $\xi_{\mathrm{ss},j}\to\xi_{\eq,j}$ ($\xi_{\mathrm{ss},j}\equiv r_j^+/(r_j^++r_j^-)$, 본 keybox net rate 고정점) 한정이다(CHARTER step4). 본 장
결론식에서는 ``배리어 낮춤''이라는 방향성 함의 용어를 피한다.
\end{keybox}
\begin{boundbox}
\textbf{Marcus bound(\AL{15}).} 선형 $-\chi_j\mathcal A_j$ 는 작은 구동력의 1차 근사다. Marcus 이론
\cite{marcus1956} 은 곡률을 주므로 $|\mathcal A_j|$ 가 클 때(역전 영역 포함) 깨진다. 본 선형식은 꼬리 영역
$V_{n,\drive}\approx U_j$ 에서 유효하며 전역 정당화가 아니다.
\end{boundbox}

\textbf{물리적 병목(속도 절단이 아니라 직렬 합성).} 전하전달이 빨라져도 다음 단계(전해질/고체 수송, 유한
활성 자리, 고체상 확산, 유한 전류)가 전체 진행을 제한할 수 있다. 이를 \emph{시간척도의 직렬 합}으로 둔다
(수치 cap 이 아니라 율속 단계 합성):
\begin{equation}
\frac{1}{k_j^{\eff}}=\frac{1}{k_j}+\frac{1}{k_{j,\lim}},
\label{eq:kphys}
\end{equation}
$k_{j,\lim}\to\infty$ 면 활성화 지배($k_j^{\eff}\to k_j$), 반대면 병목 지배로 매끄럽게 환원된다.\footnote{%
$k_{j,\lim}$ 은 수송$\cdot$유한 자리$\cdot$고체 확산$\cdot$유한 전류 등 병목의 합($1/k_{j,\lim}=\sum 1/k_{j,m}$)
이다. 독립 자료 없이 자유 파라미터로 추가하면 $k_0,\Delta G_a,\chi_j$ 와 강하게 상관되므로 물리적 근거와
독립 제약이 있을 때만 쓴다. 유한 전류를 매끄럽지 않은 clamp($\min$/step)로 강제하지 않는다(GS-4); 그런
편의 처리와 피팅 안정화는 부록 \S\ref{sec:ch1_numeric} 에서 다룬다.} 이하 $k_j$ 는 이 $k_j^{\eff}$ 를 가리킨다.

% ====================================================================
\subsection{주역 (II): 진행률 동역학과 single-mode kernel}\label{sec:dyn}
% ====================================================================
상변이는 평형 진행률을 향해 유한 속도로 완화된다(\AL{12}, \cite{onsager1931,degrootmazur1962}):
\begin{equation}
\boxed{\;\frac{\dd\xi_j}{\dd t}=k_j\,[\xi_{\eq,j}(V_n,T)-\xi_j]\;.}
\label{eq:dxidt}
\end{equation}
$\xi_j$ 가 평형을 즉시 따라가지 못하면 전압 곡선과 ICA/DVA 에 꼬리가 생긴다. 정전류 구간($|I|>0$)에서는
$\dd q/\dd t=|I|/Q_\cell$ 이므로 그 역수가 $\dd t/\dd q=Q_\cell/|I|$ 이고, 연쇄법칙
$\dd\xi_j/\dd q=(\dd\xi_j/\dd t)(\dd t/\dd q)$ 에 식~\eqref{eq:dxidt} 의 $\dd\xi_j/\dd t=k_j[\xi_{\eq,j}-\xi_j]$
를 넣으면
\begin{equation}
\frac{\dd\xi_j}{\dd q}=\frac{Q_\cell}{|I|}\,k_j\,[\xi_{\eq,j}-\xi_j].
\label{eq:dxidq}
\end{equation}

\textbf{Single-mode kernel.} 한 mode 의 지연을 $r=\xi_{\eq,j}-\xi_j$ 라 하자. peak 를 지난 영역(post-peak)
에서 목표가 거의 정지하고($\dd\xi_{\eq,j}/\dd q\simeq0$) 구동력도 완만하면($\dd\mathcal A_j/\dd q\simeq0$,
즉 $V_n$ 완만 — \emph{별개 가정 두 개}임에 유의), 식~\eqref{eq:dxidq} 는 1계 선형 ODE 가 된다.
\emph{무비약 유도}: 정의 $r=\xi_{\eq,j}-\xi_j$ 를 $q$ 로 미분하면 $\dd r/\dd q=\dd\xi_{\eq,j}/\dd q-\dd\xi_j/\dd q$
이고, post-peak 가정 $\dd\xi_{\eq,j}/\dd q\simeq0$ 으로 첫 항이 사라져 $\dd r/\dd q=-\dd\xi_j/\dd q$ 다. 여기에
식~\eqref{eq:dxidq} $\dd\xi_j/\dd q=(Q_\cell k_j/|I|)(\xi_{\eq,j}-\xi_j)=(Q_\cell k_j/|I|)\,r$ 를 넣고
$L\equiv|I|/(Q_\cell k_j)$ 로 두면 $\dd r/\dd q=-(1/L)\,r$ 다. 여기서 $L$ 은 일반적으로 $q$ 의 함수다($k_j$ 가
구동력 $\mathcal A_j$ 에 지수 의존, 식~\eqref{eq:kact}). \emph{그러나} 위 가정 (나) $\dd\mathcal A_j/\dd q\simeq0$
이면 $k_j$, 따라서 $L$ 이 이 구간에서 $q$-무관(상수)이므로 $L$ 을 적분 밖으로 뺄 수 있다 — 이것이 ``$L$ 일정''의
근거다(별개 가정이 아니라 (나)의 \emph{귀결}). 변수분리 $\dd r/r=-\dd q/L$ 를 $q_a$ 에서 $q$ 까지 적분하면
$\int_{q_a}^q \dd q'/L=(q-q_a)/L$, 즉 $\ln[r(q)/r(q_a)]=-(q-q_a)/L$ 이고 양변에 지수를 취해 지수 감쇠가 따라온다:
\begin{equation}
r(q)\simeq r(q_a)\exp\!\Big[-\frac{q-q_a}{L}\Big],\qquad \boxed{\;L=\frac{|I|}{Q_\cell\,k_j}\;}
\label{eq:single_kernel}
\end{equation}
(차원: $\mathrm A=\mathrm{C/s}$, 속도상수 $k$ 는 $\mathrm{s^{-1}}$, $Q_\cell$ 은 $\mathrm C$ 이므로
$|I|/(Q_\cell k)=(\mathrm{C/s})/(\mathrm C\cdot\mathrm{s^{-1}})=(\mathrm{C/s})/(\mathrm{C/s})=1$, 즉 무차원).
이것은 \emph{관측 꼬리 전체가 단일
지수}라는 뜻이 아니라 \textbf{하나의 mode 가 만드는 지수 kernel} 이다. 관측 꼬리는 여러 mode 의 중첩이다
(\S\ref{sec:kernel}).

\textbf{잔차에서 진행률 도함수로(무비약 연결).} 다음 절의 중첩은 잔차 $r$ 이 아니라 \emph{진행률 도함수}
$\dd\xi_j/\dd q$ 의 합이다. 둘의 연결은 식~\eqref{eq:dxidq} 에서 $\dd\xi_j/\dd q=(Q_\cell k_j/|I|)\,r=r/L$
이므로
\begin{equation}
\frac{\dd\xi_j}{\dd q}=\frac{r(q_a)}{L}\exp\!\Big[-\frac{q-q_a}{L}\Big].
\label{eq:single_dxidq}
\end{equation}
바로 이 \textbf{$1/L$ 인자}가 \S\ref{sec:kernel} kernel 의 $1/L$ 출처다 — 중첩은 각 mode 의 \emph{진행률
기여}를 합한 것이지 잔차를 합한 것이 아니다.

\textbf{저온 거동(서 단계 (2)).} $k_j=k_0e^{-(G-\chi\mathcal A)/RT}$ 가 저온서 작아져 $L$ 이 커진다 $\Rightarrow$
꼬리가 길어진다 — 관찰과 일치. 빠른 kinetics 극한($k_j\to\infty$)에서는 $L\to0$ 이며 $\dd\xi_j/\dd q\to
\dd\xi_{\eq,j}/\dd q$ 로 접근한다(작은 지연 $\times$ 큰 속도의 곱이 유한; bracket 이 단순히 0 이 되는 것이
아니라 점근 균형의 결과다).

% ====================================================================
\subsection{주역 (III): 배리어 분포 $\to$ relaxation-length spectrum}\label{sec:spectrum}
% ====================================================================
\textbf{물리.} 실제 전극은 단일 배리어 $G$ 가 아니라 입자 크기$\cdot$국소 staging 환경$\cdot$결정립계$\cdot$결함
$\cdot$응력 차이로 인한 \emph{배리어 분포} $\rho_G(G;T)$ 를 가진다(\AL{13}). 진행률을 특정 배리어 기준으로
자르지 않고, 배리어 $g$ 를 가진 집단의 진행률을 $\xi_j(g,t)$ 로 두어 분포로 평균한다($\int_0^\infty
\rho_G(g;T)\,\dd g=1$):
\begin{equation}
\xi_j(t)=\int_0^\infty \rho_G(g;T)\,\xi_j(g,t)\,\dd g,\qquad
\frac{\dd\xi_j(g,t)}{\dd t}=k_j(g)\,[\xi_{\eq,j}-\xi_j(g,t)].
\label{eq:xi_dist}
\end{equation}
\begin{flagbox}
\textbf{독립 완화 가정(평균장; \AL{13}).} 식~\eqref{eq:xi_dist} 는 배리어 $g$ 별 집단이 \emph{서로 독립적으로}
1차 완화한다고 본다(mode 간 결합 무시). 그러나 흑연 staging 은 nucleation$\cdot$상경계 이동을 동반하는
\emph{협동적 1차 상전이}다(\AL{3}, \cite{funabiki1999ea,funabiki1999jes}). 따라서 ``독립 평행 1차 완화의 분포''
근사는 mode 간 상호작용(상경계 결합$\cdot$Avrami 형 비선형 동역학)을 무시한 \emph{평균장 ansatz} 이며, 그 타당성은
FLAGGED — Ch.3(반응속도론 일반화)$\cdot$본 장 부록 \S\ref{sec:ch1_numeric}(수치)에서 검증한다.
\end{flagbox}

\textbf{배리어 $\to$ 길이 지수 매핑(꼬리가 늘어지는 근원).} 관측 꼬리의 모양은 $\rho_G$ 와 \emph{같지 않다}
— 배리어 $G$ 가 먼저 속도로, 다시 길이로 \emph{지수 변환}되기 때문이다. 식~\eqref{eq:single_kernel} 의
$L=|I|/(Q_\cell k_j)$ 에 식~\eqref{eq:kact} $k_j=k_0(T)\exp[-(G-\chi_j\mathcal A_j)/RT]$(mode 별 intrinsic
barrier $G\equiv\Delta G_{a,j}$, $\Delta G_{\eff,j}=G-\chi_j\mathcal A_j$)를 대입한다. \emph{이 닫힌 지수 매핑은
활성화 지배 극한}($k_j^{\eff}\simeq k_j$, 즉 식~\eqref{eq:kphys} 의 $k_{j,\lim}\to\infty$)\emph{을 전제}한다 — 꼬리
영역($V_{n,\drive}\approx U_j$, 저구동)에서는 활성화가 율속이라 정당하나, 병목이 유효하면
$1/k_j^{\eff}=1/k_j+1/k_{j,\lim}$ 가 지수 밖 유리식으로 들어가 $L(G)$ 가 단순 지수형이 아니다(그 경우의 수치 풀이는
부록 \S\ref{sec:ch1_numeric}). $L$ 은 $k_j$ 의 \emph{역수}에 비례하므로 $1/k_j=(1/k_0)\exp[+(G-\chi_j\mathcal A_j)/RT]$ 로
\emph{지수 부호가 반전}되어(속도가 작을수록 길이가 길다)
\begin{equation}
\boxed{\;L(G,T,V_{n,\drive})=\frac{|I|}{Q_\cell\,k_0(T)}\exp\!\Big[\frac{G-\chi_j\mathcal A_j}{RT}\Big]\;.}
\label{eq:L_of_G}
\end{equation}
양변에 로그를 취하면 $\ln L=\ln L_0+(G-\chi_j\mathcal A_j)/RT$ 이다. 즉 배리어의 표준편차 $\sigma_G$($\rho_G$ 의
표준편차)가 $\ln L$ 에서 $\sigma_G/RT$ 로 나타나, \textbf{저온일수록 $\ln L$ 의 퍼짐이 커진다} — 작은 배리어 분산이 길이에서
지수적으로 확대된다. 이것이 ``저온 긴 꼬리''의 근원이며, 고정된 배리어 분포가 온도를 통해 stretched 거동을
만든다는 메커니즘은 무질서계$\cdot$빠른 이온 전도체 동역학에서 확립돼 있다(\AL{14},
\cite{svare2000,johnston2006,lindsey1980}; 단 graphite ICA 꼬리로의 \emph{적용}은 BOUNDED, \S\ref{sec:kernel}).

\textbf{변수변환: 배리어 분포 $\to$ 완화 길이 분포(무비약, 단계별).} 우리가 만들 변수 $A_L(L)$ 은 ``완화 길이가
$L$ 인 mode 가 꼬리에 \emph{얼마나 기여}하는가''의 분포다. 이는 \emph{이미 아는} 배리어 분포 $\rho_G(G)$ 를
식~\eqref{eq:L_of_G} 의 $L(G)$ 로 \emph{좌표만 바꾼} 것 — 새 물리를 넣는 게 아니라 같은 mode 집단을 길이 축에서
다시 세는 것이다. 세 단계로 만든다.

\emph{단계 1 — 자코비안(밀도 보존).} 확률(mode 분율)은 좌표를 바꿔도 보존되므로 $A_L^{\mathrm{prob}}(L)\,\dd L
=\rho_G(G)\,\dd G$. 역변환(식~\eqref{eq:L_of_G} 을 $G$ 에 대해 풀면) $G(L)=\chi_j\mathcal A_j+RT\ln(L/L_0)$
($L_0\equiv|I|/(Q_\cell k_0)$; $G=0$ 일 때 $L=L_{\min}$, 단계 2 에서 정식화)에서 $\dd G/\dd L=RT/L>0$ 이라
일대일이므로
\begin{equation}
A_L^{\mathrm{prob}}(L)=\rho_G\big(G(L)\big)\,\Big|\frac{\dd G}{\dd L}\Big|=\rho_G\big(G(L)\big)\,\frac{RT}{L}.
\label{eq:spectrum_prob}
\end{equation}
\emph{의미}: $A_L^{\mathrm{prob}}(L)\dd L$ $=$ 완화 길이가 $[L,L+\dd L]$ 인 mode 의 \emph{분율}, 규격 보존
($\int A_L^{\mathrm{prob}}\dd L=\int\rho_G\dd G=1$). 자코비안 $RT/L$ 은 ``배리어 한 칸 $=$ 길이 몇 칸''의 환산율일
뿐 임의 가중이 아니다.

\emph{단계 2 — 지지 범위(음의 배리어 금지).} 활성화 배리어는 음수가 못 되므로($G\ge0$, \S\ref{sec:notation})
식~\eqref{eq:L_of_G} 의 단조성상 $G\ge0\Leftrightarrow L\ge L_{\min}$, $L_{\min}\equiv L_0 e^{-\chi_j\mathcal A_j/RT}$
($G=0$ 대응 최단 길이)다. 따라서 $L<L_{\min}$ 에서 $A_L^{\mathrm{prob}}=0$ — 식에 Heaviside $H(L-L_{\min})$ 로
달아 음의 배리어 영역을 잘못 적분하지 않게 한다.

\emph{단계 3 — 개수 밀도 $\to$ 용량 가중(측도 변환, 무비약).} 단계 1$\cdot$2 의 $A_L^{\mathrm{prob}}$ 는 ``각 길이에
mode 가 몇이나''(개수 밀도)다. 그러나 관측 꼬리는 \emph{용량 가중} 진행 $\Theta=Q_p^{-1}\sum_m Q_m\xi_m$
(\S\ref{sec:fiteq}; $Q_m$=mode 의 용량 기여)로 나타나므로, 길이별 기여는 \emph{개수}$\times$\emph{mode 당 용량}이다.
이산 합을 연속 적분으로 옮기면
\begin{equation*}
\Theta=Q_p^{-1}\sum_m Q_m\,\xi_m\;\xrightarrow{\text{연속극한}}\;Q_p^{-1}\!\int Q(L)\,\xi(L)\,A_L^{\mathrm{prob}}(L)\,\dd L,
\end{equation*}
즉 개수 밀도 $A_L^{\mathrm{prob}}$ 에 \emph{mode 당 용량} $Q(L)$ 이 곱해진다. 이 용량 가중을 (무차원화한)
$A_0(L)\equiv Q(L)/\bar Q$ 로 두면 $A_0$ 는 개수 밀도 $\rho_G$ 에서 \emph{용량 측도}로 가는 Radon--Nikodym
가중이다 — 입자 크기 분포$\cdot$활성 site 밀도 같은 독립 물리량에서 오며 임의 피팅 knob 이 아니다. 관측 spectrum 은
그 곱이다:
\begin{equation}
\boxed{\;A_L(L)=A_0(L)\;\underbrace{\rho_G\big(G(L)\big)\,\frac{RT}{L}\,H(L-L_{\min})}_{\displaystyle A_L^{\mathrm{prob}}(L)\ (\text{규격 보존})}\;.}
\label{eq:spectrum}
\end{equation}
$A_0$ 는 \emph{임의 피팅 knob 이 아니라} 독립 물리량(입자 크기 분포 등)에서 와야 한다. $A_0\equiv1$ 이면
$A_L=A_L^{\mathrm{prob}}$ 로 규격 보존, $A_0\ne1$ 이면 $A_L$ 은 규격화 안 된 \emph{진폭 스펙트럼}이다 — \textbf{이하
kernel 은 이 $A_L$(진폭 spectrum)을 쓴다}(관측 꼬리는 진폭 가중 중첩이므로). 자코비안 $RT/L$ 과 아래 kernel 의
$1/L$ 인자는 \emph{출처가 다른 별개}임에 유의(\S\ref{sec:kernel}).
\begin{boundbox}
\textbf{비유일성(\AL{16}).} stretched 거동에서 $\rho_G$ 로의 역매핑은 \emph{비유일}하다(재포획 등 다른
메커니즘도 같은 꼬리를 낼 수 있다 \cite{baessler1993}). 따라서 $\rho_G$ 를 관측 꼬리에서 \emph{역산하지
않고}, 독립 측정/모형한 $\rho_G$ 로부터 꼬리를 \emph{forward 예측}만 한다(\S\ref{sec:falsify}).
\end{boundbox}

% ====================================================================
\subsection{관측 꼬리 = kernel integral}\label{sec:kernel}
% ====================================================================
\textbf{$A_L$ 의 물리 의미.} $A_L(L)\,\dd L$ 은 \emph{완화 길이가 $[L,L+\dd L]$ 인 mode 들이 전체 진행에 기여하는
분율}이다 — 즉 관측 꼬리를 ``완화 길이별로 분해한 분포''다. 빠른 mode(작은 $L$)는 peak 근처에서 빨리 끝나고,
느린 mode(큰 $L$)가 꼬리를 멀리 끌고 간다. 따라서 관측 꼬리는 \emph{서로 다른 $L$ 의 지수 감쇠를 $A_L$ 로 가중
합한 것}이다: 전체 진행률의 post-peak 꼬리는 각 mode 의 진행률 기여 $\dd\xi_j/\dd q=(r(q_a;L)/L)e^{-(q-q_a)/L}$
(식~\eqref{eq:single_dxidq})를 $A_L$ 로 가중 합한 것이다. 여기서 mode 별 초기 잔차 $r(q_a;L)$ 를 post-peak
공통값 $\bar r$ 로 \emph{근사}하면 $\bar r$ 가 적분 밖으로 빠지고, 이를 진폭 규격에 재흡수($A_L\!\leftarrow\!\bar r\,A_L$)해
아래 식을 얻는다 — $r(q_a)$ 가 사라진 게 아니라 $A_L$ 의 크기 인자로 흡수된 것이다:
\begin{equation}
\boxed{\;\frac{\dd\Theta_\mathrm{tail}}{\dd q}\simeq\int_0^\infty A_L(L)\,\frac1L\,
\exp\!\Big[-\frac{q-q_a}{L}\Big]\,\dd L\;.}
\label{eq:kernel_integral}
\end{equation}
이 균일 근사는 BOUNDED 다: 빠른 mode(작은 $L$)는 $q_a$ 에서 이미 더 진행해 $r(q_a;L)$ 가 작을 수 있어 원칙적으로
$L$-의존이다. $r(q_a;L)$ 이 $L$ 의 \emph{감소}함수이면 작은-$L$ 기여를 과대평가해 \emph{꼬리 앞부분이 실제보다
가팔라지는} 쪽으로 편향되며, 이 편향은 $r(q_a;L)$ 를 직접 보유하는 Plan B(g-grid, \S\ref{sec:planB})와 대조해
정량화한다. 피적분의 $1/L$ 인자는 진행률 도함수(식~\eqref{eq:single_dxidq})에서 온 것이고, spectrum 변환의
자코비안 $RT/L$ 은 $A_L$ 안에 이미 들어 있어 둘은 \emph{별개} 출처다 — 임의 가중이 아니다. $A_L$ 은 \S\ref{sec:spectrum}
의 $H(L-L_{\min})$ 를 품으므로 적분의 실효 하한은 $L_{\min}>0$ 이고 $1/L$ 의 $L\!\to\!0$ 발산은 일어나지 않는다.
이 적분은 relaxation-rate($1/L$) 분포 위의 Laplace 형 변환이며, 단일 지수의 \emph{연속 중첩}이 일반적으로 비지수
(stretched)$\sim$power-law 꼬리를 준다(\AL{14}, \cite{johnston2006}; KWW stretched 함수와 완화시간 분포의
등가성 \cite{lindsey1980}). single-mode(식~\eqref{eq:single_kernel})는 $A_L=\delta(L-L^\ast)$ 인 특수해다
($L^\ast$=그 mode 의 실제 완화 길이 $|I|/(Q_\cell k_j)$; \S\ref{sec:spectrum} 의 prefactor 길이 $L_0$ 와 구별).
또한 식~\eqref{eq:kernel_integral} 은 목표가 멈춘($\dd\xi_\eq/\dd q\!\simeq\!0$) post-peak \emph{자유감쇠} 극한의 꼬리이며,
목표가 움직이는 일반(자기참조 포함) 경우는 동일 kernel 을 \emph{목표} $\Theta_\eq$ 에 작용시킨 \S\ref{sec:volterra}
(식~\eqref{eq:volterra})로 일반화된다(잔차가 아니라 목표에 곱함 — \S\ref{sec:volterra} 의 이중계상 경고 참조).
\begin{boundbox}
\textbf{유효범위(\AL{14}, 정직 표기).} ``고정 배리어 분포 $\to$ 지수 매핑 $\to$ 저온 긴 꼬리''의 메커니즘은
빠른 이온 전도체$\cdot$무질서 완화에서 확립돼 있다(\cite{svare2000,johnston2006,lindsey1980}). 다만 (i)
특정 $\rho_G$ 형태가 \emph{정확히 어떤} stretched 지수를 주는지는 데이터/수치(부록 \S\ref{sec:ch1_numeric})로 정하며, (ii) 위
문헌은 전도도/유전 완화 맥락이라 graphite ICA 꼬리로의 \emph{적용}이 본 작업의 기여다. 이 두 항목은
established 가 아니라 BOUNDED 로 둔다.
\end{boundbox}

% ====================================================================
\subsection{자기참조 결합과 Volterra 적분식}\label{sec:volterra}
% ====================================================================
\textbf{왜 자기참조인가(물리).} $\xi_j$ 와 $V_n$ 은 전하 보존(식~\eqref{eq:charge_balance})으로 얽혀 있다 — 진행
$\Theta$ 가 늘면 $V_n$ 이 이동하고(식~\eqref{eq:charge_balance} 미분), $V_n$ 이 이동하면 평형 목표
$\Theta_\eq(V_n)$ 가 바뀌며, 바뀐 목표가 다시 $\Theta$ 의 완화 방향을 정한다. 이 되먹임 때문에 $\Theta$ 는 \emph{자기
자신을 목표로 갖는} 적분식의 해가 된다.

\textbf{적분형 유도(single-mode 에서, 무비약).} 한 mode 의 $q$-좌표 완화는 $\dd\Theta/\dd q=(1/L)(\Theta_\eq-\Theta)$
다(식~\eqref{eq:single_kernel} 와 같은 1차 완화, rate $1/L$). 이 1계 선형 ODE 를 상수변화법(적분인자 $e^{q/L}$)으로
풀면 $\dd(\Theta e^{q/L})/\dd q=(1/L)\Theta_\eq e^{q/L}$ 이고 $0$ 에서 $q$ 까지 적분해
\begin{equation}
\Theta(q)=\underbrace{\Theta(0)\,e^{-q/L}}_{\text{초기값 자유 감쇠}}+\int_0^q \frac1L\,e^{-(q-q')/L}\,\Theta_\eq(q')\,\dd q'.
\label{eq:singlemode_integral}
\end{equation}
즉 \textbf{kernel $(1/L)e^{-(q-q')/L}$ 이 \emph{목표} $\Theta_\eq$ 에 곱해진} 형이다(상수 목표 $\Theta_\eq=\Theta_0$,
$\Theta(0)=0$ 이면 $\Theta=\Theta_0(1-e^{-q/L})$ 로 복귀 — 무결성). 이제 두 곳을 일반화한다: \textbf{(i) 자기참조} —
목표가 미지 진행에 의존, $\Theta_\eq\to\Theta_\eq(V_n(q',\Theta(q')))$; \textbf{(ii) spectrum} — 단일 $L$ 을
분포 $A_L$ 로 적분. 목표 $\Theta_\eq$ 는 \emph{모든 mode 공통}이라(식~\eqref{eq:xi_dist} 의 $\xi_{\eq,j}$ 가
$g$-무관) $L$-적분 밖으로 빠지고 적분은 kernel 에만 작용하므로, $(1/L)e^{-(q-q')/L}\to\mathcal K(q,q')=\int
A_L(1/L)e^{-(q-q')/L}\dd L$, 초기항 $\Theta(0)e^{-q/L}\to\Theta_\mathrm{init}=\int A_L\Theta_{\mathrm{init},0}(L)e^{-q/L}\dd L$
($\Theta_{\mathrm{init},0}(L)$=mode 별 초기 진행, 상수 목표 scalar $\Theta_0$ 과 구별). 그러면 인과 상한 $q$ 의 Volterra 2종이
\emph{유도}된다(``뜬금없이''가 아니라 식~\eqref{eq:singlemode_integral} 의 두 일반화):
\begin{equation}
\Theta(q)=\Theta_\mathrm{init}(q)+\int_0^q \mathcal K(q,q')\,\Theta_\eq\big(V_n(q',\Theta(q')),T\big)\,\dd q',
\qquad \mathcal K(q,q')=\int_0^\infty A_L(L)\,\tfrac1L\,e^{-(q-q')/L}\,\dd L.
\label{eq:volterra}
\end{equation}
여기서 $\Theta_\mathrm{init}(q)=\int_0^\infty A_L(L)\,\Theta_{\mathrm{init},0}(L)\,e^{-q/L}\,\dd L$ 는 peak 진입 시점의 초기 mode
진행 $\Theta_{\mathrm{init},0}(L)$ 이 외부 구동 없이 자유 감쇠하는 동차항이고, 미지 $\Theta$ 는 좌변과 목표
$\Theta_\eq(V_n(q',\Theta(q')))$ 를 통해서만 등장한다 — 즉 ``현재 진행이 현재 목표를 바꾸고, 그 목표가 다시
현재 진행을 정한다''는 되먹임이 자기참조의 정체다.
\begin{boundbox}
\textbf{이중 계상 경고.} kernel $\mathcal K$ 는 반드시 \emph{목표} $\Theta_\eq$ 에 곱한다. 만약 잔차
$[\Theta_\eq-\Theta]$ 에 곱하면 상수 목표 극한에서 정상상태가 $\Theta_0/2$ 로 어긋난다(완화를 두 번 세는
오류). 상수 목표 $\Theta_\eq=\Theta_0$, $\mathcal K=(1/L)e^{-(q-q')/L}$, $\Theta_\mathrm{init}=0$ 극한에서
식~\eqref{eq:volterra} 는 $\Theta(q)=\Theta_0(1-e^{-q/L})$ 로 single-mode ODE(식~\eqref{eq:single_kernel})
와 정확히 일치하며 $q\to\infty$ 서 목표 $\Theta_0$ 로 수렴한다(single-mode 무결성 점검). \emph{분포} $A_L$ 의 경우
같은 극한의 $q\to\infty$ 값은 $\big(\int_0^\infty A_L\dd L\big)\Theta_0$ 이므로, $A_0\equiv1$(규격 보존, \S\ref{sec:spectrum})
일 때만 $\Theta_0$ 로 수렴한다 — $A_0\ne1$(진폭 spectrum)이면 총 scale 은 $\Theta\cdot Q_p$ 정의(\S\ref{sec:fiteq})의
용량 규격에 흡수되어 절대 scale 이 재조정된다(위 single-mode 점검은 $\delta$-collapse 한정).
\end{boundbox}
% ====================================================================
\subsection{자기참조식을 닫기 (1): Plan B — 직접 g-grid (항상 보존)}\label{sec:planB}
% ====================================================================
식~\eqref{eq:volterra} 에서 \emph{값}을 얻으려면 미지 $\Theta$ 를 닫아야 한다. 두 방법을 둔다(물리 결론은 어느
쪽으로 풀든 같다): \textbf{Plan B}(본 절, 항상 보존되는 수치 core/validator)와 \textbf{Plan A}(\S\ref{sec:planA},
검증되면 쓰는 닫힌형).

\textbf{Plan B — 직접 $g$-grid 수치 적분.} 배리어를 격자 $\{g_m\}$ 로 이산화해
식~\eqref{eq:xi_dist} 를 시간 적분하며, 매 시점 전하 보존(식~\eqref{eq:charge_balance})으로 $V_n$ 을 풀어 결합한다:
\begin{equation*}
\xi_j(t)=\sum_m w_m\,\rho_G(g_m)\,\xi_j(g_m,t),\qquad \sum_m w_m\,\rho_G(g_m)=\int_0^\infty\rho_G\,\dd g=1.
\end{equation*}
즉 $w_m$ 은 격자 $\{g_m\}$ 의 quadrature 가중으로 식~\eqref{eq:xi_dist} 의 \emph{적분을 그대로 이산화}한 것이지
새 가중을 넣은 게 아니며($w_m\rho_G(g_m)$ 이 정규화 측도를 이뤄 합이 1), 비균일$\cdot$large-$L$ 편중 격자
(부록 \S\ref{sec:ch1_numeric})에서도 이 정규화는 재계산해 보존한다. 주어진 $\rho_G$ 와 완화 closure 하에서
이산화 오차만 남는 수치 기준해이며, Plan A 닫힌형을 이 기준해와 상대오차
$\varepsilon=\lVert\Theta_\mathrm A-\Theta_\mathrm B\rVert/\lVert\Theta_\mathrm B\rVert$ 로 대조하는 \emph{validator}
를 겸한다(식~\eqref{eq:closed}, 수치 스킴$\cdot$수렴 판정 상세는 부록 \S\ref{sec:ch1_numeric}). 저온$\to$큰
$L$$\to$긴 꼬리 같은 \emph{물리 결론은 (위 완화 closure 전제 하에) 닫힌형 없이도 성립}한다.

% ====================================================================
\subsection{자기참조식을 닫기 (2): Plan A — 선형화 $\to$ 닫힌형}\label{sec:planA}
% ====================================================================
Plan B 가 항상 보존되는 수치해라면, Plan A 는 \emph{데이터 피팅에 직접 쓸} 닫힌 식이다(검증 통과 시). \textbf{baseline
$\Theta^{(0)}$ 정의(먼저 명시)}: 자기참조를 끈($V_n$ 고정, 아래 $b=0$) 0차 해를 $\Theta^{(0)}(q)$ 로 둔다 —
single-mode 극한에선 $\Theta^{(0)}(q)=\Theta_0(1-e^{-(q-q_a)/L})$(식~\eqref{eq:singlemode_integral} 의 상수목표
해), 일반적으로는 Volterra(식~\eqref{eq:volterra})를 $V_n$-고정으로 평가한 해다. Plan A 는 이 $\Theta^{(0)}$ 둘레의
self-consistency 보정이다.

\emph{단계 0(선형화).} 식~\eqref{eq:volterra} 의 비선형성은 오직 $\Theta_\eq(V_n(q',\Theta(q')))$ 한 곳이다.
$\Theta^{(0)}$ 둘레에서 1차 전개하면 $\Theta_\eq(V_n(\Theta))\approx a(q')+b(q')\,\Theta(q')$, 표준 1차 Taylor
분해상 절편은 $a(q')=\Theta_\eq\big(V_n(\Theta^{(0)}(q'))\big)-b(q')\,\Theta^{(0)}(q')$(baseline 에서 평가한 목표값에서
1차 기울기$\times$baseline 을 뺀 상수항)이고, 기울기(자기참조 계수)는
\begin{equation}
b(q')=\frac{\partial\Theta_\eq}{\partial V_n}\,\frac{\partial V_n}{\partial\Theta}
=-\frac{Q_p}{C_\bg}\,\frac{\partial\Theta_\eq}{\partial V_n}\ \le\ 0,
\label{eq:selfcoupling}
\end{equation}
$\partial V_n/\partial\Theta=-Q_p/C_\bg$ 는 전하 보존(식~\eqref{eq:charge_balance})을 $Q_p\Theta\equiv\sum_j
Q_{j,\tot}\xi_j$($Q_p=\sum_j Q_{j,\tot}$=총 phase 용량), $C_\bg\equiv\partial Q_\bg/\partial V_n$(배경 chemical
capacitance, 상세 \S\ref{sec:fiteq})로 묶은 $Q_\cell q=Q_\bg(V_n)+Q_p\Theta$ 를 $\Theta$ 로 미분해 얻고($C_\bg>0$), $\partial\Theta_\eq/\partial V_n>0$(평형 진행률은 전위에 증가, 식~\eqref{eq:dxidV} 의 $\dd\xi_\eq/\dd V_n>0$
를 $\Theta_\eq=Q_p^{-1}\sum_j Q_{j,\tot}\xi_{\eq,j}$ 로 합산)이므로 $b\le0$ 이 \emph{확정}된다. $b\le0$ 의 \emph{의미}는
``진행이 늘면 목표가 줄어 진행을 \emph{되끌어당긴다}''는 음의 되먹임이다.
이로써 식~\eqref{eq:volterra} 가 미지 $\Theta$ 에 \emph{선형}이 된다: $\Theta=c+\int_0^q\mathcal K\,b\,\Theta\,\dd q'$,
$c(q)=\Theta_\mathrm{init}+\int_0^q\mathcal K\,a\,\dd q'$.
\emph{단계 1(ratio ansatz).} 사용자 PhD \cite{lee2011,son2013} 의 ratio-substitution 을 \emph{미지} $\Theta(q')$
에만 적용한다: $\Theta(q')\approx\Theta(q)\,\Theta^{(0)}(q')/\Theta^{(0)}(q)$ (진행의 \emph{모양}은 baseline 을
따르고 크기만 미지로 둔다). \emph{정당화}: 자기참조항 $\int\mathcal K b\,\Theta$ 가 $c$ 에 대해 섭동($|b|$ 작음)일
때, 형상은 0차 baseline $\Theta^{(0)}$ 이 지배하고 $b$ 는 진폭만 재조정한다 — $b\to0$ 에서 정확(식~\eqref{eq:closed}
이 $c$ 로 복귀), $|b|$ 가 커질수록 형상 왜곡이 누적되어 아래 boundbox 의 부정확 한계로 이어진다.
\emph{단계 2($\Theta(q)$ 인수분해).} 단계 1 을 단계 0 의 선형식 $\Theta=c+\int_0^q\mathcal K b\,\Theta\,\dd q'$ 의
적분에 넣으면, 앞 인자 $\Theta(q)/\Theta^{(0)}(q)$ 가 $q'$ 에 무관이라 적분 \emph{밖}으로 빠진다:
\begin{equation*}
\int_0^q\mathcal K(q,q')\,b(q')\,\Theta(q')\,\dd q'\approx\frac{\Theta(q)}{\Theta^{(0)}(q)}\int_0^q\mathcal K(q,q')\,b(q')\,\Theta^{(0)}(q')\,\dd q'.
\end{equation*}
이를 $\Theta(q)=c(q)+(\text{위 적분})$ 에 넣고 $\Theta(q)$ 항을 좌변으로 모아 분모로 묶으면:
\begin{equation}
\boxed{\;\Theta_\mathrm A(q)=\frac{c(q)}{\,1-\dfrac{1}{\Theta^{(0)}(q)}\displaystyle\int_0^q\mathcal K(q,q')\,b(q')\,\Theta^{(0)}(q')\,\dd q'\,}\;.}
\label{eq:closed}
\end{equation}
\textbf{이것이 피팅에 쓰는 닫힌형이다} — 각 항의 의미: 분자 $c$ = 자유 감쇠($\Theta_\mathrm{init}$) + baseline
목표 기여, 분모의 적분 = 자기참조 감쇠($b<0$ 이라 분모 $>1$ → 꼬리를 줄임). \emph{무결성 점검}: 자기참조가
없으면($b\to0$) 분모 $\to1$, $\Theta_\mathrm A\to c$ 로 식~\eqref{eq:volterra} 의 올바른 극한(상수 목표서
$\Theta_0(1-e^{-q/L})$, 식~\eqref{eq:single_kernel})으로 복귀한다.
\begin{boundbox}
\textbf{Fredholm$\leftrightarrow$Volterra 가정차(\AL{16}; 그대로 쓰면 안 됨).} 사용자 논문(\emph{JCP}
\textbf{147}, 144111 (2017), Eq.(32)$\to$(34), \cite{jcp2017})은 \emph{Fredholm} 2종(고정 구간, 정상상태 확률)을
ratio-substitution 으로 닫는다. graphite 의 식~\eqref{eq:volterra} 는 \emph{Volterra}(인과 상한 $q$)다. 따라서
논문 결과를 그대로 옮기지 않고, 인과성을 유지한 채 자기참조만 선형화(단계 0)해 \emph{같은} ratio 구조를 적용했다.
또 \textbf{논문 자신이 ``reaction zone 이 넓으면 ratio 근사가 부정확''이라 경고}하는데, graphite 에서 넓은
spectrum $=$ 저온 stretched tail(가장 관심 영역)이므로 closure 가 가장 필요한 곳에서 가장 부정확할 수 있다.
$\Rightarrow$ \textbf{Plan A 단독 채택 금지}: Plan B(g-grid) 대비 상대오차
$\varepsilon=\lVert\Theta_\mathrm A-\Theta_\mathrm B\rVert/\lVert\Theta_\mathrm B\rVert$ 가 운용 $T\cdot$rate 에서
작을 때만 닫힌형으로 승격한다(BOUNDED).
\end{boundbox}

% ====================================================================
\subsection{ICA 와 DVA, 그리고 피팅 가능한 닫힌 논리식}\label{sec:fiteq}
% ====================================================================
\textbf{내부 전위 미분.} 식~\eqref{eq:charge_balance} 를 $q$ 로 미분하면(등온)
$Q_\cell=\dfrac{\partial Q_\bg}{\partial V_n}\dfrac{\dd V_n}{\dd q}+\sum_j Q_{j,\tot}\dfrac{\dd\xi_j}{\dd q}$,
따라서
\begin{equation}
\frac{\dd V_n}{\dd q}=\frac{Q_\cell-\sum_j Q_{j,\tot}\,\dd\xi_j/\dd q}{\partial Q_\bg/\partial V_n}.
\label{eq:dVdq}
\end{equation}
전체 phase progress 를 $Q_p\Theta\equiv\sum_j Q_{j,\tot}\xi_j$ 로 묶는다($Q_p=\sum_j Q_{j,\tot}$=용량 scale,
$\Theta\in[0,1]$ 무차원). 이 $\Theta$ 는 \S\ref{sec:kernel} 의 $\Theta$ 와 \emph{같은 양}이며, 그 post-peak
도함수가 $\dd\Theta_\mathrm{tail}/\dd q$(식~\eqref{eq:kernel_integral})다. 외부 전하 $\dd Q=\dd Q_\ext=
Q_\cell\,\dd q$ 는 배경 저장 $\dd Q_\bg=C_\bg\,\dd V_n$($C_\bg\equiv\partial Q_\bg/\partial V_n$)과
phase progress $Q_p\,\dd\Theta$ 의 합이다:
\begin{equation}
\dd Q=C_\bg\,\dd V_n+Q_p\,\dd\Theta.
\label{eq:dQsplit}
\end{equation}
이는 식~\eqref{eq:dVdq} 양변에 $C_\bg=\partial Q_\bg/\partial V_n$ 을 곱하고 $\dd q=\dd Q/Q_\cell$ 로 다시 쓴
\emph{동일식}이다(새 가정이 아니라 같은 전하보존 미분의 다른 표현). 양변을 $\dd Q$ 로 나누면 $1=C_\bg\,(\dd V_n/\dd Q)+Q_p\,(\dd\Theta/\dd Q)$ 이고, $\dd V_n/\dd Q$ 에 대해 풀어
역수를 취하면 내부 전위 기준 ICA
\begin{equation}
\boxed{\;\frac{\dd Q}{\dd V_n}=\frac{C_\bg(V_n,T)}{\,1-Q_p\,(\dd\Theta/\dd Q)\,}\;.}
\label{eq:fiteq}
\end{equation}
식~\eqref{eq:kernel_integral} 의 꼬리 기여가 $\dd\Theta/\dd Q$ 에 들어가 ICA 꼬리로 나타난다(좌표 bridge:
외부 전하 $Q=Q_\ext=Q_\cell q$ 이므로 $\dd\Theta/\dd Q=(1/Q_\cell)\,\dd\Theta/\dd q$). \emph{적용 영역}: post-peak
구간에선 꼬리 기여가 지배해 $\dd\Theta/\dd q\simeq\dd\Theta_\mathrm{tail}/\dd q$ 이고, peak 영역에선
$\dd\Theta/\dd q\to\dd\Theta_\eq/\dd q$(빠른 kinetics 극한, \S\ref{sec:kernel} 직전 single-mode 논의)로 환원된다.
$\dd\Theta/\dd Q$ 에 넣는 $\Theta$ 는 닫힌 $\Theta_\mathrm A$(식~\eqref{eq:closed}) 또는 Plan B 수치해를 $Q$ 로
미분한 값이다. 분모
$1-Q_p\,(\dd\Theta/\dd Q)>0$ 은 단조성(식~\eqref{eq:monotone}, $C_\bg=\partial Q_\bg/\partial V_n>0$)과
\emph{동치가 아니라 별개}의 유효범위 조건이다: 전자는 배경 용량의 국소 단조성(정적 chemical capacitance
양수)이고, 분모 양수성은 식~\eqref{eq:fiteq} 의 역수 $\dd V_n/\dd Q=(1-Q_p\,\dd\Theta/\dd Q)/C_\bg$ 가
양이라는 \emph{동역학 경로 조건}($C_\bg>0$ 전제하 $\dd V_n/\dd Q>0$)이다. 두 조건이 함께 성립해야
$\dd Q/\dd V_n>0$ 이다. staging peak 에서 분모가
$0$ 에 접근하면 $\dd Q/\dd V_n$ 이 발산하는데 이것이 정상적인 ICA peak 이며, 그 너머(분모 $<0$)는 전압 해의
특이점으로 보아 피팅에서 제외한다. 측정 축으로는
apparent 전위(식~\eqref{eq:Vapp})로 환산한다: 등온 정전류에서 $\dd V_{n,\app}/\dd q=\dd V_n/\dd q+
s_I|I|\,\partial R_n/\partial q$, 그리고 $\dd Q_\ext/\dd V_{n,\app}=Q_\cell/(\dd V_{n,\app}/\dd q)$,
$\dd V_{n,\app}/\dd Q_\ext$ 가 그 역수다(DVA). \emph{전압축 특성 꼬리 길이}는 다음과 같이 환산된다: post-peak 에서
$\dd V_n/\dd q$ 를 $q_a$ 둘레 상수로 근사하면(식~\eqref{eq:single_kernel} 유도의 가정 — 꼬리에서
$\dd\mathcal A_j/\dd q\simeq0$ 인 완만 영역) $V_n-V_n(q_a)\simeq(\dd V_n/\dd q)|_{q_a}(q-q_a)$ 이고,
single-mode 꼬리의 지수 $e^{-(q-q_a)/L}$ 가 $\exp[-(V_n-V_n(q_a))/L_\varphi]$ 로 옮겨져
\begin{equation}
L_\varphi\equiv\big|\dd V_n/\dd q\big|_{q_a}\,L
\label{eq:Lphi}
\end{equation}
가 정의된다($\dd V_n/\dd q$ 가 꼬리에서 완만하다는 가정 위의 BOUNDED 근사; 기호표에 $L_\varphi$[V] 등재).

\textbf{피팅 평가 순서(deliverable).} 식~\eqref{eq:fiteq} 는 다음 순서로 계산되는 닫힌 논리식이다:
\begin{enumerate}[label=\textbf{S\arabic*.},leftmargin=2.4em]
\item 전하 보존(식~\eqref{eq:charge_balance})으로 $V_n(q,T,\{\xi_j\})$ 를 구한다.
\item 평형 목표 $\xi_{\eq,j}(V_n,T)$(식~\eqref{eq:logistic}).
\item 속도 $k_j=k_0 e^{-(G-\chi_j\mathcal A_j)/RT}$, $\mathcal A_j=s_{\phi,j}F(V_{n,\drive}-U_j)$
  (식~\eqref{eq:kact}; 병목 시 식~\eqref{eq:kphys}).
\item 배리어$\to$길이 $L(G)$(식~\eqref{eq:L_of_G}) $\to$ spectrum $A_L$(식~\eqref{eq:spectrum}).
\item 꼬리 중첩 $\dd\Theta_\mathrm{tail}/\dd q$(식~\eqref{eq:kernel_integral}); 자기참조는 \emph{본 장에서} 닫는다 —
  Plan B(g-grid 수치, 항상) 또는 Plan A 닫힌형(식~\eqref{eq:closed}, $\varepsilon$ 통과 시).
\item $\dd Q/\dd V_n$(식~\eqref{eq:fiteq}) $\to$ 측정축 $\dd Q/\dd V_{n,\app}$.
\end{enumerate}
\textbf{피팅 파라미터}: $\{U_j,w_j,Q_{j,\tot},Q_\bg\}$(평형, 저속 OCV/GITT), $\{\Delta H_a,\Delta S_a,k_0\}$
(여러 $T$ 의 Arrhenius), $\chi_j$(여러 $|I|$/전위), $\rho_G$(독립 입력 — 역산 금지). \textbf{유효범위}:
$G-\chi_j\mathcal A_j\ge0$(Marcus, \AL{15}), $V_{n,\drive}\approx U_j+O(w_j)$. (전체 spectrum$\cdot$자기참조를
쓰는 정밀 평가의 solver 스킴$\cdot$수렴 판정은 부록 \S\ref{sec:ch1_numeric}.)

\subsubsection{심플 실데이터 피팅 근사식 (single-mode tail; 본 장만으로 사용 가능)}\label{sec:ch1_simplefit}
전체 spectrum 적분(식~\eqref{eq:kernel_integral}) 없이, \emph{하나의 우세 mode} 가 꼬리를 지배하는 극한
($A_L=\delta(L-L_0)$)에서는 위 모든 단계가 \textbf{종이와 연필로 닫히는} 가장 단순한 근사식으로 줄어든다. 이것이
실데이터 1차 피팅의 출발점이다. post-peak($q>q_a$) 에서 목표가 거의 정지하면(식~\eqref{eq:single_kernel}) 진행은
\begin{equation}
\Theta_\mathrm{tail}(q)\simeq\Theta_0\Big(1-e^{-(q-q_a)/L}\Big),\qquad
\frac{\dd\Theta_\mathrm{tail}}{\dd q}\simeq\frac{\Theta_0}{L}\,e^{-(q-q_a)/L},\qquad
\boxed{\;L=\frac{|I|}{Q_\cell\,k_j}=\frac{|I|}{Q_\cell k_0}\,e^{(G-\chi_j\mathcal A_j)/RT}\;}
\label{eq:simplefit}
\end{equation}
이고, 이를 fiteq(식~\eqref{eq:fiteq})에 넣으면 ICA 꼬리는 \emph{특성 길이 $L$ 의 지수 감쇠}로 나타난다.
\textbf{데이터 피팅 절차.}
\textbf{(0) 무대 고정(식별성 전제).} \S\ref{sec:falsify} 식별성 순서대로, 먼저 저속 OCV/GITT 로
$\{U_j,w_j,Q_{j,\tot},Q_\bg\}$ 를, 펄스/전류 차단으로 $R_n$ 을 고정한다($R_n$ 과 $\chi_j$ 가 공선형이므로 $R_n$ 을
\emph{먼저} 묶지 않으면 아래가 식별 불가). 측정 $V_{n,\app}$ 는 $V_n=V_{n,\app}-s_I|I|R_n$(식~\eqref{eq:Vapp} 역산,
$R_n$ 기지)으로 내부 $V_n$ 으로 \emph{직접} 환산한다(self-consistent 루프 아님).
\textbf{(1) $L$ 추출.} 꼬리 시작점 $q_a$ 를 \emph{먼저} 고정한다(ICA peak 위치, 또는 $\dd\xi_{\eq}/\dd q$ 가 임계
이하로 떨어지는 첫 $q$) — $q_a$ 를 자유 파라미터로 두면 $L$ 과 공선형이다. 그 뒤 post-peak ICA 꼬리를
모델 $y(q)=A\,e^{-(q-q_a)/L}+(\text{baseline})$ 로 회귀해 \emph{한 숫자} $L$ 을 뽑는다 — 진폭 $A$ 는 nuisance
($\Theta_0\cdot Q_p/C_\bg$ 등을 흡수; 관심량은 $L$), 단일-mode 꼬리에선 fiteq 분모 $1-Q_p\,\dd\Theta/\dd Q\!\approx\!1$
이라 ICA 꼬리 $\approx(Q_p/C_\bg)(\Theta_0/L)e^{-(q-q_a)/L}$ 로 닫힌다(전압축이면 $L_\varphi$, 식~\eqref{eq:Lphi}).
\textbf{(2) 유효 엔탈피(prefactor$\cdot U_j(T)$ 보정).} $1/L\propto k_j=k_0(T)e^{-(G-\chi_j\mathcal A_j)/RT}$ 이고
$k_0=(k_BT/h)\kappa$(Eyring, 식~\eqref{eq:kact})에 \emph{T-선형 prefactor} 가 있어, $\ln(1/L)$ 를 그대로
$1/T$ 로 회귀하면 prefactor 의 $\ln T$ 항이 기울기를 오염시킨다(순수 $\Delta H_a$ 가 아니다). 표준 Eyring 선형화
($\ln(k/T)$ 대 $1/T$)대로 prefactor 를 보정하고, 또한 $\mathcal A_j=s_{\phi,j}F(V_{n,\drive}-U_j(T))$ 의 $U_j(T)$
T-의존 때문에 $\chi_j\mathcal A_j(T)/RT$ 가 $1/T$ 에 \emph{비선형}으로 섞이므로 이를 점별 기지값으로 먼저 빼낸다. 즉
\begin{equation*}
y(T)\equiv\ln\!\frac{1}{L\,T}-\frac{\chi_j\,\mathcal A_j(T)}{RT}\quad\text{를}\quad \frac1T\ \text{로 회귀} \;\Rightarrow\; \text{기울기}=-\frac{\Delta H_a}{R}\ (\text{순수}).
\end{equation*}
$\chi_j\mathcal A_j(T)$ 는 각 $T\cdot$측정점에서 (0)$\cdot$(3)단계 기지값으로 계산한다($\chi_j$ 는 (3)에서 먼저 얻으므로
비순환). prefactor 의 $\kappa(T)$ 약한 T-의존을 무시하는 근사는 $\Delta H_a/(RT)\gg1$($\Delta H_a\gtrsim40$ kJ/mol)
에서 BOUNDED 다. (보정 전 기울기 $-(\Delta H_a-\chi_j\mathcal A_j)/R$ 는 \emph{유효} 엔탈피로만 해석한다.)
(부록 \S\ref{sec:ch6_arrhenius} 식~\eqref{eq:ch6_arrhenius} 의 raw $\ln(1/L)$ 분해에 들어 있는 $+\chi_j\mathcal A_j/RT$ 를
여기서 빼 순수화한 것이며, 보정 전 유효 기울기 $-(\Delta H_a-\chi_j\mathcal A_j)/R$ 는 식~\eqref{eq:ch6_arrhenius_slope} 와 동일.)
\textbf{(3) $\chi_j$.} \emph{peak 근방}($\mathcal A_j=s_{\phi,j}F(V_{n,\drive}-U_j)$ 가 유한한 영역)에서 내부 전위
$V_{n,\drive}=V_n$(0단계 환산)를 가로축으로 $\partial\ln L/\partial V_{n,\drive}=-\chi_j s_{\phi,j}F/RT$ 기울기에서
$\chi_j$ 를 얻는다. \textbf{비순환 보장}: $\mathcal A_j=s_{\phi,j}F(V_{n,\drive}-U_j)$ 는 각 데이터점에서 (0)단계 환산
$V_n$ 과 (0)단계 고정 $U_j(T)$ 로 \emph{계산되는 기지값}(자유 파라미터 아님)이다. 따라서 순서는 비순환이다 — (3)에서
$\chi_j$ 를 먼저 얻고, 그 $\chi_j$ 로 (2)의 종속변수 $y(T)$ 에서 점별 보정항 $\chi_j\mathcal A_j(T)/RT$ 를 \emph{빼서} 순수
$\Delta H_a$ 를 분리한다((2)↔(3) 닭-달걀 아님). \textbf{즉 $\{L,\Delta H_a,\chi_j\}$ 가 단일 지수 꼬리 피팅
+ 식별성 순서로 나온다}(필요한 근거가 모두 본 장 안에 있다: 근사식~\eqref{eq:simplefit}, 식별성 \S\ref{sec:falsify}). 꼬리가 단일 지수에서
벗어나(stretched) 휘면 그때 비로소 spectrum(식~\eqref{eq:spectrum})$\cdot$자기참조(식~\eqref{eq:closed})가
필요하다 — \emph{단일 지수 근사의 실패 자체가 분포$\cdot$결합의 신호}다(반증, \S\ref{sec:falsify}).

\textbf{C-rate 기준식(0.2C anchor).} 전류 의존 분극을 저율 기준 곡선에 상대화하면, $0.2$C($I_{0.2C}>0$)를 anchor 로
\begin{equation}
V_{n,\app}(q;|I|)\approx V_{n,0.2C}(q)+s_I\big[\,|I|\,R_n(q,T,|I|)-I_{0.2C}R_n(q,T,I_{0.2C})\,\big],
\label{eq:02c}
\end{equation}
전류 의존이 약하면 단순형 $V_{n,\app}\approx V_{n,0.2C}+s_I(|I|-I_{0.2C})R_n(q,T)$ 다 — $0.2$C 에서 상변이 지연이
충분히 작거나 그 지연을 기준 곡선에 흡수해도 되는 경우의 실용 근사. 여기에 식~\eqref{eq:simplefit} 의 $L(|I|)$ 단축
($\partial\ln L/\partial V_{n,\drive}<0$)을 더하면 C-rate 가 커질수록 꼬리가 짧아지고 peak 겹침이 심해지는 관측이
정량 설명된다. \emph{피팅 절차와의 연결}: 식~\eqref{eq:02c} 는 여러 $|I|$ 의 $V_{n,\drive}$ 를 공통 저율 기준으로
정렬해 위 (3)단계 $\chi_j$ 회귀의 가로축 입력을 일관되게 만드는 데 쓴다.

\textbf{초기값(EMG 등 경험식).} 비선형 피팅의 \emph{초기값}으로 $\dd Q/\dd V=B(V)+\sum_j A_j\,g_\mathrm{EMG}
(V;\mu_j,\sigma_j,\lambda_j)$ 같은 peak 함수를 쓸 수 있으나 주 모델이 아니다 — 꼬리는 $k_j,\rho_G,R_n$$\cdot$전하
보존의 결합 결과이며, EMG 의 $\lambda_j$(지수 꼬리율)는 식~\eqref{eq:simplefit} 의 $1/L$ 의 경험적 대응물로만 본다.

% ====================================================================
\subsection{온도 의존 꼬리의 분해와 반증(falsification)}\label{sec:falsify}
% ====================================================================
온도에 따른 꼬리 증감을 한 항으로 단정하지 않고 \textbf{세 계층}의 결합으로 본다.
\begin{longtable}{p{0.24\textwidth} p{0.30\textwidth} p{0.36\textwidth}}
\toprule 계층 & 담당 항 & 해석 \\ \midrule \endhead
평형 열역학 & $U_j(T),w_j(T),Q_\bg(V_n,T)$ & 전이 중심 전위$\cdot$평형 폭$\cdot$배경 용량의 온도 의존 \\
동역학 지연 & $k_j(T,I),\rho_G(g;T)$ & 진행률이 평형을 못 따라가 생기는 꼬리/broadening \\
전압축 분극 & $R_n(q,T,|I|)$ & apparent 전위 이동, 기울기 변화, 분극 peak 변형 \\
\bottomrule
\end{longtable}
\textbf{전위 의존 spectrum shift.} 식~\eqref{eq:L_of_G} 에서 $\partial\ln L/\partial V_{n,\drive}
=-\chi_j s_{\phi,j}F/RT<0$(방전 방향, $\chi_j\ge0$): 구동 전위가 커지면 $L$ 이 작아져 꼬리가 짧아진다 — C-rate/전위가
꼬리를 어떻게 바꾸는지의 정량 예측이다.

\textbf{반증 규칙(forward-prediction 검정; 역산 금지).} 본 이론은 다음으로 \emph{반증 가능}하다.
\begin{enumerate}[label=(N\arabic*),leftmargin=2.4em]
\item \textbf{평형 형태}: 저속 OCV 의 near-peak 감쇠가 $\log(dQ/dV)$ 대 $(V-U)$ 면 logistic, 대 $(V-U)^2$
  면 가우시안 — 어느 쪽도 안 맞으면 평형 ansatz 재검토.
\item \textbf{Arrhenius}: $\ln k_j$(또는 $\ln(1/L)$)대 $1/T$ 가 직선이 아니면 단일 활성화 그림 기각.
\item \textbf{전위 의존}: $\partial\ln L/\partial V_{n,\drive}$ 가 위 예측과 다르면 $\chi_j$ 해석 기각.
\item \textbf{비유일성 차단}: $\rho_G$ 를 꼬리에서 역산하지 않고 독립 입력으로만 — 역산하면 비유일이라
  반증력이 없다.
\item \textbf{비퇴화 — transport$\cdot$charge-transfer 와의 분리(★ 본 이론의 가장 약한 지점).} 꼬리 길이 scaling
  $L\propto|I|$ 과 ``전위 커지면 꼬리 짧아짐''은 \emph{단순 transport 분극$\cdot$charge-transfer 도 공유}하므로,
  그것만으로 barrier-spectrum 에 귀속하면 \emph{퇴화}(오귀속)다. 구별 signature(두 조건 (a)$\cdot$(b)가 \emph{모두}
  필요하며, $k_{j,\lim}$ 이 전위무관일 때만 성립): (a) \emph{전위 의존 Arrhenius slope} — 보정 전 측정 기울기
  $\Delta H_a^{\eff}=\Delta H_a-\chi_j\mathcal A_j$(식~\eqref{eq:simplefit}; 추출은 \S\ref{sec:ch1_simplefit} 절차 (2))
  가 구동 전위에 따라 이동함. \emph{단 이것만으로는 transport 와 여전히 퇴화할 수 있다}: 고체상 확산이 율속이고
  확산계수가 점유율 의존($D=D(\xi)$, thermodynamic factor 경유)이면 transport-limited 유효 엔탈피도 전위에 따라
  이동한다(식~\eqref{eq:kphys} 의 $k_{j,\lim}$ 전위의존). 따라서 (a)는 GITT relaxation transient 의 형상
  ($\sqrt t$ Cottrell=확산 율속 vs 지수=계면 율속)으로 $k_{j,\lim}$ 의 전위의존을 독립 판정해 \emph{배제}한 뒤에만
  barrier-lowering signature 로 유효하다; (b) charge-transfer 과전압 $\eta_{\mathrm{ct}}$ 도 전위 지수의존이라 단일
  전극서 $\chi_j$ 와 공선형이므로 — \textbf{GITT/전류 차단으로 $\eta_{\mathrm{ct}}(R_{\mathrm{ct}})$ 를 독립 분리한
  뒤에만} $\chi_j$ 의 ``배리어 낮춤'' 귀속이 성립한다(단순 rate-broadening \cite{dubarry2022,fly2020} 과의 구별이
  핵심). 이 분리(확산 율속 배제 + $\eta_{\mathrm{ct}}$ 분리) 없이 단일 데이터셋에서 꼬리를 barrier-spectrum 으로
  귀속하는 것은 금지한다.
\end{enumerate}
\textbf{파라미터 제한 순서(식별성).} 모든 파라미터를 한 데이터에서 동시에 자유 피팅하면 공선형(co-linear)이라
분리 불가다. 다음 \emph{순차} 제약으로 식별성을 확보한다: \textbf{(1)} 저속 OCV/GITT(준평형, $|I|\to0$)로
$\{U_j,w_j,Q_{j,\tot},Q_\bg\}$ 를 먼저 고정 — 동역학 꼬리가 사라진 평형 무대를 잡는다. \textbf{(2)} 펄스/전류
차단으로 $R_n$(전압축 분극)을 독립 제한 — $R_n$ 과 $\chi_j$ 는 둘 다 꼬리를 늘려 공선형이므로 $R_n$ 을 먼저
묶는다. \textbf{(3)} 여러 온도의 Arrhenius($\ln(1/L)$ 대 $1/T$)로 $\{\Delta H_a,\Delta S_a,k_0\}$. \textbf{(4)}
여러 전류/전위($\partial\ln L/\partial V_{n,\drive}$)로 $\chi_j$. \textbf{(5)} $\rho_G$ 는 독립 입력(역산 금지,
N4). 이 순서를 어기면 식별성이 무너져 \emph{어떤 적합도라도} 물리 해석이 불가능하다(과적합).

% ====================================================================
\subsection{부록: 자기참조식의 수치 평가와 검증 (Plan B 상세)}\label{sec:ch1_numeric}
% ====================================================================
전체 spectrum$\cdot$자기참조를 쓰는 정밀 평가는 Plan B(직접 $g$-grid)로 푼다. 단순 꼬리는 본문
\S\ref{sec:ch1_simplefit} 의 단일 지수 근사로 충분하나, stretched 꼬리(저온)나 강한 자기참조에서는 본 수치
평가가 기준해다.
\textbf{스킴.} (i) 배리어 격자 $\{g_m\}_{m=1}^{N_g}$ 를 관심 $T\cdot$전위에서 $L(g_m)$(식~\eqref{eq:L_of_G})이
관측 윈도를 덮도록 — 특히 large-$L$(긴 꼬리) 영역 — 잡는다. (ii) mode 별 ODE
$\dd\xi_j(g_m)/\dd t=k_j(g_m)[\xi_{\eq,j}-\xi_j(g_m)]$ 를 시간 적분하되, 매 시점 전하 보존
(식~\eqref{eq:charge_balance})의 대수 제약으로 $V_n$ 을 root-find 한다 — 미분변수 $\xi_j(g_m)$ + 대수 제약
$V_n$ 의 \emph{index-1 DAE} 이며, 제약을 한 번 미분하면 $\partial Q_\bg/\partial V_n\ge\epsilon_Q>0$
(식~\eqref{eq:monotone})이 $\dd V_n/\dd t$ 를 유일하게 주므로 index 가 정확히 1 이다. (iii)
$\xi_j(t)=\sum_m w_m\rho_G(g_m)\xi_j(g_m,t)$ 로 평균.
\textbf{근거 기반 tolerance(임의 상수 아님).} root-find 종료는 전하 보존 잔차 $<\delta_q Q_\cell$($\delta_q$=좌표
분해능), 시간적분은 $\xi\in[0,1]$ 상대오차와 측정 ICA 잡음 floor 에서 잡는다 — cap/clip 이 아닌 측정$\cdot$물리
scale 기준.
\textbf{Plan A 승격 검증.} 닫힌형 $\Theta_\mathrm A$(식~\eqref{eq:closed})는 Plan B 대비
$\varepsilon=\lVert\Theta_\mathrm A-\Theta_\mathrm B\rVert/\lVert\Theta_\mathrm B\rVert$ 를 운용 $T\cdot$rate$\cdot$
tail-window 에서 측정해 $\varepsilon\le\varepsilon_{\mathrm{tol}}=\max(\varepsilon_{\mathrm{disc}},\varepsilon_{\mathrm{noise}})$
(기준해 이산화 오차 $\vee$ 측정 잡음)일 때만 승격한다. \emph{특정 수치값}은 데이터 의존이라 미리 고정하지 않는다(FLAGGED). 넓은
spectrum(저온 stretched)에서 ratio 근사가 가장 열화하므로(\cite{jcp2017} 자기명시) 이 검증이 필수다.
\textbf{격자 수렴.} $\lVert\Theta_{N_g}-\Theta_{2N_g}\rVert/\lVert\Theta_{2N_g}\rVert\le\varepsilon_{\mathrm{disc}}$
(Plan A 승격에 쓰는 이산화 오차 기준 $\varepsilon_{\mathrm{disc}}$ 와 동일값)가 될 때까지 $N_g$ 를 배가한다 —
정성적 ``수렴할 때까지''가 아니라 위 $\varepsilon_{\mathrm{tol}}$ 체인과 닫힌 정량 종료조건이다. solver 코드 구현
자체는 본 이론 범위 밖(P1; 이론식$\cdot$검증 원리만 제시).

% ====================================================================
% 부록 B — 유도식의 실데이터 피팅 워크플로 (구 Chapter 6 흡수: Ch1 자기완결)
% ====================================================================
\subsection*{\normalfont\Large 부록 B — 유도식의 실데이터 피팅 워크플로 (구 Chapter 6)}
\label{sec:ch1_fitting_practice}
\addcontentsline{toc}{subsection}{부록 B — 실데이터 피팅 워크플로}
\noindent\emph{아래 절들은 구 Chapter 6 ``피팅 실무 부록''을 본 장에 흡수한 것이다(Ch.6 해체). 본문 \S1--\S15 의 이론과 심플 근사식(\S\ref{sec:ch1_simplefit}) 위에서, 실데이터 피팅을 운영하는 식별성$\cdot$수치 해법(index-1 DAE)$\cdot$순차 제약(저속 OCV$\to$GITT$\to$다온도 Arrhenius$\to$C-rate)$\cdot$반증 워크플로를 담는다. 가정 원장 AL-60$\sim$69 를 포함한다. 발열$\cdot$히스테리시스 피팅은 각각 Ch.2$\cdot$Ch.4, Ch.5 의 정의식을 참조한다.}
\subsection{부록 B 도입 — 본문 식과 피팅 파라미터 정리}\label{sec:ch6_inherit}
% ====================================================================
본 장은 Chapter 1--5(RB 재구성본)의 다음을 \textbf{수정 없이} 입력으로 받는다(프로젝트 검수 P3-6, P5). 각 식
번호는 앞 장의 boxed 결과를 가리키며, 본 장은 그 위에 \emph{피팅 방법론} 계층을 적층한다(앞 장 본문 기본식을
고쳐 쓰지 않는다).
\begin{linkbox}
\textbf{(L1) 전하 보존식(중심 상태식, Ch1).} $Q_\cell\,q=Q_\bg(V_n,T)+\sum_{j=1}^{N_p}Q_{j,\tot}\xi_j$. $V_n$ 은
그 해(외부 입력 아님); 단조성 $\partial Q_\bg/\partial V_n\ge\epsilon_Q>0$.\\
\textbf{(L2) 진행률 동역학(Ch1$\cdot$Ch3).} $\dfrac{\dd\xi_j}{\dd t}=k_j[\xi_{\eq,j}(V_n,T)-\xi_j]$ 또는 forward/
backward 형 $r_j^+(1-\xi_j)-r_j^-\xi_j$; detailed balance $r_j^+/r_j^-=\xi_{\eq,j}/(1-\xi_{\eq,j})$;
$k_j=k_0(T)e^{-(G-\chi_j\mathcal A_j)/RT}$, $\mathcal A_j=s_{\phi,j}n_j^{\eff}F(V_{n,\drive}-U_j)$, 기본 $V_{n,\drive}=V_n$.\\
\textbf{(L3) 피팅 논리식과 평가 순서(Ch1).} $\dfrac{\dd Q}{\dd V_n}=\dfrac{C_\bg(V_n,T)}{1-Q_p\,(\dd\Theta_\mathrm{tail}/\dd Q)}$,
$C_\bg=\partial Q_\bg/\partial V_n$; 평가 순서 S1--S6 (charge-balance root $\to\xi_{\eq}\to k_j\to L(G)\to A_L\to$
kernel integral $\to\dd Q/\dd V_n\to$ 측정축 $\dd Q/\dd V_{n,\app}$).\\
\textbf{(L4) self-consistent 적분식(Ch1).} $\Theta(q)=\Theta_\mathrm{init}(q)+\int_0^q\mathcal K(q,q')\,
\Theta_\eq(V_n(q',\Theta(q')),T)\,\dd q'$, $\mathcal K(q,q')=\int_0^\infty A_L(L)\tfrac1L e^{-(q-q')/L}\dd L$ (Volterra 2종).\\
\textbf{(L5) closure 2-tier(Ch1).} \emph{직접 수치적분}(분포 $g$-grid)$=$항상 보존되는 기준 해법; \emph{Plan A}
$=$Fredholm ratio-substitution closed-form(사용자 PhD, \cite{lee2011,son2013,jcp2017})$=$검증 통과 시 가속 후보.\\
\textbf{(L6) 발열(Ch2$\cdot$Ch4).} $\dot{\mathcal Q}_\irr=\sum_j n_j^{\eff}I_j\eta_j\ge0$, $\dot{\mathcal Q}_\rev=
s^\conv|I|T\,\partial U/\partial T$ 류; calorimetry 부재 시 \textbf{forward prediction} 으로만(피팅 금지).\\
\textbf{(L7) 히스테리시스(Ch5).} signed current $I_s$$\cdot$$q_\stt$, 충전 부호 $s_{\phi,j}^b=-1$, branch target
$\xi_{\tar,j}^b=\xi_{\sstat,j}^b$, $\Delta V_\obs\ne\Delta V_\hys$, loop-area 분리 발열.\\
\textbf{(L8) 반증(Ch1).} N1 Arrhenius 직선성, N2 전위 의존 spectrum shift($\partial\ln L/\partial V_{n,\drive}<0$),
N3 저온 lineshape(평형 broadening vs disorder), N4 $\rho_G$ forward-only(역산 금지).
\end{linkbox}
\noindent\textbf{$\conv,\stt,\hys,\OCV$ 등 약어는 Ch1--5 정의 그대로}이며 본 장은 이를 재정의하지 않는다.
Chapter 1--5 본문 기본식은 수정하지 않는다(P3-6). 본 장은 이들을 \emph{데이터로 제약}하는 절차로 배치한다.

\textbf{피팅 파라미터의 전수 목록(앞 장이 남긴 것).} 본 장이 식별 대상으로 삼는 파라미터는 다음으로 닫힌다(이 외의
자유 knob 추가 금지). 각각 \emph{어느 장이 정의했고 어느 독립 실험이 1차로 제약하는지}를 \S\ref{sec:ch6_identif}
에서 매핑한다.
\begin{longtable}{p{0.20\textwidth} p{0.30\textwidth} p{0.44\textwidth}}
\toprule 파라미터 & 정의 장$\cdot$식 & 물리적 의미(피팅 시 해석) \\ \midrule \endhead
$U_j,\ w_j$ & Ch1 eq:logistic & 전이 중심 전위$\cdot$평형 전이 폭(열역학 무대) \\
$Q_{j,\tot},\ Q_\bg(V_n,T)$ & Ch1 eq:charge\_balance & 전이 용량 기여$\cdot$배경 누적 용량(전하 보존) \\
$\Delta H_{a,j},\ \Delta S_{a,j},\ k_0$ & Ch1 eq:kact (\AL{10}) & 활성화 엔탈피$\cdot$엔트로피$\cdot$prefactor(Eyring) \\
$\chi_j(\equiv\beta_j$; 조건부, 대칭 Marcus$\cdot$정상 target 극한 한정$)$ & Ch1/Ch3 (\AL{11},31) & 배리어 낮춤 민감도 $=$ transfer coefficient \\
$\rho_G(g;T)$ & Ch1 eq:spectrum (\AL{13}) & 배리어 확률밀도(\textbf{독립 입력}; 역산 금지, \AL{16}) \\
$R_n$ & Ch1 eq:Vapp / Ch3 & 음극 apparent 분극 계수($\ne R_\ct\ne R_{n,\eff}$) \\
$k_{j,\lim}$ & Ch1 eq:kphys & 수송$\cdot$확산$\cdot$유한자리 병목(독립 제약 시만) \\
$n_j^{\eff}=RT/(Fw_j)$ & Ch3 (\AL{34}) & effective 전자수(detailed balance 가 강제; 독립 자유 X) \\
$s_{\phi,j}^b,\ \xi_{\tar,j}^b,\ h_j$ & Ch5 (\AL{53}--56) & branch 부호$\cdot$target$\cdot$hysteresis state(충전 시) \\
\bottomrule
\end{longtable}

% ====================================================================
\subsection{보유 데이터와 식별 가능 / 불가 항목의 분리}\label{sec:ch6_data}
% ====================================================================
\textbf{피팅의 첫 단계는 회귀가 아니라 \emph{무엇을 제약할 수 있는지}의 분리}다. 1차 검증 기준은 사용자 보유
\textbf{음극 반쪽셀 GITT, STO--OCV 테이블, $0.1/0.2/0.5/1.0\,\mathrm{C}$ rate series} 다(\AL{63}). 이 데이터는
모델을 \emph{전부} 결정하지 않으므로, 제약 가능/불가 항을 먼저 나눠 \emph{과결정(over-fitting)} 을 차단한다.
\begin{longtable}{p{0.28\textwidth} p{0.34\textwidth} p{0.32\textwidth}}
\toprule 데이터 & 우선 제약 가능 & 직접 확정 곤란 \\ \midrule \endhead
반쪽셀 STO--OCV & $V_{n,\OCV}(q),U_j,w_j,Q_{j,\tot},Q_\bg$ 열역학 골격 & $k_j,\rho_G(g)$,\brk transport limitation \\
반쪽셀 GITT & instantaneous polarization($R_n$ 일부),\brk relaxation signature,\brk OCV 정합 & $k_j$ 단독값,\brk solid diffusion vs phase relaxation 완전 분리 \\
$0.1/0.2\,\mathrm C$ & 저율 quasi-equilibrium proxy,\brk peak 위치/폭 & strict thermodynamic equilibrium \\
$0.5/1.0\,\mathrm C$ & kinetic lag,\brk broadening,\brk transport stress test & pure frozen kinetic limit \\
다온도 series & Arrhenius 기울기($E_a$), $\Delta S_a$ & 단일 온도만으로는 $T$-의존 분리 불가 \\
온도별 수명/full-cell & 정성적 열화 trend & $U_j(T),w_j(T),k_j(T),\Delta S_j$ 정량 분리 \\
\bottomrule
\end{longtable}
\begin{boundbox}
\textbf{데이터 제약(\AL{63}).} 단일 온도 자료로 $U_j(T),w_j(T),k_j(T)$ 의 온도 의존을 분리하지 않는다 — Arrhenius
기울기는 \emph{여러} 온도가 있어야 정의된다(\S\ref{sec:ch6_arrhenius}). calorimetry 없이 Ch2/Ch4 발열 파라미터를
피팅하지 않는다 — 발열 항은 \textbf{forward prediction}(미시$=$거시 $n^{\eff}I_j\eta_j$ 로 예측만)으로 두고
calorimetry 확보 후 검증한다(\AL{67}).
\end{boundbox}
\begin{boundbox}
\textbf{등온 가정 유효 범위.} 등온 피팅은 온도 상승이 충분히 작거나 외부 온도 제어가 강할 때만 유효하다. 고율
자기발열로 $T(t)$ 가 유의하게 변하면 $U_j(T),w_j(T),k_j(T),R_n(T),k_{j,\lim}(T),\rho_G(g;T)$ 가 함께 변하므로
등온 C-rate 자료만으로 kinetic parameter 를 확정하지 않고 Ch2/Ch4 thermal coupling 으로 돌아간다(\AL{67}).
\end{boundbox}

% ====================================================================
\subsection{피팅 평가 순서 S1--S6 의 구현 (Ch1 평가 순서의 피팅 워크플로화)}\label{sec:ch6_S}
% ====================================================================
\textbf{물리.} Ch.1 의 결론식(eq:fiteq)은 평가 순서 S1--S6 으로 \emph{한 점}의 $\dd Q/\dd V_n$ 을 계산한다. 피팅은
이 계산을 파라미터의 함수로 두고, 측정 ICA 와 비교해 residual 을 최소화한다. 본 절은 S1--S6 각 단계가 어떤
\emph{수치 연산}으로 실현되는지를 명시한다 — 이 수치 연산이 곧 피팅의 inner loop(주어진 파라미터에서 모델
ICA 를 한 번 평가하는 forward 계산)이다.

\subsubsection{S1 — 전하 보존식 root 로 $V_n$ 결정 (피팅 inner-loop 의 첫 연산)}\label{sec:ch6_S1}
주어진 진행 좌표 $q$ 와 진행률 $\bm\xi=(\xi_1,\dots,\xi_{N_p})$ 에서 내부 전위 $V_n$ 은 Ch.1 전하 보존식을
$V_n$ 에 대해 푸는 \emph{root finding} 으로 얻는다(외부에서 읽는 값이 아니다, Ch1):
\begin{equation}
F_\cb^\dyn(V_n;q,\bm\xi)\;=\;Q_\bg(V_n,T)+\sum_{j=1}^{N_p}Q_{j,\tot}\xi_j-Q_\cell q\;=\;0.
\label{eq:ch6_dyn_root}
\end{equation}
$\bm\xi$ 는 이 root finding 에서 \emph{독립 입력}(시간 적분이 따로 갱신)이므로 Jacobian 은 배경 용량 기울기뿐이다:
\begin{equation}
\frac{\partial F_\cb^\dyn}{\partial V_n}=\frac{\partial Q_\bg}{\partial V_n}\;\ge\;\epsilon_Q>0.
\label{eq:ch6_dyn_deriv}
\end{equation}
Ch.1 의 단조성 floor(eq:monotone)가 유지되면 해 존재 범위에서 root 는 유일하다(중간값 정리 $+$ 단조성, Ch1).

\textbf{평형 OCV root 와의 구분(피팅에서 중요).} 평형/준평형에서는 $\xi_j=\xi_{\eq,j}(V_n)$ 이므로 root 식이
달라지고 Jacobian 에 $\xi_{\eq}$ 항이 추가된다:
\begin{equation}
F_\cb^\OCV(V_n;q)=Q_\bg(V_n)+\sum_j Q_{j,\tot}\xi_{\eq,j}(V_n)-Q_\cell q=0,\qquad
\frac{\partial F_\cb^\OCV}{\partial V_n}=\frac{\partial Q_\bg}{\partial V_n}+\sum_j Q_{j,\tot}\frac{\partial\xi_{\eq,j}}{\partial V_n}.
\label{eq:ch6_ocv_root}
\end{equation}
\textbf{두 root 의 derivative 가 다르므로(식~\eqref{eq:ch6_dyn_deriv} vs \eqref{eq:ch6_ocv_root}) 같은 Newton
Jacobian 으로 처리하면 안 된다.} 식~\eqref{eq:ch6_ocv_root} 의 $\partial F_\cb^\OCV/\partial V_n$ 은 Ch.2 의
평형 chemical capacitance(OCV 온도계수 유도에 쓴 $\dd Q/\dd V_{n,\OCV}$)와 같은 양이다 — 즉 \textbf{S1 의
평형 극한은 Ch.2 의 평형 미분용량과 자동 정합}하며, 둘이 어긋나면 피팅이 잘못된 것이다(내부 정합 점검).
\begin{verifybox}
\textbf{root-find 종료 기준은 임의 상수가 아니다(\AL{64}, GS-4$'$).} Newton 종료는 ``residual 이 충분히 작다''가
아니라 \emph{전하 단위의 물리 scale} 로 정규화한다: $|F_\cb^\dyn|\le\varepsilon_Q^{\mathrm{root}}\equiv
\delta_q\,Q_\cell$, 여기서 $\delta_q$ 는 진행 좌표 $q$ 의 측정/계산 분해능(예: $\dd q$ 격자 간격)이다. 즉 ``배경
용량 잔차가 한 격자 스텝의 외부 전하보다 작다''는 \emph{측정 분해능 기준}이며, 무근거 숫자가 아니다. 부동소수점
한계상 $\varepsilon_Q^{\mathrm{root}}$ 을 기계 epsilon 아래로 두지 않는다($\sqrt{\varepsilon_{\mathrm{mach}}}$
근방이 Newton 의 통상 floor, \cite{brenan1996}).
\end{verifybox}

\subsubsection{S2 — 평형 target $\xi_{\eq,j}$ (logistic 평가)}\label{sec:ch6_S2}
S1 의 $V_n$ 에서 Ch.1 logistic(eq:logistic) $\xi_{\eq,j}(V_n,T)=[1+\exp(-s_{\phi,j}(V_n-U_j)/w_j)]^{-1}$ 를 평가한다.
피팅에서 $U_j,w_j$ 는 S1 의 $Q_{j,\tot},Q_\bg$ 와 함께 \emph{저속 OCV} 가 1차로 제약한다(\S\ref{sec:ch6_ocvfit}).
$n_j^{\eff}=RT/(Fw_j)$ 는 $w_j$ 가 정해지면 \emph{종속}이며 독립 자유 파라미터가 아니다(Ch3 \AL{34}).

\subsubsection{S3 — 속도상수 $k_j$ (Eyring 평가, 병목 직렬 합성)}\label{sec:ch6_S3}
$\mathcal A_j=s_{\phi,j}n_j^{\eff}F(V_{n,\drive}-U_j)$(기본 $V_{n,\drive}=V_n$), $k_j^{\mathrm{act}}=k_0(T)
e^{-(G-\chi_j\mathcal A_j)/RT}$ 를 평가한다(Ch1 eq:kact). 병목이 유효하면 Ch.1 의 \emph{직렬 시간척도 합}
$1/k_j^{\eff}=1/k_j^{\mathrm{act}}+1/k_{j,\lim}$ 로 합성한다(eq:kphys). \textbf{강구동에서 $k_j^{\mathrm{act}}$ 가
커지는 것을 임의 cap 으로 자르지 않는다}(GS-4); 수치 stiffness 는 \S\ref{sec:ch6_dae} 의 implicit solver 가
흡수한다.

\subsubsection{S4 — 배리어$\to$길이 $L(G)$ 와 spectrum $A_L$}\label{sec:ch6_S4}
Ch.1 의 닫힌 지수 매핑(eq:L\_of\_G) $L(G)=\dfrac{|I|}{Q_\cell k_0(T)}\exp[(G-\chi_j\mathcal A_j)/RT]$ 과 변수변환
밀도 보존(eq:spectrum) $A_L^{\mathrm{prob}}(L)=\rho_G(G(L))\,(RT/L)\,H(L-L_{\min})$ 를 평가한다. \textbf{Heaviside
$H$ 는 Ch.1 변수변환의 자코비안 support 하한($G\ge0\Leftrightarrow L\ge L_{\min}$)에서 \emph{유도}된 정의역
경계}이며, 피팅 편의를 위한 임의 절단이 아니다(Ch1 \S spectrum; GS-4 비위반). 수치적으로는 $g$-grid 의 하한을
$L_{\min}$ 에 맞추어 음의 배리어 영역을 적분에서 제외한다.

\subsubsection{S5 — 꼬리 중첩 $\dd\Theta_\mathrm{tail}/\dd q$ (kernel integral; 자기참조 풀이는 \S\ref{sec:ch6_grid})}\label{sec:ch6_S5}
Ch.1 의 kernel integral(eq:kernel\_integral) $\dfrac{\dd\Theta_\mathrm{tail}}{\dd q}=\int_0^\infty A_L(L)\,\tfrac1L\,
e^{-(q-q_a)/L}\,\dd L$ 를 평가한다. $\xi_j$ 와 $V_n$ 이 전하 보존으로 얽혀 있어 이는 self-consistent Volterra
2종(L4)이며, 그 \emph{직접 수치 풀이가 본 장 closure 의 기준 해법}이다(\S\ref{sec:ch6_grid}--\ref{sec:ch6_closure}).

\subsubsection{S6 — $\dd Q/\dd V_n$ 과 측정축 $\dd Q/\dd V_{n,\app}$}\label{sec:ch6_S6}
Ch.1 결론식(eq:fiteq) $\dd Q/\dd V_n=C_\bg/(1-Q_p\,\dd\Theta/\dd Q)$ 를 평가하고, 측정축으로는 apparent
전위(eq:Vapp) $V_{n,\app}=V_n+s_I|I|R_n$ 로 환산한다: 등온 정전류에서 $\dd V_{n,\app}/\dd q=\dd V_n/\dd q+
s_I|I|\,\partial R_n/\partial q$, DVA 는 그 역수다. \textbf{측정 ICA 는 $\dd Q/\dd V_{n,\app}$ 와 비교}하며, 모델
출력의 $V_n$-축을 직접 데이터에 맞추지 않는다($R_n$ 이 축을 이동시키기 때문 — 식별성 \S\ref{sec:ch6_identif}).

\begin{practicebox}
\textbf{inner-loop 안정화(본문 물리식 아님, GS-4$''$).} 초기값 탐색 단계에서 $w_j>0$, $Q_{j,\tot}>0$,
$0\le\chi_j\le1$ 같은 물리 제약을 만족시키려 $\log$/$\mathrm{logit}$ 같은 bounded reparametrization 을 쓸 수
있다. 이는 \emph{최적화 변수의 변환일 뿐} 모델 물리식이 아니며, 최종 보고에서는 원 파라미터값으로 되돌려
신뢰구간과 함께 제시한다. rate 폭주를 임시 clip 하는 것도 디버깅 한정이며 최종 해석에서 제거한다(물리 병목
$k_{j,\lim}$ 으로 대체).
\end{practicebox}

% ====================================================================
\subsection{피팅 inner-loop 의 수치 엔진 (root-constrained ODE / index-1 DAE)}\label{sec:ch6_dae}
% ====================================================================
\textbf{물리.} 피팅이 한 파라미터 집합에서 모델 ICA 곡선 전체를 얻으려면 $\xi_j(t)$ 의 시간 전개를 풀어야 한다.
정전류 음극 방전에서 $\dd q/\dd t=|I|/Q_\cell$ 이고, 상태변수는 $\xi_j$, $V_n$ 은 전하 보존식이 매 시점 정하는
\emph{algebraic variable} 다. 이는 미분식 1조 $+$ 대수 제약 1조의 \textbf{index-1 DAE}(\AL{60},
\cite{brenan1996,hindmarsh2005})다:
\begin{align}
\frac{\dd\xi_j}{\dd t}&=k_j^{\eff}(V_n,q,|I|,T)\,[\xi_{\eq,j}(V_n)-\xi_j],\label{eq:ch6_ode}\\
0&=F_\cb^\dyn(V_n;q,\bm\xi).\label{eq:ch6_alg}
\end{align}
\textbf{왜 index-1 인가(무비약).} DAE 의 index 는 대수 제약을 미분해 ODE 로 환원하는 데 필요한 미분 횟수다.
식~\eqref{eq:ch6_alg} 을 시간으로 한 번 미분하면 $\dfrac{\partial Q_\bg}{\partial V_n}\dot V_n+\sum_j Q_{j,\tot}
\dot\xi_j-Q_\cell\dot q=0$ 이고, $\partial Q_\bg/\partial V_n\ge\epsilon_Q>0$(식~\eqref{eq:ch6_dyn_deriv})이라
$\dot V_n$ 에 대해 \emph{한 번}에 풀린다 — 따라서 index 가 정확히 1 이다(\AL{60}). 단조성 floor 가 이 풀림의
열쇠이므로, 피팅 중 $\partial Q_\bg/\partial V_n$ 이 0 에 접근하면 DAE 가 ill-conditioned 가 되고 이는
\emph{파라미터가 비물리}라는 신호다(reject; \S\ref{sec:ch6_S1}).

\textbf{시간 적분 절차(정전류; $q(t)$ 명시이므로 root-constrained ODE 로 구현).}
\begin{enumerate}[label=T\arabic*),leftmargin=2.4em]
\item 현재 $q(t),\xi_j(t)$ 에서 식~\eqref{eq:ch6_dyn_root} 로 $V_n(t)$ root-find(S1).
\item $V_n(t)$ 에서 $\xi_{\eq,j}$(S2), $k_j^{\mathrm{act}},k_{j,\lim},k_j^{\eff}$ 또는 $r_j^\pm$(S3) 계산.
\item 식~\eqref{eq:ch6_ode}(또는 forward/backward 형)로 $\dd\xi_j/\dd t$ 계산.
\item implicit solver(BDF 또는 Radau)로 다음 시간으로 진행(\AL{60}, \cite{hindmarsh2005}).
\item 휴지 구간에서는 $q$ 고정, $\xi_j$ relaxation 과 전하 보존 root 를 함께 갱신(Ch2/Ch5 rest; 내부 $I_j\ne0$).
\end{enumerate}
\begin{boundbox}
\textbf{stiff system 처리(cap 아님, GS-4).} 강구동에서 $k_j^{\mathrm{act}}$ 또는 $r_j^+$ 가 매우 커질 수 있다.
이를 arbitrary cap 으로 자르지 않고 \emph{stiff system} 으로 취급한다(implicit BDF/Radau 가 큰 $k_j$ 를 안정적으로
처리; \cite{brenan1996,hindmarsh2005}). 실제 rate 제한은 $k_{j,\lim}$, transport, 유한 전류, site availability
같은 물리적 병목으로 명시한다(Ch1 eq:kphys / Ch3 정합). 즉 \textbf{수치 stiffness 의 해법은 solver 선택이지
모델 변형이 아니다}.
\end{boundbox}
\begin{verifybox}
\textbf{solver tolerance 의 근거(\AL{64}, GS-4$'$).} BDF/Radau 의 상대$\cdot$절대 tolerance $(\mathrm{rtol},
\mathrm{atol})$ 는 임의 상수가 아니라 상태변수의 물리 scale 로 정한다: $\xi_j\in[0,1]$ 이므로 $\mathrm{atol}$ 은
ICA 의 관심 분해능(예: peak 높이 대비 꼬리 기여의 유의 자릿수)으로, $\mathrm{rtol}$ 은 측정 잡음 대비 over-resolve
를 피하는 수준으로 둔다. solver tolerance 를 데이터 잡음보다 훨씬 작게 두는 것은 계산 낭비이고, 훨씬 크게 두면
꼬리(tail-dominated, 작은 기여)를 놓친다 — \textbf{tolerance 는 측정 잡음 floor 와 ICA 분해능 사이}에 둔다.
\end{verifybox}

% ====================================================================
\subsection{closure: self-consistent Volterra 의 직접 수치 풀이 (기준 해법)}\label{sec:ch6_grid}
% ====================================================================
\textbf{물리.} Ch.1 의 spectrum(eq:spectrum)은 배리어 분포 $\rho_G$ 가 만드는 \emph{연속} relaxation-length 분포다.
실제 풀이는 이를 유한 격자로 이산화한다. 각 배리어 격자점 $g_m$ 의 집단이 독립 1차 완화하고, 전하 보존이 매
시점 $V_n$ 으로 결합한다:
\begin{equation}
\frac{\dd\xi_j(g_m,t)}{\dd t}=k_j^{\eff}(g_m,V_n,q,|I|)\,[\xi_{\eq,j}(V_n)-\xi_j(g_m,t)],\qquad
\xi_j(t)=\sum_{m=1}^{N_g}w_m^{(g)}\,\rho_G(g_m)\,\xi_j(g_m,t),
\label{eq:ch6_grid}
\end{equation}
$\sum_m w_m^{(g)}\rho_G(g_m)\to1$ (격자 정규화; quadrature 가중 $w_m^{(g)}$). \textbf{이 직접 $g$-grid 적분이 본
장의 기준 해법}이다 — 주어진 $\rho_G$ 와 완화 closure 하에서 \emph{이산화 오차만 남는 numerically exact reference}
이며, 모든 축약 해법은 이 해에 대해 검증한다. 이는 Ch.1 self-consistent Volterra(L4)의 직접 수치 평가다.
\begin{verifybox}
\textbf{이중 계상 점검(Ch1 boundbox 의 수치 재확인).} 상수 목표 극한($\Theta_\eq=\Theta_0$, 단일 mode
$A_L=\delta(L-L_0)$, $\Theta_\mathrm{init}=0$)에서 식~\eqref{eq:ch6_grid} 의 이산 해는 $\Theta(q)=\Theta_0
(1-e^{-q/L_0})$ 로 Ch.1 single-mode ODE(eq:single\_kernel)와 일치하고 $q\to\infty$ 서 $\Theta_0$ 로 수렴해야
한다. 만약 kernel 을 \emph{목표} 대신 \emph{잔차}에 곱하면 정상상태가 $\Theta_0/2$ 로 어긋난다 — 수치 구현이
Ch.1 Volterra(목표에 곱)와 같은지 이 극한으로 검증한다.
\end{verifybox}
\begin{boundbox}
\textbf{격자 분해능과 stretched tail(\AL{62}).} Ch.1 의 꼬리는 \emph{tail-dominated}(긴 $L$ mode 가 지배)다.
$g$-격자가 큰 $L$(작은 $k_j$, 큰 배리어) 영역을 충분히 덮지 않으면 꼬리를 놓친다. 따라서 격자 상한은 관심
온도/전위에서 $L_{\max}$ 가 관측 윈도를 덮도록 잡고, 격자 수 $N_g$ 는 $\Theta(q)$ 가 $N_g$ 에 대해 수렴할 때까지
늘린다(자기 수렴 점검; 임의 고정값 아님).
\end{boundbox}
\begin{flagbox}
\textbf{독립 완화 closure = Ch1 이 본 장에 위임한 FLAGGED 항목(미이행).} 식~\eqref{eq:ch6_grid} 의 ``각 $g_m$ 집단이
\emph{독립} 1차 완화''는 Ch.1 eq:xi\_dist 의 평균장(mean-field) ansatz 다. Ch.1 은 이를 flagbox 로 정직 표기하며
``흑연 staging 은 nucleation$\cdot$상경계 이동을 동반하는 \emph{협동적} 1차 상전이라 mode 간 결합(Avrami)을 무시한
평균장이며, 그 타당성은 FLAGGED — Ch.3$\cdot$본 부록(§\ref{sec:ch6_reduction})에서 검증''이라 위임했다. 그러나 위임된 \emph{협동적
staging 검증 자체는 본 작업에서 미이행}이다(독립 완화의 정량 타당성은 operando/calorimetry 측정 + 협동 모형 비교를
요하며 본 이론 범위 밖). 따라서 본 절 $g$-grid 의 ``numerically exact reference''는 그 closure \emph{하에서만}
exact 이며 \emph{exactness 는 이산화 오차에 한하고 closure 자체에는 미치지 않는다} — FLAGGED 로 승계한다.
\end{flagbox}

% ====================================================================
\subsection{closure 축약: 허용되는 수학적 가속과 그 검증}\label{sec:ch6_reduction}
% ====================================================================
직접 격자(\S\ref{sec:ch6_grid})는 비용이 크므로 피팅 반복에서 가속이 필요할 수 있으나, 물리 의미를 훼손하지
않고 \emph{기준 해 대비 검증}돼야 한다. 다음 축약은 수학적 근거가 있을 때 허용한다.
\begin{enumerate}[label=\textbf{(\alph*)},leftmargin=2.4em]
\item \textbf{Adaptive quadrature.} $\rho_G(g)$ 가 매끄럽고 support 제한($L\ge L_{\min}$)/tail 빠른 감소면, 모델을
  바꾸지 않고 적분 정확도만 확보하는 수치 적분법이므로 허용(\AL{65}).
\item \textbf{Moment matching.} 분포 폭이 작고 kernel 이 매끄러우면 $k_j(g)$/memory kernel 을 저차 moment 로 근사.
  단 moment closure 가 \emph{tail-dominated} kinetics 를 놓치지 않는지 기준 해와 비교한다(Ch.1 stretched tail 주의).
\item \textbf{Prony/rational approximation.} 연속 relaxation spectrum 을 유한 exponential mode 로 근사(\AL{65}):
\begin{equation}
K_j(t)=\int_0^\infty\rho_G(g)\,k_j(g)\,e^{-k_j(g)t}\,\dd g\;\approx\;\sum_{\ell=1}^{N_r}a_\ell\,e^{-t/\tau_\ell},
\qquad a_\ell\ge0,\ \tau_\ell>0.
\label{eq:ch6_prony}
\end{equation}
\end{enumerate}
\textbf{$a_\ell\ge0,\tau_\ell>0$ 의 근거(무비약).} $K_j(t)$ 는 \emph{양의} 분포 $\rho_G k_j$ 위의 지수 감쇠 평균이라
완전 단조(completely monotone) 함수다; 그 유한 mode 근사가 물리적 완화이려면 각 mode 가 \emph{음이 아닌} 진폭과
\emph{양의} 시간상수를 가져야 한다(음수 가중/허수 time constant 는 진동/발산을 낳아 비물리). 따라서 부호 제약은
피팅 편의가 아니라 \emph{완전 단조성의 보존}이다(\AL{65}; KWW stretched 함수와 완화시간 분포의 등가성
\cite{lindsey1980}, 지수 감쇠의 연속합 \cite{johnston2006}). 정전류에서 시간$\leftrightarrow$좌표 대응은
$t-t'\leftrightarrow(q-q')Q_\cell/|I|$, $\tau_\ell\leftrightarrow L_\ell$ 이므로 식~\eqref{eq:ch6_prony} 는 Ch.1
relaxation-length spectrum 의 유한 mode 표현이다.
\begin{boundbox}
\textbf{축약의 한계(되먹임).} $V_n$ 이 전하 보존으로 $\xi_j$ 와 결합돼 고율/큰 lag 에서 $\xi_j\to V_n\to\xi_{\eq,j}
\to\dd\xi_j/\dd t$ 되먹임이 강해진다. Prony/rational/moment 축약은 \emph{준정상$\cdot$약결합} 근사로만 쓰고, 직접
$g$-grid DAE 대비 오차를 \S\ref{sec:ch6_closure} 의 절차로 측정한다. 되먹임 오차가 허용 범위를 넘으면 기준
해법으로 회귀한다.
\end{boundbox}

% ====================================================================
\subsection{Plan A (Fredholm ratio-substitution closed-form) 의 위치와 ★ 검증 tolerance 의 정직한 정의}\label{sec:ch6_closure}
% ====================================================================
Ch.1 의 \textbf{Plan A}(사용자 PhD ratio-substitution closed-form, \cite{lee2011,son2013,jcp2017})는 문헌적$\cdot$
수학적 기반이 있는 \emph{가속 후보}다(\AL{61}). 본 장은 이를 유지하되, \textbf{직접 $g$-grid 기준 해를 대체하는
기본 모델이 아니라 검증 대상}으로 둔다. Ch.1 의 승격 조건 M1--M5 를 본 장에서 수치적으로 집행한다.
\begin{enumerate}[label=M\arabic*),leftmargin=2.4em]
\item \textbf{kernel mapping.} 흑연 배리어 분포 문제로의 kernel mapping 이 명확하고 적분 안팎 function class 가 같다.
\item \textbf{function class.} ratio-substitution 이 가정하는 분포 class 가 흑연 $\rho_G$ 와 양립한다.
\item \textbf{수렴성.} ratio truncation/Neumann series 의 수렴이 수치적으로 확인된다.
\item \textbf{$\delta$ baseline.} $\rho_G(g)\to\delta(g-g_0)$ 극한에서 단일 배리어 해(Ch.1 single-mode
  $r\propto e^{-(q-q_a)/L_0}$)로 환원된다.
\item \textbf{정량 오차.} 기준 해 대비 상대오차가 검증 tolerance 이하다(아래 ★ 정의).
\end{enumerate}

\subsubsection{★ 검증 tolerance $\varepsilon_\tol$ 의 정직한 재정의 (방향-3: FLAGGED 허위 정정)}\label{sec:ch6_epsilon}
\textbf{문제(Phase A 적발).} consolidated\_ch6 의 M5 는 정량 오차 조건을
\[
\varepsilon=\frac{\lVert\Theta_{\mathrm A}-\Theta_{\mathrm B}\rVert}{\lVert\Theta_{\mathrm B}\rVert}\le\varepsilon_\tol
\]
로 쓰면서 $\varepsilon_\tol$ 을 \emph{established 임계값처럼} 사용했다. 그러나 (i) 그 \emph{값}의 근거가 본문에
없고, (ii) 문서 전체에 정직 표기용 flagbox 가 0건이었다. 이는 GS-3$\cdot$GS-4 위반(근거 없는 수치 임계값)이다.
본 절은 이를 \textbf{근거 있는 수치 tolerance 로 재서술}하되, 정당화되지 않는 특정 숫자는 FLAGGED 로 정직 표기한다.

\textbf{무엇을 비교하는가(정의 명확화).} $\Theta_{\mathrm B}$ 는 직접 $g$-grid 기준 해(\S\ref{sec:ch6_grid}),
$\Theta_{\mathrm A}$ 는 Plan A closed-form 해다. norm $\lVert\cdot\rVert$ 은 \emph{관심 tail-window} 위의 이산
$L^2$ 노름으로 명시한다: $\lVert f\rVert^2=\sum_{i\in\mathcal W}(f(q_i))^2\,\Delta q$, $\mathcal W=$ post-peak 꼬리
구간. window 를 명시하지 않으면 peak 영역의 큰 값이 분모를 지배해 꼬리 오차를 가린다(\textbf{꼬리가 본 이론의
관심사}이므로 window 가 필수다).

\textbf{tolerance 의 \emph{근거 있는} 결정(임의 상수 금지).} $\varepsilon_\tol$ 은 ``작으면 좋다''가 아니라 두 물리
scale 의 \emph{큰 쪽}으로 둔다(\AL{66}):
\begin{equation}
\varepsilon_\tol\;=\;\max\!\Big(\underbrace{\varepsilon_{\mathrm{disc}}}_{\text{기준 해 자체의 이산화 오차}},\
\underbrace{\varepsilon_\noise}_{\text{측정 잡음으로 인한 ICA 불확실도}}\Big).
\label{eq:ch6_epsilon}
\end{equation}
\emph{한 줄씩 근거}: (a) Plan A 를 기준 해 $\Theta_{\mathrm B}$ 보다 더 정확하게 만들 수는 없다 — $\Theta_{\mathrm B}$
자신이 격자 이산화 오차 $\varepsilon_{\mathrm{disc}}$ 를 품기 때문이다. 따라서 $\varepsilon_\tol<\varepsilon_{\mathrm{disc}}$
를 요구하는 것은 무의미하다(기준 해의 정밀도를 넘는 일치를 강요). $\varepsilon_{\mathrm{disc}}$ 는 \S\ref{sec:ch6_grid}
의 격자 자기수렴(격자 배가 시 $\Theta_{\mathrm B}$ 변화)으로 \emph{측정}한다. (b) 모델 예측을 데이터 잡음보다 더
정밀하게 일치시키는 것은 관측 가능한 차이가 아니다 — ICA 의 잡음 유발 불확실도 $\varepsilon_\noise$ 는 측정
전압/전류 잡음을 $\dd Q/\dd V$ 로 전파해 추정한다($dQ/dV$ 는 미분이라 잡음이 증폭됨, \cite{dubarry2022,fly2020}
\AL{63}). 따라서 \textbf{Plan A 가 기준 해와 ``측정으로 구별 불가능한 만큼'' 일치하면 충분}하며, 그 경계가
식~\eqref{eq:ch6_epsilon} 다. 두 scale 중 \emph{큰 쪽}을 쓰는 이유는, 더 작은 쪽까지 맞추려는 것이 위 (a)(b)
어느 한쪽에서 무의미해지기 때문이다.
\begin{flagbox}
\textbf{$\varepsilon_\tol$ 의 \emph{특정 숫자값}은 미확정(\AL{66}).} 식~\eqref{eq:ch6_epsilon} 은 $\varepsilon_\tol$
을 \emph{측정 가능한 두 scale} 로 환원했으나, $\varepsilon_{\mathrm{disc}},\varepsilon_\noise$ 의 \emph{수치}는
사용자 보유 데이터의 실제 잡음 수준$\cdot$격자 수렴 거동에서 산출돼야 한다 — 본 이론 문서가 ``$\varepsilon_\tol=10^{-2}$''
같은 \emph{고정 숫자}를 established 로 제시하지 않는다. consolidated\_ch6 가 $\varepsilon_\tol$ 을 무근거 임계값처럼
쓴 것을 본 절이 \emph{측정 기반 정의}로 교체했고, 구체 값은 데이터 확보 후 결정하는 \textbf{미확정 항}으로 정직
표기한다(GS-3$\cdot$GS-4$'$ 준수). $L^2$ vs $L^\infty$ norm 선택, window 경계 정의도 데이터 의존이라 함께 미확정이다.
\end{flagbox}
\begin{enumerate}[label=F\arabic*),leftmargin=2.4em,start=6]
\item 실패 시 fallback 순서: 직접 grid(\S\ref{sec:ch6_grid}) $\to$ adaptive quadrature $\to$ Prony/rational
  (\S\ref{sec:ch6_reduction}). 즉 \emph{기준 해는 언제나 직접 grid} 이며 Plan A 는 통과 시에만 가속용으로 끼운다.
\end{enumerate}
\begin{boundbox}
\textbf{Plan A 의 broad-kernel 한계(★ \AL{62}, jcp2017 자기명시).} 사용자 논문(\cite{jcp2017}, Eq.~34 직후)은
``ratio-substitution 근사의 정확도는 reaction zone(분포)이 매우 넓어질 때 나빠진다''고 \emph{스스로 경고}한다.
graphite 의 넓은 spectrum($=$저온 stretched tail, \emph{관심 영역})은 바로 그 broad-kernel 영역이라 Plan A 가
\textbf{가장 부정확할 수 있다}. 또한 사용자 논문은 \emph{Fredholm}(고정 구간) ratio closure 이고 graphite 문제는
\emph{Volterra}(인과 상한 $q$, Ch.1 L4)라 자기참조 선형화가 선행돼야 한다. 따라서 \textbf{Plan A 단독 채택 금지;
항상 직접 grid 대비 식~\eqref{eq:ch6_epsilon} 검증.} graphite transition 개수$\cdot$분포 폭$\cdot$hysteresis branch
offset 을 Fredholm correction factor 안에 흡수하지 않는다(식별성 보존).
\end{boundbox}

% ====================================================================
\subsection{파라미터 식별성(identifiability)과 공선형 차단}\label{sec:ch6_identif}
% ====================================================================
\textbf{물리.} 피팅의 핵심 위험은 \emph{여러 파라미터가 같은 곡선 변화를 만들어 서로 흉내 내는} 공선형
(co-linearity)이다. 흑연 ICA 에서 가장 위험한 두 쌍을 먼저 닫고, 일반 규칙을 제시한다.

\subsubsection{$R_n$ vs $\chi_j$ — 분극과 mobility 가속의 공선형}\label{sec:ch6_RnChi}
\textbf{둘 다 꼬리를 늘린다.} 분극 $R_n$ 은 apparent 전위축을 이동시켜(eq:Vapp) ICA peak 를 넓히고, mobility 가속
$\chi_j$ 는 $L=|I|/(Q_\cell k_j)$ 를 통해 꼬리 길이를 직접 바꾼다(Ch.1 ideal-width 형
$\partial\ln L/\partial V_{n,\drive}=-\chi_j s_{\phi,j}F/RT$; 일반형은 $-\chi_j s_{\phi,j}n_j^{\eff}F/RT
=-\chi_j s_{\phi,j}/w_j$, $n_j^{\eff}=RT/Fw_j$). 한 데이터에서 둘을 동시에 자유 피팅하면 공선형이라 분리 불가다(Ch1 \S falsify, Ch3
\AL{39}).

\textbf{분리 전략(순차 제약).} \emph{독립 실험으로 $R_n$ 을 먼저 고정}한다: 펄스/전류 차단(GITT 의 순간 step)은
$|I|\to0$ 직후 전압 도약으로 \emph{순간} 분극을 주며, 이는 mobility 완화(느린 spectrum)와 시간척도가 분리된다.
$R_n$ 을 이렇게 제한한 \emph{뒤에} 같은 모델로 $\chi_j$ 를 본다. 순서를 바꾸면(먼저 $\chi_j$ 자유화) 분극이
$\chi_j$ 로 흡수돼 가짜 가속을 얻는다.
\begin{keybox}
\textbf{세 저항은 다른 물리다(Ch1$\cdot$Ch3$\cdot$Ch4 계승).} apparent voltage-axis 계수 $R_n$ $\ne$ charge-transfer
저항 $R_\ct=RT/(FA_\eff j_{0,n})$(Ch3 small-signal) $\ne$ heat-generation 저항 $R_{n,\eff}$(Ch4 발열). 피팅에서
이 셋을 한 ``$R$''로 합치면 식별성과 발열 예측이 모두 무너진다 — \textbf{$R_n\ne R_\ct\ne R_{n,\eff}$} 를
파라미터 단계에서 분리 유지한다.
\end{keybox}

\subsubsection{$R_n/R_\ct/R_{n,\eff}$ 의 실험적 분리 (Ch3 계승)}\label{sec:ch6_threeR}
$R_\ct$ 는 \emph{소진폭 EIS/Tafel}(small-signal 선형 영역, \cite{bardfaulkner})에서, $R_n$ 은 \emph{ICA 전압축
이동}(apparent), $R_{n,\eff}$ 는 \emph{calorimetry}(발열률 $\div I^2$)에서 각각 독립 제약한다. 한 실험이 셋을 모두
주지 않으므로, 같은 데이터에서 동시에 자유화하지 않는다(\AL{39},\AL{47},\AL{68}).

\subsubsection{일반 식별성 규칙 — 동시 자유화 금지 집합}\label{sec:ch6_noFree}
다음을 \emph{같은} 자료에서 동시에 자유화하면 식별성이 무너진다(Ch1--3 누적 경고의 통합):
\begin{equation}
\{\,R_n,\ k_j^{\eff},\ w_j,\ \rho_G(g),\ Q_\bg,\ Q_{j,\tot},\ k_{j,\lim}\,\}.
\label{eq:ch6_no_free}
\end{equation}
\textbf{순차 제약 사슬(deliverable).} 식~\eqref{eq:ch6_no_free} 의 공선형은 \emph{독립 실험을 순서대로 거는} 것으로만
풀린다. 본 장의 권고 순서는 \S\ref{sec:ch6_ocvfit}--\ref{sec:ch6_crate} 에서 단계별로 실현한다:
\[
\underbrace{\text{저속 OCV}}_{U_j,w_j,Q_{j,\tot},Q_\bg}\ \to\ \underbrace{\text{GITT}}_{R_n,\ \text{완화 signature}}\ \to\
\underbrace{\text{다온도 Arrhenius}}_{\Delta H_a,\Delta S_a,k_0}\ \to\ \underbrace{\text{C-rate series}}_{\chi_j,\ k_{j,\lim}}.
\]
$\rho_G$ 는 어느 단계에서도 관측 꼬리에서 \emph{역산하지 않고}(비유일, \AL{16}) 독립 입력으로만 둔다.

% ====================================================================
\subsection{S1 실현: 저속 OCV 로 열역학 골격 제약}\label{sec:ch6_ocvfit}
% ====================================================================
저속(준평형) OCV/STO--OCV 테이블을 평형 음극 전위 $V_{n,\OCV}(q)$ 의 실측 proxy 로 쓴다. \textbf{이 자료만으로
$U_j,w_j,Q_{j,\tot},Q_\bg$ 가 유일 결정되지는 않으므로} 다음 물리 제약 하에서 피팅한다(과적합 차단):
\begin{enumerate}[label=O\arabic*),leftmargin=2.4em]
\item 총용량 정합: $\sum_j Q_{j,\tot}+\Delta Q_\bg=Q_\cell\cdot\Delta q_{\mathrm{full}}$(Ch1 전하 보존의 적분형).
\item $Q_\bg(V_n)$ 단조성 $\partial Q_\bg/\partial V_n\ge\epsilon_Q>0$ 및 smoothness(Ch1 eq:monotone).
\item transition 개수 $N_p$ 를 필요 이상 늘리지 않음(parsimony; 관측 $N_p\approx3\sim4$, \AL{3}).
\item $Q_{j,\tot}>0$(방전 방향 용량 기여; Ch1).
\item $w_j$ 는 \emph{평형} transition width 이며 kinetic 배리어 분포 $\rho_G(g)$ 와 혼동 금지(Ch1 구분; 단
  $n_j^{\eff}w_j=RT/F$ 정합, Ch3 \AL{34}).
\item 잔차를 임의 background wiggle 로 모두 흡수하지 않음(GS-4; 잔차는 다음 단계가 설명할 신호일 수 있음).
\end{enumerate}
\textbf{평형 형태의 판별(Ch1 N1 의 실현).} 저속 OCV 의 near-peak 감쇠를 $\log(\dd Q/\dd V)$ 대 $(V-U)$(직선이면
logistic, Ch1 eq:logistic) 또는 대 $(V-U)^2$(직선이면 Gaussian-cumulative)로 그려 평형 isotherm 형태를
\emph{판별}한다. \textbf{logistic 이외 형태(erf 등)는 흑연에 대한 열역학적 유도가 없는 FLAGGED ansatz}(Ch1
flagbox)이므로, 데이터가 logistic 을 기각하면 그때 BOUNDED ansatz 로만 도입한다.
\begin{boundbox}
\textbf{저속이 평형은 아니다.} $0.1/0.2\,\mathrm C$ 는 \emph{quasi-equilibrium proxy} 이지 strict thermodynamic
equilibrium 이 아니다(\AL{63}). peak 위치/폭에 잔존 kinetic broadening 이 섞일 수 있으므로, OCV 골격은 GITT
relaxation 끝점(\S\ref{sec:ch6_gittfit})과 교차 검증한다.
\end{boundbox}

% ====================================================================
\subsection{S1--S2 보강: GITT 로 분극과 완화 signature 제약}\label{sec:ch6_gittfit}
% ====================================================================
GITT(galvanostatic intermittent titration)는 짧은 정전류 펄스 $+$ 긴 휴지의 반복으로 \emph{순간 응답}과 \emph{완화
응답}을 함께 준다. \textbf{GITT relaxation 을 곧바로 $k_j$ 의 직접 측정값으로 해석하지 않는다} — 여러 과정이
섞이기 때문이다:
\begin{equation}
\text{GITT relaxation}=\underbrace{\text{solid diffusion}}_{}+\underbrace{\text{electrolyte relaxation}}_{}+
\underbrace{\text{phase/staging relaxation}}_{\text{Ch1 }k_j}+\underbrace{\text{OCV curvature}}_{}+\underbrace{\text{measurement filtering}}_{}.
\label{eq:ch6_gitt_decomp}
\end{equation}
\textbf{권장 절차(순차 분리).}
\begin{enumerate}[label=G\arabic*),leftmargin=2.4em]
\item 펄스 직후 \emph{순간} step 으로 ohmic$\cdot$fast polarization 제약($R_n$ 의 일부; Ch3 $\eta_\ct$ 분리 시작).
\item 장시간 완화 끝점을 저속 OCV(\S\ref{sec:ch6_ocvfit})와 비교해 평형 골격 검증(O 제약과 정합 확인).
\item relaxation curve 를 단일 $k_j$ 로 단정하지 않고 single- vs multi-time-constant 후보를 비교한다(Ch.1
  spectrum 의 다중 $L$ — 여러 시간상수가 보이면 분포 $\rho_G$ 의 직접 증거).
\item C-rate series(\S\ref{sec:ch6_crate})와 결합해 $k_j$, diffusion limitation, transport limitation 을 분리.
\end{enumerate}
\textbf{$\chi_j$ vs $\eta_\ct$ 공선형(Ch1/Ch3 N 의 실현).} GITT 의 순간 step($\eta_\ct$, charge-transfer)과 느린
relaxation(barrier spectrum)을 분리한 \emph{후에만} $\chi_j(=\beta_j$; 조건부, 대칭 Marcus$\cdot$정상 target 극한 한정$)$ 귀속이 성립한다. 즉 \textbf{$R_n$ 과 $\chi_j$
는 GITT 의 시간척도 분리로 비로소 독립 제약}된다(\S\ref{sec:ch6_RnChi} 의 분리 전략을 데이터로 실현).

% ====================================================================
\subsection{S3 핵심: 다온도 Arrhenius 로 활성화 파라미터 제약}\label{sec:ch6_arrhenius}
% ====================================================================
\textbf{물리.} 활성화 파라미터 $\Delta H_{a,j},\Delta S_{a,j},k_0$ 는 \emph{단일 온도} 데이터로는 분리되지 않는다 —
$k_j=k_0 e^{-(\Delta H_a-T\Delta S_a-\chi_j\mathcal A_j)/RT}$ 에서 한 온도는 한 방정식만 주기 때문이다. 여러 온도가
있어야 Arrhenius 직선의 \emph{기울기와 절편}이 따로 정해진다.

\textbf{Arrhenius 분해(무비약).} 꼬리 길이 $L=|I|/(Q_\cell k_j)$ 의 로그를 취하면(Ch1 eq:L\_of\_G)
\begin{equation}
\ln\frac1L\;=\;\ln k_j-\ln\frac{|I|}{Q_\cell}\;=\;\ln k_0-\frac{\Delta H_{a,j}}{RT}+\frac{\Delta S_{a,j}}{R}
+\frac{\chi_j\mathcal A_j}{RT}-\ln\frac{|I|}{Q_\cell}.
\label{eq:ch6_arrhenius}
\end{equation}
\emph{한 줄씩}: $\Delta G_{a,j}=\Delta H_{a,j}-T\Delta S_{a,j}$(Ch1 \AL{10})를 $k_j=k_0 e^{-(\Delta G_{a,j}-\chi_j
\mathcal A_j)/RT}$ 에 넣고 로그를 취했다. \textbf{$\mathcal A_j$ 를 상수로 고정하고}(같은 동작점) $\ln(1/L)$ 을 $1/T$ 에
대해 그리면, $\chi_j\mathcal A_j/(RT)=(\chi_j\mathcal A_j/R)\cdot(1/T)$ \emph{도} $1/T$ 에 선형이라 기울기에 들어간다.
따라서 기울기는 순수 $-\Delta H_{a,j}/R$ 가 \emph{아니라}
\begin{equation}
\frac{\dd\ln(1/L)}{\dd(1/T)}\Big|_{\mathcal A_j}=-\frac{\Delta H_{a,j}-\chi_j\mathcal A_j}{R}=-\frac{\Delta H_{a,j}^{\eff}}{R}
\label{eq:ch6_arrhenius_slope}
\end{equation}
즉 \emph{유효} 활성화 엔탈피 $\Delta H_{a,j}^{\eff}=\Delta H_{a,j}-\chi_j\mathcal A_j$ 를 준다. \textbf{순수 $\Delta H_a$ 를
얻으려면} 구동력을 없앤 극한 $\mathcal A_j\to0$(꼬리 영역 $V_{n,\drive}\to U_j$, 무구동)로 외삽하거나, $\chi_j$ 가 독립
제약된 뒤 $\chi_j\mathcal A_j/R$ 만큼 보정한다. $1/T\to0$ 외삽 절편 $=\ln k_0+\Delta S_{a,j}/R-\ln(|I|/Q_\cell)$
(prefactor$\cdot$엔트로피 묶음; $\ln(|I|/Q_\cell)$ 상수항 포함)다. $k_0$ 와 $\Delta S_a$ 는 절편에서 공선형이므로
$k_0=(k_BT/h)\kappa$ 의 transition-state 절대값을 \emph{인용}(Ch1 \AL{10}, \cite{eyring1935})해 $\Delta S_a$ 를
분리하되, $\kappa(T)$ 미지로 $\ln\kappa$ 가 $\Delta S_a$ 로 흡수되는 잔여 공선형은 BOUNDED 다(\AL{69};
\cite{bardfaulkner}). $\Delta S_a$ 를 reaction entropy(Ch2 \AL{24}, 별개 물리)와 혼동하지 않는다.
\begin{boundbox}
\textbf{$\mathcal A_j$ 고정의 의미(★ 정정).} 흔한 오해와 달리, $\chi_j\mathcal A_j/(RT)$ 항은 $\mathcal A_j$ 를
\emph{고정해도} $1/(RT)=(1/R)(1/T)$ 를 통해 기울기에 기여한다 — 따라서 고정 $\mathcal A_j$ Arrhenius 기울기는 순수
$-\Delta H_a/R$ 가 아니라 \emph{유효} $-(\Delta H_a-\chi_j\mathcal A_j)/R$ 다(식~\eqref{eq:ch6_arrhenius_slope}).
전위를 \emph{바꾸면} $\mathcal A_j$ 자체가 변해 기울기를 더 휘게 하므로, 깨끗한 분리는 한 동작점 기울기에서
$\chi_j\mathcal A_j/R$ 를 보정하거나 $\mathcal A_j\to0$ 외삽으로 한다. 이 $\chi_j\mathcal A_j$ 혼입이 N1
반증(\S\ref{sec:ch6_falsify})의 대상이다. (참고: $\mathcal A_j$ 를 일반형 $n_j^\eff F(V-U)=(RT/w_j)(V-U)$ 로 쓰면
$\chi_j\mathcal A_j/RT=\chi_j(V-U)/w_j$ 로 explicit $1/T$ 가 상쇄되나, 그땐 $w_j(T)$ 의 온도 표류가 같은 자리에 남아
— 어느 convention 이든 순수 $\Delta H_a$ 분리는 $\mathcal A_j\to0$/보정이 필요하다.)
\end{boundbox}
\textbf{Arrhenius 직선성 = 반증 가능성(Ch1 N2).} $\ln(1/L)$ 대 $1/T$ 가 \emph{직선이 아니면} 단일 활성화 그림이
기각된다 — 분포 $\rho_G$ 의 폭이 온도에 따라 유효 기울기를 휘게 하거나(stretched, Ch1 \AL{14}), transport 가
지배하는 경우다. 직선성 자체가 모델의 검정이다.

% ====================================================================
\subsection{S3--S5 검증: C-rate series 로 전위 의존 가속과 꼬리 제약}\label{sec:ch6_crate}
% ====================================================================
$0.1/0.2\,\mathrm C$ 는 quasi-equilibrium proxy(peak 위치/폭), $1.0\,\mathrm C$ 는 kinetic-lag/transport stress
test 다(strict frozen limit 아님). \textbf{검증 순서(순차 제약의 마지막 단계).}
\begin{enumerate}[label=C\arabic*),leftmargin=2.4em]
\item 저속 OCV 로 평형 골격($U_j,w_j,Q_{j,\tot},Q_\bg$) 제약(\S\ref{sec:ch6_ocvfit}).
\item GITT 로 fast polarization($R_n$)$\cdot$relaxation signature 제약(\S\ref{sec:ch6_gittfit}).
\item 다온도로 활성화 파라미터($\Delta H_a,\Delta S_a,k_0$) 제약(\S\ref{sec:ch6_arrhenius}).
\item $0.1/0.2\,\mathrm C$ 에서 ICA/DVA peak 위치$\cdot$폭 검증(S1--S2; 평형 골격 재확인).
\item $0.5/1.0\,\mathrm C$ 에서 \textbf{꼬리$\cdot$broadening$\cdot$shift} 로 $\chi_j$ 제약(S3--S6, kernel integral).
  전위 의존 $\partial\ln L/\partial V_{n,\drive}<0$ 가 C-rate 증가 시 꼬리 단축으로 나타나는지 확인(\AL{62}).
\item 잔차를 $\rho_G(g),k_{j,\lim},R_n$, transport 중 \emph{어느 항으로} 설명 가능한지 독립 제약과 함께 판단
  (한 항으로 모두 흡수 금지; 식~\eqref{eq:ch6_no_free}).
\end{enumerate}
\textbf{이 순서가 식별성을 보장하는 이유.} 각 단계가 \emph{앞 단계에서 고정된 파라미터} 위에서 \emph{새 파라미터
하나(군)} 만 자유화하므로, 식~\eqref{eq:ch6_no_free} 의 공선형 집합이 한 번에 풀리지 않고 차례로 닫힌다 — 이것이
``동시 자유 피팅'' 과 ``순차 식별'' 의 결정적 차이다.

% ====================================================================
\subsection{반증(falsification): forward-only 예측의 데이터 검정}\label{sec:ch6_falsify}
% ====================================================================
\textbf{물리.} Ch.1--5 는 모두 \emph{forward-only}(파라미터$\to$예측) 모델이다. 피팅이 ``곡선을 맞췄다''로 끝나면
어떤 모델도 충분한 자유도로 맞출 수 있어 의미가 없다. \textbf{모델의 가치는 \emph{반증 가능성}}이며, 본 절은 Ch.1
의 반증 규칙 N1--N4 를 \emph{데이터 시험}으로 실현한다.
\begin{loopbox}
\textbf{도입 관찰의 반증 가능 시험(Ch1 N1--N4 의 데이터 실현).} 도입 관찰 ``저온$\cdot$고율서 긴 겹침 꼬리''를
다음으로 시험한다(어느 하나라도 어긋나면 해당 귀속을 \emph{기각}).
\begin{itemize}[leftmargin=1.6em]
\item \textbf{(N1) Arrhenius 기울기 vs $E_D$.} 꼬리 길이의 $\ln(1/L)$ 대 $1/T$ 기울기(식~\eqref{eq:ch6_arrhenius})가
  독립 측정 확산 활성화 $E_D$ 와 \emph{일치}하고 \emph{전위 의존이 없으면}, 꼬리는 barrier-spectrum 이 아니라
  transport 지배다 $\Rightarrow$ barrier-spectrum 귀속 \emph{기각}. C-rate$\times$온도 교차표로 측정.
\item \textbf{(N2) 전위 의존 spectrum shift.} 전위/$|I|$ 증가 시 spectrum 이 short-$L$ 로 이동하는지
  ($\partial\ln L/\partial V_{n,\drive}<0$, Ch1 eq:L\_of\_G). 이동이 없으면 $\chi_j$ 항 \emph{기각}(가속의 전위
  의존이 모델의 핵심 예측).
\item \textbf{(N3) 저온 lineshape.} 저속 OCV near-peak lineshape 이 저온에서 \emph{안 좁아지면}, 평형 broadening
  ($w=RT/F$ 는 저온서 좁아짐, Ch1 eq:dxidV)으로 설명 안 되는 heterogeneous disorder 기여가 존재 $\Rightarrow$
  kinetic 단독 귀속 기각, disorder 와 분리(Ch1 flagbox).
\item \textbf{(N4) $\rho_G$ forward-only.} $\rho_G(g)$ 를 관측 꼬리에서 \emph{역산하지 않고}(비유일, \AL{16})
  독립 근거(입자 크기 분포$\cdot$EIS$\cdot$계산)로 주어 forward 예측만 한다 — 역산하면 비유일이라 반증력이 없다.
\end{itemize}
이로써 \textbf{도입 관찰 $\to$ 인과 사슬(Ch1) $\to$ 발열/속도론/히스테리시스(Ch2--5) $\to$ 피팅식 $\to$ 반증
가능한 데이터 시험(본 부록)} 으로 한 바퀴가 닫힌다.
\end{loopbox}
\begin{boundbox}
\textbf{$\rho_G$ 역산 금지의 피팅 함의(\AL{16}).} 피팅 도구는 잔차를 줄이려 $\rho_G$ 를 자유화하고 싶어 하지만,
stretched 거동에서 $\rho_G$ 로의 역매핑은 비유일(재포획 등 다른 메커니즘도 같은 꼬리)이다. 따라서 $\rho_G$ 를
피팅 변수로 풀면 반증력이 사라진다 — 독립 입력으로 고정하고, 그 입력이 데이터를 \emph{예측}하는지만 본다.
\end{boundbox}

% ====================================================================
\subsection{calorimetry 부재 시 발열 항의 forward-prediction 제한}\label{sec:ch6_heat}
% ====================================================================
\textbf{Ch2/Ch4 의 발열 항을 ICA/전압 데이터만으로 피팅하지 않는다}(\AL{67}). 비가역 발열
$\dot{\mathcal Q}_\irr=\sum_j n_j^{\eff}I_j\eta_j\ge0$(Ch4)과 가역 발열 $\dot{\mathcal Q}_\rev$(Ch2 Bernardi
energy balance \cite{bernardi1985}; 다공성 전극 열모델 \cite{thomasnewman2003})은 \emph{전류$\cdot$과전압$\cdot$
entropy coefficient} 로부터 \textbf{forward 로 예측}만 하고, 그 예측을 \emph{calorimetry} 가 확보됐을 때
검증한다(\AL{67}). 이유: ICA/전압 데이터는 발열을 직접 보지 않으므로, 발열 파라미터를
전압 잔차에 피팅하면 다른 전압 효과(분극 등)와 공선형이 되어 가짜 값을 얻는다(식~\eqref{eq:ch6_no_free} 의
정신을 발열로 확장).
\begin{boundbox}
\textbf{발열의 ICA 되먹임(등온 한계의 신호).} 등온 가정이 깨지면(고율 자기발열) Ch.2 의 되먹임
($\dot{\mathcal Q}\to T\uparrow\to k_j\uparrow\to L\downarrow$ 꼬리 단축, 동시에 $w=RT/F\uparrow$ 로 평형 peak
폭 증가)이 ICA tail 에 두 부호로 작용한다. 따라서 고율 ICA tail 을 단일 부호로 단정하지 않고, 등온이 의심되면
Ch2/Ch4 thermal coupling 으로 돌아간다(\AL{67}).
\end{boundbox}

% ====================================================================
\subsection{충전$\cdot$양극$\cdot$full-cell 로의 확장 조건 (음극 방전 검증 후)}\label{sec:ch6_extension}
% ====================================================================
\begin{enumerate}[label=E\arabic*),leftmargin=2.4em]
\item \textbf{충전 branch.} Ch.5 의 signed current $I_s$ 와 metastable branch target $\xi_{\tar,j}^b=\xi_{\sstat,j}^b$,
  유도된 부호 $s_{\phi,j}^b=-1$ 을 사용한다(\AL{53}--56). branch switching 시 $\xi_j,h_j,V_n$ continuity(재초기화 X),
  branch 별 root finding(\S\ref{sec:ch6_S1})으로 매 시점 $V_n$ 결정.
\item \textbf{히스테리시스 분리 검증.} $\Delta V_\obs\ne\Delta V_\hys$(Ch5 \AL{58}): 관측 loop 폭에서 polarization
  ($R_n^b$; $|I_s|\to0$ 외삽으로 소멸)$\cdot$kinetic lag$\cdot$transport 를 빼고 남는 잔존 분지 차만 true
  hysteresis 후보로 둔다. empirical branch fit 은 metastable free-energy/domain memory 분리 안 되면 \emph{validation
  candidate} 로만 취급(주 결론서 제외, Ch5 forward-only).
\item \textbf{양극.} 양극에 별도 전하 보존식$\cdot$STO--OCV$\cdot$GITT 를 같은 절차로 적용.
\item \textbf{full-cell.} $V_\cell=V_p-V_n-\eta_\loss$ 또는 apparent electrode voltage 표현; loss/polarization
  이중 계상 금지(Ch5). full-cell ICA/DVA 는 양극/음극/branch signed loss/transport/polarization 이 섞인 관측
  신호라 음극 단독과 직접 같지 않다.
\item \textbf{순서.} \textbf{음극 방전이 먼저 식별$\cdot$검증되지 않으면 full-cell 조합 해석으로 넘어가지 않는다.}
\end{enumerate}

% ====================================================================
\subsection{금지되는 피팅 편의와 실무 보조의 격리}\label{sec:ch6_forbidden}
% ====================================================================
\textbf{금지(본문 물리식으로 쓰지 않음, GS-4).}
\begin{enumerate}[label=(\arabic*),leftmargin=2.4em]
\item 강구동 rate 폭주를 물리 의미 없는 cap/clip 으로 자르기(대신 stiff solver $+$ 물리 병목 $k_{j,\lim}$).
\item 잔차를 arbitrary background correction(임의 wiggle)으로 모두 흡수(O6).
\item branch hysteresis 를 단순 charge/discharge offset 으로만 피팅(Ch5 metastable 분리 없이).
\item negative quadrature weight/비물리 time constant 로 relaxation spectrum 맞추기(완전 단조성 위반, \S\ref{sec:ch6_reduction}).
\item Fredholm/Padé(Plan A) correction factor 를 불분명한 fitting knob 으로 사용(\S\ref{sec:ch6_closure}).
\item $\rho_G$ 를 관측 꼬리에서 역산(비유일, N4).
\item 근거 없는 수치 임계값($\varepsilon_\tol$ 등)을 established 로 사용(\S\ref{sec:ch6_epsilon} 의 측정 기반 정의로 대체).
\end{enumerate}
\begin{practicebox}
\textbf{실무 보조(본문 물리식 아님, GS-4$''$).} 초기값 탐색$\cdot$디버깅에서 bounded transform, 임시 clipped rate,
smoothed background 같은 수치 보조를 \emph{임시} 사용할 수 있으나 본 모델 물리식이 아니다. 최종 해석$\cdot$논문용
식$\cdot$검증 결과에서는 제거하고 물리적 제한항(병목 $k_{j,\lim}$) 또는 검증된 축약법으로 대체한다. 피팅 보조가
결과에 남았는지 여부는 \S\ref{sec:ch6_selfcheck} 검수표가 점검한다.
\end{practicebox}

% ====================================================================
\subsection{부록 B 자체 검수표 (구 Chapter 6)}\label{sec:ch6_selfcheck}
% ====================================================================
{\small
\begin{longtable}{p{0.74\textwidth} p{0.16\textwidth}}
\toprule 검수 항목 & 결과 \\ \midrule \endhead
Ch1--5 본문 수정 없이 부록만 적층(P3-6,\brk P5) & 통과(\S\ref{sec:ch6_inherit}) \\
★ 피팅 실무가 주역, 수치해법은 inner-loop solver 로 종속(사용자 5-30) & 통과(\S\ref{sec:ch6_S},\ref{sec:ch6_dae}) \\
피팅 평가 순서 S1--S6 을 피팅 워크플로로 실현 & 통과(\S\ref{sec:ch6_S}) \\
피팅 파라미터 전수 목록 $+$ 정의 장$\cdot$1차 제약 실험 매핑 & 통과(\S\ref{sec:ch6_inherit},\ref{sec:ch6_identif}) \\
★ 식별성: $R_n$ vs $\chi_j$ 공선형 분리(독립 실험 순차) & 통과(\S\ref{sec:ch6_RnChi}) \\
★ $R_n\ne R_\ct\ne R_{n,\eff}$ 분리 유지(Ch1$\cdot$3$\cdot$4 계승) & 통과(\S\ref{sec:ch6_threeR} keybox) \\
동시 자유화 금지 집합 $+$ 순차 제약 사슬(OCV$\to$GITT$\to$Arrhenius$\to$C-rate) & 통과(\S\ref{sec:ch6_noFree}) \\
dynamic root vs OCV root 의 derivative 구분(식~\eqref{eq:ch6_dyn_deriv} vs \eqref{eq:ch6_ocv_root}) & 통과(\S\ref{sec:ch6_S1}) \\
OCV root 분모 $=$ Ch2 평형 chemical capacitance 정합 & 통과(\S\ref{sec:ch6_S1}) \\
전하 보존식 해 존재 조건을 residual 로 흡수 X(비물리 파라미터 reject) & 통과(\S\ref{sec:ch6_S1},\ref{sec:ch6_dae}) \\
index-1 DAE \emph{왜 index-1 인지} 유도(단조성 floor) & 통과(\S\ref{sec:ch6_dae}) \\
강구동 rate 를 cap 으로 자르지 않고 stiff$+$물리 병목으로 처리 & 통과(\S\ref{sec:ch6_dae} boundbox) \\
직접 $g$-grid 를 기준 해법으로; 모든 축약은 대비 검증 & 통과(\S\ref{sec:ch6_grid}) \\
이중 계상(목표 vs 잔차) 수치 점검 극한 & 통과(\S\ref{sec:ch6_grid} verifybox) \\
Prony/rational $a_\ell\ge0,\tau_\ell>0$ 을 완전 단조성으로 유도(편의 아님) & 통과(\S\ref{sec:ch6_reduction}) \\
★ FLAGGED 허위 정정: $\varepsilon_\tol$ 을 측정 기반(이산화 오차$\vee$잡음)으로 재정의 & 통과(\S\ref{sec:ch6_epsilon} eq~\eqref{eq:ch6_epsilon}) \\
★ $\varepsilon_\tol$ 특정 숫자값을 flagbox 로 정직 표기(established 사용 X) & 통과(\S\ref{sec:ch6_epsilon} flagbox) \\
★ root-find/solver tolerance 도 물리 scale$\cdot$잡음으로 근거화(GS-4$'$) & 통과(verifybox \S\ref{sec:ch6_S1},\ref{sec:ch6_dae}) \\
Plan A(Fredholm)를 M1--M5$\cdot\varepsilon$ 가속 후보로 제한;\brk broad-kernel/Volterra 한계 명시 & 통과(\S\ref{sec:ch6_closure}) \\
저속 OCV 골격(O1--O6); $w_j$ vs $\rho_G(g)$ 구분, $n^{\eff}w_j=RT/F$ 정합 & 통과(\S\ref{sec:ch6_ocvfit}) \\
GITT relaxation 을 $k_j$ 직접 측정으로 단정 X; $\chi$ vs $\eta_\ct$ 분리 & 통과(\S\ref{sec:ch6_gittfit}) \\
★ 다온도 Arrhenius 로 $\Delta H_a$ 기울기 식별; $\mathcal A_j$ 고정 조건 명시 & 통과(\S\ref{sec:ch6_arrhenius}) \\
$0.1\mathrm C/1.0\mathrm C$ 를 strict limit 아닌 proxy/stress test 로 & 통과(\S\ref{sec:ch6_crate}) \\
★ 도입 관찰 $\to$ 반증 N1--N4 데이터 시험으로 loop 폐합 & 통과(\S\ref{sec:ch6_falsify} loopbox) \\
calorimetry 없이 발열 파라미터 피팅 X(forward prediction) & 통과(\S\ref{sec:ch6_heat}) \\
피팅 편의(cap/clip/transform)를 본문서 배제,\brk practicebox 격리 & 통과(\S\ref{sec:ch6_S6},\ref{sec:ch6_forbidden}) \\
충전/양극/full-cell 확장을 음극 방전 검증 이후로 보류 & 통과(\S\ref{sec:ch6_extension}) \\
모든 가정 AL-60$\sim$69 인용; 무근거 established 0;\brk FLAGGED 정직 표기 & 통과(\S\ref{sec:ch6_al}) \\
``자명/clearly/쉽게/obviously'' 본문 0건(G-noleap) & 통과 \\
\bottomrule
\end{longtable}
}

% ====================================================================
\subsection{부록 B 가정 원장 (구 Chapter 6, Assumption Ledger, AL-60$\sim$69)}\label{sec:ch6_al}
% ====================================================================
본 장이 사용한 가정과 그 tier$\cdot$근거$\cdot$검증 DOI 다(통합 master 와 정합; AL-64$\sim$69 가 본 장 신규 등재).
{\small
\begin{longtable}{p{0.05\textwidth} p{0.49\textwidth} p{0.12\textwidth} p{0.26\textwidth}}
\toprule AL & 가정 진술 & tier & 근거(검증 DOI) \\ \midrule \endhead
60 & index-1 DAE(root-constrained ODE),\brk BDF/Radau 시간 적분; 단조성 floor 가 index-1 보장 & GROUNDED & brenan1996(10.1137/\brk 1.9781611971224),\brk hindmarsh\brk 2005(10.1145/\brk 1089014.1089020) \\
61 & Plan A closure: Fredholm 2종 ratio-substitution closed-form(사용자 PhD) & BOUNDED & jcp2017(10.1063/\brk 1.5000882),\brk lee2011(10.1063/\brk 1.3565476),\brk son2013(10.1063/\brk 1.4802584) \\
62 & ★ ratio ansatz 의 broad-kernel 열화 $=$ stretched-tail(저온) 영역서 부정확 $\to$ 직접 grid validator 필수 & BOUNDED & jcp2017 Eq.34 자기명시(10.1063/\brk 1.5000882) \\
63 & ICA($dQ/dV$) rate-dependency$\cdot$미분 잡음 증폭$\cdot$best-practice & GROUNDED & dubarry2022(10.3389/\brk fenrg.2022.1023555),\brk fly2020(10.1016/\brk j.est.2020.101329) \\
64 & ★(신규) root-find/solver 종료 tolerance 는 물리 scale($\delta_q Q_\cell$, $\xi\in[0,1]$)$\cdot$기계 epsilon floor 로 정함(임의 상수 아님) & BOUNDED & brenan1996(10.1137/\brk 1.9781611971224); 수치해석 표준 \\
65 & ★(신규) 분포 적분 축약(adaptive quadrature/moment/Prony)은 기준 grid 대비 검증 시만;\brk Prony $a_\ell\ge0,\tau_\ell>0$ 은 완전 단조성 보존 & BOUNDED & lindsey1980(10.1063/\brk 1.440530,\brk KWW$\leftrightarrow$분포),\brk johnston2006(10.1103/\brk PhysRevB.74.184430, 연속합) \\
66 & ★(신규) closure 검증 tolerance $\varepsilon_\tol=\max(\varepsilon_{\mathrm{disc}},\varepsilon_\noise)$; \emph{특정 숫자값은 데이터 의존 미확정} & FLAGGED(값) / BOUNDED(정의) & dubarry2022(10.3389/\brk fenrg.2022.1023555),\brk fly2020(10.1016/\brk j.est.2020.101329); 값은 사용자 데이터서 산출 \\
67 & ★(신규) calorimetry 부재 시 Ch2/Ch4 발열 항은 forward prediction 으로만(전압 데이터에 발열 파라미터 피팅 X) & BOUNDED & bernardi1985(10.1149/\brk 1.2113792),\brk thomasnewman\brk 2003(10.1149/\brk 1.1531194) \\
68 & ★(신규) 식별성: $\{R_n,k_j^{\eff},w_j,\rho_G,Q_\bg,Q_{j,\tot},k_{j,\lim}\}$ 동시 자유화 금지; 순차 제약(OCV$\to$GITT$\to$Arrhenius$\to$C-rate)으로만 식별 & BOUNDED & bard\brk faulkner(ISBN 978-0-\brk 471-\brk 04372-0),\brk fly2020(10.1016/\brk j.est.2020.101329) \\
69 & ★(신규) Arrhenius 분해: $\mathcal A_j$ 고정 기울기 $=-(\Delta H_a-\chi_j\mathcal A_j)/R$ = \emph{유효} 엔탈피($\chi\mathcal A/RT$ 가 $1/T$-선형 기여); 순수 $\Delta H_a$ 는 $\mathcal A_j\to0$ 외삽/보정; $k_0,\Delta S_a$ 절편 공선형($\kappa$ 미지 잔여) & BOUNDED & eyring1935(10.1063/\brk 1.1749604),\brk bard\brk faulkner(ISBN 978-0-\brk 471-\brk 04372-0) \\
\bottomrule
\end{longtable}
}
\textbf{tier 요약.} GROUNDED: AL-60,63,69(분해). BOUNDED: AL-61,62,64,65,67,68,69(분리). \textbf{FLAGGED: AL-66 의
$\varepsilon_\tol$ \emph{특정 숫자값}}(정의는 측정 기반으로 BOUNDED, 값만 데이터 확보 후 확정). 본 장은 무근거
established 수치 임계값을 본문에 두지 않는다(consolidated\_ch6 의 $\varepsilon_\tol$ 허위를 \S\ref{sec:ch6_epsilon}
에서 정정).

% ====================================================================
\subsection{부록 B 의 결론 (구 Chapter 6)}\label{sec:ch6_concl}
% ====================================================================
첫째, \textbf{본 부록(구 Chapter 6)은 유도식(Ch1--5)을 ICA/DVA 데이터에 \emph{피팅하는 방법론}이다}. 피팅 평가 순서 S1--S6
(\S\ref{sec:ch6_S})이 한 점의 $\dd Q/\dd V_n$ 을 주는 forward 계산이고, 이를 측정 ICA 와 비교해 파라미터를
식별한다. \textbf{수치 해법(전하 보존식 root-find, index-1 DAE 시간 적분, self-consistent Volterra 의 직접 grid
풀이)은 이 피팅의 inner-loop solver 로 \emph{종속}}되며 주역이 아니다(사용자 5-30 의도 반영).

둘째, \textbf{피팅의 본질은 회귀가 아니라 순차 식별}이다(\S\ref{sec:ch6_identif}). $R_n$ 과 $\chi_j$ 처럼 같은
꼬리를 만드는 공선형 파라미터는 한 데이터에서 분리되지 않으므로, 저속 OCV$\to$GITT$\to$다온도 Arrhenius$\to$
C-rate series 를 순서대로 걸어 식~\eqref{eq:ch6_no_free} 의 동시 자유화 금지 집합을 차례로 닫는다. $R_n\ne R_\ct
\ne R_{n,\eff}$ 의 세 저항 분리(Ch1$\cdot$3$\cdot$4 계승)가 식별성과 발열 예측을 동시에 지킨다.

셋째, \textbf{전하 보존식 root 의 dynamic/OCV 두 형태는 derivative 가 다르므로 분리}하며(식~\eqref{eq:ch6_dyn_deriv}
vs \eqref{eq:ch6_ocv_root}), 후자의 분모는 Ch.2 평형 chemical capacitance 와 정합한다. index-1 DAE 는 단조성
floor 가 index 를 1 로 만든다는 \emph{유도}로 정당화되며, 강구동 stiffness 는 cap 이 아니라 implicit solver 와
물리 병목으로 처리한다.

넷째, \textbf{closure 의 기준 해법은 직접 $g$-grid}(\S\ref{sec:ch6_grid})이고, adaptive quadrature/moment matching/
Prony-rational/Plan A(Fredholm)는 모두 기준 해 대비 검증될 때만 쓴다. Plan A 는 문헌적 가속 후보이되 흑연의 넓은
spectrum(저온 stretched tail)에서 가장 부정확할 수 있어(\AL{62}, jcp2017 자기명시) 항상 직접 grid 가 validator 다.

다섯째, \textbf{★ consolidated\_ch6 의 $\varepsilon_\tol$ 허위(근거 없는 수렴 임계값, flagbox 0건)를 정정}했다
(\S\ref{sec:ch6_epsilon}). $\varepsilon_\tol$ 을 ``작으면 좋다''가 아니라 기준 해의 이산화 오차와 측정 잡음으로
인한 ICA 불확실도의 \emph{큰 쪽}(식~\eqref{eq:ch6_epsilon})으로 \emph{측정 기반} 재정의하고, 그 \emph{특정 숫자값}은
데이터 의존이라 FLAGGED 로 정직 표기했다. root-find$\cdot$solver tolerance 도 같은 원리로 물리 scale$\cdot$잡음
floor 에서 근거화했다(\AL{64}, GS-4$'$).

여섯째, \textbf{반증(forward-only)이 모델의 가치}다(\S\ref{sec:ch6_falsify}). Arrhenius 직선성(N1), 전위 의존
spectrum shift(N2), 저온 lineshape(N3), $\rho_G$ 역산 금지(N4)로 도입 관찰을 \emph{기각 가능}하게 시험한다.
calorimetry 부재 시 발열 항은 피팅하지 않고 forward prediction 으로만 둔다(\AL{67}).

일곱째, \textbf{도입 관찰(저온$\cdot$고율 긴 겹침 꼬리)} $\to$ 인과 사슬(Ch1) $\to$ 발열/속도론/히스테리시스
(Ch2--5) $\to$ 피팅식과 순차 식별$\to$ 반증 가능 데이터 시험(본 부록 N1--N4)으로 \emph{한 바퀴가 닫힌다}. 음극
방전 모델이 먼저 식별$\cdot$검증된 뒤에야 충전 branch$\cdot$양극$\cdot$full-cell 로 확장한다.

\subsection*{Chapter 2--5 로의 전달}
본 장이 확정한 상태변수($V_n,\xi_j,k_j,\xi_{\eq,j}$)와 식 위에서: Ch.2 는 \emph{전지 가역 반응열}을 평형
전위 온도계수로 잇고, Ch.3 은 $k_j$ 를 forward/backward 반응속도론(Level B)으로 확대하며, Ch.4 는 그 위에서
발열을, Ch.5 는 충방전 히스테리시스를 다룬다. \emph{(기존 Ch.6 의 피팅 실무$\cdot$수치 내용은 본 장
\S\ref{sec:ch1_simplefit}(심플 근사식)$\cdot$\S\ref{sec:planB}$\cdot$\S\ref{sec:planA}(Plan B/A closure)$\cdot$
\S\ref{sec:ch1_numeric}(수치 평가) 및 \emph{부록 B}(\S\ref{sec:ch1_fitting_practice}: 식별성$\cdot$수치 해법$\cdot$순차 제약$\cdot$반증 워크플로)로 흡수했다 — Ch1 자기완결. Ch.6 은 해체되었다.)}



% ====================================================================
% ===== Chapter 2 body (원본 graphite_ica_ch2_rebuilt.tex, byte-verbatim; heading 1단계 demote) =====
% ====================================================================
\clearpage
\section{Chapter 2: 전지 가역 반응열 + 비가역 소산열 (Bernardi 에너지 balance)}\label{chap:ch2}

% ====================================================================
\subsection*{서 — 인과 사슬에서 ``발열''이 들어설 자리}
\addcontentsline{toc}{subsection}{서 — 발열의 자리}
% ====================================================================
Chapter 1 은 하나의 관찰(``저온$\cdot$고율에서 ICA peak 의 꼬리가 길어져 다음 peak 와 겹친다'')을 풀어,
\textbf{열역학이 무대}(내부 전위 $V_n$ 과 평형 target $\xi_{\eq,j}$)를 깔고 \textbf{동역학이 주역}(진행률 $\xi_j$
의 유한속도 완화)이 되어 꼬리를 만든다는 단일 인과 사슬을 세웠다. Chapter 2 는 \emph{똑같은 상태변수}
$\{V_n,\xi_j,k_j,\xi_{\eq,j},U_j,w_j\}$ 위에서 ``그 과정이 어떤 열을 주고받는가''를 잇는다.

발열은 ICA/DVA 곡선 그 자체가 아니다. 그것은 \textbf{(i)} 비등온 조건에서 $V_n,\xi_j,T$ 가 되먹임으로 결합되어
ICA tail 의 온도 의존을 바꾸고, \textbf{(ii)} 안전$\cdot$열관리$\cdot$열화 해석으로 이어지는 연결 고리다. 본 장의
물리 초점은 \textbf{전지 가역 반응열}이다: 평형 전위의 온도계수 $\partial U/\partial T$ 가 reaction entropy 를
주며, 그것이 Bernardi 에너지 balance 의 가역열 $q_\rev$ 로 직결된다. 비가역열 $q_\irr$ 은 과전압 소산으로
항상 $\ge0$(2법칙)이다.

\begin{quote}\itshape
한 문장 요지: 가역 열은 열역학(평형 entropy, 전위 온도계수)이 전담하고, 비가역 열은 동역학(flux$\times$과전압,
소산)이 전담한다 — 이는 Chapter 1 의 ``열역학=방향(target), 동역학=속도(mobility)'' 분업(keystone Level A)의
\textbf{열적 대응}이다. 그래서 두 열을 한 경험항으로 섞지 않는다.
\end{quote}

\noindent\textbf{본 장이 확정하지 \emph{않는} 것.} 완전한 Butler--Volmer, exchange current 와 계면 과전압의
분리, forward/backward rate 와 detailed balance 기반 entropy production 의 \emph{정량 분해}, branch 의존 가역열
(히스테리시스)은 Chapter 3--5 로 넘긴다. 본 장은 \textbf{물리적 연결 장}이며 발열 모델의 폐쇄가 아니다. 또한 본
장은 Chapter 1 과 같이 \textbf{방전(탈리튬화) 방향}을 기본으로 하되, 가역열 부호는 피팅 파라미터가 아니라
convention bookkeeping factor 로 분리한다(\S\ref{sec:ch2_decomp}).

\begin{groundingbox}
\textbf{지배 표준(GS, Ch1 계승).} \textbf{(GS-1)} 모든 식은 물리적 의미를 prose 로 먼저 말한 뒤 수식화하고
중간 단계를 생략하지 않는다(학부 무비약). \textbf{(GS-2)} 결론은 출판 수준 엄밀성. \textbf{(GS-3)} 모든 가정은
통합 Assumption Ledger(\AL{\#}: 논문/교과서 근거 + 검증 DOI)를 인용하며, 근거가 없으면 FLAGGED 로 표기하고
established 로 쓰지 않는다. \textbf{(GS-4)} 임의 clip/cap/softplus/step-function 정의식과 근거 없는 수치
임계값을 물리 모델로 쓰지 않는다. \textbf{부호 규약.} 방전 기준($s_{\phi,j}=+1$), heat-generation positive
convention, 전류 크기 $|I|>0$. 가역열의 방향 부호는 convention factor $s^{\conv}$ 로 분리해 명시 후 고정한다.
\textbf{위계.} 본 장은 \emph{ideal-width 특수형}($n_j^{\eff}=RT/(Fw_j)=1$)을 기본으로 쓰며, 일반형
$n_j^{\eff}\ne1$ 의 비가역열 합은 Chapter 4 로 forward-ref 한다(\S\ref{sec:ch2_irr}; RB\_CHARTER step 3).
\end{groundingbox}

% ====================================================================
\subsection{Chapter 1 에서 전달받는 기준식}\label{sec:ch2_inherit}
% ====================================================================
본 장은 Chapter 1(RB 재구성본)의 다음을 \textbf{수정 없이} 입력으로 받는다(프로젝트 검수 항목 P3-6). 각 식 번호는
Chapter 1 의 boxed 결과를 가리킨다.
\begin{linkbox}
\textbf{(L1) 전하 보존식(무대).} $Q_\cell\,q=Q_\bg(V_n,T)+\sum_{j=1}^{N_p}Q_{j,\tot}\,\xi_j$. 주어진
$(q,T,\{\xi_j\})$ 에서 이 식을 $V_n$ 에 대해 푼 것이 내부 전위이며, 평형 특수해($\xi_j=\xi_{\eq,j}$)가
$V_n=V_{n,\OCV}(q,T)$ 다. 단조성 floor $\partial Q_\bg/\partial V_n\ge\epsilon_Q>0$ 로 해가 유일.\\
\textbf{(L2) 진행률 동역학(주역).} $\dfrac{\dd\xi_j}{\dd t}=k_j(V_n,q,T,I)\,[\xi_{\eq,j}(V_n,T)-\xi_j]$,
$k_j=k_j^{\eff}$(활성화+병목 직렬 합성). 정전류 구간서만 $\dfrac{\dd\xi_j}{\dd t}=\dfrac{|I|}{Q_\cell}\dfrac{\dd\xi_j}{\dd q}$.\\
\textbf{(L3) 평형 진행률(logistic).} $\xi_{\eq,j}(V_n,T)=[1+\exp(-s_{\phi,j}(V_n-U_j(T))/w_j(T))]^{-1}$,
$\partial\xi_{\eq,j}/\partial V_n=s_{\phi,j}\,\xi_{\eq,j}(1-\xi_{\eq,j})/w_j$. ideal 폭 $w_j=RT/F$.\\
\textbf{(L4) 세 전위 구분.} 내부 $V_n$(보존식 해), apparent $V_{n,\app}=V_n+s_I|I|R_n$(분극 포함, 평형
전위 아님), 구동 $V_{n,\drive}$(속도식 구동, 기본 $=V_n$), 평형 $V_{n,\OCV}$(평형 보존식 해).\\
\textbf{(L5) Level A 분업.} $\chi_j\mathcal A_j$ ($\mathcal A_j=s_{\phi,j}F(V_{n,\drive}-U_j)$, $n_j^{\eff}=1$
특수형)는 \emph{방향성 없는} mobility 가속이고, 평형 target $\xi_{\eq,j}$ 의 방향은 열역학이 전담한다.
여기서 $\chi_j\in[0,1]$ 는 무차원 mobility 가속 인자(방향성 없음)로, Ch1 의 $k_j\propto\exp(-(G-\chi_j\mathcal A_j)/RT)$
지수에 들어가는 보정이다. 방향성 barrier lowering($\beta_j$)은 Ch3 Level B.\\
\textbf{(L6) 집계 진행률.} $Q_p\Theta\equiv\sum_jQ_{j,\tot}\xi_j$ ($\Theta\in[0,1]$ 무차원,
$Q_p=\sum_jQ_{j,\tot}$). 따라서 $Q_\cell q=Q_\bg+Q_p\Theta$.
\end{linkbox}

\textbf{신규 기호(Ch1 위에 추가).} 본 장에서만 새로 도입하는 기호를 모은다(Ch1 기호는 위 (L1)--(L6) 과 동일물).
\begin{longtable}{>{\raggedright\arraybackslash}p{0.17\textwidth} >{\raggedright\arraybackslash}p{0.16\textwidth} >{\raggedright\arraybackslash}p{0.59\textwidth}}
\toprule 기호 & 단위 & 의미 \\ \midrule \endhead
$\dot{\mathcal Q}$ & W & 열률(시간율); 첨자 $\src$(열원)$\cdot$$\rev$(가역)$\cdot$$\irr$(비가역)$\cdot$$\loss$(손실) \\
$q_\rev,\ q_\irr$ & W & 가역(entropic, 부호가변)$\cdot$비가역(소산, $\ge0$) 열률(per-transition 또는 합) \\
$\Delta S_{\rxn,j}$ & J/(mol K) & transition $j$ 의 \emph{reaction} entropy(평형량;\brk products$-$reactants) \\
$\Delta S_{a,j}$ & J/(mol K) & \emph{activation} entropy(전이상태;\brk Ch1 $\Delta G_{a,j}=\Delta H_{a,j}-T\Delta S_{a,j}$). $\Delta S_{\rxn}$ 와 \textbf{별개} \\
$\Pi_j$ & V & Peltier-type 계수 $=T\,\partial U_j/\partial T$ \\
$\eta_j$ & V & 과전압(국소 평형 이탈) $=V_{n,\drive}-V_{\eq,j}(\xi_j)$ \\
$V_{\eq,j}(\xi_j)$ & V & $\xi_j$ 가 현재값일 때 진행이 멈추는 국소 평형 전위(logistic 역함수) \\
$J_{n,j}$ & mol/s & transition $j$ 의 진행률 mole flux $=Q_{j,\tot}(\dd\xi_j/\dd t)/F$ \\
$\sigma_j$ & W/(mol K) & per-mode near-equilibrium entropy production rate \\
$g''_j$ & J/mol & 자유에너지 곡률 $\partial^2 G_j/\partial\xi_j^2|_{\eq}>0$ (thermodynamic factor; 안정성) \\
$D_q$ & C/V & 평형 chemical capacitance $\partial Q_\bg/\partial V_n+\sum_jQ_{j,\tot}\partial\xi_{\eq,j}/\partial V_n$ \\
$h_\bg$ & V/K & background \textbf{전위} 온도계수 $(\partial V_{\bg,\eq}/\partial T)_{Q_\bg}$ [V/K] — entropy coefficient(J/(mol\,K)) \emph{아님}; $n=1$ 특수형서 $\Delta S_\bg=-F\,h_\bg$ \\
$C_{\mathrm{th}}$ & J/K & 음극 control volume 유효 열용량 \\
$h,\ A,\ T_\amb$ & W/(m$^2$K),\brk m$^2$,\brk K & 대류계수$\cdot$표면적$\cdot$주위 온도 \\
$s^{\conv}$ & --- & convention bookkeeping factor $\in\{+1,-1\}$ (피팅 X) \\
$\chi_j$ & --- & 무차원 mobility 가속 인자(방향성 없음;\brk Ch1 $k_j\propto\exp(-(G-\chi_j\mathcal A_j)/RT)$ 지수 보정) \\
\bottomrule
\end{longtable}
$F=96485\ \mathrm{C/mol}$(Faraday), $R=8.314\ \mathrm{J/(mol\,K)}$(기체상수). 열률 dot 표기는 시간율 [W].

% ====================================================================
\subsection{평형 전위의 온도계수 $(\partial V_{n,\OCV}/\partial T)_q$}\label{sec:ch2_ocvT}
% ====================================================================
\textbf{물리.} 가역 열은 평형 상태함수의 온도 의존에서 나온다. Chapter 1 의 전하 보존식이 $V_n$ 을 결정하므로,
가역 열의 기초가 되는 전위 온도계수도 \emph{외부 lookup 이 아니라 전하 보존식에서 유도}해야 한다(프로젝트 검수
P3-2). 그래야 Chapter 1 의 ``OCV 는 보존식의 평형 특수해''라는 중심식 흐름이 열 계층에서도 유지된다. 평형
($\xi_j=\xi_{\eq,j}$)에서 식 (L1) 은
\begin{equation}
Q_\cell\,q=Q_\bg(V_n,T)+\sum_{j=1}^{N_p}Q_{j,\tot}\,\xi_{\eq,j}(V_n,T),\qquad V_n=V_{n,\OCV}(q,T).
\label{eq:ch2_eqbal}
\end{equation}

\textbf{기준 용량 가정(\AL{22}).} 본 유도에서 $Q_\cell,Q_{j,\tot}$ 은 온도 무관 기준 용량 매개변수로 둔다.
실제 accessible capacity / transition allocation 의 온도 의존은 별도 용량 보정 또는 후속 열화 모델에서
다룬다(BOUNDED). $Q_{j,\tot}(T)$ 를 허용하면 아래 고정-$q$ 온도미분 분자에 $\sum_j(\partial Q_{j,\tot}/\partial T)\,\xi_{\eq,j}$
항이 추가된다(\S\ref{sec:ch2_ocvT_general}).

\subsubsection{고정 $q$ 에서의 온도 미분 (무비약)}\label{sec:ch2_ocvT_fix}
$q$ 를 고정하고 식~\eqref{eq:ch2_eqbal} 의 양변을 $T$ 로 전미분한다. 좌변 $Q_\cell q$ 는 $q$ 고정이면 상수이므로
그 $T$-미분은 $0$ 이다. 우변에서 $\xi_{\eq,j}$ 는 \emph{두 경로}로 $T$ 에 의존한다: (가) $V_{n,\OCV}(T)$ 를 통한
간접 의존(연쇄법칙)과 (나) $U_j(T),w_j(T)$ 를 통한 $V_n$ 고정 명시적 의존. $Q_\bg(V_n,T)$ 도 같은 두 경로다.
전미분을 한 줄씩 펼치면(이변수 함수의 전미분 $\dd f=\partial_{V_n}f\,\dd V_n+\partial_T f\,\dd T$, 양변을 $\dd T$
로 나누고 $q$ 고정)
\begin{equation}
0=\frac{\partial Q_\bg}{\partial V_n}\left(\frac{\partial V_{n,\OCV}}{\partial T}\right)_q+\frac{\partial Q_\bg}{\partial T}
+\sum_{j}Q_{j,\tot}\!\left[\frac{\partial\xi_{\eq,j}}{\partial V_n}\left(\frac{\partial V_{n,\OCV}}{\partial T}\right)_q
+\left(\frac{\partial\xi_{\eq,j}}{\partial T}\right)_{V_n}\right].
\label{eq:ch2_implicitT_raw}
\end{equation}
$(\partial V_{n,\OCV}/\partial T)_q$ 가 들어간 항을 묶으면(분배법칙)
\begin{equation}
0=\underbrace{\left[\frac{\partial Q_\bg}{\partial V_n}+\sum_{j}Q_{j,\tot}\frac{\partial\xi_{\eq,j}}{\partial V_n}\right]}_{\equiv\,D_q\ (\text{평형 chemical capacitance})}
\left(\frac{\partial V_{n,\OCV}}{\partial T}\right)_q
+\frac{\partial Q_\bg}{\partial T}+\sum_{j}Q_{j,\tot}\left(\frac{\partial\xi_{\eq,j}}{\partial T}\right)_{V_n}.
\label{eq:ch2_implicitT}
\end{equation}
$D_q$ 가 \emph{평형 chemical capacitance} 인 것은 단위가 $\mathrm{C/V}$ 이고($\partial Q_\bg/\partial V_n$ 의
단위 + $Q_{j,\tot}\cdot\partial\xi_{\eq,j}/\partial V_n$ 의 단위 $=\mathrm C/\mathrm V$), $V_n$ 변화에 대한
저장 전하의 응답을 재기 때문이다. $D_q>0$ 은 Chapter 1 의 단조성 floor 에서 보장된다: $\partial Q_\bg/\partial V_n\ge\epsilon_Q>0$
(식 (L1)) 이고 $\partial\xi_{\eq,j}/\partial V_n=s_{\phi,j}\xi_{\eq,j}(1-\xi_{\eq,j})/w_j\ge0$ (방전
$s_{\phi,j}=+1$, $0\le\xi_{\eq,j}\le1$, $w_j>0$). 두 항이 모두 비음이고 첫 항이 양의 하한을 가지므로 $D_q\ge\epsilon_Q>0$.
따라서 식~\eqref{eq:ch2_implicitT} 을 $(\partial V_{n,\OCV}/\partial T)_q$ 에 대해 풀 때 $0$ 으로 나누는 특이점이
없고(implicit function theorem 의 비특이 조건),
\begin{boundbox}
이하 logistic 분자$\cdot$분모는 방전 $s_{\phi,j}=+1$ 특수형으로 닫는다; 일반 부호는
$\xi_{\eq}(1-\xi_{\eq})/w\to s_{\phi,j}\xi_{\eq}(1-\xi_{\eq})/w$ 로 복원(충전 transition 혼재 시 $D_q>0$$\cdot$부호 재검토 필요).
\end{boundbox}
\begin{equation}
\boxed{\;\left(\frac{\partial V_{n,\OCV}}{\partial T}\right)_q
=-\,\frac{\dfrac{\partial Q_\bg}{\partial T}+\displaystyle\sum_{j}Q_{j,\tot}\left(\dfrac{\partial\xi_{\eq,j}}{\partial T}\right)_{V_n}}
{\dfrac{\partial Q_\bg}{\partial V_n}+\displaystyle\sum_{j}Q_{j,\tot}\dfrac{\partial\xi_{\eq,j}}{\partial V_n}}\;.}
\label{eq:ch2_dVocvdT}
\end{equation}
\textbf{차원 검산.} 분자 단위 $=\partial Q_\bg/\partial T$ 의 $\mathrm{C/K}$, 분모 $=D_q$ 의 $\mathrm{C/V}$,
따라서 전체 $=(\mathrm{C/K})/(\mathrm{C/V})=\mathrm{V/K}$ — 전위 온도계수의 올바른 단위다. \textbf{이 식은 평형
상태함수의 온도계수이며 apparent 전위 경로 미분이 아니다}(P3-1): 즉 $(\partial V_{n,\OCV}/\partial T)_q\ne
\dd V_{n,\app}/\dd T$. 후자에는 $R_n(q,T,|I|)$ 의 온도 의존(비가역 분극)이 섞인다(\S\ref{sec:ch2_revOCV}).

\subsubsection{logistic 평형 진행률의 온도 항 (무비약)}\label{sec:ch2_ocvT_logistic}
식~\eqref{eq:ch2_dVocvdT} 의 분자$\cdot$분모를 닫으려면 $\partial\xi_{\eq,j}/\partial V_n$ 과
$(\partial\xi_{\eq,j}/\partial T)_{V_n}$ 가 필요하다. Chapter 1 logistic(식 (L3), 방전 $s_{\phi,j}=+1$)
$\xi_{\eq,j}=[1+\exp(-(V_n-U_j(T))/w_j(T))]^{-1}$ 에서 출발한다. 보조변수 $z_j\equiv(V_n-U_j(T))/w_j(T)$ 로 두면
$\xi_{\eq,j}=(1+e^{-z_j})^{-1}$ 이고, 표준 logistic 항등식 $\dd\xi_{\eq,j}/\dd z_j=\xi_{\eq,j}(1-\xi_{\eq,j})$
(Chapter 1 \S 평형 진행률에서 유도)를 쓴다.

\emph{(i) $V_n$ 미분.} $z_j$ 의 $V_n$ 의존은 $\partial z_j/\partial V_n=1/w_j$ 뿐이므로 연쇄법칙으로
\begin{equation}
\frac{\partial\xi_{\eq,j}}{\partial V_n}=\frac{\dd\xi_{\eq,j}}{\dd z_j}\frac{\partial z_j}{\partial V_n}
=\frac{\xi_{\eq,j}(1-\xi_{\eq,j})}{w_j(T)},
\label{eq:ch2_dxidV}
\end{equation}
이는 Chapter 1 식 (L3) 의 $s_{\phi,j}=+1$ 특수형과 일치한다(검산).

\emph{(ii) $V_n$ 고정 $T$ 미분.} $z_j$ 는 $T$ 에 두 경로($U_j(T)$ 와 $w_j(T)$)로 의존한다. 몫의 미분
$z_j=(V_n-U_j)/w_j$ 를 $T$ 로 미분하면($V_n$ 고정), $U'_j\equiv\dd U_j/\dd T$, $w'_j\equiv\dd w_j/\dd T$ 로
두고
\begin{equation}
\left(\frac{\partial z_j}{\partial T}\right)_{V_n}=\frac{-U'_j\,w_j-(V_n-U_j)\,w'_j}{w_j^2}
=-\frac{U'_j}{w_j}-\frac{(V_n-U_j)\,w'_j}{w_j^2},
\label{eq:ch2_dzdT}
\end{equation}
(첫 등호는 몫미분 $\dd(u/v)=(u'v-uv')/v^2$ 에서 $u=V_n-U_j$, $u'=-U'_j$, $v=w_j$, $v'=w'_j$). 다시 연쇄법칙
$\partial\xi_{\eq,j}/\partial T=(\dd\xi_{\eq,j}/\dd z_j)(\partial z_j/\partial T)$ 로
\begin{equation}
\left(\frac{\partial\xi_{\eq,j}}{\partial T}\right)_{V_n}
=\xi_{\eq,j}(1-\xi_{\eq,j})\left[-\frac{U'_j(T)}{w_j(T)}-\frac{V_n-U_j(T)}{w_j(T)^2}\,w'_j(T)\right].
\label{eq:ch2_dxidT}
\end{equation}
\textbf{물리 해석.} 식~\eqref{eq:ch2_dVocvdT} 에 식~\eqref{eq:ch2_dxidV},\eqref{eq:ch2_dxidT} 를 넣으면
$(\partial V_{n,\OCV}/\partial T)_q$ 가 $Q_\bg$, $U_j$, $w_j$, $\xi_{\eq,j}$ 의 온도 의존이 결합된 결과로
드러난다. 이 결합을 단일 경험 보정항으로 흡수하면 물리 해석성이 약해지므로, 본 장은 결합을 명시적으로 둔다.
\begin{boundbox}
\textbf{ideal width 의 함의(\AL{23}).} ideal limit $w_j=RT/F$ 이면 $w'_j=R/F>0$. 식~\eqref{eq:ch2_dxidT} 의
첫 항 $-U'_j/w_j$ 는 reaction-entropy 기원(방전 $s_{\phi,j}=+1$ 에서; $U'_j=\partial U_j/\partial T$, \S\ref{sec:ch2_entropy})이고, 둘째
항 $-(V_n-U_j)w'_j/w_j^2$ 는 $(V_n-U_j)$ 의 부호에 따라 transition 의 양쪽에서 부호가 바뀐다(중심 $V_n=U_j$
에서 $0$). 단 $w'_j>0$ 은 ideal $w_j=RT/F$ 한정($w'_j=R/F$)이며, non-ideal $w_j(T)$ 에서는 $w'_j$ 부호가
미정이라 둘째 항 부호반전 결론도 ideal 범위 내 한정이다. 실제 graphite 의 $U_j(T)$(entropy profile)는 비단조라 $(\partial V_{n,\OCV}/\partial T)_q$ 가
SOC 에 따라 부호 반전할 수 있다 \cite{reynier2004,thomasnewman2003}(\AL{23}); 이는 측정된 anode entropy
곡선과 정합 검증할 대상이다(BOUNDED — 본 식은 logistic ideal 모델 내에서만 정확).
\end{boundbox}

\subsubsection{기준 용량의 온도 의존(일반형, forward-ref)}\label{sec:ch2_ocvT_general}
\AL{22} 의 기준 용량 가정을 풀어 $Q_{j,\tot}=Q_{j,\tot}(T)$ 를 허용하면, 식~\eqref{eq:ch2_implicitT} 의 명시적
$T$-의존 항에 $\sum_j(\partial Q_{j,\tot}/\partial T)\,\xi_{\eq,j}$ 가 추가되어
\begin{equation}
\left(\frac{\partial V_{n,\OCV}}{\partial T}\right)_q
=-\frac{\dfrac{\partial Q_\bg}{\partial T}+\displaystyle\sum_j\!\left[Q_{j,\tot}\!\left(\dfrac{\partial\xi_{\eq,j}}{\partial T}\right)_{V_n}\!\!+\dfrac{\partial Q_{j,\tot}}{\partial T}\,\xi_{\eq,j}\right]}{D_q}
\label{eq:ch2_dVocvdT_gen}
\end{equation}
가 된다(곱미분 $\partial[Q_{j,\tot}\xi_{\eq,j}]/\partial T|_{V_n}=Q_{j,\tot}\partial_T\xi_{\eq,j}+\xi_{\eq,j}\partial_T Q_{j,\tot}$
에서 둘째 항이 추가됨). 본 장 본문은 식~\eqref{eq:ch2_dVocvdT}(온도 무관 용량) 을 기본으로 쓰고, 용량 보정은
열화 모델로 분리한다.

% ====================================================================
\subsection{전지 가역 반응열의 토대: $\partial U/\partial T=$ reaction entropy}\label{sec:ch2_entropy}
% ====================================================================
가역 반응열을 쓰려면 평형 전위의 온도계수가 \emph{무엇인지}를 열역학 1법칙 수준에서 먼저 못박아야 한다. 그
정체는 reaction entropy 다 — 그리고 그것은 Chapter 1 Eyring rate 의 \emph{activation entropy} 와 별개다(혼동
금지). 이 절이 본 장 가역열 유도의 출발점이다.

\subsubsection{$\partial U_j/\partial T=\Delta S_{\rxn,j}/(s_{\phi,j}F)$ (무비약)}\label{sec:ch2_dUdT}
\textbf{물리.} transition $j$ 의 intercalation half-reaction 의 reaction Gibbs energy $\Delta G_{\rxn,j}$ 는
그 평형 전위 $U_j(T)$ 와 전기화학 평형으로 연결된다. 한 전자 이동(또는 $n_j$ 전자, 본 장 ideal 특수형
$n_j^{\eff}=1$)에 대해 \cite{bardfaulkner,newman}(\AL{21})
\begin{equation}
\Delta G_{\rxn,j}=-\,s_{\phi,j}F\,U_j(T),
\label{eq:ch2_dGrxn}
\end{equation}
($s_{\phi,j}=\pm1$ 방향 지시자, Chapter 1 의 $\mathcal A_j=s_{\phi,j}F(V_{n,\drive}-U_j)$ 와 \emph{동일 부호
convention}; 방전 $s_{\phi,j}=+1$). 여기서 prefactor 는 $s_{\phi,j}F=(\text{부호 }s_{\phi,j})\times(n^{\eff}=1)\times F$
로, 일반형은 $s_{\phi,j}\,n_j^{\eff}\,F$ 다(본 장 ideal 특수형 $n_j^{\eff}=1$). \textbf{주의(혼동 금지):} $\Delta G_{\rxn,j}$(평형 reaction Gibbs energy)
와 Chapter 1 의 $\mathcal A_j$(비평형 driving force, 중심 $U_j$ 기준)는 \emph{별개 양}이다 — $\Delta G_{\rxn,j}$
는 전위 $U_j$ 자체에 비례하고, $\mathcal A_j$ 는 구동 전위가 $U_j$ 에서 벗어난 정도에 비례한다($\Delta G_{\rxn,j}\ne-\mathcal A_j$).

표준 열역학 항등식 $(\partial\Delta G/\partial T)_p=-\Delta S$ 를 식~\eqref{eq:ch2_dGrxn} 에 적용한다. 좌변은
$\partial\Delta G_{\rxn,j}/\partial T=-s_{\phi,j}F\,\partial U_j/\partial T$($s_{\phi,j},F$ 온도 무관), 우변은
$-\Delta S_{\rxn,j}$ 이므로 $-s_{\phi,j}F\,\partial U_j/\partial T=-\Delta S_{\rxn,j}$. 양변을 $-s_{\phi,j}F$
로 나누면
\begin{equation}
\boxed{\;\frac{\partial U_j}{\partial T}=\frac{\Delta S_{\rxn,j}}{s_{\phi,j}F}.\;}
\label{eq:ch2_entropy_coeff}
\end{equation}
\textbf{차원 검산.} 우변 $=[\mathrm{J/(mol\,K)}]/[\mathrm{C/mol}]=\mathrm{J/(C\,K)}=\mathrm{V/K}$($\mathrm{J/C}=\mathrm V$)
— 전위 온도계수의 단위와 일치. 즉 평형 전위의 온도의존이 곧 \emph{reaction entropy} 를 준다. 이는
potentiometric entropy 측정(준평형 OCV 를 여러 $T$ 에서)으로 \emph{ICA 와 독립하게} 얻는 양이며, 흑연
$\mathrm{Li_xC_6}$ 의 $\partial U/\partial T$ 가 staging 에 따라 부호$\cdot$크기가 달라짐이 실측돼 있다
\cite{reynier2004,thomasnewman2003}(\AL{21}).

\begin{linkbox}
\textbf{Chapter 1 과의 접점(전달식).} 식~\eqref{eq:ch2_entropy_coeff} 의 $U_j(T)$ 는 Chapter 1 의 평형
target $\xi_{\eq,j}(V_n,T)$ 의 \emph{중심 전위}이자 driving force $\mathcal A_j$ 의 기준점이다. 따라서
$\partial U_j/\partial T$ 는 Chapter 1 의 $U_j(T)$ 를 \emph{독립 측정으로 제약}한다 — Chapter 1 에서 피팅
파라미터였던 $U_j(T)$ 의 온도 의존이 여기서 calorimetry/potentiometry 와 묶인다. 또한 식~\eqref{eq:ch2_dxidT}
의 첫 항 $-U'_j/w_j$ 가 바로 이 $\partial U_j/\partial T$ 이므로, $(\partial V_{n,\OCV}/\partial T)_q$
(식~\eqref{eq:ch2_dVocvdT})와 transition reaction entropy 가 한 뿌리($U_j(T)$)에서 나옴이 확인된다.
\end{linkbox}

\subsubsection{reaction entropy $\ne$ activation entropy (혼동 금지)}\label{sec:ch2_two_entropy}
Chapter 1 의 Eyring rate \cite{eyring1935}(\AL{10}) $k_j=(k_BT/h)\kappa\exp(-\Delta G_{a,j}/RT)$ 에
$\Delta G_{a,j}=\Delta H_{a,j}-T\Delta S_{a,j}$ 를 대입한다. 지수 안에서 $-\Delta G_{a,j}/RT=
-(\Delta H_{a,j}-T\Delta S_{a,j})/RT=-\Delta H_{a,j}/RT+\Delta S_{a,j}/R$ 로 분해되고, 지수의 합은 지수의
곱이므로
\begin{equation}
k_j=\frac{k_BT}{h}\kappa\,\exp\!\Big(\frac{\Delta S_{a,j}}{R}\Big)\exp\!\Big(-\frac{\Delta H_{a,j}}{RT}\Big),
\label{eq:ch2_eyring_entropy}
\end{equation}
즉 \emph{activation} entropy $\Delta S_{a,j}$ 는 rate 의 \emph{prefactor}($e^{\Delta S_{a,j}/R}$ 온도무관 인자;
$k_BT/h$ 의 약한 $T$ 의존은 별도 존재)에서 나온다 — 전이상태와 reactant 의 entropy 차이다. 반면 $\Delta S_{\rxn,j}$(식~\eqref{eq:ch2_entropy_coeff})
은 \emph{평형}(products$-$reactants) entropy 다.
\begin{boundbox}
\textbf{두 entropy 의 별개성(\AL{24}).} $\Delta S_{a,j}$ 와 $\Delta S_{\rxn,j}$ 은 일반적으로 \emph{다르며}
서로 대입할 수 없다(BOUNDED). 둘은 전이상태의 entropy 를 통해서만(BEP 류 가정 하에서) 부분적으로 연결될 수
있으나, 본 장은 그런 연결을 \emph{가정하지 않는다}. 따라서 $\partial U/\partial T$ 로 $\Delta S_{a,j}$ 를
추정하거나 그 역을 하는 것을 금지한다: $\partial U/\partial T$ 는 reaction entropy 만, rate-prefactor 는
activation entropy 만 준다. (이 구분이 무너지면 Chapter 1 의 spectrum 파라미터 $\Delta S_a$ 와 본 장의 가역열
$\Delta S_\rxn$ 이 오염되어 식별이 불가능해진다.)
\end{boundbox}

% ====================================================================
\subsection{열원 항의 물리적 분해와 control-volume 열수지}\label{sec:ch2_decomp}
% ====================================================================
음극 control volume 내부 열원을 \emph{물리 기원별}로 분해한다(피팅 편의 분해가 아니라 entropy/affinity/overpotential
기원 분해; \cite{bernardi1985,rao1997}, \AL{20}):
\begin{equation}
\dot{\mathcal Q}_{\src,n}=\dot{\mathcal Q}_{\rev,n}+\dot{\mathcal Q}_{\irr,n}+\dot{\mathcal Q}_{\rlx,n}^{\cand}.
\label{eq:ch2_srcdecomp}
\end{equation}
(각 항 [W].) $\dot{\mathcal Q}_{\rlx,n}^{\cand}$ 의 비가역분은 휴지서 식~\eqref{eq:ch2_qirrkin_ext} 로 $\ge0$ 닫히고
($\dot{\mathcal Q}_{\irr,n}$ 과 별개 아님), 가역 entropic 분만 미정이다.
$\dot{\mathcal Q}_{\rev,n}$(가역, 열역학 기원)은 흡열일 수 있으므로 $\dot{\mathcal Q}_{\src,n}$ 이 항상 양수일
필요는 없다(가역열의 부호가변성). $\dot{\mathcal Q}_{\irr,n}$(비가역, 동역학 기원)은 $\ge0$. $\dot{\mathcal Q}_{\rlx,n}^{\cand}$
은 외부 전류가 없거나 작아도 내부 $\xi_j$ relaxation 으로 발생 가능한 \emph{후보} 항이다(\S\ref{sec:ch2_rest};
확정 분해는 Ch3--4). 외부 열손실까지 포함한 control-volume 1법칙 열수지는
\begin{equation}
C_{\mathrm{th}}\frac{\dd T}{\dd t}=\dot{\mathcal Q}_{\src,n}-\dot{\mathcal Q}_{\loss},\qquad
\dot{\mathcal Q}_{\loss}=hA(T-T_\amb),
\label{eq:ch2_balance}
\end{equation}
(좌변 $\mathrm{(J/K)(K/s)=W}$, $\dot{\mathcal Q}_{\loss}$ 단위 $\mathrm{(W/(m^2K))\,m^2\,K=W}$ — 차원 정합).
이것이 Bernardi--Pawlikowski--Newman 의 2항 balance($q_\rev+q_\irr$)에 본 장 후보 확장항 $\dot{\mathcal Q}_{\rlx}^{\cand}$
를 더한 형태다($\dot{\mathcal Q}_{\rlx}^{\cand}$ 는 hypothesis — \texttt{\textbackslash cite} 없음) \cite{bernardi1985,rao1997}.
\begin{longtable}{>{\raggedright\arraybackslash}p{0.27\textwidth} >{\raggedright\arraybackslash}p{0.61\textwidth}}
\toprule 항 & 물리적 의미 \\ \midrule \endhead
$\dot{\mathcal Q}_{\rev,n}$ & 평형 entropy coefficient($\partial V_{n,\OCV}/\partial T$) 또는 transition entropy($\Delta S_{\rxn,j}$)에 따른 가역 열(열역학 기원, 부호가변) \\
$\dot{\mathcal Q}_{\irr,n}$ & overpotential$\cdot$reaction affinity$\cdot$transport limitation$\cdot$internal dissipation 에 의한 비가역 열(동역학 기원, $\ge0$) \\
$\dot{\mathcal Q}_{\rlx,n}^{\cand}$ & 외부 전류가 작아도 내부 $\xi_j$ relaxation 으로 발생 가능한 후보 열 (부호: sign-open(가역분)) \\
\bottomrule
\end{longtable}
\begin{keybox}
\textbf{Ch1 분업의 열적 대응(Level A 계승).} $\dot{\mathcal Q}_{\rev}$ 는 Chapter 1 의 ``방향=열역학(target/entropy)''에,
$\dot{\mathcal Q}_{\irr}$ 는 ``속도=동역학(mobility/affinity)''에 대응한다. Chapter 1 keystone(Level A)에서
$\chi_j\mathcal A_j$ 는 방향성 없는 mobility 가속이고 방향(target)은 열역학이 전담했다 — 그 분업이 열에서는
\emph{가역열=entropy(방향) / 비가역열=소산(속도)} 으로 나타난다. 이 대응이 ``정상 target 극한
$\xi_{\mathrm{ss},j}\to\xi_{\eq,j}$''에서 Level B(Ch3)와 일치한다는 한정은 RB\_CHARTER step 4 와 같다
(가역/비가역의 \emph{정량} 분해는 forward/backward rate 비대칭이 필요하므로 Ch3--4).
\end{keybox}

% ====================================================================
\subsection{가역 열 (1): 평형 전위 온도계수 기반 (Bernardi)}\label{sec:ch2_revOCV}
% ====================================================================
\textbf{물리 + 무비약 유도(\AL{20}).} 평형 준정적 경로에서 한 transition 의 reaction rate 는
$\dot n_j=I_j/(s_{\phi,j}F)$ [mol/s] 다(Faraday 법칙: 전류 $I_j$ 가 $s_{\phi,j}F$ 쿨롱당 1몰 반응을 진행).
entropic(가역) 열률은 reversible 경로의 $T\,\Delta S$ 흐름 $q_\rev=T\,\Delta S_{\rxn,j}\,\dot n_j$ 이고, 여기에
식~\eqref{eq:ch2_entropy_coeff} 의 $\Delta S_{\rxn,j}=s_{\phi,j}F\,\partial U_j/\partial T$ 를 대입하면
$s_{\phi,j}F$ 와 $\dot n_j$ 의 $1/(s_{\phi,j}F)$ 가 상쇄되어
\begin{equation}
\boxed{\;q_{\rev,j}=T\,\Delta S_{\rxn,j}\,\frac{I_j}{s_{\phi,j}F}=I_j\,T\,\frac{\partial U_j}{\partial T}
=I_j\,\Pi_j,\qquad \Pi_j\equiv T\,\frac{\partial U_j}{\partial T}.\;}
\label{eq:ch2_qrev}
\end{equation}
\textbf{차원 검산.} $I_j\,T\,(\partial U_j/\partial T)=\mathrm A\cdot\mathrm K\cdot(\mathrm{V/K})=\mathrm{A\,V}=\mathrm W$
— 열률의 올바른 단위. $\Pi_j=T\,\partial U_j/\partial T$ 가 Peltier-type 계수[V]이며 $q_{\rev,j}=I_j\Pi_j$ 는
전류$\times$Peltier 전위 형이다. \textbf{부호.} $I_j$ 가 부호 있는 전류이므로 $q_{\rev,j}$ 는 부호가변이다:
$\partial U_j/\partial T>0$ 이고 $I_j>0$(한 방향)이면 흡열, 전류 방향이 바뀌면 부호 반전. 이 방향 부호는 피팅
파라미터가 아니라 convention bookkeeping 이다.

정전류 $|I|=|\sum_j I_j|$ 이고 $(\partial V_{n,\OCV}/\partial T)_q$ 가 식~\eqref{eq:ch2_dVocvdT} 로 모든 transition
기여를 집계하므로, $\sum_j I_j\,T\,\partial U_j/\partial T=s_{\rev,n}^{\conv}|I|\,T(\partial V_{n,\OCV}/\partial T)_q$
(\textbf{(1)} 부호전류 $\to s^{\conv}|I|$ 규약, \textbf{(2)} 집계는 식~\eqref{eq:ch2_dVocvdT}). 따라서 음극
control volume 의 reduced form(전류 크기 $|I|$ 와 convention factor 로 표기)은
\begin{equation}
\dot{\mathcal Q}_{\rev,n}^{\OCV}=s_{\rev,n}^{\conv}\,|I|\,T\left(\frac{\partial V_{n,\OCV}}{\partial T}\right)_q,
\label{eq:ch2_revOCV}
\end{equation}
$s_{\rev,n}^{\conv}\in\{+1,-1\}$ 는 음극 전위 부호$\cdot$전류 방향$\cdot$heat-generation-positive 규약을 명시한
뒤 고정하는 bookkeeping factor 다(피팅 X; RB\_CHARTER step 1 의 $q_\rev$ 부호 convention). 여기서
$(\partial V_{n,\OCV}/\partial T)_q$ 는 \S\ref{sec:ch2_ocvT} 의 식~\eqref{eq:ch2_dVocvdT} 로 \emph{전하 보존식
에서 유도}된 값이지 외부 lookup 이 아니다. 차원: $|I|\,T\,(\partial V/\partial T)=\mathrm A\cdot\mathrm K\cdot
\mathrm{(V/K)}=\mathrm W$ \checkmark.

\textbf{금지되는 연결(P3-1).}
\begin{equation}
\dot{\mathcal Q}_{\rev,n}\ \not\propto\ |I|\,T\,\frac{\dd V_{n,\app}}{\dd T}.
\label{eq:ch2_forbidden}
\end{equation}
$V_{n,\app}=V_n+s_I|I|R_n$ 는 평형 전위가 아니므로 그 경로 미분은 reversible entropy coefficient 가 아니다.
apparent 전위의 온도 변화에는 $R_n(q,T,|I|)$ 의 온도 의존(비가역 분극)이 섞여 있어, 이를 가역열로 쓰면 비가역
소산을 가역 entropy 로 오계상한다.
\begin{boundbox}
\textbf{다중 transition 합(\AL{20}).} 여러 transition 이 동시에 진행하면 $I_j$ = transition $j$ 의 partial
current($\sum_j I_j=I$)로 읽고 $\dot{\mathcal Q}_{\rev}=\sum_j I_j\,T\,\partial U_j/\partial T$ 로 합산한다.
단일 우세 transition 극한에서 식~\eqref{eq:ch2_qrev} 로 환원된다. partial current 의 정량 배분은 forward/backward
rate(Ch3)가 필요하므로 본 장에서는 합산 형식만 둔다(BOUNDED).
\end{boundbox}

% ====================================================================
\subsection{가역 열 (2): transition entropy basis ($\xi_j$ 기반)}\label{sec:ch2_revxi}
% ====================================================================
Chapter 1 의 핵심 내부 상태변수는 $\xi_j$ 다. 따라서 상변이/staging transition 의 entropy 변화는 \emph{진행률
시간율} $\dd\xi_j/\dd t$ 를 통해 열 항으로도 표현할 수 있다. 이는 식~\eqref{eq:ch2_revOCV} 의 OCV 계수 방식과
\emph{같은 entropy 정보}를 더 미시적인 수준에서 쓴 것이다(\AL{21}).

\textbf{무비약 유도.} transition $j$ 의 진행률 mole flux 는 $J_{n,j}=Q_{j,\tot}(\dd\xi_j/\dd t)/F$ [mol/s]
다($Q_{j,\tot}\,\dd\xi_j$ 가 진행에 든 전하[C], $F$ 로 나눠 몰수, 시간율). 그 reversible 열률은 $T\,\Delta S_j\,J_{n,j}$
이므로
\begin{equation}
\dot{\mathcal Q}_{\rev,\xi}=\sum_{j=1}^{N_p}s_{\rev,\xi,j}^{\conv}\,T\,\Delta S_j\,J_{n,j}
=\sum_{j=1}^{N_p}s_{\rev,\xi,j}^{\conv}\,T\,\Delta S_j\,\frac{Q_{j,\tot}}{F}\,\frac{\dd\xi_j}{\dd t}.
\label{eq:ch2_revxi}
\end{equation}
$\Delta S_j\equiv\Delta S_{\rxn,j}$ = $j$번째 effective transition 의 방전 방향 진행 mol electron(또는 mol
Li-equivalent) 1몰당 entropy change [J/(mol K)]. \textbf{차원 검산(한 줄씩).}
\begin{equation}
T\,\Delta S_j\,\frac{Q_{j,\tot}}{F}\,\frac{\dd\xi_j}{\dd t}\sim
\mathrm K\cdot\frac{\mathrm J}{\mathrm{mol\,K}}\cdot\frac{\mathrm C}{\mathrm{C\,mol^{-1}}}\cdot\mathrm s^{-1}
=\mathrm K\cdot\frac{\mathrm J}{\mathrm{mol\,K}}\cdot\mathrm{mol}\cdot\mathrm s^{-1}=\mathrm W.
\label{eq:ch2_revxi_dim}
\end{equation}
($Q_{j,\tot}/F$ 의 단위 $\mathrm C/(\mathrm{C\,mol^{-1}})=\mathrm{mol}$; $\xi_j$ 무차원이라 $\dd\xi_j/\dd t$ 는
$\mathrm s^{-1}$.) \textbf{좌표 변환 주의.} 정전류 방전 구간서만 (L2) 의 $\dd\xi_j/\dd t=(|I|/Q_\cell)\dd\xi_j/\dd q$
로 $q$-좌표식으로 바꿀 수 있다:
\begin{equation}
\frac{\dd\xi_j}{\dd t}=\frac{|I|}{Q_\cell}\frac{\dd\xi_j}{\dd q}\qquad(\text{정전류 한정}).
\label{eq:ch2_revxi_q}
\end{equation}
휴지/가변전류 구간에서는 $|I|/Q_\cell$ 이 시간 상수가 아니므로 이 변환을 쓰지 않고 $\dd\xi_j/\dd t$ 를 직접
쓴다(\S\ref{sec:ch2_rest}).

\subsubsection{OCV 계수 방식과 entropy basis 방식의 정합 (double counting 금지)}\label{sec:ch2_double}
$\dot{\mathcal Q}_{\rev,n}^{\OCV}$(식~\eqref{eq:ch2_revOCV})와 $\dot{\mathcal Q}_{\rev,\xi}$(식~\eqref{eq:ch2_revxi})
는 \emph{같은 평형 entropy 정보}를 서로 다른 수준에서 표현한다. 따라서 둘을 독립 발열항으로 동시에 자유롭게 더하면
double counting 이다(P3 검수 항목). 물리 우선 원칙은 둘 중 하나를 선택하거나, 동시에 쓸 때 \textbf{정합 제약}을
부과하는 것이다:
\begin{enumerate}[label=\Alph*),leftmargin=2.0em]
\item 전체 평형 저장 반응을 $(\partial V_{n,\OCV}/\partial T)_q$ 하나로 대표(식~\eqref{eq:ch2_revOCV}).
\item $Q_\bg$ 와 각 $\xi_j$ transition 에 별도 entropy coefficient($h_\bg,\Delta S_j$)를 부여하되, 이들이
  OCV 온도계수를 재현하도록 제약(식~\eqref{eq:ch2_revxi}, \eqref{eq:ch2_revbg}).
\end{enumerate}
\textbf{정합 제약(평형 경로에서).} 평형 준정적 방전($\xi_j=\xi_{\eq,j}$, $V_n=V_{n,\OCV}$)에서 두 표현이 같은
가역 열률을 주어야 한다:
\begin{equation}
\boxed{\;s_{\rev,n}^{\conv}|I|\,T\left(\frac{\partial V_{n,\OCV}}{\partial T}\right)_q
=s_\bg^{\conv}\,T\,h_\bg\,\dot Q_\bg^{\prog}
+\sum_{j}s_{\rev,\xi,j}^{\conv}\,T\,\Delta S_j\,\frac{Q_{j,\tot}}{F}\,\frac{\dd\xi_{\eq,j}}{\dd t}\;.}
\label{eq:ch2_consistency}
\end{equation}
\begin{boundbox}
이 정합은 \emph{reaction-entropy 성분($\propto\partial U_j/\partial T$)에 한정}해 성립한다. OCV 방식 좌변은 추가로
(i) logistic width 의 온도의존 $-(V_n-U_j)w'_j/w_j^2$ 항(식~\eqref{eq:ch2_dxidT})과 (ii) 분모 $D_q$ 정규화를
포함하므로, 두 방식은 \emph{부분적으로만} 같은 정보다 — width$\cdot D_q$ 기여는 OCV 방식 고유다. 따라서
double-counting 차단은 reaction-entropy 성분에 대해 부과되며, 등호는 등온$\cdot$준정적 + 단일 우세 transition
(또는 $D_q$ 정규화 명시) 하에서 닫힌다(\AL{25}).
\end{boundbox}
이 제약을 두면 B 방식의 $\{h_\bg,\Delta S_j\}$ 가 A 방식의 $(\partial V_{n,\OCV}/\partial T)_q$ 와 정합되어
double counting 이 제거된다. 제약 없이 둘을 더하면 동일 entropy 가 두 번 계상된다. \textbf{정합성 점검(극한):}
background 가 없고($h_\bg=0$) 단일 transition($N_p=1$, $\xi_{\eq,1}$ 만) 이며 모든 $s^{\conv}=+1$ 인 극한에서,
식~\eqref{eq:ch2_consistency} 우변 $=T\,\Delta S_1\,(Q_{1,\tot}/F)\,\dd\xi_{\eq,1}/\dd t$ 이고 좌변에 평형
경로의 $|I|=Q_\cell\,\dd q/\dd t$ 와 $(\partial V_{n,\OCV}/\partial T)_q$ 를 넣으면 두 표현이 같은 $\Delta S_1$
정보(식~\eqref{eq:ch2_entropy_coeff})로 환원됨이 확인된다(\AL{25}). \emph{부호 규약 정합 주의}: 좌$\cdot$우변의
$s^{\conv}$($s_{\rev,n}^{\conv},s_\bg^{\conv},s_{\rev,\xi,j}^{\conv}$)는 모두 \emph{동일한 heat-positive 규약}
(발열을 양으로)에서 유도돼야 본 제약이 부호 모순 없이 닫힌다 — 서로 다른 bookkeeping 규약으로 독립 도입하면 같은
reaction entropy 가 좌우에서 반대 부호로 들어와 제약이 깨진다.

% ====================================================================
\subsection{배경 용량과 background entropy coefficient}\label{sec:ch2_bg}
% ====================================================================
$Q_\bg$ 는 잔류 background capacity 저장소이지 entropy 자체가 아니다. background 영역의 reversible heat 를
쓰려면 별도 entropy coefficient 가 필요하다(\AL{26}):
\begin{equation}
h_\bg(V_n,T)\equiv\left(\frac{\partial V_{\bg,\eq}}{\partial T}\right)_{Q_\bg},
\label{eq:ch2_hbg}
\end{equation}
($V_{\bg,\eq}$ = background 저장의 국소 평형 전위; 단위 V/K). background reversible heat 후보는
\begin{equation}
\dot{\mathcal Q}_{\rev,\bg}^{\cand}=s_\bg^{\conv}\,T\,h_\bg(V_n,T)\,\dot Q_\bg^{\prog},
\label{eq:ch2_revbg}
\end{equation}
$\dot Q_\bg^{\prog}$ 는 실제 background 저장/방출 \emph{진행} 율(단위 C/s)이다.

\textbf{total derivative 와 progress 의 구분(무비약, 중요).} 전하 보존식 (L1) 을 시간 미분한다. $Q_\bg$ 는
$(V_n,T)$ 의 함수이므로 전미분 연쇄법칙으로
\begin{equation}
Q_\cell\frac{\dd q}{\dd t}=\frac{\partial Q_\bg}{\partial V_n}\frac{\dd V_n}{\dd t}+\frac{\partial Q_\bg}{\partial T}\frac{\dd T}{\dd t}
+\sum_jQ_{j,\tot}\frac{\dd\xi_j}{\dd t}.
\label{eq:ch2_dQdt}
\end{equation}
비등온 조건에서 $\dd Q_\bg/\dd t$ 전체를 곧바로 reaction progress 로 해석하면 안 된다 — 온도 변화에 의한 항이
섞이기 때문이다. 개념적으로 분리하면
\begin{equation}
\dot Q_\bg^{\mathrm{tot}}=\dot Q_\bg^{\prog}+\dot Q_\bg^{\coord},\qquad
\dot Q_\bg^{\prog}=\frac{\partial Q_\bg}{\partial V_n}\frac{\dd V_n}{\dd t},\qquad
\dot Q_\bg^{\coord}=\frac{\partial Q_\bg}{\partial T}\frac{\dd T}{\dd t},
\label{eq:ch2_bgsplit}
\end{equation}
여기서 $\dot Q_\bg^{\prog}$(전위 구동 progress)만 heat-generating 반응 진행으로 식~\eqref{eq:ch2_revbg} 에
들어가고, $\dot Q_\bg^{\coord}$(온도 변화에 따른 capacity-coordinate shift)는 그대로 reaction rate 로 넣지
않는다. 이 구분을 빠뜨리면 단순 온도 drift 가 가역 발열로 오계상된다.

\textbf{단위 비대칭의 해소.} bg 항은 전위형 $T\,h_\bg\,\dot Q_\bg^{\prog}$ [W]이고 transition 항은 entropy형
$T\,\Delta S_j\,(Q_{j,\tot}/F)\,\dd\xi_j/\dd t$ [W]로, 둘 다 [W]이나 $h_\bg$ 가 이미 $\partial V/\partial T$
전위계수라 $/F$ 가 흡수되어 있다(transition 항의 $\Delta S_j/F$ 에 대응). 또한 $h_\bg$(온도 경로)와
$\dot Q_\bg^{\prog}$($V_n$ 경로)의 좌표 연결은 background 평형 항등식
$(\partial Q_\bg/\partial T)_{V_n}=-(\partial Q_\bg/\partial V_n)\,h_\bg$ 로 주어진다.

% ====================================================================
\subsection{비가역 열: 과전압 $\cdot$ flux$\times$force $\cdot$ $I\eta$}\label{sec:ch2_irr}
% ====================================================================
가역 열은 평형 entropy 가 주지만, 실제 전류가 흐르면 평형에서 벗어난 만큼 \emph{소산}이 생긴다. 그 비가역 열은
비평형 열역학의 flux$\times$force(entropy production)에서 출발하며, 2법칙으로 항상 $\ge0$ 이다(\AL{27},
\cite{degrootmazur1962,onsager1931}).

\subsubsection{공액 force = 과전압 $\eta_j$ (무비약)}\label{sec:ch2_fluxforce}
진행률 flux $J_{n,j}=Q_{j,\tot}(\dd\xi_j/\dd t)/F$ [mol/s] 에 \emph{공액(conjugate)인 열역학적 force 는
국소 평형으로부터의 이탈, 즉 과전압} $\eta_j$ 다(\cite{newman,bardfaulkner}, \AL{27}). 그 정의와 entropy
production 은
\begin{equation}
T\dot S_{\irr,n}=\sum_j J_{n,j}\,F\eta_j\ \ge\ 0,\qquad
\eta_j\equiv V_{n,\drive}-V_{\eq,j}(\xi_j),\quad
V_{\eq,j}(\xi_j)=U_j(T)+w_j(T)\ln\frac{\xi_j}{1-\xi_j}.
\label{eq:ch2_fluxforce}
\end{equation}
\textbf{$V_{\eq,j}(\xi_j)$ 의 유도(logistic 역함수).} Chapter 1 logistic(식 (L3), $s_{\phi,j}=+1$)
$\xi_{\eq,j}=[1+e^{-(V_n-U_j)/w_j}]^{-1}$ 을 ``$\xi_j$ 가 현재값일 때 그것을 평형으로 만드는 전위''에 대해
역으로 푼다. $\xi_j=[1+e^{-(V_{\eq,j}-U_j)/w_j}]^{-1}$ 로 두고 $\xi_j$ 에 대해 정리하면 $1/\xi_j-1=e^{-(V_{\eq,j}-U_j)/w_j}$,
즉 $(1-\xi_j)/\xi_j=e^{-(V_{\eq,j}-U_j)/w_j}$. 양변에 자연로그를 취하면 $\ln[(1-\xi_j)/\xi_j]=-(V_{\eq,j}-U_j)/w_j$,
부호를 뒤집어 $\ln[\xi_j/(1-\xi_j)]=(V_{\eq,j}-U_j)/w_j$, 따라서 $V_{\eq,j}(\xi_j)=U_j+w_j\ln[\xi_j/(1-\xi_j)]$
— 식~\eqref{eq:ch2_fluxforce} 의 셋째 식이다(중간 단계 전부 명시).

\textbf{부호 정합($\ge0$ 의 무비약 확인; 기본형 $V_{n,\drive}=V_n$ 전제).} \emph{먼저 전제를 명시한다}: flux 의
평형 target 은 Chapter 1 (L2) 의 $\xi_{\eq,j}(V_n)$(\emph{내부} 전위 기준)이고, force 의 평형은 $V_{\eq,j}(\xi_j)$
를 통한 $\xi_{\eq,j}(V_{n,\drive})$(\emph{구동} 전위 기준)이다. 두 평형이 같은 집합이 되는 것은 RB\_CHARTER step 2
의 기본형 $V_{n,\drive}=V_n$($\eta_{n,\drive}=0$)에서이며, 아래 항별 동부호는 \emph{이 전제 하에서} 닫힌다.
$V_{\eq,j}$ 는 $\xi_j$ 의 증가함수다($\ln[\xi_j/(1-\xi_j)]$ 가 증가).
$\xi_j<\xi_{\eq,j}$ 이면 (정의상) $V_{\eq,j}(\xi_j)<V_{n,\drive}$ 이라 $\eta_j=V_{n,\drive}-V_{\eq,j}>0$
이고, 동시에 (L2) 에서 $\dd\xi_j/\dd t=k_j(\xi_{\eq,j}-\xi_j)>0$ 이라 $J_{n,j}>0$ — 두 양이 \emph{같은 부호}라
$J_{n,j}F\eta_j>0$. 반대로 $\xi_j>\xi_{\eq,j}$ 이면 둘 다 음이라 곱은 여전히 양. 평형($\xi_j=\xi_{\eq,j}$)에서는
$\eta_j\to0$ 과 $\dd\xi_j/\dd t\to0$ 이 \emph{동시에} 일어나 entropy production 이 사라진다. 따라서
\emph{기본형에서} $T\dot S_{\irr,n}=\sum_j J_{n,j}F\eta_j\ge0$ 이 항별로 보장된다(2법칙 정합). \emph{강구동
$V_{n,\drive}\ne V_n$ 일반형}에서는 flux 와 force 의 평형 기준이 갈려 항별 동부호가 자동 보장되지 않으며, 그 경우의
양정치는 forward/backward rate 분해(Chapter 3)와 $n^\eff$ 일반형(Chapter 4)에서 닫는다(BOUNDED, \AL{28}).
\begin{boundbox}
\textbf{rate-affinity 와 heat-force 의 구분(\AL{28}, 중요).} Chapter 1 의 \emph{rate-구동} affinity
$\mathcal A_j=s_{\phi,j}F(V_{n,\drive}-U_j)$ 는 전이 \emph{중심} $U_j$ 기준이라 평형에서도 0 이 아니다(BEP
장벽 낮춤용 — 평형에서도 forward/backward 가 같은 유한 속도로 진행). 반면 \emph{heat-구동}(dissipation)
force 는 국소 평형 이탈 $\eta_j$ 다. 둘의 관계는 식~\eqref{eq:ch2_fluxforce} 의 $V_{\eq,j}$ 를 대입해
\begin{equation*}
\frac{\mathcal A_j}{F}=V_{n,\drive}-U_j=\underbrace{(V_{n,\drive}-V_{\eq,j}(\xi_j))}_{=\eta_j}+\underbrace{(V_{\eq,j}(\xi_j)-U_j)}_{=w_j\ln[\xi_j/(1-\xi_j)]},
\qquad\text{즉}\quad\frac{\mathcal A_j}{F}=\eta_j+w_j\ln\frac{\xi_j}{1-\xi_j}.
\end{equation*}
둘째 항(configurational 가역 혼합 entropy)은 \emph{비가역 열이 아니라} 가역 entropy basis(\S\ref{sec:ch2_revxi})
로 간다. 즉 \textbf{rate 에는 $\mathcal A_j$, dissipation 에는 $\eta_j$} 를 쓴다 — 이 둘을 혼동하면 가역 혼합
entropy 가 비가역 열로 오계상된다(BOUNDED — 본 분해는 logistic 평형 모델 내에서 정확).
\end{boundbox}

\subsubsection{전류$\times$과전압 축약형 $q_\irr=|I|\eta$ (CHARTER 부호 양정치)}\label{sec:ch2_Ieta}
flux$\times$force(식~\eqref{eq:ch2_fluxforce})를 전류와 과전압으로 축약하면, 일관된 부호 규약에서 비가역 열률은
\begin{equation}
\boxed{\;q_{\irr}=|I|\,\eta\ \ge\ 0,\qquad \eta=\sum_j \frac{J_{n,j}F}{I}\,\eta_j\ \ (\ge0),\;}
\label{eq:ch2_qirr}
\end{equation}
즉 $q_\irr$ 은 \emph{전류 크기 $|I|$ 와 과전압 크기 $\eta$ 의 곱}으로, 부호가 이미 정렬된 양정치 형이다
(RB\_CHARTER step 1: $q_\irr=|I|\eta$ 형, ``$I\eta$ 단독'' 금지 — 단일 transition 항별 부호는 BOUNDED 이나
flux 와 force 가 식~\eqref{eq:ch2_fluxforce} 에서 같은 부호임이 보장되어 합 $\ge0$ 이 2법칙으로 성립). $\eta$ 의
정의에 절댓값이 없어도 각 항 $J_{n,j}F\eta_j\ge0$ (식~\eqref{eq:ch2_fluxforce} 동부호, \S\ref{sec:ch2_fluxforce})
이므로 $\ge0$ 이 이미 닫혀, 절댓값 이중정당화는 불필요하다. 과전압
$\eta$ 는 charge-transfer($\eta_\kin$, Ch1 의 effective barrier 를 넘는 활성화), ohmic($\eta_\ohm$,
$R_{n,\eff}$ 분극; $R_n$ 아님, \S\ref{sec:ch2_Rn}), transport($\eta_\conc$) 기여의 합으로 분해되며 각각 $\ge0$ 이다.
transition별(식~\eqref{eq:ch2_qirr})$\leftrightarrow$메커니즘별(본 식) 분해는 동일 $\eta$ 의 두 관점이다.
\begin{boundbox}
\textbf{본 장은 $n^{\eff}=1$ 특수형(\AL{20}, RB\_CHARTER step 3).} 식~\eqref{eq:ch2_qirr} 의 $q_\irr$ 합은
ideal-width($n_j^{\eff}=RT/(Fw_j)=1$)를 전제한 특수형이다. effective width 일반형에서는 transition별 비가역
열률이 $\dot{\mathcal Q}_{\irr}=\sum_j n_j^{\eff}\,I_j\,\eta_j$ 로 $n_j^{\eff}$ 가중을 받으며, 이 \emph{일반형}
과 그 부호 양정치 보존($\sum_j n_j^{\eff}I_j\eta_j\ge0$)의 완전 유도는 \textbf{Chapter 4} 로 넘긴다(forward-ref).
본 장은 Chapter 4 일반형의 ideal-width 특수해로 정합한다(모순이 아니라 일반/특수 위계).
\end{boundbox}

\subsubsection{kinetic-lag 의 entropy production $\sigma_j\propto k_j r_j^2$ (near-equilibrium)}\label{sec:ch2_sigma}
Chapter 1 의 relaxation $\dd\xi_j/\dd t=k_j r_j$, $r_j=\xi_{\eq,j}-\xi_j$ 를 비평형 열역학의 미시 형태로 본다.
국소 자유에너지 $G_j(\xi_j)$ 를 평형점 $\xi_{\eq,j}$ 근처에서 Taylor 전개한다: $G_j(\xi_j)=G_j(\xi_{\eq,j})
+\tfrac12 g''_j(\xi_j-\xi_{\eq,j})^2+\cdots$ (평형점이라 1차항 $\partial G_j/\partial\xi_j|_{\eq}=0$,
$g''_j\equiv\partial^2 G_j/\partial\xi_j^2|_{\eq}$). 복원력(affinity)은 자유에너지 기울기의 음수이므로 1차
(선형) 근사에서
\begin{equation}
\mathcal X_j\equiv-\frac{\partial G_j}{\partial\xi_j}\Big|_{\xi_j}\simeq -g''_j(\xi_j-\xi_{\eq,j})=g''_j\,(\xi_{\eq,j}-\xi_j)=g''_j\,r_j,
\qquad g''_j>0,
\label{eq:ch2_affinity}
\end{equation}
($g''_j>0$ 은 평형 안정성 = regular-solution/lattice-gas 의 thermodynamic factor; Chapter 1 평형 진행률의
오목성과 동형, \cite{bazant2013}, \AL{29}). flux 는 $J_j=\dd\xi_j/\dd t=k_j r_j$다(\emph{무차원} 진행률
시간율 $J_j$ [1/s]이며, mole flux $J_{n,j}=Q_{j,\tot}(\dd\xi_j/\dd t)/F$ [mol/s]와 단위가 다름 — 혼동 금지). linear irreversible
thermodynamics 의 entropy production(flux$\times$force$/T$)은 \cite{degrootmazur1962,onsager1931}
\begin{equation}
\boxed{\;\sigma_j=\frac{J_j\,\mathcal X_j}{T}=\frac{g''_j\,k_j}{T}\,r_j^{\,2}\ \ge 0,\;}
\label{eq:ch2_sigma}
\end{equation}
즉 entropy production 은 lag $r_j$ 의 \emph{제곱}에 비례하여 항상 $\ge0$(2법칙; $g''_j>0$, $k_j>0$,
$r_j^2\ge0$). \textbf{차원 주의(G-dim).} $g''_j=\partial^2 G_j/\partial\xi_j^2$ 는 \emph{몰당} 자유에너지 곡률
[J/mol]이라 $\sigma_j=g''_j k_j r_j^2/T$ 는 \emph{intensive} [W/(mol$\cdot$K)]이고 $T\sigma_j=g''_j k_j r_j^2$
는 [W/mol]이다. 따라서 이것은 per-mode 비가역 열률 \emph{밀도}이며, control-volume 에너지 balance
(식~\eqref{eq:ch2_balance}, [W])의 형제항($q_\rev,q_\irr$ — 모두 $J_{n,j}=(Q_{j,\tot}/F)\dot\xi_j$ 의 몰 인자를
품은 extensive [W])과 차원을 맞추려면 그 mode 의 몰 함량 $(Q_{j,\tot}/F)$ [mol]을 곱해
\begin{equation}
\dot{\mathcal Q}_{\irr,\kin,j}=\frac{Q_{j,\tot}}{F}\,T\sigma_j=\frac{Q_{j,\tot}}{F}\,g''_j k_j r_j^2\quad[\mathrm W]
\label{eq:ch2_qirrkin_ext}
\end{equation}
로 환산한다(W/mol $\to$ W; Chapter 1 의 ``밀도는 적분변수의 역수 차원'' 교훈과 같은 bookkeeping).
\begin{boundbox}
\textbf{유효범위(\AL{29}).} 식~\eqref{eq:ch2_affinity}--\eqref{eq:ch2_sigma} 는 평형 근처 1차(선형) 전개다.
tail 영역($\xi_j\approx\xi_{\eq,j}$, 즉 $r_j$ 작음)에서 유효하며 큰 $r_j$ 의 전역 정당화가 아니다(Chapter 1
의 near-equilibrium relaxation 과 동일 bound). 큰 $r_j$ 에서는 $g'''$ 등 고차 보정이 필요하다.
\end{boundbox}

\subsubsection{$R_n$ 과 heat-generation resistance 의 구분 (P3-1)}\label{sec:ch2_Rn}
Chapter 1 의 $R_n$ 은 \emph{apparent 전압축 이동 계수}다: $V_{n,\app}-V_n=s_I|I|R_n(q,T,|I|)$. 이를 곧바로
실제 dissipated-heat resistance 로 동일시하지 않는다. 물리적 발열 저항으로 해석하려면 $R_n$ 이 ohmic /
charge-transfer / transport / film / concentration polarization 중 무엇을 대표하는지 독립 실험 또는 Chapter 3
kinetics 로 제한되어야 한다. 작은 과전압(선형 응답, $\eta\simeq R_{n,\eff}|I|$)에서만 reduced form
\begin{equation}
\dot{\mathcal Q}_{\irr,n}\approx I^2R_{n,\eff},\qquad R_{n,\eff}\ \ne\ R_n\ (\text{원칙적으로}),
\label{eq:ch2_i2r}
\end{equation}
($q_\irr=|I|\eta=|I|\cdot R_{n,\eff}|I|=I^2 R_{n,\eff}$, $I^2=|I|^2$). $R_{n,\eff}$ 는 열 발생에 대응하는
물리적 effective resistance 이며, apparent shift 계수 $R_n$ 과 일반적으로 다르다.
\begin{boundbox}
\textbf{식별성 경고(\AL{28}).} 전압 자료만으로 $R_n,R_{n,\eff},R_{\mathrm{ct}}$ 를 동시에 자유 피팅하면
apparent shift 와 실제 dissipated heat 가 혼동된다. 논문용 물리 모델에서는 $R_n$ 을 heat resistance 로 직접
쓰지 말고, pulse / current-interrupt / impedance / calorimetry 로 제한된 경우에만 $I^2R$ 형태를 쓴다(GS-4).
\end{boundbox}

% ====================================================================
\subsection{휴지 구간 relaxation 과 열 발생 가능성}\label{sec:ch2_rest}
% ====================================================================
휴지 구간에서는 $I=0,\ \dd q/\dd t=0$ 이지만 내부 진행률은 (L2) 로 계속 완화된다:
$\dd\xi_j/\dd t=k_j^{\eff}(\xi_{\eq,j}-\xi_j)$. 따라서 전하 보존식은 매 시간점에서 다시 풀린다($q$ 는 고정,
$\xi_j(t)$ 변화에 맞춰 $V_n(t)$ 가 따라감):
\begin{equation}
Q_\cell\,q_{\mathrm{rest}}=Q_\bg(V_n(t),T)+\sum_jQ_{j,\tot}\xi_j(t).
\label{eq:ch2_rest_balance}
\end{equation}
휴지 중 외부 전류에 의한 $I^2R$ 열(식~\eqref{eq:ch2_i2r})은 $I=0$ 이라 사라진다. 그러나 $\xi_j$ relaxation 이
있으면 entropy exchange 후보와 internal dissipation 후보가 남는다:
\begin{equation}
\dot{\mathcal Q}_{\xi,\rlx}^{\cand}=\sum_j s_{\rev,\xi,j}^{\conv}\,T\,\Delta S_j\,\frac{Q_{j,\tot}}{F}\,\frac{\dd\xi_j}{\dd t}
\qquad(\text{정전류 좌표 변환 금지: } \dd\xi_j/\dd t \text{ 직접 사용}).
\label{eq:ch2_relax}
\end{equation}
이는 entropy basis heat exchange \emph{후보}일 뿐이며 확정된 가역 열 또는 확정된 비가역 열로 단정하지 않는다.
다만 휴지($I=0$)서도 $r_j=\xi_{\eq,j}-\xi_j\ne0$, $\dd\xi_j/\dd t=k_j r_j\ne0$ 이므로 비가역 floor 는 식~\eqref{eq:ch2_qirrkin_ext}
로 정량 확정된다: $q_{\irr,\rlx,j}=(Q_{j,\tot}/F)\,g''_j k_j r_j^2\ge0$ (near-eq, \AL{29}).
\begin{flagbox}
\textbf{왜 후보인가(\AL{29}).} 휴지 relaxation 의 가역/비가역 분리는 forward/backward rate 의 비대칭(detailed
balance 로부터의 이탈)에 달려 있다. Chapter 1 Level A 의 $k_j=r_j^++r_j^-$ 는 \emph{합}만 주고 비대칭(entropy
production 의 정량값)은 주지 않으므로, $\dot{\mathcal Q}_{\xi,\rlx}$ 의 가역/비가역 \emph{정량} 분해는 Chapter
3--4 의 forward/backward rate 와 free-energy/entropy-production 구조가 있어야 확정된다. 비가역 floor 는 near-eq
에서 정량 확정($\ge0$, 식~\eqref{eq:ch2_qirrkin_ext})이며, Ch3--4 가 닫는 것은 (i) 가역 entropic 성분 배분과
(ii) 비선형/큰-$r_j$ 일반화다.
\end{flagbox}

% ====================================================================
\subsection{온도 되먹임 구조와 계산 순서}\label{sec:ch2_feedback}
% ====================================================================
온도는 Chapter 1 의 다음 항들에 되먹임된다(각 항이 $T$ 에 의존):
\begin{equation}
Q_\bg(V_n,T),\ \xi_{\eq,j}(V_n,T),\ U_j(T),\ w_j(T),\ \rho_G(g;T),
\label{eq:ch2_feedT1}
\end{equation}
\begin{equation}
k_j^{\eff}(V_n,q,T,I),\ k_{j,\lim}(V_n,q,T,I),\ R_n(q,T,|I|),\ V_{n,\OCV}(q,T).
\label{eq:ch2_feedT2}
\end{equation}
따라서 비등온 계산에서 $V_n,\xi_j,T$ 를 독립 후처리하지 않고 같은 시간 적분 루프에서 갱신한다(P3-3 순환
의존성). 계산 순서:
\begin{enumerate}[label=\arabic*),leftmargin=2.0em]
\item 외부 전류에서 $q(t)$ 업데이트.
\item 현재 $q,T,\{\xi_j\}$ 에서 전하 보존식 (L1) 으로 $V_n$ root-find.
  (이 (2)는 $\xi_{\eq,j}(V_n,T)$ 의 $V_n$ 의존으로 implicit 비선형 방정식 — P3-3 분류 $=$ 수치해법(부록 B);
  단조성 floor $\partial Q_\bg/\partial V_n\ge\epsilon_Q>0$, $D_q\ge\epsilon_Q>0$ 로 유일$\cdot$수렴이며 논리공백$\cdot$물리충돌이 아니다.)
\item $V_n,T$ 에서 $\xi_{\eq,j}$, 활성화 $k_j^{\mathrm{act}}$, 필요 시 $k_{j,\lim}$ 및 직렬 합성 $k_j^{\eff}$ 계산.
\item $\dd\xi_j/\dd t$ (L2) 계산.
\item 물리적으로 정의된 과전압$\cdot$affinity$\cdot$entropy coefficient 로 열원 항(식~\eqref{eq:ch2_srcdecomp}) 계산:
  가역 $\dot{\mathcal Q}_{\rev}$(식~\eqref{eq:ch2_revOCV} 또는 \eqref{eq:ch2_revxi}, 정합 제약 식~\eqref{eq:ch2_consistency}),
  비가역 $q_\irr=|I|\eta$(식~\eqref{eq:ch2_qirr}). (후보항 $\dot{\mathcal Q}_{\rlx}^{\cand}$ 는 본 장 정량 미확정 $\to$ src 합에서 분리 표기.)
\item 열수지식(식~\eqref{eq:ch2_balance})으로 $T(t)$ 업데이트.
\end{enumerate}
(수치 적분기$\cdot$root-find solver 구현과 수렴 판정은 본 이론 장 범위 밖이며 Ch1 부록 B 로 분리한다; GS-4 의
이론식/피팅 실무 분리.)
\begin{boundbox}
\textbf{되먹임이 ICA tail 에 미치는 영향.} 발열 $\to T\uparrow\to k_j\uparrow$ ($L=|I|/(Q_\cell k_j)\downarrow$,
Chapter 1 spectrum shift) $\to$ 꼬리 단축; 동시에 ideal $w_j=RT/F\uparrow$ 로 평형 peak 폭 증가(Chapter 1
식 $\partial\xi_{\eq,j}/\partial V_n\propto1/w_j$). 두 효과가 경쟁하므로 비등온 ICA tail 을 단일 부호로 단정
하지 않는다(Chapter 1 의 온도 tail 3계층 분해와 정합). 이것이 본 장 발열이 ICA tail 로 되먹임되는 핵심 고리다.
\end{boundbox}

% ====================================================================
\subsection{비가역열(kinetic-lag 성분 $q_{\irr,\kin}$) tail = ICA tail 의 열적 거울 (consistency 가설)}\label{sec:ch2_thermal_tail}
% ====================================================================
\textbf{물리.} Chapter 1 은 ICA tail 을 relaxation-length spectrum 의 kernel integral 로 설명했다. 같은
relaxation 이 비가역 열(kinetic-lag 성분 $q_{\irr,\kin}$, $\sigma_j$ 채널)을 만들므로(식~\eqref{eq:ch2_sigma}), \emph{같은 kinetic 정보}가 비가역열 tail 에도
나타나야 한다 — 이것이 ICA tail 의 ``열적 거울''이며, calorimetry 로 Chapter 1 의 kinetic 귀속을 \emph{독립
modality 로 반증}할 수 있게 한다.

\textbf{무비약 유도(per-mode).} post-peak tail 에서 $k_j\approx\mathrm{const}$, $\xi_{\eq,j}\approx\mathrm{const}$
근사 하에 $r_j$ 는 단일 지수다: Chapter 1 post-peak 에서 단일 mode 의 잔차는 $r_j(q)=r_j(q_a)\exp[-(q-q_a)/L]$
($L=|I|/(Q_\cell k_j)$, Chapter 1 single-mode kernel). 이를 식~\eqref{eq:ch2_sigma} 의 per-mode 비가역
열률에 넣는다. $r_j^2=r_j(q_a)^2\exp[-2(q-q_a)/L]$ 이므로
\begin{equation}
\frac{\dd q_{\irr,\kin,j}}{\dd q}\ \propto\ g''_j\,k_j\,r_j(q)^2
= g''_j\,k_j\,r_j(q_a)^2\,\exp\!\Big[-\frac{2(q-q_a)}{L}\Big].
\label{eq:ch2_thermal_single}
\end{equation}
\textbf{핵심 비교(두 kernel 의 차이).} Chapter 1 의 ICA progress kernel 은 $\dd\xi_j/\dd q=r_j/L\propto(1/L)e^{-(q-q_a)/L}$
로 (i) kinematic 인자 $1/L$ 을 갖고 (ii) 감쇠 지수가 $-(q-q_a)/L$ 이다. 반면 비가역열 kernel(식~\eqref{eq:ch2_thermal_single})
은 $r_j^2$ 때문에 (i) kinematic $1/L$ 인자가 \emph{없고} (ii) 감쇠 지수가 \emph{2배} $-2(q-q_a)/L$ 다(arc-length
척도가 절반 $L/2$, 단일 mode 기준 2배 빠른 감쇠). spectrum 으로 중첩하면
\begin{equation}
\boxed{\;\frac{\dd q_{\irr,\kin}}{\dd q}\simeq\int_0^\infty A_L^{(2)}(L;T,V_{n,\drive})\,
\exp\!\Big[-\frac{2(q-q_a)}{L}\Big]\,\dd L,\;}
\label{eq:ch2_thermal_kernel}
\end{equation}
여기서 $A_L^{(2)}$ 는 Chapter 1 spectrum 의 \emph{지지집합} $\{L\}$ 을 그대로 쓰되 per-mode weight $g''_j k_j r_j(q_a)^2$
(kinematic $1/L$ 부재 포함)을 실은 \textbf{재가중} 분포다 — Chapter 1 의 $A_L$ 자체와 같지 않다.
\begin{boundbox}
\textbf{spectrum 일반화 한계(\AL{29}).} ``decay length 비 $=2$''(균일 $L/2$)는 \emph{단일 우세 mode}
($A_L\to\delta(L-L_0)$, Chapter 1 single-mode 극한)에서만 엄밀하다. Chapter 1 이 명시한 stretched/비지수
spectrum(broad $\rho_G$, kernel mixture)에서는 두 관측 tail 이 단순히 ``2배 빠른'' 관계가 아니라, per-mode
arc-length doubling $q-q_a\to2(q-q_a)$ 이 \emph{재가중}과 함께 나타난다(예: power-law tail 이면 같은 멱지수의
상수배 재척도이지 ``2배 감쇠''가 아님). 따라서 정량 검정은 ``정확히 2배''가 아니라 \emph{동일 $\{L\}$ 지지집합
으로부터의 forward prediction 정합}(Chapter 1 forward-only 원칙)으로 수행한다(BOUNDED).
\end{boundbox}

\textbf{$T$ 거동(Chapter 1 stretched tail 의 열적 거울).} 저온이면 $L$ 이 커지고(Chapter 1: $k_j\propto e^{-(G-\chi_j\mathcal A_j)/RT}$
가 작아짐) relaxation 이 느려, post-peak 같은 arc-length 에서 \emph{잔여} lag $r_j$ 가 (고온 대비) 더 크게
유지되어 식~\eqref{eq:ch2_thermal_kernel} 의 비가역열 tail 이 더 길고 더 크다. 고온이면 $L$ 이 작아져 비가역열
tail 이 짧아진다. 이는 near-equilibrium tail($r_j$ 작음, 식~\eqref{eq:ch2_affinity} 선형 전개) 내의 \emph{상대적}
잔여 lag 비교다.
\begin{flagbox}
\textbf{상태(\AL{29}, novel).} ``저온 긴 ICA tail $\leftrightarrow$ 저온 큰 비가역열 tail''의 대응은 novel
hypothesis 다. consistency 로만 제시하며, \emph{확정}은 calorimetry cross-check($q_\rev$$\cdot$mixing 항과
ohmic($I^2R$)$\cdot$transport 성분을 먼저 차감한 뒤 비가역열 tail 이 Chapter 1 spectrum $\{L\}$ 로부터의
forward prediction 과 정합하는지) 통과를
전제로 한다(Chapter 1 의 forward-only 반증 원칙의 calorimetry modality 확장). 불일치는 kinetic-lag dissipation
그림(따라서 Chapter 1 의 kinetic-spectrum 귀속)을 기각/재검토하게 한다. \emph{독립성 한정}: 이 검정은 $r_j(q_a)$
를 post-peak 공통값으로 근사하는 단계를 Chapter 1 ICA kernel 과 \emph{공유}하므로(동일 근사 재사용), 완전 독립
반증이 아니라 \emph{준독립}(semi-independent) cross-modality 정합 검정으로 한정된다.
\end{flagbox}

% ====================================================================
\subsection{식별성 및 실험 요구조건}\label{sec:ch2_ident}
% ====================================================================
물리 우선 모델에서도 식별성은 사라지지 않는다. 해결책은 임의 clipping 이 아니라 독립 실험과 물리적 prior 다.
\begin{longtable}{>{\raggedright\arraybackslash}p{0.36\textwidth} >{\raggedright\arraybackslash}p{0.52\textwidth}}
\toprule 상관 항 & 주의점 \\ \midrule \endhead
$R_n$ 와 $R_{n,\eff}$ & apparent voltage shift 와 실제 dissipated heat 혼동(식~\eqref{eq:ch2_i2r}) \\
$(\partial V_{n,\OCV}/\partial T)_q$ 와 $\Delta S_j$ & 둘 다 reversible heat 설명 $\to$ double counting(식~\eqref{eq:ch2_consistency} 제약 필요). 저율 OCV series 가 식~\eqref{eq:ch2_consistency} 좌변 $(\partial V_{\OCV}/\partial T)_q$ 을 고정해 B 방식 $\{h_\bg,\Delta S_j\}$ 합을 제약 \\
$\Delta S_{\rxn}$ 와 $\Delta S_a$ & reaction entropy(평형, $\partial U/\partial T$)와 activation entropy(rate prefactor) 별개 — 대입 금지(\AL{24}) \\
$k_j^{\mathrm{act}}$ 와 $k_{j,\lim}$ & 강구동 가속과 물리 병목 분리(Ch1 계승) \\
$Q_\bg$ 와 $h_\bg$ & 배경 용량과 background entropy coefficient 를 독립 자료 없이 분리 곤란 $\to$ 저속 OCV series 의 SOC 분해 + 식~\eqref{eq:ch2_bgsplit}(\AL{26}) 로 해소 \\
$hA$ 와 내부 열원 항 & 표면 온도만 있으면 heat generation 과 heat loss 가 상호 보상 \\
\bottomrule
\end{longtable}
필요 실험 제약:
\begin{enumerate}[label=\arabic*),leftmargin=2.0em]
\item 저율 OCV/quasi-OCV 온도 series: $(\partial V_{n,\OCV}/\partial T)_q$, $U_j(T)$, $w_j(T)$, $\partial U/\partial T$(reaction entropy) 제한.
\item rate series: $k_j^{\mathrm{act}}$, $k_{j,\lim}$, transport limitation, tail broadening 분리.
\item pulse/current-interrupt/impedance: ohmic 및 charge-transfer 성분 제한($R_n$ vs $R_{n,\eff}$).
\item 휴지 relaxation 자료: $\xi_j$ relaxation 과 thermal relaxation 분리(\S\ref{sec:ch2_rest}).
\item calorimetry/온도 측정: heat source scale, $q_\rev$ vs $q_\irr$ 분리, heat loss($hA$) 제한, 비가역열 tail forward-prediction 검정(\S\ref{sec:ch2_thermal_tail}).
\end{enumerate}
\textbf{한계.} calorimetry 자료가 부재하면 $q_\rev$ vs $q_\irr$(가역/비가역) 분해가 전압 자료만으로는 닫히지
않으며, $\partial U_j/\partial T$ 분해도 AL-23 의 SOC 분리(staging별 entropy profile)를 전제로만 위계화된다.

% ====================================================================
\subsection{Chapter 2 의 가정 원장 (AL-20$\sim$29)}\label{sec:ch2_al}
% ====================================================================
본 장에서 사용한 가정을 통합 ledger(RB\_AL\_MASTER) Ch2 그룹(AL-20$\sim$29)으로 정리한다. AL-20,21 은 기등재분
(Foundation), AL-22$\sim$29 는 본 장 S2 에서 채운 신규 정의다. tier = GROUNDED/BOUNDED/FLAGGED.
{\footnotesize
\setlength{\tabcolsep}{3pt}
\begin{longtable}{>{\raggedright\arraybackslash}p{0.05\textwidth} >{\raggedright\arraybackslash}p{0.355\textwidth} >{\raggedright\arraybackslash}p{0.155\textwidth} >{\raggedright\arraybackslash}p{0.195\textwidth} >{\raggedright\arraybackslash}p{0.13\textwidth}}
\toprule AL\# & 가정 진술 & tier & 근거 cite (DOI) & 사용 식 \\ \midrule \endhead
AL-20 & 전지 에너지 balance: $q_\rev=s^\conv|I|T\,\partial U/\partial T$, $q_\irr=|I|\eta$ (Bernardi--Pawlikowski--Newman) & GROUNDED & bernardi1985,\brk rao1997 (10.1149/\brk 1.2113792;\brk 10.1149/\brk 1.1837884) & \eqref{eq:ch2_srcdecomp},\eqref{eq:ch2_balance},\eqref{eq:ch2_qrev},\eqref{eq:ch2_revOCV},\eqref{eq:ch2_qirr} \\
AL-21 & 평형 전위 온도계수 = reaction entropy $\partial U/\partial T=\Delta S_\rxn/(s_\phi F)$;\brk graphite $\mathrm{Li_xC_6}$ 실측 & GROUNDED & reynier2004,\brk thomasnewman\brk 2003 (10.1149/\brk 1.1646152;\brk 10.1149/\brk 1.1531194) & \eqref{eq:ch2_entropy_coeff},\eqref{eq:ch2_revxi} \\
AL-22 & 기준 용량 $Q_\cell,Q_{j,\tot}$ 온도 무관 매개변수($Q_{j,\tot}(T)$ 보정은 일반형으로 분리) & BOUNDED & thomasnewman\brk 2003 (10.1149/\brk 1.1531194) & \eqref{eq:ch2_eqbal},\eqref{eq:ch2_dVocvdT},\eqref{eq:ch2_dVocvdT_gen} \\
AL-23 & graphite $U_j(T)$ entropy profile 비단조 → $(\partial V_{n,\OCV}/\partial T)_q$ SOC 부호 반전 가능 & BOUNDED & reynier2004,\brk thomasnewman\brk 2003 (10.1149/\brk 1.1646152;\brk 10.1149/\brk 1.1531194) & \eqref{eq:ch2_dxidT} boundbox \\
AL-24 & reaction entropy $\Delta S_\rxn$($\partial U/\partial T$) $\ne$ activation entropy $\Delta S_a$(Eyring prefactor); 대입 금지 & BOUNDED & eyring1935,\brk bard\brk faulkner (10.1063/\brk 1.1749604;\brk ISBN 0-471-\brk 04372-0) & \eqref{eq:ch2_eyring_entropy} boundbox \\
AL-25 & OCV 계수 가역열 $=$ transition entropy basis 가역열 (정합 제약,\brk double counting 금지) & GROUNDED & bernardi1985,\brk thomasnewman\brk 2003 (10.1149/\brk 1.2113792;\brk 10.1149/\brk 1.1531194) & \eqref{eq:ch2_consistency} \\
AL-26 & background entropy coefficient $h_\bg=(\partial V_{\bg,\eq}/\partial T)$;\brk progress vs coordinate-shift 분리 & BOUNDED & rao1997,\brk thomasnewman\brk 2003 (10.1149/\brk 1.1837884;\brk 10.1149/\brk 1.1531194) & \eqref{eq:ch2_hbg},\eqref{eq:ch2_revbg},\eqref{eq:ch2_bgsplit} \\
AL-27 & 비가역열 = flux$\times$force(entropy production); 공액 force = 과전압 $\eta_j$;\brk 2법칙 $\ge0$ & GROUNDED & degrootmazur\brk 1962 (book,\brk ISBN),\brk onsager1931 (10.1103/\brk PhysRev.37.405) & \eqref{eq:ch2_fluxforce},\eqref{eq:ch2_qirr} \\
AL-28 & rate-affinity $\mathcal A_j$(중심 $U_j$) $\ne$ heat-force $\eta_j$(국소 평형 이탈); $R_n\ne R_{n,\eff}$ & BOUNDED & newman,\brk bard\brk faulkner (ISBN 0-471-\brk 47756-3;\brk 0-471-\brk 04372-0) & \eqref{eq:ch2_fluxforce} boundbox,\eqref{eq:ch2_i2r} \\
AL-29 & near-eq entropy production $\sigma_j=(g''_j/T)k_j r_j^2\ge0$ ($g''_j>0$ 안정성); 비가역열 tail = ICA tail 열적 거울 (per-mode arc-length doubling;\brk spectrum 재척도) & BOUNDED(σ) / FLAGGED(열적 tail novel) & degrootmazur\brk 1962,\brk onsager1931,\brk bazant2013 (10.1103/\brk PhysRev.37.405;\brk 10.1021/\brk ar300145c) & \eqref{eq:ch2_affinity},\eqref{eq:ch2_sigma},\eqref{eq:ch2_thermal_kernel},\eqref{eq:ch2_relax} \\
\bottomrule
\end{longtable}
}
\textbf{AGP.} 모든 본문 가정식은 AL-\# 를 인용한다. FLAGGED(AL-29 의 열적 tail novel)는 established 로 쓰지
않고 consistency+falsification 로만 제시한다. BOUNDED(AL-22,23,24,26,28,29-$\sigma$)는 유효범위를 병기한다.
임의 clip/cap/softplus/step 정의식 0(GS-4). solver/numerical-validation 구현 0(Ch1 부록 B 로 분리).

% ====================================================================
\subsection{Chapter 2 자체 검수표}\label{sec:ch2_selfcheck}
% ====================================================================
{\small
\begin{longtable}{>{\raggedright\arraybackslash}p{0.74\textwidth} >{\raggedright\arraybackslash}p{0.16\textwidth}}
\toprule 검수 항목 & 결과 \\ \midrule \endhead
Chapter 1(RB) 상태변수/기호를 수정 없이 입력으로 받았는가(P3-6) & 통과(\S\ref{sec:ch2_inherit}) \\
발열 항이 entropy/affinity/overpotential 물리에서 출발하는가(피팅 편의 X) & 통과(\S\ref{sec:ch2_decomp}) \\
$\partial U/\partial T=\Delta S_\rxn/(s_\phi F)$ 무비약 유도 + reaction$\ne$activation entropy 구분 & 통과(\S\ref{sec:ch2_entropy}, \AL{21},\AL{24}) \\
가역열 $q_\rev=s^\conv|I|T(\partial V_{n,\OCV}/\partial T)_q=I_j\Pi_j$ Bernardi balance, 부호$\cdot$차원 일치 & 통과(식~\eqref{eq:ch2_qrev},\eqref{eq:ch2_revOCV}) \\
비가역열 $q_\irr=|I|\eta\ge0$ (CHARTER 부호 양정치, ``$I\eta$ 단독'' 회피) & 통과(식~\eqref{eq:ch2_qirr}) \\
비가역열 공액 force = 과전압 $\eta_j=V_{n,\drive}-V_{\eq,j}(\xi_j)$,\brk rate-affinity $\mathcal A_j$ 와 구분 & 통과(식~\eqref{eq:ch2_fluxforce}, \AL{28}) \\
평형서 $\eta_j\to0,\dot\xi_j\to0$ 동시 $\Rightarrow$ entropy production$\to0$; 가역 혼합 entropy 오계상 회피 & 통과(\S\ref{sec:ch2_fluxforce}) \\
$V_n,V_{n,\OCV},V_{n,\app},V_{n,\drive}$ 구분 유지; $V_{n,\drive}$ 기본 $=V_n$ (CHARTER step2) & 통과(\S\ref{sec:ch2_inherit} L4) \\
$(\partial V_{n,\OCV}/\partial T)_q$ 를 평형 전하 보존식에서 유도(외부 lookup X,\brk P3-2); $D_q>0$ Ch1 floor & 통과(식~\eqref{eq:ch2_dVocvdT}) \\
고정-$q$ 온도미분서 chain($V_{n,\OCV}(T)$)과 명시 $T$ 의존 분리 무비약 & 통과(식~\eqref{eq:ch2_implicitT}) \\
$V_{n,\app}$ 온도 미분을 reversible heat 로 쓰지 않음 명시 & 통과(식~\eqref{eq:ch2_forbidden}) \\
상변이 heat-rate 입력을 $\dd\xi_j/\dd t$ 로 유지, $\dd\xi_j/\dd q$ 는 정전류 한정 & 통과(식~\eqref{eq:ch2_revxi_q}) \\
$\Delta S_j$ mol electron/Li-equiv 기준 + 차원 W 확인 & 통과(식~\eqref{eq:ch2_revxi_dim}) \\
OCV entropy 방식과 transition entropy basis double counting 정합 제약 & 통과(식~\eqref{eq:ch2_consistency}) \\
$n^\eff=1$ 특수형 명시 + Ch4 일반형 $\sum_j n^\eff_j I_j\eta_j$ forward-ref (CHARTER step3) & 통과(\S\ref{sec:ch2_Ieta}) \\
near-eq entropy production $\sigma_j=(g''_j/T)k_j r_j^2\ge0$ (2법칙) & 통과(식~\eqref{eq:ch2_sigma}) \\
$R_n,R_{n,\eff},R_{\mathrm{ct}}$ 구분, $I^2R$ 는 선형응답 한정 & 통과(식~\eqref{eq:ch2_i2r}) \\
$Q_\bg$ total derivative 와 progress rate 구분 & 통과(식~\eqref{eq:ch2_bgsplit}) \\
휴지 relaxation heat 를 확정항 아닌 후보로, 가역/비가역 정량 분해는 Ch3--4 위임 & 통과(\S\ref{sec:ch2_rest}) \\
온도 되먹임 순환(P3-3)을 같은 적분 루프로 처리, 계산 순서 명시 & 통과(\S\ref{sec:ch2_feedback}) \\
비가역열 tail = ICA tail 열적 거울 (FLAGGED novel,\brk forward-prediction 검정) & 통과(\S\ref{sec:ch2_thermal_tail}) \\
keystone Level A 열적 대응 + Level B 일치는 정상 target 극한 한정 (CHARTER step4) & 통과(\S\ref{sec:ch2_decomp}) \\
가역 열 부호를 피팅 X,\brk convention factor $s^{\conv}$ 로 분리 & 통과(\S\ref{sec:ch2_revOCV}) \\
임의 clipping/stabilizer 를 기본 열원 항으로 쓰지 않음(GS-4) & 통과 \\
모든 가정 AL-20$\sim$29 인용; 미확정 FLAGGED;\brk BOUNDED 유효범위 병기 & 통과(\S\ref{sec:ch2_al}) \\
Ch3--5 이월 항목 분리 & 통과(\S\ref{sec:ch2_transmit}) \\
\bottomrule
\end{longtable}
}

% ====================================================================
\subsection{후속 장으로 넘기는 항목}\label{sec:ch2_transmit}
% ====================================================================
\begin{enumerate}[label=\arabic*),leftmargin=2.0em]
\item 완전한 Butler--Volmer / generalized BV; exchange current density 와 계면 overpotential 의 분리 (Ch3).
\item forward/backward rate 와 detailed balance 기반 entropy production (Ch3--4) — 휴지 relaxation heat
  (식~\eqref{eq:ch2_relax})의 가역/비가역 \emph{정량} 확정.
\item 비가역열 일반형 $\dot{\mathcal Q}_{\irr}=\sum_j n_j^{\eff}I_j\eta_j$ 와 그 부호 양정치 보존의 완전 유도
  (Ch4; 본 장은 $n^\eff=1$ 특수형).
\item microscopic free-energy surface 기반 relaxation heat (Ch4); full-cell heat balance 와 양극/음극 발열
  분리 (Ch4--5).
\item 충전 branch, hysteresis, branch-dependent reversible heat($s_{\phi,j}^b$ 부호 유도) (Ch5).
\item 온도 되먹임 시간 적분기$\cdot$root-find solver$\cdot$수렴 판정 (Ch1 부록 B).
\end{enumerate}

% ====================================================================
\subsection{Chapter 2 의 결론}\label{sec:ch2_concl}
% ====================================================================
첫째, Chapter 1 전하 보존식은 평형에서 $(\partial V_{n,\OCV}/\partial T)_q$ 를 유도한다(식~\eqref{eq:ch2_dVocvdT}).
이는 가역 열과 연결되는 상태함수 온도계수이며 apparent 전위 $V_{n,\app}$ 경로 미분과 구분된다.

둘째, 그 온도계수의 정체는 reaction entropy 다: $\partial U_j/\partial T=\Delta S_{\rxn,j}/(s_{\phi,j}F)$
(식~\eqref{eq:ch2_entropy_coeff}). 이는 Chapter 1 Eyring rate 의 activation entropy $\Delta S_{a,j}$ 와
\emph{별개}이며 상호 대입을 금지한다(식~\eqref{eq:ch2_eyring_entropy}).

셋째, 전지 가역 반응열은 Bernardi 에너지 balance 의 $q_\rev=s^{\conv}|I|T(\partial V_{n,\OCV}/\partial T)_q
=I_j\Pi_j$ (식~\eqref{eq:ch2_qrev},\eqref{eq:ch2_revOCV})로, 부호가변(entropic)이다. transition entropy basis
$\dot{\mathcal Q}_{\rev,\xi}$ (식~\eqref{eq:ch2_revxi})는 같은 entropy 정보의 미시 표현이라 정합 제약
(식~\eqref{eq:ch2_consistency})으로 묶어 double counting 을 막는다.

넷째, 비가역 열은 flux$\times$\emph{과전압} $\sum_jJ_{n,j}F\eta_j$, 축약형 $q_\irr=|I|\eta\ge0$
(식~\eqref{eq:ch2_qirr}; CHARTER 부호 양정치), 또는 물리적으로 제한된 $I^2R_{n,\eff}$ 다. 공액 force 는 국소
평형 이탈 $\eta_j=V_{n,\drive}-V_{\eq,j}(\xi_j)$ 이며 전이 중심 기준 rate-affinity $\mathcal A_j$ 와 구분된다
(가역 혼합 entropy 의 비가역 오계상 방지). 본 장은 ideal-width $n^\eff=1$ 특수형이며 일반형은 Chapter 4 다.

다섯째, kinetic-lag 의 near-eq entropy production $\sigma_j=(g''_j/T)k_j r_j^2\ge0$ (식~\eqref{eq:ch2_sigma})
이 비가역열 tail 을 만들며, 이것이 Chapter 1 ICA tail 의 \emph{열적 거울}(per-mode arc-length doubling, FLAGGED
novel)이라는 consistency 가설을 준다 — calorimetry 가 Chapter 1 의 kinetic 귀속을 독립 modality 로 반증할
handle 이다.

여섯째, Chapter 2 는 발열 모델 폐쇄가 아니라 물리적 연결 장이다. 상세 과전압, exchange current,
forward/backward kinetics 와 entropy production 의 정량 분해, 충방전 hysteresis heat 는 Chapter 3--5 에서
다룬다. 이로써 ``열역학=가역(방향/entropy), 동역학=비가역(속도/affinity)'' 분업이 Chapter 1 의 인과 사슬과
정합되게 유지된다.

\subsection*{Chapter 3 으로의 전달}
본 장이 확정한 열 계층($q_\rev=I_j\Pi_j$, $q_\irr=|I|\eta$, $\sigma_j\ge0$)과 상태변수($V_n,\xi_j,k_j,U_j,
\eta_j$) 위에서: Chapter 3 은 Chapter 1 의 mobility 가속 $k_j$ 를 forward/backward 반응속도론(Level B,
$\beta_j\equiv\chi_j$; 조건부 동일물, 대칭 Marcus$\cdot$정상 target 극한 한정)으로 확대하여 본 장이 후보로 남긴 휴지 relaxation heat 의 가역/비가역 \emph{정량} 분해와
detailed balance 기반 entropy production 을 닫는다. Chapter 4 는 그 위에서 비가역열 일반형
$\sum_j n_j^{\eff}I_j\eta_j$ 와 full-cell 발열을, Chapter 5 는 충방전 히스테리시스의 branch 의존 가역열을 다룬다.



% ====================================================================
% ===== Chapter 3 body (원본 graphite_ica_ch3_rebuilt.tex, byte-verbatim; heading 1단계 demote) =====
% ====================================================================
\clearpage
\section{Chapter 3: 반응속도론 일반화 (Level B) — forward/backward kinetics $\cdot$ detailed balance}\label{chap:ch3}

% ====================================================================
\subsection*{서 — 인과 사슬에서 ``반응속도론''의 자리}
\addcontentsline{toc}{subsection}{서 — 반응속도론의 자리}
% ====================================================================
Chapter 1 은 ICA 꼬리의 주역을 \emph{동역학}(진행률 $\xi_j$ 의 유한속도 완화)으로 지목하되, 그 동역학을
\textbf{Level A}(평형 target $\xi_{\eq,j}$ 을 향한 mobility $k_j$ 완화)로 \emph{coarse-grained} 하게 두었다.
Level A 의 핵심 진술은 Chapter 1 의 완화식 $\dd\xi_j/\dd t=k_j(\xi_{\eq,j}-\xi_j)$ 와 그 keystone(``$k_j$=속도,
$\xi_{\eq,j}$=방향''; $\chi_j\mathcal A_j$ 는 \emph{방향성 없는} mobility 가속)이었다. Chapter 3 은 이 한 단계를
\textbf{확대(zoom-in)}한다: ``그 mobility 는 미시적으로 어떤 forward/backward transition kinetics 에서
나오는가, 그리고 그 미시식이 Chapter 1 식을 \emph{정확히} 회복하는가''.

본 장은 이 확대를 \textbf{다섯 단계}로 푼다. 각 단계는 학부 수준(전기화학 평형 $\Delta G=-nF E$, Nernst,
Eyring rate, 미적분 Taylor 1차$\cdot$지수$\cdot$로그)으로 따라갈 수 있게 중간을 생략하지 않는다.
\begin{enumerate}[label=\textbf{(\arabic*)},leftmargin=2.2em]
\item \textbf{진행은 forward 와 backward 의 차다.} 가장 일반적인 형 $\dd\xi_j/\dd t=r_j^+(1-\xi_j)-r_j^-\xi_j$
  에서 출발하면(\S\ref{sec:ch3_db}), 평형(net flux $0$)이 곧 \emph{detailed balance} $r_j^+/r_j^-=\xi_{\eq,j}/(1-\xi_{\eq,j})$
  를 준다. Chapter 1 logistic 을 대입하면 $\ln(r_j^+/r_j^-)=(V_n-U_j)/w_j$ 가 무비약으로 따라온다(\AL{30}).
\item \textbf{Level A 는 그 forward/backward 의 \emph{대칭} 축약이다.} split $r_j^+=k_j\xi_{\eq,j}$,
  $r_j^-=k_j(1-\xi_{\eq,j})$ 가 Chapter 1 식을 \emph{대수적 항등}으로 복원한다(\S\ref{sec:ch3_consistency}).
  여기서 $k_j$ 는 forward-only rate 가 아니라 mobility(=$r_j^++r_j^-$)이고, 방향(target)은 열역학이 전담한다.
\item \textbf{Level B 는 방향성을 넣는다.} forward 장벽만 $-\beta_j\mathcal A_j$, backward 는 $+(1-\beta_j)\mathcal A_j$
  로 \emph{비대칭}하게 낮추면(\S\ref{sec:ch3_levelB}) rate ratio 가 $\mathcal A_j$ 에 의존한다(방향성 O). 이때
  Chapter 1 의 $\chi_j$ 가 transfer coefficient $\beta_j$ 와 \emph{동일물}임을 명시한다($\chi_j\equiv\beta_j$; \emph{대칭 Marcus 1차$\cdot$정상 target 극한 한정} — 일반은 별개).
\item \textbf{detailed-balance 정합이 $n_j^{\eff}$ 를 고정한다.} Level B 의 ratio 가 Chapter 1 평형(단계 (1))과
  같으려면 $n_j^{\eff}=RT/(Fw_j)$ 가 \emph{강제}된다(\S\ref{sec:ch3_neff}). 이 정합 하에서 비평형 정상 target
  $\xi_{\sstat,j}$ 가 평형 target $\xi_{\eq,j}$ 로 환원된다 — \textbf{Level A 는 Level B 의 detailed-balance
  극한}이라는 keystone 의 \emph{한정} 명제다(CHARTER step4).
\item \textbf{관측 저항은 세 층이다.} Butler--Volmer 형(\S\ref{sec:ch3_bv})과 그 small-signal 한계에서
  charge-transfer $R_\ct$ 가 나오고, 이는 Chapter 1 의 apparent shift $R_n$ 과도 Chapter 2 의 heat-generation
  $R_{n,\eff}$ 와도 다르다 — $R_n\ne R_\ct\ne R_{n,\eff}$(\AL{31}, \S\ref{sec:ch3_Rn}).
\end{enumerate}

\noindent\textbf{한 문장 요지.} \emph{Chapter 1 의 relaxation ODE 는 forward/backward transition kinetics 의
coarse-grained 축약형이며}(Level A = Level B 의 detailed-balance 극한), 강구동에서 forward rate 증가는
물리적으로 허용하되 그 제한은 임의 cap 이 아니라 transport$\cdot$site$\cdot$diffusion$\cdot$current-conservation
병목으로 도입한다. 본 장은 발열식의 최종 부호와 일반형(Ch.4)$\cdot$충전 branch/hysteresis(Ch.5)를 닫지 않는다.
Chapter 1$\cdot$2 와 같이 \textbf{방전(탈리튬화) 방향}($s_{\phi,j}=+1$)을 기본으로 한다.

\begin{groundingbox}
\textbf{지배 표준(GS, Ch1/Ch2 계승) 및 keystone.} \textbf{(GS-1)} 모든 식은 물리적 의미를 prose 로 먼저 말한 뒤
수식화하고 중간 단계를 생략하지 않는다(학부 무비약). \textbf{(GS-2)} 결론은 출판 수준 엄밀성. \textbf{(GS-3)} 모든
가정은 통합 Assumption Ledger(\AL{\#}: 논문/교과서 근거 + 검증 DOI)를 인용하며, 근거가 없으면 FLAGGED 로
표기하고 established 로 쓰지 않는다. \textbf{(GS-4)} 임의 clip/cap/softplus/step-function 정의식과 근거 없는
수치 임계값을 물리 모델로 쓰지 않는다(속도 제한은 물리적 병목으로). \textbf{Keystone(Ch1 계승).} Level A:
$\chi_j\mathcal A_j$ 는 \emph{방향성 없는} mobility 가속(target=열역학 전담, ratio 불변). Level B(본 장):
forward 장벽만 $-\beta_j\mathcal A_j$, backward 는 $+(1-\beta_j)\mathcal A_j\Rightarrow$ rate ratio 가
$\mathcal A_j$ 에 의존(방향성 O). 본 장 transfer coefficient $\beta_j$ 는 Chapter 1 의 $\chi_j$ 와 \textbf{동일물}
이다($\chi_j\equiv\beta_j$; \emph{대칭 Marcus 1차$\cdot$정상 target 극한 한정} — 일반은 별개). 방전 기준 $s_{\phi,j}=+1$, $|I|>0$.
\end{groundingbox}

% ====================================================================
\subsection{Chapter 1--2 에서 전달받는 고정 기준식}\label{sec:ch3_inherit}
% ====================================================================
본 장은 Chapter 1$\cdot$2(RB 재구성본)의 다음을 \textbf{수정 없이} 입력으로 받는다(프로젝트 검수 P3-6). 각 식
번호는 앞 장의 boxed 결과를 가리키며, 본 장은 그 위에 미시 kinetics 계층을 \emph{적층}한다.
\begin{linkbox}
\textbf{(L1) 전하 보존식(무대).} $Q_\cell\,q=Q_\bg(V_n,T)+\sum_{j=1}^{N_p}Q_{j,\tot}\xi_j$. $V_n$ 은 그 해(외부
입력 아님); 평형 특수해가 $V_{n,\OCV}(q,T)$. 단조성 floor $\partial Q_\bg/\partial V_n\ge\epsilon_Q>0$ 로 해 유일.\\
\textbf{(L2) relaxation ODE(Level A 주역).} $\dfrac{\dd\xi_j}{\dd t}=k_j^{\mathrm{phys}}(V_n,q,T,I)\,[\xi_{\eq,j}(V_n,T)-\xi_j]$,
$k_j^{\mathrm{phys}}=k_j^{\eff}$(활성화+병목 직렬 합성). 정전류 구간서만 $\dd\xi_j/\dd t=(|I|/Q_\cell)\,\dd\xi_j/\dd q$.\\
\textbf{(L3) 활성화 계층.} $\Delta G_{\eff,j}=\Delta G_{a,j}(T)-\chi_j\mathcal A_j$,
$k_j^{\mathrm{act}}=\nu_j(T)\,e^{-\Delta G_{\eff,j}/RT}$, $1/k_j^{\mathrm{phys}}=1/k_j^{\mathrm{act}}+1/k_{j,\lim}$.
$\Delta G_{\eff,j}<0$ 은 비물리 영역이 아니라 강구동 accelerated regime.\\
\textbf{(L4) 평형 진행률(logistic).} $\xi_{\eq,j}(V_n,T)=[1+\exp(-s_{\phi,j}(V_n-U_j(T))/w_j(T))]^{-1}$, ideal
폭 $w_j=RT/F$(이 (L4)는 effective width $w_j$ 만 전달하며, $n_j^{\eff}=RT/(Fw_j)$ 의 \emph{정의}는 \S\ref{sec:ch3_neff}
에서 detailed-balance 정합으로 1회 도입한다).\\
\textbf{(L5) 네 전위 구분.} 내부 $V_n$(보존식 해), apparent $V_{n,\app}=V_n+s_I|I|R_n$(분극 포함, 평형 전위
아님), 구동 $V_{n,\drive}$(속도식 구동, 기본 $=V_n$), 평형 $V_{n,\OCV}$(평형 보존식 해).\\
\textbf{(L6) Ch2 경계.} 가역 열 = entropy($\partial U/\partial T$, 방향); 비가역 열 = flux$\times$과전압
$\eta_j=V_{n,\drive}-V_{\eq,j}(\xi_j)$($\ge0$, 속도); rate-affinity $\mathcal A_j$(중심 $U_j$ 기준)와
heat-force $\eta_j$(국소 평형 이탈) 구분; $R_n\ne R_{n,\eff}$. 본 장은 $n_j^{\eff}=1$ 특수형 한계를 일반화한다.\\
\textbf{(L7) spectrum.} 각 effective transition 내부에 activation barrier 분포 $\rho_j(g;T)$ 존재(Ch1
relaxation-length spectrum 의 근원). $\rho_j(g;T)$ 는 \textbf{Chapter 1 의 barrier 분포 $\rho_G(g;T)$ 를
transition $j$ 로 인덱싱한 동일물}이다(이후 장은 $\rho_j$ 로 통일).
\end{linkbox}
임의 clipping/softplus/bounded transform 을 기본 반응속도식으로 쓰지 않는다(GS-4); 강구동 forward rate 증가는
허용하고 제한은 물리적 병목으로 둔다.

\textbf{신규 기호(Ch1/Ch2 위에 추가).} 본 장에서만 새로 도입하는 기호를 모은다(앞 장 기호는 (L1)--(L7) 과 동일물).
\begin{longtable}{p{0.16\textwidth} p{0.13\textwidth} >{\raggedright\arraybackslash}p{0.63\textwidth}}
\toprule 기호 & 단위 & 의미 \\ \midrule \endhead
$r_j^+,\ r_j^-$ & 1/s & transition $j$ 의 forward(방전 방향)$\cdot$backward(역방향) rate \\
$k_j^{\mathrm{phys}}$ & 1/s & Level A mobility $=r_j^++r_j^-$ (평형 접근 속도; (L2) 의 $k_j$ 와 동일물). 단 이 합$=$mobility 항등은 \emph{detailed-balance split}(Level A, \S\ref{sec:ch3_split}) 하에서만 성립하며, 일반 Level B $r_j^\pm$ 에서는 합 $r_j^++r_j^-$ 이 mobility 와 다를 수 있다(식~\eqref{eq:ch3_ksum} 의 비대칭 $\mathcal A_j$-응답) \\
$\xi_{\sstat,j}$ & --- & 비평형 정상(stationary) 진행률 target $=r_j^+/(r_j^++r_j^-)$ \\
$\beta_j$ & --- & transfer coefficient ($0\le\beta_j\le1$); \textbf{$\equiv\chi_j$} (Ch1 배리어 낮춤 민감도; \emph{대칭 Marcus 1차$\cdot$정상 target 극한 한정}, 일반은 별개) \\
$\Delta G_j^{+\ddagger},\Delta G_j^{-\ddagger}$ & J/mol & forward/backward intrinsic activation free energy (구동력 $0$ 기준) \\
$\Delta H_j^{\pm\ddagger}$ & J/mol & forward/backward 활성화 엔탈피 (Eyring; 다온도 $\ln(r/T)$ vs $1/T$ 기울기 $=-\Delta H^\ddagger/R$, \S\ref{sec:ch3_ident}) \\
$\Delta S_j^{\pm\ddagger}$ & J/(mol\,K) & forward/backward 활성화 엔트로피 (Eyring; 절편 $=\Delta S^\ddagger/R+\ln(k_B/h)+\ln a$) \\
$E_{\nu,j}$ & J/mol & prefactor Arrhenius 활성화 에너지(lumped; $E_{\nu,j}\approx\Delta H^\ddagger+RT$, 식~\eqref{eq:ch3_T}). $\Delta H_{a,j}^\ddagger$ 와 동일 $1/T$ 기울기로 합만 식별 \\
$\nu_j^+,\nu_j^-$ & 1/s & forward/backward Eyring prefactor ($=(k_BT/h)\kappa^\pm$) \\
$a_j^+,a_j^-$ & --- & site availability/activity/phase accessibility factor (forward/backward) \\
$C_j(\xi_j,T)$ & --- & rate-ratio 의 configurational baseline (prefactor$\cdot$intrinsic barrier$\cdot$activity 비) \\
$n_j^{\eff}$ & --- & effective electron number $=RT/(Fw_j)$ (ideal $w_j=RT/F$ 서 $1$) \\
$\eta_n$ & V & 계면 charge-transfer 과전압 $=\phi_s-\phi_e-U_n$ (Butler--Volmer 입력) \\
$\phi_s,\phi_e$ & V & 고체상$\cdot$전해질상 전위 \\
$U_n$ & V & local equilibrium potential (과전압 기준) \\
$\eta_{n,\drive}$ & V & reduced kinetic driving correction $=V_{n,\drive}-V_n$ (기본 $0$); $\eta_n$ 과 별개 \\
$\tilde\eta_j$ & V & affinity 전위 정규화 $=\mathcal A_j/(n_j^\eff F)$, \emph{중심} $U_j$ 기준(식~\eqref{eq:ch3_Aj} 둘째식);\brk L6 heat-force $\eta_j$(\emph{국소} 평형 $V_{\eq,j}(\xi_j)$ 기준)와 별개 \\
$j_n,\ j_{0,n}$ & A/m$^2$ & 음극 반응 전류밀도,\brk exchange current density \\
$\alpha_n$ & --- & Butler--Volmer cathodic transfer coefficient (계면;\brk anodic 측 $=1-\alpha_n$). Level B $\beta_j$ 와는 \emph{상보} 관계 $\alpha_n=1-\beta_j$(동일 \emph{값} 아님; \S\ref{sec:ch3_bv}) \\
$A_\eff$ & m$^2$ & 유효 반응 면적 (전류밀도$\to$전류 환산) \\
$R_\ct,\ R_{\ct,n}$ & $\Omega$ & charge-transfer resistance (small-signal); $R_n,R_{n,\eff}$ 와 별개 \\
$I_j$ & A & transition $j$ lumped current $=Q_{j,\tot}\,\dd\xi_j/\dd t$ \\
$\mathcal A_n,\mathcal A_{\resid}$ & J/mol & generalized affinity$\cdot$residual nonideality affinity \\
$\theta$ & --- & 표면 점유율(surface coverage; $j_{0,n}$ 의 인자) \\
\bottomrule
\end{longtable}
$F=96485\ \mathrm{C/mol}$, $R=8.314\ \mathrm{J/(mol\,K)}$, $\mathrm{J/C}=\mathrm V$. rate 의 단위는
일관되게 $\mathrm{s^{-1}}$ 이다 (mobility, forward, backward 모두 진행률 $\xi_j$ 의 시간율 기여이므로).

% ====================================================================
\subsection{전위 $\cdot$ 과전압 $\cdot$ 구동력의 계층화 (CHARTER step1$\cdot$2)}\label{sec:ch3_potentials}
% ====================================================================
\textbf{물리.} Chapter 1$\cdot$2 가 분리한 네 전위(L5) 위에, 본 장은 계면 charge-transfer 의 전위 변수
$\phi_s,\phi_e,U_n,\eta_n$ 를 추가한다. 이들은 \emph{계면}의 미시 변수이고, transition 구동력 $\mathcal A_j$ 는
Chapter 1 의 \emph{reduced(transition-level)} 구동력이다 — 둘을 임의로 동일시하지 않는 것이 본 절의 핵심이다.
{\setlength{\tabcolsep}{4pt}
\begin{longtable}{p{0.14\textwidth} >{\raggedright\arraybackslash}p{0.41\textwidth} >{\raggedright\arraybackslash}p{0.37\textwidth}}
\toprule 기호 & 의미 & 본 장 역할 \\ \midrule \endhead
$V_n$ & 전하 보존식 해(내부 음극 전위) & thermodynamic state coordinate \\
$V_{n,\OCV}$ & 평형 전하 보존식 해 & 준평형 기준 \\
$V_{n,\app}$ & 분극+실험 전압축 이동 포함 apparent 전위 & 관측량(직접 과전압 아님) \\
$V_{n,\drive}$ & Ch1 속도식 구동 전위(기본 $=V_n$) & reduced driving coordinate \\
$\phi_s,\phi_e$ & 고체상/전해질상 전위 & 계면 kinetics 변수 \\
$U_n$ & local equilibrium potential & 계면 overpotential 기준 \\
$\eta_n$ & 계면 charge-transfer 과전압 & Butler--Volmer 입력 \\
\bottomrule
\end{longtable}}
\textbf{계면 과전압(정의).} 표준 계면 charge-transfer 과전압은 고체상$\cdot$전해질상 전위차에서 국소 평형
전위를 뺀 것이다 \cite{bardfaulkner,newman}(\AL{31}):
\begin{equation}
\eta_n=\phi_s-\phi_e-U_n.
\label{eq:ch3_eta_standard}
\end{equation}
이는 \emph{관측} apparent 전위 $V_{n,\app}$ 와 같지 않다 — $V_{n,\app}=V_n+s_I|I|R_n$ 에는 ohmic drop,
concentration polarization, surface film, fitting residual 이 섞이기 때문이다(\S\ref{sec:ch3_Rn}; $R_n$ 분해).

\textbf{구동력(CHARTER step1 — $s_{\phi,j}$ 명시형 기준).} Chapter 1 transition 구동력(reduced affinity)을
\emph{본 장에서 기준형으로 동결}한다. 부호 인자 $s_{\phi,j}$ 와 effective electron number $n_j^{\eff}$ 를 명시한
일반형은(\AL{31}; RB\_CHARTER step1 의 $\mathcal A_j=s_{\phi,j}n_j^\eff F(V_\drive-U_j)$ 명시형 기준장)
\begin{equation}
\boxed{\;\mathcal A_j=s_{\phi,j}\,n_j^{\eff}\,F\,[V_{n,\drive}-U_j(T)]\;,\qquad \tilde\eta_j\equiv\frac{\mathcal A_j}{n_j^{\eff}F}=s_{\phi,j}[V_{n,\drive}-U_j(T)]\;.}
\label{eq:ch3_Aj}
\end{equation}
$s_{\phi,j}\in\{+1,-1\}$ 는 $\mathcal A_j>0$ 이면 방전 방향 진행이 촉진되도록 선택한다(방전 기본 $s_{\phi,j}=+1$). 이
$s_{\phi,j}$ 는 (L4) logistic 부호 인자와 \emph{동일물}이다($\mathcal A_j>0\Leftrightarrow V_{n,\drive}>U_j\Leftrightarrow\xi_{\eq,j}>1/2$ 정렬).
\textbf{차원 검산.} $n_j^{\eff}$(무차원)$\cdot F[\mathrm{C/mol}]\cdot[\mathrm V]=\mathrm{C\cdot V/mol}=\mathrm{J/mol}$
이므로 $\mathcal A_j$ 는 몰당 자유에너지 구동력이다. $n_j^{\eff}$ 는 \S\ref{sec:ch3_neff} 에서 detailed balance
정합으로 \emph{고정}되며, ideal limit $n_j^{\eff}=1$ 이면 Chapter 1 원형 $\mathcal A_j=s_{\phi,j}F(V_{n,\drive}-U_j)$
로 환원된다(본 장이 일반화하는 대상). 식~\eqref{eq:ch3_Aj} 둘째 식은 $\tilde\eta_j$ 가 $\mathcal A_j$ 를 $n_j^\eff F$
로 정규화한 \emph{전위 단위}[V] 구동력임을 보인다(Ch2 의 dissipation force $\eta_j$ 와 같은 \emph{단위}이지만 \emph{별개
기호}다). 둘의 관계는 $\tilde\eta_j=s_{\phi,j}[V_{n,\drive}-U_j]$(\emph{중심} $U_j$ 기준)와 (L6) heat-force
$\eta_j=V_{n,\drive}-V_{\eq,j}(\xi_j)$(\emph{현재} $\xi_j$ 의 \emph{국소} 평형 기준)로, $s_{\phi,j}=+1$ 기본형에서
$\tilde\eta_j-\eta_j=V_{\eq,j}(\xi_j)-U_j=w_j\ln[\xi_j/(1-\xi_j)]$(\S\ref{sec:ch3_db} 검산에서 재확인). rate 의
중심-기준 구동력에는 $\tilde\eta_j$(또는 $\mathcal A_j$), dissipation 에는 국소-평형 $\eta_j$ 를 쓴다.

\textbf{구동 전위(CHARTER step2 — $V_{n,\drive}$ 정의장).} 속도식이 쓰는 구동 전위와 내부 전위의 가장 보수적
연결을 \emph{본 장에서 정의로 동결}한다(RB\_CHARTER step2 의 $V_{n,\drive}$ 정의장):
\begin{equation}
\boxed{\;V_{n,\drive}=V_n+\eta_{n,\drive}\;,\qquad \text{기본형}\ \eta_{n,\drive}=0\Rightarrow V_{n,\drive}=V_n\;,\qquad \eta_{n,\drive}\ne\eta_n\;.}
\label{eq:ch3_drive}
\end{equation}
$\eta_{n,\drive}$ 는 Chapter 1 의 reduced kinetic driving correction 이며, 기본 전개에서는 $0$(즉
$V_{n,\drive}=V_n$, $\eta_{n,\drive}=0$)으로 둔다. 이는 RB\_CHARTER step2 의 ``$\eta_j$ 의 기준 전위를 한
문서 내 단일($V_{n,\drive}$)로 유지, 기본 $V_{n,\drive}=V_n$''와 정합한다.
\begin{boundbox}
\textbf{$\eta_{n,\drive}$ 와 계면 $\eta_n$ 의 임의 동일시 금지(\AL{32}).} 둘은 \emph{서로 다른 층위}다:
$\eta_{n,\drive}$ 는 transition-level reduced 구동력 보정(중심 $U_j$ 기준 affinity 의 일부), $\eta_n$ 은
계면 charge-transfer 과전압(국소 평형 $U_n$ 기준). 둘을 연결하려면 $V_n,\phi_s,\phi_e,U_n$ 의 대응과 reference
electrode convention 이 별도로 필요하므로, 본 장 기본 전개는 $\eta_{n,\drive}=0$ 으로 두고 강구동 일반화
($\eta_{n,\drive}\ne0$)는 BOUNDED 로 표기한다. 기본형에서 $V_{n,\drive}=V_n$ 이므로 아래 detailed balance
정합(\S\ref{sec:ch3_neff})과 keystone 검증(\S\ref{sec:ch3_consistency})은 \emph{이 전제 하에서} 닫힌다.
\end{boundbox}

% ====================================================================
\subsection{Forward/backward transition kinetics 와 detailed balance}\label{sec:ch3_db}
% ====================================================================
\textbf{물리.} 상변이 진행률 $\xi_j$ 는 가장 일반적으로 \emph{앞으로 가는 transition}(미점유$\to$점유, 방전
방향)과 \emph{뒤로 가는 transition}(점유$\to$미점유)의 \emph{차}로 변한다. forward 는 아직 진행 안 한 fraction
$(1-\xi_j)$ 에, backward 는 이미 진행한 fraction $\xi_j$ 에 비례한다(2-상태 mass-action; \cite{eyring1935,degrootmazur1962},
\AL{30}):
\begin{equation}
\boxed{\;\frac{\dd\xi_j}{\dd t}=r_j^+\,(1-\xi_j)-r_j^-\,\xi_j\;.}
\label{eq:ch3_fb}
\end{equation}
$r_j^+$=방전 방향 transition rate[1/s], $r_j^-$=역방향[1/s]. \textbf{차원 검산.} $r_j^\pm[\mathrm{s^{-1}}]\times$
무차원 fraction $=\mathrm{s^{-1}}$ 이고 좌변 $\dd\xi_j/\dd t$ 도 $\mathrm{s^{-1}}$($\xi_j$ 무차원) — 정합.
\textbf{극한 검산.} $r_j^-=0$(역반응 없음)이면 $\dd\xi_j/\dd t=r_j^+(1-\xi_j)$ 로 단조 포화($\xi_j\to1$);
$r_j^+=0$ 이면 $\dd\xi_j/\dd t=-r_j^-\xi_j$ 로 단조 소멸($\xi_j\to0$) — 부호 직관과 일치.

\textbf{평형 = detailed balance(무비약).} 평형은 net flux 가 $0$ 인 정상점이다. 식~\eqref{eq:ch3_fb} 에서
$\dd\xi_j/\dd t=0$ 과 $\xi_j=\xi_{\eq,j}$ 를 놓으면
\begin{equation}
r_j^+(1-\xi_{\eq,j})=r_j^-\xi_{\eq,j}
\ \Longrightarrow\
\frac{r_j^+}{r_j^-}=\frac{\xi_{\eq,j}}{1-\xi_{\eq,j}}\quad(\text{local detailed balance}).
\label{eq:ch3_db}
\end{equation}
(둘째 등호: 양변을 $r_j^-(1-\xi_{\eq,j})$ 로 나눔; \emph{열린구간} $0<\xi_{\eq,j}<1$ 에서 분모 $\ne0$.) 즉 \emph{rate 의 비가
평형 점유율의 odds 와 같다}는 것이 detailed balance 의 내용이다. 경계 $r_j^-\to0$ 은 비를 발산시키므로, 이 경우
ratio 형 대신 \emph{곱형} detailed balance $r_j^+(1-\xi_{\eq,j})=r_j^-\xi_{\eq,j}$(식~\eqref{eq:ch3_db} 의 좌식)로
극한 $\xi_{\eq,j}\to1$ 을 닫는다($r_j^+=0\Rightarrow\xi_{\eq,j}\to0$ 동형).

\textbf{Chapter 1 logistic 대입(keystone 1차 결과; CHARTER step3 의 $w_j$ scale).} Chapter 1 평형 진행률
(L4, $s_{\phi,j}=+1$) $\xi_{\eq,j}=[1+\exp(-(V_n-U_j)/w_j)]^{-1}$ 의 odds 를 계산한다. $\xi_{\eq,j}=E/(1+E)$,
$E\equiv\exp[(V_n-U_j)/w_j]$ 로 두면 $1-\xi_{\eq,j}=1/(1+E)$ 이므로 $\xi_{\eq,j}/(1-\xi_{\eq,j})=E$. 이를
식~\eqref{eq:ch3_db} 에 넣고 양변에 자연로그를 취하면(중간 단계 전부 명시: $\ln E=(V_n-U_j)/w_j$)
\begin{equation}
\boxed{\;\ln\frac{r_j^+}{r_j^-}=\ln\frac{\xi_{\eq,j}}{1-\xi_{\eq,j}}=\frac{V_n-U_j(T)}{w_j(T)}\;.}
\label{eq:ch3_db_width}
\end{equation}
이것이 \AL{30}(forward/backward + detailed balance)의 boxed 결과다(\emph{$s_{\phi,j}=+1$ 방전 기본형}; 이 식은
\emph{평형점}에서만 성립하고, 비평형 일반 ratio $\ln(r_j^+/r_j^-)=C_j+\mathcal A_j/RT$ 는 \S\ref{sec:ch3_levelB}
식~\eqref{eq:ch3_ratio_B} 에서 도입한다). \textbf{차원 검산.} 좌변은 rate 비의
로그(무차원); 우변은 $[\mathrm V]/[\mathrm V]$(전위/폭) 무차원 — 정합. \textbf{이 식은 $w_j$ 를 $RT/F$ 에 강제로
묶지 않는다}: $w_j$ 는 nonideality$\cdot$domain heterogeneity$\cdot$finite width 를 포함하는 effective
equilibrium width 이며(Chapter 1 (L4) 의 effective width 와 동일물), 그 effective 성격이 \S\ref{sec:ch3_neff}
에서 $n_j^{\eff}=RT/(Fw_j)$ 로 정량화된다(CHARTER step3).
\begin{boundbox}
\textbf{rate-affinity $\mathcal A_j$ 와 detailed-balance 우변의 관계(\AL{30}; Ch2 (L6) 정합).} 식~\eqref{eq:ch3_db_width}
우변 $(V_n-U_j)/w_j$ 는 \emph{중심} $U_j$ 기준이라 식~\eqref{eq:ch3_Aj} 의 $\tilde\eta_j=s_{\phi,j}(V_{n,\drive}-U_j)$
와 같은 ``중심 기준''이고($V_{n,\drive}=V_n$ 기본형서 $\tilde\eta_j=V_n-U_j$, $s_{\phi,j}=+1$), Ch2 의 dissipation
force(국소 평형 $V_{\eq,j}(\xi_j)$ 기준)와는 다르다. 즉 detailed balance 의 ln 비는 \emph{평형 odds}(중심 기준)
를 정하고, Ch2 의 $\eta_j=V_{n,\drive}-V_{\eq,j}(\xi_j)$ 는 \emph{현재 상태의} 평형 이탈을 잰다. 둘의 차이가
$V_{\eq,j}(\xi_j)-U_j=w_j\ln[\xi_j/(1-\xi_j)]$(Chapter 2 의 configurational 가역 entropy 항)이다 — rate 에는
$\mathcal A_j$(중심), dissipation 에는 $\eta_j^{\mathrm{Ch2}}$(국소 평형)를 쓰라는 Ch2 (L6)$\cdot$\AL{28} 의
구분이 본 장에서도 보존된다(BOUNDED — logistic 평형 모델 내에서 정확).
\end{boundbox}

% ====================================================================
\subsection{Chapter 1 relaxation ODE 와의 정합 (Level A = Level B 의 detailed-balance 극한)}\label{sec:ch3_consistency}
% ====================================================================
\subsubsection{consistent split (Level A 회복; 무비약 항등)}\label{sec:ch3_split}
\textbf{물리.} Chapter 1 의 $\dd\xi_j/\dd t=k_j^{\mathrm{phys}}(\xi_{\eq,j}-\xi_j)$ 는 식~\eqref{eq:ch3_fb} 의
\emph{특수 축약형}이다. 이것을 보장하는 forward/backward split 을 명시한다(\AL{30}):
\begin{equation}
r_j^+=k_j^{\mathrm{phys}}\,\xi_{\eq,j},\qquad r_j^-=k_j^{\mathrm{phys}}\,(1-\xi_{\eq,j}).
\label{eq:ch3_split}
\end{equation}
\textbf{대입(무비약, 한 줄씩).} 이 split 을 식~\eqref{eq:ch3_fb} 에 넣고 전개하면
\begin{align}
\frac{\dd\xi_j}{\dd t}
&=k_j^{\mathrm{phys}}\xi_{\eq,j}(1-\xi_j)-k_j^{\mathrm{phys}}(1-\xi_{\eq,j})\xi_j \notag\\
&=k_j^{\mathrm{phys}}\big[\xi_{\eq,j}-\xi_{\eq,j}\xi_j-\xi_j+\xi_{\eq,j}\xi_j\big]
=k_j^{\mathrm{phys}}\big[\xi_{\eq,j}-\xi_j\big].
\label{eq:ch3_recover}
\end{align}
(둘째 줄: 첫 괄호 $\xi_{\eq,j}(1-\xi_j)=\xi_{\eq,j}-\xi_{\eq,j}\xi_j$, 둘째 괄호 $(1-\xi_{\eq,j})\xi_j=\xi_j-\xi_{\eq,j}\xi_j$;
$-\xi_{\eq,j}\xi_j$ 와 $+\xi_{\eq,j}\xi_j$ 가 \emph{정확히 상쇄}되어 Chapter 1 식이 복원된다.) 이 split 은
detailed balance(식~\eqref{eq:ch3_db})를 \emph{자동으로} 만족하며($r_j^+/r_j^-=\xi_{\eq,j}/(1-\xi_{\eq,j})$),
$\xi_{\eq,j}$ 를 전류 의존 target 으로 바꾸지 않는다.
\begin{boundbox}
\textbf{keystone 닫힘조건(★ split$=$일반 rate 평형극한의 필요충분조건; \AL{30}).} 위 대칭 split 이 Level B
일반 rate(식~\eqref{eq:ch3_rplus},\allowbreak\eqref{eq:ch3_rminus})의 \emph{평형 극한}과 동일물이 되는 필요충분조건은,
평형에서 $r_j^+/r_j^-$ 가 \emph{$\xi_j$-무관}(즉 $C_j(\xi_j,T)$ 의 $\xi_j$-의존이 detailed balance 와 상충하지
않음)이며, 이는 \S\ref{sec:ch3_neff} 의 $n_j^{\eff}$ 고정과 \emph{동일 closure}다. 이 조건이 깨지면(예: $a_j^\pm$ 의
강한 $\xi_j$ 의존) split 은 detailed balance 를 만족해도 일반 rate 의 정상점과 어긋날 수 있다(\S\ref{sec:ch3_neff}
verifybox 의 (ii) 와 동일 입도).
\end{boundbox}

\begin{keybox}
\textbf{Chapter 1 의 $k_j^{\mathrm{phys}}$ 는 forward-only rate 가 아니라 mobility 다 (Level A).} 식~\eqref{eq:ch3_split}
에서 $r_j^++r_j^-=k_j^{\mathrm{phys}}[\xi_{\eq,j}+(1-\xi_{\eq,j})]=k_j^{\mathrm{phys}}$ 이므로 Chapter 1 의
$k_j^{\mathrm{phys}}$ 는 정확히 \emph{forward 와 backward 의 합}(평형 접근 속도, mobility)이다 — Chapter 1
keystone 의 ``$k_j=r_j^++r_j^-$''(Ch1 §forward/backward 유도 결과를 (L2) 적층 입력으로 받음)와 동일하다(기본형
$\eta_{n,\drive}=0$, 식~\eqref{eq:ch3_drive} 전제). 그리고 Chapter 1 의 $\chi_j\mathcal A_j$
가 $k_j^{\mathrm{phys}}$ \emph{전체}에 들어가면 $r_j^+,r_j^-$ 가 \emph{같은 배율}로 커져 비
$r_j^+/r_j^-=\xi_{\eq,j}/(1-\xi_{\eq,j})$ 가 \emph{불변}이다 — 방향(target)은 열역학(\S\ref{sec:ch3_db})이
전담하고 $\chi_j\mathcal A_j$ 는 \emph{방향 없는} 속도 scale 일 뿐이다(\textbf{Level A=방향성 X}). \emph{방향을
가진} 배리어 낮춤(forward 만 $-\beta\mathcal A$, backward 는 $+(1-\beta)\mathcal A$ 라 비가 바뀜)은
\textbf{Level B}(\S\ref{sec:ch3_levelB})에서 도입하며, 그때 $\chi_j$ 가 transfer coefficient $\beta_j$ 와
\emph{동일물}이 된다($\chi_j\equiv\beta_j$; 조건부 동일물 — 대칭 Marcus$\cdot$정상 target 극한 한정).
\end{keybox}

\subsubsection{nonequilibrium stationary target $\xi_{\sstat,j}$}\label{sec:ch3_xiss}
\textbf{물리.} 전류/과전압이 forward/backward 비 \emph{자체}를 바꾸면(Level B), 진행이 멈추는 새 정상점이
생긴다. 식~\eqref{eq:ch3_fb} 에서 $\dd\xi_j/\dd t=0$ 을 \emph{일반} $r_j^\pm$ 에 대해 풀면(평형을 가정하지 않고)
$r_j^+(1-\xi_{\sstat,j})=r_j^-\xi_{\sstat,j}\Rightarrow r_j^+=(r_j^++r_j^-)\xi_{\sstat,j}$ 이므로
\begin{equation}
\boxed{\;\xi_{\sstat,j}(V_n,T,I)=\frac{r_j^+}{r_j^++r_j^-}\;.}
\label{eq:ch3_xiss}
\end{equation}
\textbf{검산.} $r_j^-=0$ 이면 $\xi_{\sstat,j}=1$(완전 진행), $r_j^+=0$ 이면 $\xi_{\sstat,j}=0$ — 극한 직관과 일치;
$0\le\xi_{\sstat,j}\le1$(두 비음 rate 의 비). 이는 true thermodynamic equilibrium $\xi_{\eq,j}$ 가 아니라
\emph{nonequilibrium stationary progress} 다. detailed balance(식~\eqref{eq:ch3_db})가 성립하면, \emph{split 을
가정하지 않고} 일반형 $\xi_{\sstat,j}=r_j^+/(r_j^++r_j^-)$ 에 detailed balance($r_j^+/r_j^-=\xi_{\eq,j}/(1-\xi_{\eq,j})$)를
직접 대입하여
\begin{equation*}
\xi_{\sstat,j}=\frac{r_j^+}{r_j^++r_j^-}=\frac{r_j^+/r_j^-}{r_j^+/r_j^-+1}
=\frac{\xi_{\eq,j}/(1-\xi_{\eq,j})}{\xi_{\eq,j}/(1-\xi_{\eq,j})+1}=\xi_{\eq,j}
\end{equation*}
로 환원된다(순환 없음). 이때 consistent split(식~\eqref{eq:ch3_split})은 이 환원을 달성하는 \emph{한 표현}임이
사후 확인된다. Chapter 5
의 branch target 도 이 $\xi_{\sstat,j}$ 계열이며 true equilibrium 과 구분한다(\S\ref{sec:ch3_transmit}).

% ====================================================================
\subsection{Level B: 명시적 활성화 rate 와 directional barrier lowering}\label{sec:ch3_levelB}
% ====================================================================
\textbf{물리.} 이제 mobility 를 ``속도 한 덩어리''로 두지 않고, forward/backward 각각의 \emph{명시적} Eyring
rate 로 세운다. Chapter 1 의 ``배리어 낮춤''을 \emph{방향성 있게} 분해하는 것이 본 절이다. transfer coefficient
$0\le\beta_j\le1$ 가 구동력 $\mathcal A_j$ 를 forward 와 backward 에 \emph{비대칭}으로 배분한다 — forward 장벽은
$\beta_j\mathcal A_j$ 만큼 낮아지고, backward 장벽은 $(1-\beta_j)\mathcal A_j$ 만큼 \emph{높아진다}
(\cite{evanspolanyi1938,bardfaulkner,bazant2013}, \AL{31}). Eyring rate $r=\nu\,a\,e^{-\Delta G^\ddagger/RT}$ 에
이 비대칭 장벽을 넣으면
\begin{equation}
r_j^+=\nu_j^+(T)\,a_j^+(\xi_j,V_n,T)\,\exp\!\Big[-\frac{\Delta G_j^{+\ddagger}(T)-\beta_j\mathcal A_j}{RT}\Big],
\label{eq:ch3_rplus}
\end{equation}
\begin{equation}
r_j^-=\nu_j^-(T)\,a_j^-(\xi_j,V_n,T)\,\exp\!\Big[-\frac{\Delta G_j^{-\ddagger}(T)+(1-\beta_j)\mathcal A_j}{RT}\Big].
\label{eq:ch3_rminus}
\end{equation}
$a_j^\pm$=site availability/activity/phase accessibility (무차원), $\nu_j^\pm$=Eyring prefactor[1/s].
\textbf{차원 검산.} 지수 인자 $\Delta G^\ddagger,\beta_j\mathcal A_j$ 모두 $\mathrm{J/mol}$, $RT$ 도
$\mathrm{J/mol}$ 이라 지수는 무차원; $\nu_j^\pm[\mathrm{s^{-1}}]\times a_j^\pm$(무차원)$=\mathrm{s^{-1}}$ — rate
단위 정합. \textbf{forward 가속의 부호.} $\mathcal A_j>0$(방전 구동)일 때 forward 지수의 $-(-\beta_j\mathcal A_j)/RT
=+\beta_j\mathcal A_j/RT>0$ 이라 $r_j^+$ 증가, backward 지수 $-(1-\beta_j)\mathcal A_j/RT<0$ 이라 $r_j^-$ 감소 —
방전 방향으로 net 진행이 촉진(방향성 O).

\textbf{rate ratio 가 $\mathcal A_j$ 에 의존한다(무비약).} 두 식의 비의 로그를 취한다. 지수의 차가 로그의 차이므로
\begin{equation}
\ln\frac{r_j^+}{r_j^-}=\underbrace{\ln\frac{\nu_j^+a_j^+}{\nu_j^-a_j^-}-\frac{\Delta G_j^{+\ddagger}-\Delta G_j^{-\ddagger}}{RT}}_{\equiv\,C_j(\xi_j,T)}
+\frac{\beta_j\mathcal A_j+(1-\beta_j)\mathcal A_j}{RT}.
\label{eq:ch3_ratio_raw}
\end{equation}
$\mathcal A_j$ 항의 계수가 $\beta_j+(1-\beta_j)=1$ 로 \emph{합쳐지므로}(forward 의 $+\beta_j\mathcal A_j$ 와
backward 장벽 증가가 비에서 $-[-(1-\beta_j)\mathcal A_j]=+(1-\beta_j)\mathcal A_j$ 로 더해짐),
\begin{equation}
\boxed{\;\ln\frac{r_j^+}{r_j^-}=C_j(\xi_j,T)+\frac{\mathcal A_j}{RT}\;.}
\label{eq:ch3_ratio_B}
\end{equation}
\textbf{핵심 관찰($\beta_j$ 의 상쇄).} \emph{ratio 의 $\mathcal A_j$-기울기는 $1/RT$ 로 $\beta_j$ 에 무관}하다
($\beta_j$ 가 식~\eqref{eq:ch3_ratio_raw}$\to$\eqref{eq:ch3_ratio_B} 에서 상쇄). 즉 $\beta_j$ 는 ratio(=방향)를
정하지 않고, \emph{sum} $r_j^++r_j^-$(=mobility)의 $\mathcal A_j$-응답 \emph{비대칭}만 결정한다. mobility 의
명시형은(\AL{31})
\begin{equation}
k_j^{\mathrm{phys}}=r_j^++r_j^-
=\nu_j^+a_j^+\,e^{-\Delta G_j^{+\ddagger}/RT}\,e^{\beta_j\mathcal A_j/RT}
+\nu_j^-a_j^-\,e^{-\Delta G_j^{-\ddagger}/RT}\,e^{-(1-\beta_j)\mathcal A_j/RT}.
\label{eq:ch3_ksum}
\end{equation}
\textbf{$\beta_j$ 의 역할(극한 검산).} $\mathcal A_j=0$(구동력 없음)이면 두 지수가 $1$ 이라 $k_j^{\mathrm{phys}}$ 가
$\beta_j$ 무관($=\nu_j^+a_j^+e^{-\Delta G^{+\ddagger}/RT}+\nu_j^-a_j^-e^{-\Delta G^{-\ddagger}/RT}$); $\mathcal A_j>0$
이면 forward 항이 $e^{\beta_j\mathcal A_j/RT}$ 로 커지고 backward 항이 $e^{-(1-\beta_j)\mathcal A_j/RT}$ 로 작아져,
$\beta_j$ 가 클수록 mobility 의 $\mathcal A_j$-증가가 가팔라진다. $\beta_j=1/2$(대칭)는 Chapter 1 의 대칭 Marcus
극한과 같다(거기서 ``공통 배율''로 평균낸 그 대칭).
\begin{boundbox}
\textbf{Marcus bound 계승(\AL{33}).} forward/backward 장벽의 \emph{선형} 이동($\mp\beta_j\mathcal A_j$,
$\mp(1-\beta_j)\mathcal A_j$)은 Chapter 1 과 같이 작은 구동력 1차 근사다. $|\mathcal A_j|\sim\lambda$(재구성
에너지)에서 Marcus 곡률(역전 영역)로 깨진다 \cite{marcus1956}. 본 식은 꼬리 영역 $V_{n,\drive}\approx U_j$ 에서
유효하며 전역 정당화가 아니다(BOUNDED; Chapter 1 \AL{15} 의 Level B 대응).
\end{boundbox}

\subsubsection{detailed-balance 정합 제약: $n_j^{\eff}=RT/(Fw_j)$ (CHARTER step3 정의장)}\label{sec:ch3_neff}
\textbf{물리.} Level B 의 rate ratio(식~\eqref{eq:ch3_ratio_B})가 Chapter 1 의 평형(식~\eqref{eq:ch3_db_width})과
\emph{모든 $V_n$ 에서} 일치해야 한다(detailed balance 는 평형에서 성립). 이 정합이 effective electron number
$n_j^{\eff}$ 를 \emph{고정}한다. 기본형 $V_{n,\drive}=V_n$($\eta_{n,\drive}=0$, 식~\eqref{eq:ch3_drive})에서
\begin{equation}
C_j+\frac{\mathcal A_j}{RT}\Big|_{V_{n,\drive}=V_n}=\frac{V_n-U_j}{w_j}.
\label{eq:ch3_dbmatch}
\end{equation}
$\mathcal A_j=s_{\phi,j}n_j^{\eff}F(V_n-U_j)$(식~\eqref{eq:ch3_Aj}, $s_{\phi,j}=+1$)를 넣으면 좌변의 $V_n$-의존
항은 $n_j^{\eff}F(V_n-U_j)/(RT)$ 다. (\emph{전제}: $U_j,w_j,n_j^{\eff}$ 는 $V_n$-무의존인 온도 함수이므로 아래
$V_n$-미분에서 상수로 취급한다.) 우변과 \emph{$V_n$ 전 구간}에서 같으려면 (i) $V_n$ 에 \emph{선형}인 항의
기울기가 같고 (ii) 상수항이 같아야 한다. (i) 기울기 정합: $V_n$ 으로 미분하면 $C_j$ 가 $V_n$-무의존이라는
가정(아래 boundbox) 하에 좌변 기울기 $=n_j^{\eff}F/(RT)$, 우변 기울기 $=1/w_j$ 이므로
\begin{equation}
\boxed{\;n_j^{\eff}=\frac{RT}{F\,w_j(T)}\;,\qquad C_j\big|_{\text{midpoint }V_n=U_j}=0\;.}
\label{eq:ch3_neff}
\end{equation}
(ii) 상수항: $V_n=U_j$(중심)에서 우변 $=0$ 이고 좌변 $\mathcal A_j$ 항 $=0$ 이므로 $C_j|_{\text{midpoint}}=0$.
\textbf{차원 검산.} $RT[\mathrm{J/mol}]/(F[\mathrm{C/mol}]\,w_j[\mathrm V])=\mathrm{J/(C\cdot V)}=\mathrm{V/V}=1$
(무차원) — $n_j^{\eff}$ 가 전자수(무차원)로 올바르다. \textbf{함의.} effective width $w_j$ 와 effective electron
number $n_j^{\eff}$ 는 독립이 아니라 $n_j^{\eff}w_j=RT/F$ 로 묶인다(ideal $w_j=RT/F\Rightarrow n_j^{\eff}=1$).
\textbf{이 제약이 없으면} BV-형 barrier lowering($F\eta/RT$ scale)과 logistic 평형 폭 $w_j$ 가 불일치한다 —
모든 prefactor/barrier/transfer coefficient 를 독립 자유 피팅하면 thermodynamic consistency 가 깨진다(P3 검수).
이것이 RB\_CHARTER step3 의 ``$n_j^{\eff}=RT/(Fw_j)$ 정의를 도입 장(Ch3)에 1회''에 해당하는 \emph{정의장}이다.
\begin{boundbox}
\textbf{정합 가정 명시($C_j$ 의 $V_n$-무의존; \AL{34}).} 식~\eqref{eq:ch3_neff} 의 기울기 정합은 $C_j$ 가 $V_n$
에 무의존(활성도비 $a_j^+/a_j^-$ 의 전위 의존을 $\mathcal A_j$ 항으로 모두 흡수)이라는 가정에 의존한다. 일반적으로
ln 비의 전 $V_n$ 의존은 $C_j$(configurational)와 $\mathcal A_j/RT$(potential)로 분할되며 그 partition 은 모델
선택이다. 본 장은 \emph{전체 ln 비가 detailed balance(식~\eqref{eq:ch3_db_width}) $=(V_n-U_j)/w_j$ 와 같아야
한다}는 \textbf{model-independent 제약}을 우선하고, 전위 의존을 $\mathcal A_j$ 에 귀속하는 표준 partition 에서
$n_j^{\eff}=RT/(Fw_j)$ 를 얻는다($s_{\phi,j}=+1$). $a_j^\pm$ 가 강한 $V_n$ 의존을 가지면 effective width 가
재정의되므로, partition 을 명시한 뒤 $n_j^{\eff}$ 를 고정한다(BOUNDED).
\end{boundbox}
\begin{boundbox}
\textbf{Chapter 1 과의 reconcile(\AL{34}; Ch1 식 형태 불변).} Chapter 1 본문은 $\mathcal A_j=s_{\phi,j}F(V_{n,\drive}-U_j)$
(즉 $n_j^{\eff}=1$)에 \emph{effective} $w_j$ 를 함께 썼다. 위 제약상 이는 $w_j=RT/F$(ideal)에서만 정확히
정합한다. 따라서 nonideal effective width($w_j\ne RT/F$)를 쓸 때는 \textbf{Chapter 1 의 $\mathcal A_j$ 의 $F$ 도
$n_j^{\eff}F=RT/w_j$ 로 읽어야} 정합이 유지된다 — Chapter 1 식의 \emph{형태}를 바꾸지 않고 affinity scale 의
\emph{해석}만 정밀화한 것이며(Chapter 1 의 logistic baseline FLAG 와 같은 층위의 단서), 앞 장과 충돌하지 않는다.
\end{boundbox}

\begin{verifybox}
\textbf{Level A $=$ Level B 의 detailed-balance 극한 (CHARTER step4 keystone 검증).} 식~\eqref{eq:ch3_xiss} 에
식~\eqref{eq:ch3_ratio_B} 를 대입한다. $\xi_{\sstat,j}=r_j^+/(r_j^++r_j^-)=1/(1+r_j^-/r_j^+)$ 이고, 여기서
$r_j^-/r_j^+=\exp(-\ln(r_j^+/r_j^-))$($\because\ \ln(r_j^-/r_j^+)=-\ln(r_j^+/r_j^-)$)이므로
$\xi_{\sstat,j}=[1+e^{-\ln(r_j^+/r_j^-)}]^{-1}=[1+e^{-(C_j+\mathcal A_j/RT)}]^{-1}$. 정합 제약(식~\eqref{eq:ch3_neff}; $C_j$ 무의존 가정, $V_{n,\drive}=V_n$)
하에서 $C_j+\mathcal A_j/RT=(V_n-U_j)/w_j$ 이므로 $\xi_{\sstat,j}=[1+e^{-(V_n-U_j)/w_j}]^{-1}=\xi_{\eq,j}$. 즉
\textbf{Chapter 1 의 평형 target 은 Level B 의 detailed-balance \emph{stationary 극한}}이며, 두 계층은 충돌하지
않는다(P3-6 통과). \textbf{한정(★ CHARTER step4 — 무조건 ``$A=B$'' 금지).} 일치하는 것은 \emph{정상 target}
$\xi_{\sstat,j}\to\xi_{\eq,j}$ 이고, 그 조건은 \emph{둘}이다: (i) detailed balance($C_j$ 의 $V_n$-무의존 partition),
\emph{그리고} (ii) ratio $r_j^+/r_j^-$(즉 $C_j$)가 $\xi_j$\emph{-무관}이어야 한다. (여기서 무관 대상은 \emph{현재}
$\xi_j$(점유율)이며, ratio 의 $\xi_{\eq,j}(=V_n$ 함수$)$-의존은 detailed balance 가 \emph{요구}하는 것이라
모순이 아니다.) (ii)가 있어야 $\xi_{\sstat,j}=
r_j^+/(r_j^++r_j^-)$ 가 음함수(implicit fixed point)가 아닌 \emph{닫힌형}으로 평가되어 $\xi_{\eq,j}$ 와 일치한다
— 이는 Chapter 1 keystone 의 ``$r^\pm$ 를 $\xi_j$-무관 국소 상수로 본다''는 가정과 \emph{동일 입도}이며,
식~\eqref{eq:ch3_split} 의 대칭 split 이 이를 만족한다($a_j^\pm$ 가 강한 $\xi_j$ 의존을 가지면 detailed balance 가
성립해도 $\xi_{\sstat,j}\ne\xi_{\eq,j}$ 가능). Level A 는 forward/backward \emph{대칭}(공통 배율, ratio 불변,
방향성 X) 특수극한이고, Level B 는 \emph{비대칭}($\beta_j$ 가 sum 의 $\mathcal A_j$-응답을 비대칭으로, 방향성 O)
일반형이다. branch 비대칭 $\xi_{\sstat,j}\ne\xi_{\eq,j}$ 은 $C_j$ 가 detailed balance 를 이탈하거나 $r_j^+/r_j^-$ 가
$\xi_j$-의존일 때(독립 landscape, metastable branch) 발생하며 Chapter 5 로 넘긴다. 이 한정은
Chapter 1 keybox 의 ``정상 target 극한 $\xi_{\sstat,j}\to\xi_{\eq,j}$ 한정''(좁은 표현)과 \emph{동일 입도}다.
\end{verifybox}

\subsubsection{강구동 accelerated regime 과 물리적 병목(cap 금지)}\label{sec:ch3_limiter}
강구동에서 forward 유효 장벽 $\Delta G_j^{+\ddagger}-\beta_j\mathcal A_j$ 가 음수가 될 수 있다. 이를 $0$ 에 잘라
붙이지 않는다(GS-4); 이는 비물리 영역이 아니라 단순 activation 식이 \emph{매우 큰} forward rate 를 예측하는
accelerated regime 이다. 실제 제한은 \emph{직렬 병목}(시간척도의 합)으로 나타난다 — 수치 cap 이 아니라 율속
단계 합성이며, Chapter 1 의 $1/k_j^{\eff}=1/k_j+1/k_{j,\lim}$ 를 \emph{(강구동 branch 용) forward-only
template} 으로 한정 일반화한 것이다(\AL{35}):
\begin{equation}
\frac{1}{r_{j,\obs}^+}=\frac{1}{r_j^+}+\frac{1}{r_{j,\trans}^+}+\frac{1}{r_{j,\site}^+}+\frac{1}{r_{j,\curr}^+}.
\label{eq:ch3_forward_limiter}
\end{equation}
$r_{j,\trans}^+$=전해질/고체 수송, $r_{j,\site}^+$=유한 활성 자리, $r_{j,\curr}^+$=유한 외부 전류(전류 보존,
\S\ref{sec:ch3_current}) 병목. $r_j^+\to\infty$(활성화 지배)면 $r_{j,\obs}^+\to(\sum_{k\ne\mathrm{act}}1/r_k^+)^{-1}$
(나머지 병목들의 \emph{조화 합}; reciprocal 합이므로 단순 $\min$ 이 아님)으로 매끄럽게 포화하고, 한 병목이
지배적이면 그 \emph{최소} 병목 rate 로 근사된다; 병목이 모두 $\to\infty$ 면 $r_{j,\obs}^+\to r_j^+$ 로 환원된다. \textbf{병목의 비대칭 주의}: forward
에만 작용하는 병목은 $r_j^+/r_j^-$ 를 바꿔 detailed balance 를 이탈시킬 수 있으므로(\S\ref{sec:ch3_ident}),
병목이 forward/backward/net 중 어디에 작용하는지 본문에서 반드시 명시한다(비대칭 $\Rightarrow\xi_{\sstat}\ne\xi_{\eq}$, Ch5).
\textbf{기본형(detailed balance 보존)}에서는 같은 직렬-병목 합성을 mobility($r^++r^-$ 전체) 또는 net flux 에
\emph{대칭}으로 적용해 $r_j^+/r_j^-$ 비를 보존한다. 즉 기본형 병목은 mobility 수준의 대칭 직렬식
\begin{equation}
\frac{1}{k_{j,\obs}^{\mathrm{phys}}}=\frac{1}{k_j^{\mathrm{phys}}}+\frac{1}{k_{j,\lim}}
\qquad(\text{mobility 전체에 적용}\Rightarrow r_j^+/r_j^-\ \text{불변})
\label{eq:ch3_mobility_limiter}
\end{equation}
로 쓰며, 이것이 (L3) 의 $1/k_j^{\mathrm{phys}}=1/k_j^{\mathrm{act}}+1/k_{j,\lim}$ 와 동형이다. 반면
식~\eqref{eq:ch3_forward_limiter} 의 forward-only 형은 강구동 branch(Chapter 5)에서 detailed balance 이탈을
\emph{의도적으로} 도입하는 template 이며, 기본 전개의 기본값이 아니다.
\begin{practicebox}
강구동서 explicit rate 가 발산적으로 커지면 초기 fitting/solver debugging 에서 임시 cap 을 \emph{계산 보조}로
쓸 수 있으나, 이는 \emph{본문 물리식이 아니며} 최종 해석에서는 식~\eqref{eq:ch3_forward_limiter} 의 transport/site/
diffusion/finite-current 병목으로 대체한다(GS-4; solver 구현은 Ch1 부록 B).
\end{practicebox}

% ====================================================================
\subsection{Butler--Volmer 계층과 $R_n\ne R_\ct\ne R_{n,\eff}$}\label{sec:ch3_bv}
% ====================================================================
\textbf{물리.} 계면 charge-transfer 가 지배적이면, 음극 반응 전류밀도는 forward/backward 의 \emph{차}를
charge-transfer 과전압 $\eta_n$ 으로 쓴 Butler--Volmer(BV) 형이다. 이는 \S\ref{sec:ch3_levelB} 의 transition-level
$r_j^\pm$ 를 \emph{계면 charge-transfer} 수준에서 다시 쓴 것이며, BV 의 cathodic $\alpha_n$ 은 Level B forward
$\beta_j$ 와 \emph{상보} 관계 $\alpha_n=1-\beta_j$(anodic 측 대응; 동일 \emph{값} 아닌 상보 \emph{관계})에 있다
(\cite{bardfaulkner,newman}, \AL{31}):
\begin{equation}
\boxed{\;j_n=j_{0,n}(T,\theta)\Big[\exp\!\Big(\frac{(1-\alpha_n)F\eta_n}{RT}\Big)-\exp\!\Big(-\frac{\alpha_n F\eta_n}{RT}\Big)\Big]\;.}
\label{eq:ch3_bv}
\end{equation}
$\eta_n>0$ 이 방전 방향 anodic flux 를 키우는 convention 이며, $j_{0,n}$=exchange current density[A/m$^2$].
\textbf{BV 가 forward$-$backward 차임의 무비약 확인.} 식~\eqref{eq:ch3_rplus},\allowbreak\eqref{eq:ch3_rminus} 형에서
net rate $\propto r_j^+-r_j^-$ 를 계면 변수로 옮기면, forward 지수는 $+(1-\alpha_n)F\eta_n/RT$(anodic),
backward 지수는 $-\alpha_n F\eta_n/RT$(cathodic)로 갈리고, 평형($\eta_n=0$)에서 두 지수가 $1$ 이라 $j_n=0$
(net flux 소멸 = exchange 평형) — 부호$\cdot$극한 정합. 두 지수의 transfer coefficient 합 $(1-\alpha_n)+\alpha_n=1$
은 Level B 의 $\beta_j+(1-\beta_j)=1$(식~\eqref{eq:ch3_ratio_raw})과 같은 구조다.

\textbf{small-signal 극한과 $R_\ct$(무비약 Taylor).} 작은 과전압 $|F\eta_n/RT|\ll1$ 에서 각 지수를 1차 Taylor
전개한다: $e^{x}\simeq1+x$. 식~\eqref{eq:ch3_bv} 에 넣으면
\begin{equation}
j_n\simeq j_{0,n}\Big[\big(1+\tfrac{(1-\alpha_n)F\eta_n}{RT}\big)-\big(1-\tfrac{\alpha_n F\eta_n}{RT}\big)\Big]
=j_{0,n}\,\frac{[(1-\alpha_n)+\alpha_n]F\eta_n}{RT}=j_{0,n}\,\frac{F\eta_n}{RT},
\label{eq:ch3_bv_linear}
\end{equation}
(상수 $1$ 이 상쇄, transfer coefficient 합 $=1$). 즉 small-signal 에서 $j_n$ 은 $\eta_n$ 에 \emph{선형}이고
$\alpha_n$ 무관이다(detailed balance 의 ratio 가 $\beta_j$ 무관인 것과 같은 상쇄). 전류 $I_n=A_\eff j_n$ 로
환산하면 $\partial I_n/\partial\eta_n|_0=A_\eff\,\partial j_n/\partial\eta_n=A_\eff\,j_{0,n}F/RT$(식~\eqref{eq:ch3_bv_linear}
의 $\eta_n$-미분)이고, 그 \emph{역수}가 charge-transfer resistance($\eta_n/I_n$ 의 small-signal 값)다:
\begin{equation}
\boxed{\;R_{\ct,n}=\frac{\partial\eta_n}{\partial I_n}\Big|_{\eta_n\to0}=\frac{RT}{F\,A_\eff\,j_{0,n}}\;.}
\label{eq:ch3_rct}
\end{equation}
\textbf{차원 검산.} $RT[\mathrm{J/mol}]/(F[\mathrm{C/mol}]\,A_\eff[\mathrm{m^2}]\,j_{0,n}[\mathrm{A/m^2}])
=\mathrm{J}/(\mathrm{C}\cdot\mathrm{A})$ 이고, $\mathrm{A}=\mathrm{C/s}$ 이므로 $\mathrm{C\cdot A}=\mathrm{C^2/s}$,
따라서 $\mathrm J/(\mathrm{C^2/s})=\mathrm{J\,s/C^2}=(\mathrm{J/C})/(\mathrm{C/s})=\mathrm{V/A}=\Omega$
— charge-transfer 저항의 올바른 단위.
\begin{keybox}
\textbf{세 저항의 계층 $R_n\ne R_\ct\ne R_{n,\eff}$ (\AL{31}).} 세 저항은 \emph{서로 다른 물리}를 잰다:
\begin{equation}
\boxed{\;R_n\ne R_{\ct,n},\qquad R_n\ne R_{n,\eff},\qquad R_{\ct,n}\ne R_{n,\eff}\;.}
\label{eq:ch3_Rn_not_equal}
\end{equation}
(i) $R_n$(Chapter 1) = apparent \emph{voltage-axis shift} 계수 $V_{n,\app}-V_n=s_I|I|R_n$ — 관측 전위 이동.
(ii) $R_{\ct,n}$(본 절, 식~\eqref{eq:ch3_rct}) = 계면 \emph{charge-transfer} small-signal 저항 — exchange
current $j_{0,n}$ 의 역수 scale. (iii) $R_{n,\eff}$(Chapter 2 의 $q_\irr=I^2R_{n,\eff}$ 형) = \emph{heat-generation}
effective 저항 $q_\irr=I^2R_{n,\eff}$ — 실제 dissipated heat. 이 셋을 한 데이터로 동시에 자유 피팅하면 apparent
shift$\cdot$charge-transfer$\cdot$dissipated heat 가 혼동된다(\S\ref{sec:ch3_Rn} 식별성).
\end{keybox}
\begin{boundbox}
\textbf{BV $\leftrightarrow$ Level B 대응(층위 구분; \AL{32}).} 식~\eqref{eq:ch3_bv} 의 cathodic $\alpha_n$ 은
식~\eqref{eq:ch3_rplus} 의 forward $\beta_j$ 와 \emph{상보} 관계 $\alpha_n=1-\beta_j$(동일 \emph{값}이 아닌 anodic 측
상보 \emph{관계})이며, BV 는 \emph{계면 charge-transfer}
과전압 $\eta_n$(식~\eqref{eq:ch3_eta_standard})을 쓰고 Level B 는 \emph{transition} 구동력 $\mathcal A_j$
(식~\eqref{eq:ch3_Aj})을 쓴다. $\eta_{n,\drive}\ne\eta_n$(식~\eqref{eq:ch3_drive})이므로 $\mathcal A_j/(n_j^\eff F)$
와 $\eta_n$ 을 임의로 동일시하지 않는다 — 둘의 연결은 reference electrode convention 과 $V_n,\phi_s,\phi_e,U_n$
대응이 필요하다(BOUNDED).
\end{boundbox}

% ====================================================================
\subsection{Generalized affinity 와 double counting 금지}\label{sec:ch3_genaff}
% ====================================================================
\textbf{물리.} 비이상 열역학$\cdot$staging free energy$\cdot$surface contribution 을 포함하려면 단순 $F\eta_n$
대신 \emph{generalized affinity} 를 쓴다. 단, 같은 자유에너지 항을 두 곳에 넣는 double counting 을 막는 것이
핵심이다(\AL{36}):
\begin{equation}
\mathcal A_n=F\eta_n+\mathcal A_{\resid}.
\label{eq:ch3_genaff}
\end{equation}
\textbf{차원.} $F\eta_n=\mathrm{(C/mol)\cdot V}=\mathrm{J/mol}$, $\mathcal A_{\resid}$ 도 $\mathrm{J/mol}$ — 정합.
$U_n$ 을 local thermodynamic equilibrium potential 로 정의하면 activity/chemical potential/staging free energy
\emph{일부가 이미 $U_n$ 안에} 있다(그것이 $\eta_n=\phi_s-\phi_e-U_n$ 의 $U_n$). \textbf{operational 정의}:
$U_n\equiv$ 측정 OCV at $I\to0$(reference 전극 고정). 따라서 $\mathcal A_{\resid}$ 는 \emph{이 reference 에서 제외된
항만}으로 — $U_n$ 에 \emph{포함되지 않은} residual nonideality, phase-boundary driving force, surface contribution
으로만 — 해석한다. 귀속 불변식: staging $\mu$ 가 이미 $U_n$ 에 포함되면 $\partial\mathcal A_{\resid}/\partial(\text{staging }\mu)=0$
(동일 항을 $\mathcal A_{\resid}$ 가 다시 세지 않음).
\begin{boundbox}
\textbf{double counting 금지(\AL{36}).} 같은 자유에너지 항을 $U_n$ 과 $\mathcal A_{\resid}$ 에 동시에 넣지
않는다. 예: $U_n$ 이 이미 graphite staging chemical potential 을 포함하면, 동일 staging free-energy 를
$\mathcal A_{\resid}$ 에 다시 더하는 것은 double counting 이다(Chapter 2 의 ``$\mathcal A_j/F=\eta_j+w_j\ln[\xi_j/(1-\xi_j)]$
분해에서 configurational 항을 가역 entropy 로''와 같은 bookkeeping 원칙). 어느 항이 어디 귀속되는지 명시한 뒤
사용한다(BOUNDED).
\end{boundbox}

% ====================================================================
\subsection{전류 보존과 transition current 분해}\label{sec:ch3_current}
% ====================================================================
\textbf{물리.} forward rate $r_j^+$ 가 아무리 커도 외부 전류가 한정되면, 각 transition 이 담당하는 전류의 합이
전하 보존식과 외부 전류를 \emph{동시에} 만족해야 한다. 이것이 강구동 saturation 의 \emph{물리적} 근거(임의 cap
아님)다. 전하 보존식 (L1) 을 시간 미분하면(Chapter 2 식과 동형, \AL{37})
\begin{equation}
Q_\cell\frac{\dd q}{\dd t}=\frac{\partial Q_\bg}{\partial V_n}\frac{\dd V_n}{\dd t}+\frac{\partial Q_\bg}{\partial T}\frac{\dd T}{\dd t}+\sum_j Q_{j,\tot}\frac{\dd\xi_j}{\dd t}.
\label{eq:ch3_current_balance}
\end{equation}
등온 방전 정전류 구간($\dd T/\dd t=0$, $Q_\cell\,\dd q/\dd t=|I|$)에서 좌변 $=|I|$ 이고, 우변을 background
progress 전류와 transition 전류로 묶으면
\begin{equation}
|I|=I_\bg^{\mathrm{prog}}+\sum_j I_j,\qquad I_\bg^{\mathrm{prog}}=\frac{\partial Q_\bg}{\partial V_n}\frac{\dd V_n}{\dd t},\qquad I_j=Q_{j,\tot}\frac{\dd\xi_j}{\dd t}.
\label{eq:ch3_current_decomp}
\end{equation}
\textbf{부호 보조정리($\sum_j I_j\le|I|$ 의 성립 조건).} 방전 정전류$\cdot$$s_{\phi,j}=+1$ 단일방향 가정 하
모든 transition 이 같은 방향으로 진행하므로 $\dd\xi_j/\dd t\ge0\Rightarrow I_j=Q_{j,\tot}\,\dd\xi_j/\dd t\ge0$ 이다.
또 $I_\bg^{\mathrm{prog}}=(\partial Q_\bg/\partial V_n)(\dd V_n/\dd t)$ 의 부호를 평가한 뒤(단조성 floor
$\partial Q_\bg/\partial V_n\ge\epsilon_Q>0$, (L1)), $I_\bg^{\mathrm{prog}}\ge0$ 인 영역에서
$\sum_j I_j=|I|-I_\bg^{\mathrm{prog}}\le|I|$ 가 성립한다(부등호의 유일 성립조건은 이 두 \emph{비음수성}이다).
\textbf{차원 검산.} $I_j=Q_{j,\tot}[\mathrm C]\cdot\dd\xi_j/\dd t[\mathrm{s^{-1}}]=\mathrm{C/s}=\mathrm A$;
$I_\bg^{\mathrm{prog}}=(\partial Q_\bg/\partial V_n)[\mathrm{C/V}]\cdot\dd V_n/\dd t[\mathrm{V/s}]=\mathrm{C/s}=\mathrm A$
— 모두 전류. 비등온서는 $(\partial Q_\bg/\partial T)(\dd T/\dd t)$ 가 추가되며 Chapter 2 처럼 coordinate shift
와 progress 를 구분한다(온도 drift 를 반응 전류로 오계상 금지).
\textbf{전류 보존 = 강구동 rate 제한의 물리적 근거.} 식~\eqref{eq:ch3_current_decomp} 의 $\sum_j I_j\le|I|$
제약이 식~\eqref{eq:ch3_forward_limiter} 의 $r_{j,\curr}^+$ 병목의 출처다 — $r_j^+$ 가 커도 transition current
합은 외부 전류를 넘을 수 없다(임의 cap 이 아니라 current conservation + coupled transport 에서 유도).

% ====================================================================
\subsection{Chapter 1 의 $R_n$ 의 전기화학적 분해}\label{sec:ch3_Rn}
% ====================================================================
\textbf{물리.} Chapter 1 의 apparent 전위 $V_{n,\app}=V_n+s_I|I|R_n(q,T,|I|)$ 의 $R_n$ 은 전기화학적으로 여러
성분의 \emph{apparent aggregate} 일 수 있다. 이는 해석용 분해이며 모든 항 동시 자유 피팅을 뜻하지 않는다(\AL{38}):
\begin{equation}
s_I|I|R_n\approx\eta_{\ct,n}+\eta_{\ohm,n}+\eta_{\conc,n}+\eta_{\film,n}+\eta_{\resid,n}.
\label{eq:ch3_Rn_decomp}
\end{equation}
\begin{boundbox}
\textbf{부호$\cdot$식별성(\AL{38}).} 각 $\eta_{\bullet,n}\ge0$ 으로 둔다(방전 convention 의 공통 부호 $s_I$ 가
좌변에 흡수되어 우변 항들은 비음수 분극 기여). 본 5항 분해는 $T\times|I|\times f$(주파수/시간척도) 다채널
신호(EIS/GITT/pulse)로 각 항을 \emph{분리}하는 것을 전제하며, 단일 $I$--$V$ 곡선만으로는 $\eta_{\ct,n}$ 만
식별 가능하고 나머지는 lumped residual 에 흡수된다(BOUNDED).
\end{boundbox}
\begin{longtable}{p{0.20\textwidth} >{\raggedright\arraybackslash}p{0.42\textwidth} >{\raggedright\arraybackslash}p{0.26\textwidth}}
\toprule 항 & 물리적 의미 & 분리 신호 \\ \midrule \endhead
$\eta_{\ct,n}$ & charge-transfer overpotential (식~\eqref{eq:ch3_eta_standard},\allowbreak\eqref{eq:ch3_rct}; $R_{\ct,n}$ 의 전압강하) & EIS 중주파 charge-transfer arc \\
$\eta_{\ohm,n}$ & electrolyte/solid ohmic drop & current-interrupt 순간 강하 \\
$\eta_{\conc,n}$ & concentration polarization / transport limitation & relaxation tail / 저주파 \\
$\eta_{\film,n}$ & SEI/surface film contribution & 고주파 RC arc \\
$\eta_{\resid,n}$ & reduced model 에서 아직 분리 못 한 residual & (lumped; 미분리) \\
\bottomrule
\end{longtable}
\textbf{Ch1 falsification 연결.} $\chi_j(=\beta_j$; 조건부 동일물, 대칭 Marcus$\cdot$정상 target 극한 한정$)$ barrier-lowering 기울기(Chapter 1 spectrum shift
$\partial\ln L/\partial V_{n,\drive}=-\chi_j s_{\phi,j}F/RT$)와 $\eta_{\ct,n}$ 의 Tafel 기울기는 단일 전극에서
\emph{co-linear} 일 수 있으므로, 다중 $T\times|I|$ + pulse/GITT 로 $\eta_{\ct}$(따라서 $R_{\ct,n}$)를 독립
분리한 \emph{후}에만 $\chi_j$ 귀속이 성립한다(Chapter 1 의 ``$R_n$ 과 $\chi_j$ 동시 자유 피팅 금지''와 같은
식별성 원칙). 이 분해의 $R_n$ 은 식~\eqref{eq:ch3_Rn_not_equal} 의 $R_{\ct,n},R_{n,\eff}$ 와 여전히 구분된다.

% ====================================================================
\subsection{C-rate 효과의 전기화학적 재분해}\label{sec:ch3_crate}
% ====================================================================
\textbf{물리.} Chapter 1 의 정전류 좌표식을 본 장 mobility $k_j^{\mathrm{phys}}=r_j^++r_j^-$ 로 다시 쓰면
\begin{equation}
\frac{\dd\xi_j}{\dd q}=\frac{Q_\cell}{|I|}\,k_j^{\mathrm{phys}}\,[\xi_{\eq,j}-\xi_j],
\label{eq:ch3_crate}
\end{equation}
(Chapter 1 식~$\dd\xi_j/\dd q=(Q_\cell k_j/|I|)(\xi_{\eq,j}-\xi_j)$ 와 동일; $k_j=k_j^{\mathrm{phys}}=r_j^++r_j^-$).
고율(C-rate) 효과는 단순 $k_j/|I|$ 경쟁이 아니라 다음 다섯 요소의 결합이다(\AL{39}):
\begin{longtable}{p{0.22\textwidth} >{\raggedright\arraybackslash}p{0.26\textwidth} >{\raggedright\arraybackslash}p{0.40\textwidth}}
\toprule 요소 & 수식 위치 ($\mathcal F$ 인자, 식~\eqref{eq:ch3_tail_function}) & 의미 \\ \midrule \endhead
체류 시간 & $1/|I|$ ($=Q_\cell/|I|$ scale) & 고율일수록 같은 $\dd q$ 에서 반응 시간 짧음 \\
평형 target lag & $\xi_{\eq,j}-\xi_j$ & state 가 target 못 따라가면 mismatch 증가 \\
활성화 mobility & $k_j^{\mathrm{act}}$ ($\mathcal A_j,w_j(T)$ 통해) & 전위 구동력+activation barrier 로 rate 증가 \\
물리적 병목 & $k_{j,\lim}$ & transport/site/diffusion/current 제한(식~\eqref{eq:ch3_forward_limiter}) \\
분극 & $R_n$ ($V_{n,\drive},\eta_n$ co-linear) & 구동력+apparent shift 변화 \\
\bottomrule
\end{longtable}
따라서 고율 tail/broadening 은 이 요소들의 \emph{의존성 schematic}(정량 모델이 아니라 어떤 변수에 의존하는지의
구조적 나열; BOUNDED)이다:
\begin{equation}
\mathrm{tail/broadening}\sim\mathcal F\!\Big(\tfrac{1}{|I|},\,\mathcal A_j,\,k_j^{\mathrm{act}},\,k_{j,\lim},\,R_n,\,
w_j(T),\,(\xi_{\eq,j}-\xi_j),\,\rho_j(g;T)\Big).
\label{eq:ch3_tail_function}
\end{equation}
여기서 $\mathcal A_j,R_n,V_{n,\drive}$ 는 $V_{n,\drive}$ 를 통해 서로 종속(co-linear)이므로 독립 인자가 아니다
(\S\ref{sec:ch3_ident} 식별성; 단일 신호로 동시 분리 불가). 또 tail 길이 $L$ 은 체류시간 $\propto1/|I|$ 를 통해
$|I|$ 와 부호 정합한다($\partial L/\partial|I|>0$: 고율일수록 lag 누적으로 tail 길어짐).
강구동서 $k_j^{\mathrm{act}}$ 가 증가해도 외부 전류 + transport limitation 이 전체 관측 tail 을 제한할 수 있다
(Chapter 1 spectrum shift 와 정합; $r_{j,\curr}^+$ 병목, \S\ref{sec:ch3_current}). \textbf{귀속 순위.} 따라서
1차 C-rate 효과는 체류시간 $1/|I|$ 이고, 강구동 율속은 병목 $k_{j,\lim}/r_{j,\curr}^+$ 이지 활성화 $k_j^{\mathrm{act}}$
가 아니다($k_j^{\mathrm{act}}$ 는 강구동서 오히려 증가; 관측 tail 제한의 주역은 병목).

% ====================================================================
\subsection{온도 의존 kinetics 와 식별성}\label{sec:ch3_ident}
% ====================================================================
\textbf{물리.} prefactor$\cdot$barrier$\cdot$exchange current$\cdot$transport 가 모두 온도 의존을 가질 수 있다.
대표적으로 Arrhenius 형(\AL{39}):
\begin{equation}
\nu_j(T)=\nu_{j,\mathrm{ref}}\exp\!\Big[-\frac{E_{\nu,j}}{R}\Big(\frac1T-\frac1{T_{\mathrm{ref}}}\Big)\Big],\quad
j_{0,n}(T,\theta)=j_{0,n}^{\mathrm{ref}}(\theta)\exp\!\Big[-\frac{E_{a,j_0}}{R}\Big(\frac1T-\frac1{T_{\mathrm{ref}}}\Big)\Big].
\label{eq:ch3_T}
\end{equation}
\textbf{차원 검산.} 지수 $E/R\cdot(1/T)=\mathrm{(J/mol)}/\mathrm{(J/(mol\,K))}\cdot\mathrm{K^{-1}}$ 무차원 — 정합.

\textbf{활성화 $\Delta H^\ddagger\cdot\Delta S^\ddagger$ 의 다온도 분리(Eyring 1차).} lumped Arrhenius $E_{\nu,j}$ 는
prefactor 의 $\propto T$ 인자($\nu=(k_BT/h)\kappa$, 식~\eqref{eq:ch3_rplus} 근원)를 은닉한다. Eyring rate
$r=a\,(k_BT/h)\,e^{-\Delta G^\ddagger/RT}$($\Delta G^\ddagger=\Delta H^\ddagger-T\Delta S^\ddagger$)에서 $T$
선형 인자를 좌변으로 옮겨 1차 분해하면
\begin{equation}
\boxed{\;\ln\frac{\nu_j^\pm}{T}=\ln\frac{k_B}{h}+\frac{\Delta S_j^{\pm\ddagger}}{R}-\frac{\Delta H_j^{\pm\ddagger}}{RT}\;.}
\label{eq:ch3_eyring}
\end{equation}
다온도에서 $\ln(r/T)$ 대 $1/T$ 직선의 \emph{기울기} $=-\Delta H^\ddagger/R$, \emph{절편}
$=\Delta S^\ddagger/R+\ln(k_B/h)+\ln a$ 이므로, 기울기에서 $\Delta H^\ddagger$, 절편에서
$\Delta S^\ddagger+$prefactor 를 분리한다(순수 Arrhenius $\ln r$ vs $1/T$ 는 $T$ 인자 은닉으로 둘을 혼합).
lumped 관계는 $E_{\nu,j}\approx\Delta H^\ddagger+RT$(Eyring$\to$Arrhenius 1차 보정)다.

\textbf{공선형 쌍(다온도 fit 으로 풀리지 않는 것).} (i) $\{\nu_j^\pm,a_j^\pm\}$ 는 $\ln(\nu a)$ 로만 진입하므로
\emph{곱}만 식별된다(개별 분리 불가). (ii) $\{E_{\nu,j},\Delta H_{a,j}^\ddagger\}$ 는 동일 $1/T$ 기울기로
들어가 \emph{합}만 식별된다. (iii) $\beta(=\alpha)$ 는 \emph{온도} fit 으로 풀리지 않으며, 분리하려면 \emph{등온}
Tafel 스윕이 필요하다(forward/backward 비대칭은 $\eta$-기울기로만 드러남). Arrhenius fit 절차 자체는 Ch1 부록 B
참조. 그러나 다음 항들은 강하게 상관된다(물리적 식별성 문제):
\begin{equation}
j_{0,n}(T,\theta)\leftrightarrow\nu_j(T)\leftrightarrow\Delta G_{a,j}(T)\leftrightarrow\chi_j(=\beta_j)\leftrightarrow k_{j,\lim}\leftrightarrow R_n\leftrightarrow\rho_j(g;T).
\label{eq:ch3_ident_chain}
\end{equation}
같은 ICA/rate 자료만으로 이 모두를 독립 결정할 수 없다 $\Rightarrow$ 독립 실험$\cdot$사전 물리 제약$\cdot$단계적
검증이 필요하다(Chapter 1$\cdot$2 의 forward-only$\cdot$식별성 원칙 계승). 위 사슬의 $\chi_j(=\beta_j)$ 는
조건부 동일물(대칭 Marcus$\cdot$정상 target 극한 한정)임에 유의한다.
\begin{boundbox}
\textbf{병목 도입 층위 명시(\AL{35}; 비대칭$\Rightarrow$detailed balance 충돌).} Coarse-grained(Level A)에서
$k_{j,\lim}$ 은 $\xi_j\to\xi_{\eq,j}$ relaxation mobility 의 병목이다(대칭, ratio 불변; 식~\eqref{eq:ch3_mobility_limiter}).
explicit(Level B)에서는
제한이 $r_j^+$, $r_j^-$, 또는 net flux 에 작용할 수 있다 — transport/finite-current 가 forward/backward 에
\emph{비대칭}으로 들어가면 $r_j^+/r_j^-$ 가 바뀌어 detailed balance(식~\eqref{eq:ch3_db})와 충돌하고
$\xi_{\sstat,j}\ne\xi_{\eq,j}$ 가 된다(branch, Chapter 5). 따라서 제한항이 mobility 전체/forward flux/net current
중 \emph{어디}에 작용하는지 반드시 명시한다(BOUNDED).
\end{boundbox}

% ====================================================================
\subsection{Chapter 4 $\cdot$ Chapter 5 로 전달되는 항목}\label{sec:ch3_transmit}
% ====================================================================
본 장이 확정한 미시 kinetics 변수와 그 후속 역할을 정리한다.
\begin{longtable}{p{0.30\textwidth} >{\raggedright\arraybackslash}p{0.58\textwidth}}
\toprule 항 & 후속 역할 \\ \midrule \endhead
$\eta_n,\mathcal A_n,\mathcal A_j,\eta_j$ & Ch4 비가역 열/entropy production driving force(dissipation 은 $\eta_j$,\brk Ch2 정합) \\
$r_j^+,r_j^-$ & Ch4 transition-level dissipation + relaxation entropy production;\brk Ch5 branch 비대칭 \\
$I_j=Q_{j,\tot}\,\dd\xi_j/\dd t$ & 상변이 heat-rate / transition current(Ch4 일반형 $\sum_j n_j^\eff I_j\eta_j$) \\
$n_j^{\eff}=RT/(Fw_j)$ & Ch4 비가역열 일반형 가중(CHARTER step3 일반형);\brk Ch2 특수형($n^\eff=1$)의 모장 \\
$R_{\ct,n},R_{n,\eff}$ & Ch4 small-signal heat / dissipative loss 분해 \\
$k_j^{\mathrm{act}},k_{j,\lim},\beta_j$ & Ch4 rate acceleration vs 병목;\brk Ch5 charge/discharge 비대칭 \\
$\xi_{\sstat,j}$ vs $\xi_{\eq,j}$ & Ch5 branch-dependent target / hysteresis state($C_j$ 의 detailed-balance 이탈) \\
\bottomrule
\end{longtable}
Chapter 5 로 넘기는 항목: 충전 branch 부호 convention($s_{\phi,j}^b$), branch-dependent target/hysteresis
state, charge/discharge asymmetric kinetics, path-dependent $U_j^{\mathrm{ch}},U_j^{\mathrm{dis}}$, full-cell
voltage decomposition, charge/discharge heat asymmetry. \textbf{본 장이 닫지 \emph{않는} 것}: 발열식의 최종
부호와 비가역열 일반형의 양정치 보존($\sum_j n_j^\eff I_j\eta_j\ge0$; Ch4), 충전 branch 부호 유도(Ch5).

% ====================================================================
\subsection{Chapter 3 의 가정 원장 (AL-30$\sim$39)}\label{sec:ch3_al}
% ====================================================================
본 장에서 사용한 가정을 통합 ledger(RB\_AL\_MASTER) Ch3 그룹(AL-30$\sim$39)으로 정리한다. AL-30,31 은 기등재분
(Foundation), AL-32$\sim$39 는 본 장 S2 에서 채운 신규 정의다(Phase A 적발 dangling AL-K1$\sim$K3 의 실제 정의
부여 포함: K1$\to$AL-30, K2$\to$AL-31, K3$\to$AL-31/32). tier = GROUNDED/BOUNDED/FLAGGED.
{\footnotesize
\setlength{\tabcolsep}{3pt}
\begin{longtable}{>{\raggedright\arraybackslash}p{0.045\textwidth} >{\raggedright\arraybackslash}p{0.335\textwidth} >{\raggedright\arraybackslash}p{0.135\textwidth} >{\raggedright\arraybackslash}p{0.205\textwidth} >{\raggedright\arraybackslash}p{0.155\textwidth}}
\toprule AL\# & 가정 진술 & tier & 근거 cite (DOI/ISBN) & 사용 식 \\ \midrule \endhead
AL-30 & forward/backward kinetics(식~\eqref{eq:ch3_fb}) + detailed balance(식~\eqref{eq:ch3_db_width});\brk consistent split(식~\eqref{eq:ch3_split}) → Ch1 ODE 복원 & GROUNDED & bard\brk faulkner,\brk newman,\brk degrootmazur\brk 1962\,(book),\brk onsager1931 (ISBN 978-0-\brk 471-\brk 04372-0;\brk\ 978-0-\brk 471-\brk 47756-3;\brk\ onsager1931 DOI 10.1103/\brk PhysRev.37.405) & \eqref{eq:ch3_fb},\allowbreak\eqref{eq:ch3_db},\allowbreak\eqref{eq:ch3_db_width},\allowbreak\eqref{eq:ch3_split},\allowbreak\eqref{eq:ch3_recover} \\
AL-31 & Butler--Volmer(식~\eqref{eq:ch3_bv});\brk small-signal $R_\ct{=}RT/(FA_\eff j_0)$; 계층 $R_n{\ne}R_\ct{\ne}R_{n,\eff}$;\brk Level B 비대칭 장벽($\beta_j{\equiv}\chi_j$, 식~\eqref{eq:ch3_rplus}--\eqref{eq:ch3_rminus}) & GROUNDED & bard\brk faulkner,\brk bazant2013,\brk evanspolanyi\brk 1938,\brk eyring1935 (ISBN 978-0-\brk 471-\brk 04372-0;\brk\ DOI 10.1021/\brk ar300145c;\brk\ 10.1039/\brk TF9383400011;\brk\ 10.1063/\brk 1.1749604) & \eqref{eq:ch3_Aj},\allowbreak\eqref{eq:ch3_rplus},\allowbreak\eqref{eq:ch3_rminus},\allowbreak\eqref{eq:ch3_ratio_B},\allowbreak\eqref{eq:ch3_ksum},\allowbreak\eqref{eq:ch3_bv},\allowbreak\eqref{eq:ch3_rct},\allowbreak\eqref{eq:ch3_Rn_not_equal} \\
AL-32 & 계면 과전압 $\eta_n=\phi_s-\phi_e-U_n$ $\ne$ reduced $\eta_{n,\drive}=V_{n,\drive}-V_n$ ($\ne V_{n,\app}$); 임의 동일시 금지(reference-electrode convention 필요) & BOUNDED & bard\brk faulkner,\brk newman (ISBN 978-0-\brk 471-\brk 04372-0;\brk 978-0-\brk 471-\brk 47756-3) & \eqref{eq:ch3_eta_standard},\allowbreak\eqref{eq:ch3_drive} drive boundbox, \eqref{eq:ch3_bv} BV boundbox \\
AL-33 & forward/backward 장벽 선형 이동($\mp\beta_j\mathcal A_j$ 등)은 소구동력 1차; $|\mathcal A_j|\sim\lambda$ 서 Marcus 곡률로 깨짐(역전 영역) & BOUNDED & marcus1956,\brk bazant2013 (10.1063/\brk 1.1742723;\brk 10.1021/\brk ar300145c) & \eqref{eq:ch3_rplus},\allowbreak\eqref{eq:ch3_rminus} Marcus boundbox \\
AL-34 & detailed-balance 정합 제약 $n_j^{\eff}{=}RT/(Fw_j)$(식~\eqref{eq:ch3_neff}); $C_j$ 의 $V_n$-무의존 partition 명시;\brk Ch1 $\mathcal A_j$ 의 $F{\to}n^\eff F$ 재해석(식 형태 불변;\brk partition 가정은 BOUNDED) & GROUNDED/\brk BOUNDED & bazant2013,\brk newman (DOI 10.1021/\brk ar300145c;\brk\ ISBN 978-0-\brk 471-\brk 47756-3) & \eqref{eq:ch3_dbmatch},\allowbreak\eqref{eq:ch3_neff},\allowbreak\eqref{eq:ch3_xiss} 검증 \\
AL-35 & 강구동 forward rate 제한 = 직렬 병목(식~\eqref{eq:ch3_forward_limiter}, 수치 cap 아님); 비대칭 병목 → detailed balance 이탈($\xi_\sstat\ne\xi_\eq$; 직렬합성 GROUNDED, 비대칭 BOUNDED) & GROUNDED/\brk BOUNDED & newman,\brk degrootmazur\brk 1962\,(book),\brk onsager1931 (ISBN 978-0-\brk 471-\brk 47756-3;\brk onsager1931 DOI 10.1103/\brk PhysRev.37.405) & \eqref{eq:ch3_forward_limiter}, \S\ref{sec:ch3_ident} boundbox \\
AL-36 & generalized affinity $\mathcal A_n=F\eta_n+\mathcal A_{\resid}$; $U_n$ 에 포함된 자유에너지를 $\mathcal A_{\resid}$ 에 중복 계상 금지(double counting) & BOUNDED & bazant2013,\brk newman (10.1021/\brk ar300145c;\brk ISBN 978-0-\brk 471-\brk 47756-3) & \eqref{eq:ch3_genaff} double-counting boundbox \\
AL-37 & 전류 보존 $|I|=I_\bg^{\mathrm{prog}}+\sum_j I_j$(식~\eqref{eq:ch3_current_decomp}, 전하보존식 시간미분);\brk coordinate/progress 분리(Ch2 정합) & GROUNDED & doyle1993,\brk newman (DOI 10.1149/\brk 1.2221597;\brk ISBN 978-0-\brk 471-\brk 47756-3) & \eqref{eq:ch3_current_balance},\allowbreak\eqref{eq:ch3_current_decomp} \\
AL-38 & $R_n$ apparent aggregate 분해(식~\eqref{eq:ch3_Rn_decomp}, 해석용; 동시 자유 피팅 X); $\chi_j$ vs Tafel co-linearity & BOUNDED & bard\brk faulkner,\brk newman (ISBN 978-0-\brk 471-\brk 04372-0;\brk 978-0-\brk 471-\brk 47756-3) & \eqref{eq:ch3_Rn_decomp} \\
AL-39 & kinetics 온도 의존(Arrhenius $\nu_j,j_{0,n}$); $j_0\!\leftrightarrow\!\nu\!\leftrightarrow\!\Delta G_a\!\leftrightarrow\!\chi\!\leftrightarrow\!k_{j,\lim}\!\leftrightarrow\!R_n\!\leftrightarrow\!\rho_j$ 식별성(독립 실험 필요) & BOUNDED & bard\brk faulkner,\brk bazant2013,\brk fly2020 (ISBN 978-0-\brk 471-\brk 04372-0;\brk 10.1021/\brk ar300145c;\brk 10.1016/\brk j.est.2020.101329) & \eqref{eq:ch3_crate},\allowbreak\eqref{eq:ch3_tail_function},\allowbreak\eqref{eq:ch3_T},\allowbreak\eqref{eq:ch3_ident_chain} \\
\bottomrule
\end{longtable}
}
\textbf{AGP(Assumption Grounding Policy).} 모든 본문 가정식은 AL-\# 를 인용한다. 본 장에 FLAGGED 항목은 없다
— 모든 식이 GROUNDED(반응속도론$\cdot$BV$\cdot$전하보존 확립) 또는 BOUNDED(유효범위 병기: partition 선택,
Marcus 곡률, 비대칭 병목, double-counting, 식별성)다. 임의 clip/cap/softplus/step 정의식 0(GS-4; 강구동 제한은
식~\eqref{eq:ch3_forward_limiter} 물리 병목, 수치 cap 은 practicebox 로 격하). solver/numerical 구현 0(Ch1 부록 B 분리).

% ====================================================================
\subsection{Chapter 3 자체 검수표}\label{sec:ch3_selfcheck}
% ====================================================================
{\small
\begin{longtable}{>{\raggedright\arraybackslash}p{0.74\textwidth} >{\raggedright\arraybackslash}p{0.16\textwidth}}
\toprule 검수 항목 & 결과 \\ \midrule \endhead
Chapter 1--2(RB) 상태변수/기호를 수정 없이 입력으로 받음(P3-6) & 통과(\S\ref{sec:ch3_inherit}) \\
반응속도론이 전하 보존식을 대체하지 않음 & 통과(\S\ref{sec:ch3_current}) \\
forward/backward kinetics 를 중심 구조로 둠 & 통과(식~\eqref{eq:ch3_fb}) \\
Ch1 relaxation ODE = forward/backward 축약(consistent split 무비약 항등 유도) & 통과(식~\eqref{eq:ch3_recover}) \\
$\xi_{\eq,j}$ 를 전류 의존 target 으로 바꾸지 않음; $\xi_{\sstat,j}$ 별도 정의 & 통과(식~\eqref{eq:ch3_xiss}) \\
★ Level A = Level B 의 detailed-balance 극한($\xi_\sstat\to\xi_\eq$) 검증 + 한정어(CHARTER step4) & 통과(verifybox) \\
★ Level A=방향성 X(ratio 불변) / Level B=방향성 O($\beta_j$ 비대칭) 구분 & 통과(keybox \S\ref{sec:ch3_split}) \\
★ $\beta_j$ 가 ratio 의 $\mathcal A$-기울기서 상쇄,\brk sum(mobility) 비대칭만 결정(무비약) & 통과(식~\eqref{eq:ch3_ratio_B},\allowbreak\eqref{eq:ch3_ksum}) \\
★ $\chi_j\equiv\beta_j$ 조건부 동일물 명시(대칭 Marcus$\cdot$정상 target 극한 한정; CHARTER step4) & 통과(\S 서,\brk keybox,\brk groundingbox) \\
★ detailed-balance 정합 제약 $n_j^{\eff}=RT/(Fw_j)$ 도출(BV-폭 불일치 방지;\brk CHARTER step3 정의장) & 통과(식~\eqref{eq:ch3_neff}) \\
$n_j^{\eff}$ 도출의 $C_j$ 무의존 partition 가정 명시;\brk model-independent 제약 우선 & 통과(\AL{34} boundbox) \\
Ch1 의 $\mathcal A_j(n{=}1)$+effective $w_j$ reconcile($F\to n^\eff F=RT/w_j$, 식 형태 불변) & 통과(\AL{34} reconcile boundbox) \\
★ $\mathcal A_j=s_{\phi,j}n_j^\eff F(V_\drive-U_j)$ 명시형 기준(CHARTER step1) & 통과(식~\eqref{eq:ch3_Aj} boxed) \\
★ $V_{n,\drive}=V_n+\eta_{n,\drive}$ 정의 + 기본 $V_{n,\drive}=V_n$(CHARTER step2) & 통과(식~\eqref{eq:ch3_drive} boxed) \\
$\Delta G_{\eff}<0$ 을 강구동 accelerated regime 으로(cap X) & 통과(\S\ref{sec:ch3_limiter}) \\
강구동 제한을 transport/site/diffusion/current 병목(직렬); 수치 cap 은 practicebox 격하 & 통과(식~\eqref{eq:ch3_forward_limiter}) \\
BV forward exponential 증가 허용;\brk net=forward$-$backward;\brk small-signal $R_\ct$ Taylor 유도 & 통과(식~\eqref{eq:ch3_bv},\allowbreak\eqref{eq:ch3_bv_linear},\allowbreak\eqref{eq:ch3_rct}) \\
$\alpha_n=1-\beta_j$ 상보 관계(동일 값 아님), $\eta_n\ne\eta_{n,\drive}\ne V_{n,\app}$ 구분 & 통과(\S\ref{sec:ch3_bv} BV boundbox) \\
$R_n\ne R_\ct\ne R_{n,\eff}$ 계층(세 물리 구분) & 통과(식~\eqref{eq:ch3_Rn_not_equal} keybox) \\
$V_{n,\app},V_{n,\drive},\eta_{n,\drive},\eta_n,U_n$ 구분 & 통과(\S\ref{sec:ch3_potentials}) \\
generalized affinity chemical-potential double counting 금지 & 통과(\S\ref{sec:ch3_genaff}) \\
$\chi_j$ vs $\eta_\ct$ Tafel co-linearity 경고(독립 분리 후 귀속) & 통과(\S\ref{sec:ch3_Rn}) \\
전류 보존 = 강구동 rate 제한의 물리적 근거($r_{j,\curr}^+$) & 통과(식~\eqref{eq:ch3_current_decomp}) \\
비대칭 병목 $\Rightarrow$ detailed balance 충돌 가능 명시 & 통과(\S\ref{sec:ch3_ident}) \\
모든 식 차원$\cdot$부호$\cdot$극한 검산 명시 & 통과(\S\ref{sec:ch3_db},\ref{sec:ch3_levelB},\ref{sec:ch3_bv} 등) \\
``자명/clearly/쉽게/obviously'' 본문 0건(G-noleap) & 통과 \\
모든 가정 AL-30$\sim$39 인용;\brk FLAGGED 0;\brk BOUNDED 유효범위 병기 & 통과(\S\ref{sec:ch3_al}) \\
Ch4/Ch5 이월 항목 분리($n^\eff$ 일반형 → Ch4) & 통과(\S\ref{sec:ch3_transmit}) \\
\bottomrule
\end{longtable}
}

% ====================================================================
\subsection{Chapter 3 의 결론}\label{sec:ch3_concl}
% ====================================================================
첫째, Chapter 1 의 relaxation ODE 는 forward/backward transition kinetics 의 \emph{축약형}이다. 기본 split
$r_j^+=k_j^{\mathrm{phys}}\xi_{\eq,j}$, $r_j^-=k_j^{\mathrm{phys}}(1-\xi_{\eq,j})$ 는 true equilibrium target 을
유지하며 Chapter 1 식을 \emph{대수적 항등}으로 복원한다(식~\eqref{eq:ch3_recover}). 이때 $k_j^{\mathrm{phys}}$ 는
forward-only rate 가 아니라 mobility $=r_j^++r_j^-$ 다.

둘째, directional barrier lowering(Level B)에서 transfer coefficient $\beta_j(\equiv\chi_j$; 조건부 동일물, 대칭 Marcus$\cdot$정상 target 극한 한정$)$ 는 rate ratio 의
$\mathcal A_j$-기울기에서 \emph{상쇄}되고($\beta+(1-\beta)=1$) mobility(sum)의 $\mathcal A_j$-응답 \emph{비대칭}만
결정한다(식~\eqref{eq:ch3_ratio_B},\allowbreak\eqref{eq:ch3_ksum}). \textbf{Level A=방향성 없는 대칭 가속(ratio 불변),
Level B=방향성 있는 비대칭} 이라는 구분이 명확히 닫힌다.

셋째, Level B 의 ratio 가 Chapter 1 평형(detailed balance, 식~\eqref{eq:ch3_db_width})과 정합하려면
$n_j^{\eff}=RT/(Fw_j)$ 가 강제된다($n_j^{\eff}w_j=RT/F$, 식~\eqref{eq:ch3_neff}). 이 정합 하에서 비평형 정상 target
$\xi_{\sstat,j}$ 가 평형 target $\xi_{\eq,j}$ 로 환원된다 — \textbf{Level A 는 Level B 의 detailed-balance 극한}이며
일치는 \emph{정상 target} $\xi_{\sstat,j}\to\xi_{\eq,j}$ \emph{한정}이다(★ 무조건 ``$A=B$'' 아님; CHARTER step4,
verifybox). 모든 prefactor/barrier/transfer coefficient 를 독립 자유 피팅하면 thermodynamic consistency 가 깨진다.

넷째, 강구동에서 forward rate 증가는 물리적으로 허용하되 softplus/$\max$ cap 으로 절단하지 않는다(GS-4). 제한은
transport/site/diffusion/finite-current/current-conservation 병목(식~\eqref{eq:ch3_forward_limiter},\allowbreak\eqref{eq:ch3_current_decomp})
으로 표현하며, 비대칭 병목은 detailed balance 를 이탈시켜 $\xi_{\sstat}\ne\xi_{\eq}$(branch, Ch5)를 만든다.

다섯째, Butler--Volmer 의 small-signal 한계가 charge-transfer 저항 $R_{\ct,n}=RT/(FA_\eff j_{0,n})$ 을 주며
(식~\eqref{eq:ch3_rct}), 이는 Chapter 1 의 apparent voltage-axis 계수 $R_n$ 과도 Chapter 2 의 heat-generation
$R_{n,\eff}$ 와도 다르다 — \textbf{$R_n\ne R_\ct\ne R_{n,\eff}$}(식~\eqref{eq:ch3_Rn_not_equal}). 이 구분은
Chapter 4 발열식과 Chapter 5 hysteresis 해석에서도 유지된다.

여섯째, 본 장의 $r_j^\pm,\eta_n,\eta_j,\mathcal A_j,I_j,n_j^{\eff},k_j^{\mathrm{act}},k_{j,\lim},\beta_j,\xi_{\sstat,j}$ 는
Chapter 4 의 entropy production/발열 일반형($\sum_j n_j^{\eff}I_j\eta_j$)과 Chapter 5 의 branch/hysteresis 로
전달된다. 이로써 Chapter 1 의 ``열역학=방향, 동역학=속도'' 분업이 미시 forward/backward 수준에서 \emph{정량화}
되되 앞 장과 충돌하지 않는다.

\subsection*{Chapter 4 로의 전달}
본 장이 확정한 미시 kinetics($r_j^\pm$, detailed balance, $n_j^{\eff}=RT/(Fw_j)$, $\xi_{\sstat,j}$, $R_\ct$)와
상태변수 위에서: Chapter 4 는 forward/backward rate 비대칭으로 transition-level entropy production 과 비가역열
일반형 $\dot{\mathcal Q}_\irr=\sum_j n_j^{\eff}I_j\eta_j$ 의 양정치 보존을 닫고(Ch2 가 후보로 남긴 휴지
relaxation heat 의 가역/비가역 \emph{정량} 분해 포함), Chapter 5 는 $C_j$ 의 detailed-balance 이탈로 충방전
히스테리시스의 branch 의존 target 과 부호를 다룬다.



% ====================================================================
% ===== Chapter 4 body (원본 graphite_ica_ch4_rebuilt.tex, byte-verbatim; heading 1단계 demote) =====
% ====================================================================
\clearpage
\section{Chapter 4: entropy-production 발열 정량화 ($\sum_j n_j^{\eff}I_j\eta_j$)}\label{chap:ch4}

% ====================================================================
\subsection*{서 — 인과 사슬에서 ``발열 정량화''의 자리}
\addcontentsline{toc}{subsection}{서 — 발열 정량화의 자리}
% ====================================================================
Chapter 1 은 ICA 꼬리의 주역을 \emph{동역학}(진행률 $\xi_j$ 의 유한속도 완화)으로 지목했고, Chapter 2 는 그
위에서 발열을 \emph{가역}(평형 entropy, 전위 온도계수)과 \emph{비가역}(과전압 소산, flux$\times$force) 기원으로
\emph{연결}했다. Chapter 3 은 그 비가역 소산의 동역학 원천인 mobility $k_j$ 를 \emph{미시} forward/backward
transition rate $r_j^\pm$(Level B)로 \emph{확대}하고, detailed balance 정합으로 effective electron number
$n_j^{\eff}=RT/(Fw_j)$ 를 고정했다. Chapter 4 는 이 셋을 합쳐 \textbf{하나의 정량 명제}를 닫는다.
\begin{quote}\itshape
Chapter 3 의 forward/backward rate $r_j^\pm$ 로부터 각 transition 의 \textbf{entropy production} 을 비평형
열역학의 표준형(network/master-equation)으로 세우면, 그것은 거시 비가역 발열 \emph{일반형}
$\dot{\mathcal Q}_\irr=\sum_j n_j^{\eff}I_j\eta_j\ge0$ 으로 \textbf{정확히 환원}된다. Chapter 2 가 쓴
$\sum_j I_j\eta_j$ 는 이 일반형의 ideal-width($n_j^{\eff}=1$) 특수해다.
\end{quote}
이 명제를 \textbf{다섯 단계}로 푼다. 각 단계는 학부 수준(비평형 열역학 flux$\times$force, Faraday 법칙,
지수$\cdot$로그 대수, 2법칙 $\dot S_\irr\ge0$)으로 따라갈 수 있게 중간을 생략하지 않는다.
\begin{enumerate}[label=\textbf{(\arabic*)},leftmargin=2.2em]
\item \textbf{비가역 열의 출발점은 flux$\times$force(entropy production)다.} 임의 heat correction 이 아니라
  conjugate flux $J_\alpha$ 와 thermodynamic force $X_\alpha$ 의 곱 $T\dot S_\irr=\sum_\alpha J_\alpha X_\alpha\ge0$
  이 2법칙이 보장하는 기본형이다(\S\ref{sec:ch4_irr}, \AL{40}).
\item \textbf{각 transition 의 entropy production 은 network 형으로 쓴다.} Chapter 3 의 forward/backward
  flux $\mathcal J_j^\pm$ 로 $\dot{\mathcal Q}_{\irr,j}^{\tr}=\tfrac{Q_{j,\tot}}{F}RT(\mathcal J_j^+-\mathcal J_j^-)
  \ln(\mathcal J_j^+/\mathcal J_j^-)\ge0$ 이 따라온다 — ``두 양수의 (차)$\times$(로그비)는 항상 음이 아니다''라는
  부등식이 2법칙을 직접 준다(\S\ref{sec:ch4_tr}, \AL{41},\AL{43}).
\item \textbf{미시 entropy production 이 거시 $n_j^{\eff}I_j\eta_j$ 로 환원된다.} detailed balance(Ch3)와
  $n_j^{\eff}=RT/(Fw_j)$(Ch3)를 대입하면 transition chemical affinity 가 $\mathcal A_j^{\chem}=n_j^{\eff}F\eta_j$
  로 정리되어 $\dot{\mathcal Q}_{\irr,j}^{\tr}=n_j^{\eff}I_j\eta_j$ 가 된다(\S\ref{sec:ch4_tr} verifybox,
  \AL{44}). 여기서 $\eta_j=V_{n,\drive}-V_{\eq,j}(\xi_j)$ 의 기준 전위를 \emph{단일}로 못박는다(CHARTER step2).
\item \textbf{일반형의 부호 양정치를 닫는다.} $\dot{\mathcal Q}_\irr=\sum_j n_j^{\eff}I_j\eta_j\ge0$ 이며
  $\sum_j I_j\eta_j$(Ch2)는 그 $n^{\eff}=1$ 특수해다(\S\ref{sec:ch4_general}, \AL{45}; CHARTER step3).
\item \textbf{가역 열$\cdot$transport$\cdot$휴지 relaxation 을 같은 원칙으로 정렬한다.} 가역 열은 평형 전위
  온도계수 또는 transition entropy basis(택일, 정합 제약), transport 는 flux$\times$force, 휴지 relaxation 은
  외부 $I=0$ 라도 내부 transition flux 의 dissipation 으로 남는다(\S\ref{sec:ch4_revOCV}--\ref{sec:ch4_rest}).
  Bernardi 가역열($q_\rev=s^\conv|I|T\,\partial U/\partial T$)이 미시 reaction entropy 의 합과 정합함을
  재확인한다.
\end{enumerate}

\noindent\textbf{한 문장 요지.} \emph{발열은 새 피팅항이 아니라, 이미 깔린 인과 사슬
($V_n\to\xi_{\eq,j}\to k_j/r_j^\pm\to\dd\xi_j/\dd t$) 위에서 계산되는 에너지 계층이며, 비가역 열은
transition-level entropy production 의 거시 표현 $\sum_j n_j^{\eff}I_j\eta_j\ge0$ 으로 닫힌다.} 본 장은
Chapter 1--3 과 같이 \textbf{방전(탈리튬화) 방향}($s_{\phi,j}=+1$)을 기본으로 한다. 충전 branch/hysteresis
heat 의 branch 의존 부호(Ch5)와 수치 해법(Ch1 부록 B)은 닫지 않는다.

\begin{groundingbox}
\textbf{지배 표준(GS, Ch1--3 계승) 및 발열 근거 원칙.} \textbf{(GS-1)} 모든 식은 물리적 의미를 prose 로 먼저
말한 뒤 수식화하고 중간 단계를 생략하지 않는다(학부 무비약). \textbf{(GS-2)} 결론은 출판 수준 엄밀성.
\textbf{(GS-3)} 모든 가정은 통합 Assumption Ledger(\AL{\#}: 논문/교과서 근거 + 검증 DOI)를 인용하며, 근거가
없으면 FLAGGED 로 표기하고 established 로 쓰지 않는다. \textbf{(GS-4)} 임의 clip/cap/softplus/step-function
정의식과 근거 없는 수치 임계값을 물리 모델로 쓰지 않는다(속도 제한은 물리적 병목으로, log 발산은 domain guard
로). \textbf{발열 근거 원칙.} 모든 발열 항은 (i) entropy production, (ii) conjugate flux$\times$force, (iii)
평형 전위 온도계수, (iv) transition entropy$\times$진행률, (v) transport/finite-current dissipative loss, (vi)
자유에너지 감소를 동반하는 internal relaxation 중 하나의 \emph{물리} 근거를 가져야 한다. 임의 heat correction /
capped heat / empirical gain / residual heat 는 기본식에 넣지 않는다(실무 팁 분리). 방전 기준 $s_{\phi,j}=+1$,
$|I|>0$, heat-generation-positive convention.
\end{groundingbox}

% ====================================================================
\subsection{Chapter 1--3 에서 전달받는 고정 기준식}\label{sec:ch4_inherit}
% ====================================================================
본 장은 Chapter 1$\cdot$2$\cdot$3(RB 재구성본)의 다음을 \textbf{수정 없이} 입력으로 받는다(프로젝트 검수 P3-6).
각 식 번호는 앞 장의 boxed 결과를 가리키며, 본 장은 그 위에 entropy-production 계층을 \emph{적층}한다(전하
보존식을 발열로 대체하지 않는다).
\begin{linkbox}
\textbf{(L1) 전하 보존식(중심 상태식, 발열로 대체 안 됨).} $Q_\cell\,q=Q_\bg(V_n,T)+\sum_{j=1}^{N_p}Q_{j,\tot}\xi_j$.
$V_n$ 은 그 해(외부 입력 아님); 평형 특수해가 $V_{n,\OCV}(q,T)$. 단조성 floor $\partial Q_\bg/\partial V_n\ge\epsilon_Q>0$
로 해 유일.\\
\textbf{(L2) kinetics(Ch3 미시 + Ch1 축약).} $\dfrac{\dd\xi_j}{\dd t}=r_j^+(1-\xi_j)-r_j^-\xi_j$, 축약형
$=k_j^{\mathrm{phys}}(\xi_{\eq,j}-\xi_j)$ ($k_j^{\mathrm{phys}}=r_j^++r_j^-$, mobility); 정전류 구간서만
$\dfrac{\dd\xi_j}{\dd t}=\dfrac{|I|}{Q_\cell}\dfrac{\dd\xi_j}{\dd q}$.\\
\textbf{(L3) Ch2 dissipation force(비가역 열 공액 force).} $\eta_j=V_{n,\drive}-V_{\eq,j}(\xi_j)$,
$V_{\eq,j}(\xi_j)=U_j(T)+w_j(T)\ln[\xi_j/(1-\xi_j)]$ — 현재 $\xi_j$ 의 \emph{국소 평형} 기준(중심 affinity
$\mathcal A_j$ 와 구분). \textbf{기준 전위는 $V_{n,\drive}$ 단일}(CHARTER step2; \S\ref{sec:ch4_eta_single}).\\
\textbf{(L4) Ch3 미시 정합.} detailed balance $r_j^+/r_j^-=\xi_{\eq,j}/(1-\xi_{\eq,j})$, 그 ln 형
$\ln[\xi_{\eq,j}/(1-\xi_{\eq,j})]=(V_n-U_j)/w_j$; effective electron number $n_j^{\eff}=RT/(Fw_j)$ ($\mathcal A_j
=s_{\phi,j}n_j^{\eff}F(V_{n,\drive}-U_j)$, $s_{\phi,j}=+1$ 방전 기본).\\
\textbf{(L5) 네 전위 구분.} 내부 $V_n$(보존식 해), apparent $V_{n,\app}=V_n+s_I|I|R_n$(분극 포함, 평형 전위
아님), 구동 $V_{n,\drive}=V_n+\eta_{n,\drive}$(기본 $\eta_{n,\drive}=0\Rightarrow V_{n,\drive}=V_n$), 평형
$V_{n,\OCV}$(평형 보존식 해).\\
\textbf{(L6) Ch2 가역열 경계.} $q_{\rev,j}=s_{\phi,j}\,I_j\,T\,\partial U_j/\partial T=s_{\phi,j}I_j\Pi_j$
(방전 $s_{\phi,j}=+1$ 서 $I_j\Pi_j$; $\Pi_j=T\partial U_j/\partial T$, Peltier-type); reduced $\dot{\mathcal Q}_{\rev,n}^{\OCV}=s_{\rev,n}^{\conv}|I|T(\partial V_{n,\OCV}/\partial T)_q$;
$\partial U_j/\partial T=\Delta S_{\rxn,j}/(s_{\phi,j}F)$ (reaction entropy, activation entropy $\Delta S_{a,j}$
와 별개).\\
\textbf{(L7) Ch2 비가역열 특수형($n^{\eff}=1$).} flux$\times$force $T\dot S_{\irr,n}=\sum_j J_{n,j}F\eta_j\ge0$,
축약 $q_\irr=|I|\eta\ge0$. \textbf{본 장이 일반화하는 대상} — Ch2 가 forward-ref 한 일반형 $\sum_j n_j^{\eff}I_j\eta_j$
를 본 장에서 닫는다(\S\ref{sec:ch4_general}).
\end{linkbox}
임의 heat correction / capped heat / empirical residual 을 기본 발열식으로 쓰지 않는다(GS-4); 강구동 제한과
log 발산은 물리적 병목과 numerical domain guard 로 둔다.

\textbf{Ch1--3 상속 기호(본 장 첫 출현).} 위 (L1)--(L7) 에 함께 쓰이는 상속 기호: $N_p$ $=$ 동시 진행 transition
(상변이) 개수(식 (L1) 합 상한), $s_I\in\{+1,-1\}$ $=$ apparent voltage-axis 전류 부호 convention(식 (L5)
$V_{n,\app}=V_n+s_I|I|R_n$, 피팅 X), $\rho_j(g;T)$ $=$ transition $j$ 의 quenched barrier 분포 밀도(Chapter 1
$g$-spectrum, $[\mathrm{mol/J}]$ 형). \textbf{신규 기호(Ch1--3 위에 추가).} 본 장에서만 새로 도입하는 기호를 모은다(앞 장 기호는 (L1)--(L7) 과 동일물).
\begin{longtable}{p{0.16\textwidth} p{0.13\textwidth} >{\raggedright\arraybackslash}p{0.63\textwidth}}
\toprule 기호 & 단위 & 의미 \\ \midrule \endhead
$\dot S_\irr$ & W/K & entropy production rate (비가역 entropy 생성률;\brk 2법칙 $\ge0$) \\
$J_\alpha,\ X_\alpha$ & (쌍) & generalized conjugate flux $\cdot$ thermodynamic force (곱 $J_\alpha X_\alpha$ 가 W) \\
$J_n,\ \mathcal A_n$ & mol/s,\brk J/mol & 음극 molar reaction flux $\cdot$ generalized reaction affinity \\
$\mathcal J_j^+,\ \mathcal J_j^-$ & 1/s & transition $j$ 의 forward $=r_j^+(1-\xi_j)$ $\cdot$ backward $=r_j^-\xi_j$ unidirectional flux \\
$\mathcal A_j^{\chem}$ & J/mol & transition $j$ 의 chemical affinity $=RT\ln(\mathcal J_j^+/\mathcal J_j^-)$ (network thermo) \\
$I_j$ & A & transition $j$ lumped current $=Q_{j,\tot}\,\dd\xi_j/\dd t$ (Ch3 동일물) \\
$\dot n_j$ & mol/s & transition $j$ molar rate $=(Q_{j,\tot}/F)\,\dd\xi_j/\dd t$ \\
$\dot{\mathcal Q}_{\irr,j}^{\tr}$ & W & transition $j$ 의 transition-level 비가역 발열률 (entropy production $\times T$) \\
$\dot{\mathcal Q}_{\src,n}$ & W & 음극 control volume 내부 총 열원 (가역 흡열 가능) \\
$\dot{\mathcal Q}_{\irr,n},\dot{\mathcal Q}_{\rev,n}$ & W & 계면 반응 비가역 발열 $\cdot$ 가역 발열 (음극 합) \\
$\dot{\mathcal Q}_{\rlx,n},\dot{\mathcal Q}_{\transp,n}$ & W & 내부 relaxation 발열 $\cdot$ transport/병목 dissipative loss \\
$\Delta S_j(\equiv\Delta S_{\rxn,j})$ & J/(mol K) & transition $j$ 의 reaction entropy (1 mol electron/Li-equiv 당) \\
$h_\bg$ & V/K & background reversible-heat potential coefficient $=(\partial V_{\bg,\eq}/\partial T)$ (전위형 V/K,\brk entropy J/(mol K) 아님;\brk Ch2 동일물) \\
$R_{\transp,n},R_{n,\eff}$ & $\Omega$ & transport effective resistance $\cdot$ heat-generation effective resistance \\
$J_\ell,\ X_\ell$ & (쌍) & transport flux $\cdot$ 그 공액 force (chemical/electrolyte-potential/concentration gradient) \\
$C_{\mathrm{th}}$ & J/K & 음극 control volume 유효 열용량 (Ch2 동일물) \\
$s^{\conv},\ s_\bg^\conv$ & --- & convention bookkeeping factor $\in\{+1,-1\}$ (피팅 X;\brk Ch2 동일물) \\
\bottomrule
\end{longtable}
$F=96485\ \mathrm{C/mol}$(Faraday), $R=8.314\ \mathrm{J/(mol\,K)}$(기체상수), $\mathrm{J/C}=\mathrm V$. 열률 dot
표기는 시간율 [W]. unidirectional flux $\mathcal J_j^\pm$ 와 net flux $\dd\xi_j/\dd t=\mathcal J_j^+-\mathcal J_j^-$
의 단위는 $\mathrm{s^{-1}}$($\xi_j$ 무차원).

\textbf{직접 입력은 시간 미분량.} 모든 발열률$\cdot$entropy production 의 직접 입력은 시간 미분량 $\dd\xi_j/\dd t$
또는 시간당 flux 다. transition 별 전하/molar 진행률은(Ch3 (L2) kinetics 의 $I_j$ 정의와 동일 — 본 장 재기입)
\begin{equation}
I_j=Q_{j,\tot}\,\frac{\dd\xi_j}{\dd t},\qquad
\dot n_j=\frac{Q_{j,\tot}}{F}\,\frac{\dd\xi_j}{\dd t}\ [\mathrm{mol\,s^{-1}}].
\label{eq:ch4_Ij}
\end{equation}
\textbf{차원 검산.} $I_j=Q_{j,\tot}[\mathrm C]\cdot\dd\xi_j/\dd t[\mathrm{s^{-1}}]=\mathrm{C/s}=\mathrm A$;
$\dot n_j=(Q_{j,\tot}/F)[\mathrm{C}/(\mathrm{C\,mol^{-1}})]\cdot[\mathrm{s^{-1}}]=\mathrm{mol/s}$ — 정합. 면적
정규화가 필요하면 유효 반응면적 $A_{\eff}$ 로 area-flux $J_j=\dot n_j/A_{\eff}$ 를 \emph{별도} 정의한다(lumped
molar rate $\dot n_j$ 와 area-flux $J_j$ 혼동 금지).

% ====================================================================
\subsection{$\eta_j$ 의 기준 전위 단일화 (CHARTER step2 — 이중정의 해소)}\label{sec:ch4_eta_single}
% ====================================================================
본 절은 \textbf{본 장의 가장 먼저 못박아야 할 규약}이다. Chapter 2 는 dissipation force(과전압)를
$\eta_j=V_{n,\drive}-V_{\eq,j}(\xi_j)$ 로(구동 전위 기준) 정의했고, Chapter 3 은 중심 기준 affinity 의
전위 단위 형 $\eta_j=s_{\phi,j}(V_{n,\drive}-U_j)$ 를 별도로 썼다. 본 장의 \emph{발열} 유도(특히 미시=거시
정합, \S\ref{sec:ch4_tr})에서는 ``비가역 열 공액 force''를 쓰므로, 그것을 \textbf{Chapter 2 의 dissipation
force 하나로 단일화}한다(\AL{42}; RB\_CHARTER step2).
\begin{keybox}
\textbf{본 장 비가역 열 공액 force 의 단일 정의(CHARTER step2).} 비가역 발열의 공액 force 는 \emph{현재}
$\xi_j$ 의 국소 평형으로부터의 이탈, 즉
\begin{equation}
\boxed{\;\eta_j\;\equiv\;V_{n,\drive}-V_{\eq,j}(\xi_j)\;,\qquad V_{\eq,j}(\xi_j)=U_j(T)+w_j(T)\ln\frac{\xi_j}{1-\xi_j}\;,}
\label{eq:ch4_eta_def}
\end{equation}
이며 \textbf{기준 전위는 $V_{n,\drive}$ 단일}이다($V_n$ 과 $V_{n,\drive}$ 혼용 0). 본 정의는 방전 부호 생략형
이며, 일반(충전 포함) 부호는 $\eta_j$ 가 아니라 $\mathcal A_j$ 의 scale 인자 $s_{\phi,j}$ 로 처리한다($\eta_j$ 자체는
부호 인자 없음; \S\ref{sec:ch4_tr} 식~\eqref{eq:ch4_affinity_eta}). 기본 전개에서는 RB\_CHARTER
step2 의 \textbf{기본형 $V_{n,\drive}=V_n$}($\eta_{n,\drive}=0$, 식 (L5))을 쓰며, 이 전제는 미시=거시 대입
(\S\ref{sec:ch4_tr})에서 \emph{명시적으로} 사용된다. 강구동($\eta_{n,\drive}\ne0$, 따라서 $V_{n,\drive}\ne V_n$)
일반화는 BOUNDED 로 표기하고 유효범위를 병기한다.
\end{keybox}
\textbf{왜 단일화가 필요한가(이중정의의 위험).} 만약 $\eta_j$ 를 어떤 식에서는 $V_n$ 기준
($V_n-V_{\eq,j}(\xi_j)$), 다른 식에서는 $V_{n,\drive}$ 기준($V_{n,\drive}-V_{\eq,j}(\xi_j)$)으로 섞어 쓰면,
$V_{n,\drive}\ne V_n$ 인 강구동에서 \emph{같은 기호 $\eta_j$ 가 서로 다른 값}을 가리켜 비가역 발열의 부호와
크기가 식마다 달라진다 — 이것이 Phase A 가 CRITICAL 로 적발한 ``$\eta$ 이중정의''다. 본 장은 \emph{정의}
\eqref{eq:ch4_eta_def} 를 $V_{n,\drive}$ 기준으로 동결하고, 기본형에서 $V_{n,\drive}=V_n$ 임을 사용 시점마다
명시해 모순을 차단한다(Chapter 2 \S 비가역 열 부호 정합 + Chapter 3 \S$V_{n,\drive}$ 정의장과 동일 입도).

\textbf{$V_{\eq,j}(\xi_j)$ 의 유도(logistic 역함수; Ch2 재확인, 무비약).} Chapter 1 logistic(식 (L4),
$s_{\phi,j}=+1$) $\xi_{\eq,j}=[1+e^{-(V_n-U_j)/w_j}]^{-1}$ 을 ``$\xi_j$ 가 현재값일 때 그것을 평형으로 만드는
전위''에 대해 역으로 푼다. $\xi_j=[1+e^{-(V_{\eq,j}-U_j)/w_j}]^{-1}$ 로 두고 정리하면 $1/\xi_j-1=
e^{-(V_{\eq,j}-U_j)/w_j}$, 즉 $(1-\xi_j)/\xi_j=e^{-(V_{\eq,j}-U_j)/w_j}$. 양변에 자연로그를 취하면
$\ln[(1-\xi_j)/\xi_j]=-(V_{\eq,j}-U_j)/w_j$, 부호를 뒤집어 $\ln[\xi_j/(1-\xi_j)]=(V_{\eq,j}-U_j)/w_j$, 따라서
$V_{\eq,j}(\xi_j)=U_j+w_j\ln[\xi_j/(1-\xi_j)]$ — 식~\eqref{eq:ch4_eta_def} 의 둘째 식이다(중간 단계 전부 명시;
Chapter 2 의 동일 유도를 본 장 기준 전위 단일화 위에서 재확인).
\begin{boundbox}
\textbf{rate-affinity $\mathcal A_j$ 와 heat-force $\eta_j$ 의 구분(\AL{42}; Ch2$\cdot$Ch3 계승).} 식 (L4) 의
$\mathcal A_j=s_{\phi,j}n_j^{\eff}F(V_{n,\drive}-U_j)$ 는 \emph{중심} $U_j$ 기준이라 평형에서도 0 이 아니다
(rate 구동용). 본 절 $\eta_j$(식~\eqref{eq:ch4_eta_def})는 \emph{현재} $\xi_j$ 의 국소 평형 $V_{\eq,j}(\xi_j)$
기준이라 평형($\xi_j=\xi_{\eq,j}$)에서 0 이 된다(dissipation 용). 둘의 차이는 식~\eqref{eq:ch4_eta_def} 를
$\mathcal A_j$ 의 전위 단위 형 $\mathcal A_j/(n_j^{\eff}F)=s_{\phi,j}(V_{n,\drive}-U_j)$ 와 비교하면(아래 분해
$\mathcal A_j/(n_j^{\eff}F)-\eta_j=V_{\eq,j}(\xi_j)-U_j$ 는 $V_{n,\drive}$ \emph{임의}에서 성립하는 \emph{항등식}
이며 기본형 $V_{n,\drive}=V_n$ 에 의존하지 않는다 — 진짜 기본형 의존은 미시=거시 닫힘인 \S\ref{sec:ch4_tr}
3단계에만 있다; $s_{\phi,j}=+1$ 표기만 방전 한정)
\[
\frac{\mathcal A_j}{n_j^{\eff}F}=(V_{n,\drive}-U_j)
=\underbrace{(V_{n,\drive}-V_{\eq,j}(\xi_j))}_{=\eta_j}
+\underbrace{(V_{\eq,j}(\xi_j)-U_j)}_{=w_j\ln[\xi_j/(1-\xi_j)]},
\]
즉 $\mathcal A_j/(n_j^{\eff}F)=\eta_j+w_j\ln[\xi_j/(1-\xi_j)]$ 이며, 둘째 항(configurational 가역 혼합 entropy)은
\emph{가역} entropy basis(\S\ref{sec:ch4_revxi})로 가고 \emph{비가역} dissipation 에는 $\eta_j$ 만 쓴다 — 이
구분을 혼동하면 가역 혼합 entropy 가 비가역 열로 오계상된다(BOUNDED — logistic 평형 모델 내에서 정확).
\end{boundbox}

% ====================================================================
\subsection{내부 열원 항의 물리적 분해 (Ch2 3-분해의 정련)}\label{sec:ch4_decomp}
% ====================================================================
\textbf{물리.} 음극 control volume 내부 열원을 \emph{물리 기원별}로 분해한다(피팅 편의 분해가 아니라
entropy-production / 평형 entropy / dissipative-loss 기원 분해; 근거는 항별로 귀속한다: control-volume 1법칙
에너지수지 \emph{형식}은 \cite{bernardi1985,rao1997}; 비가역 항의 network entropy-production 형
$\sum_j n_j^{\eff}I_j\eta_j$ 는 \cite{degrootmazur1962,schnakenberg1976}(\S\ref{sec:ch4_irr},\ref{sec:ch4_tr});
$\dot{\mathcal Q}_{\rlx,n}$ 의 자유에너지 relaxation 분해는 본 장 신규(Ch2 후보 확정); \AL{46}):
\begin{equation}
\dot{\mathcal Q}_{\src,n}=\dot{\mathcal Q}_{\irr,n}+\dot{\mathcal Q}_{\rev,n}+\dot{\mathcal Q}_{\rlx,n}+\dot{\mathcal Q}_{\transp,n}.
\label{eq:ch4_decomp}
\end{equation}
\begin{longtable}{p{0.28\textwidth} >{\raggedright\arraybackslash}p{0.60\textwidth}}
\toprule 항 & 물리적 의미 \\ \midrule \endhead
$\dot{\mathcal Q}_{\irr,n}$ & 계면 반응/transition 진행의 비가역 entropy production ($\sum_j n_j^{\eff}I_j\eta_j$, $\ge0$) \\
$\dot{\mathcal Q}_{\rev,n}$ & 평형 entropy coefficient 또는 transition entropy 가역 열 (열역학 기원, 부호가변) \\
$\dot{\mathcal Q}_{\rlx,n}$ & 외부 전류와 무관하게 내부 상태가 자유에너지를 낮추며 relax 할 때의 열 (확정엔 자유에너지 함수 필요) \\
$\dot{\mathcal Q}_{\transp,n}$ & solid diffusion/electrolyte/film/finite-current limitation 의 dissipative loss ($\ge0$) \\
\bottomrule
\end{longtable}
\textbf{Chapter 2 3-분해와의 관계(정련, 이중계상 금지; \AL{46}).} 본 4-분해는 Chapter 2 의 3-분해
$\dot{\mathcal Q}_{\src,n}=\dot{\mathcal Q}_{\irr,n}+\dot{\mathcal Q}_{\rev,n}+\dot{\mathcal Q}_{\rlx,n}^{\cand}$
를 \emph{정련}한 것이다 — Chapter 2 에서 $\dot{\mathcal Q}_{\irr,n}$ 에 묶여 있던 transport/diffusion
dissipation 을 본 장에서 $\dot{\mathcal Q}_{\transp,n}$ 으로 \emph{분리}했고, 후보로만 남았던 relaxation heat
$\dot{\mathcal Q}_{\rlx,n}^{\cand}$ 을 본 장에서 분해한다(\S\ref{sec:ch4_rest}). 여기서 \emph{비가역 floor
$\ge0$ 는 entropy production 으로 식으로 확정}되나, \emph{가역/비가역 정량 분리 비율}은 detailed-balance 정합
정도 $+$ 독립 calorimetry 전제 하에서만 닫힌다(\S\ref{sec:ch4_rest} flagbox; ``확정''은 floor 의 확정이지
분리 비율의 확정이 아니다). 또 control-volume 합과 transition 합이 같은 양임을 명시한다:
$\dot{\mathcal Q}_{\rlx,n}\equiv\dot{\mathcal Q}_{\xi,\rlx}$(\S\ref{sec:ch4_rest} 식~\eqref{eq:ch4_rest_split}).
즉 기호 관계는
\begin{equation}
\dot{\mathcal Q}_{\irr,n}^{\mathrm{Ch2}}\;\supseteq\;\dot{\mathcal Q}_{\irr,n}^{\mathrm{Ch4}}+\dot{\mathcal Q}_{\transp,n},
\label{eq:ch4_ch2refine}
\end{equation}
이며, 본 장 분할(계면 반응 EP $+$ transport dissipation)이 Ch2 의 $\dot{\mathcal Q}_{\irr,n}^{\mathrm{Ch2}}$ 를
\emph{완전 분할}하면(잔차 $=0$) 등호로 닫힌다. 두 장의 항을 함께 쓸 때 transport dissipation 을 \emph{이중
계상하지 않는다}(본 장 $\dot{\mathcal Q}_{\irr,n}$
은 계면 반응/transition entropy production 만 포함). 가역 열이 흡열일 수 있으므로 $\dot{\mathcal Q}_{\src,n}$
은 항상 양수 ``생성량''이 아니라 internal heat source term 이다. 외부 손실 포함 control-volume 1법칙 열수지는
(Chapter 2 식과 동형)
\begin{equation}
C_{\mathrm{th}}\frac{\dd T}{\dd t}=\dot{\mathcal Q}_{\src,n}-\dot{\mathcal Q}_{\loss},\qquad
\dot{\mathcal Q}_{\loss}=hA(T-T_\amb),
\label{eq:ch4_balance}
\end{equation}
(좌변 $\mathrm{(J/K)(K/s)=W}$ — 차원 정합). $\dot{\mathcal Q}_{\loss}$ 의 상세 열전달 모델(대류계수 $h$, 표면적
$A$, 주위 온도 $T_\amb$)은 본 장에서 닫지 않는다(Chapter 2 의 control-volume 형 계승).

% ====================================================================
\subsection{비가역 열의 출발점: flux--force 와 entropy production}\label{sec:ch4_irr}
% ====================================================================
\textbf{물리.} 비가역 열은 임의 heat correction 이 아니라 \emph{비평형 열역학의 entropy production} 에서 나온다.
가장 기본적 출발점은 conjugate flux $J_\alpha$ 와 그에 공액인 thermodynamic force $X_\alpha$ 의 곱이며, 2법칙이
그 합의 비음성을 보장한다(\cite{degrootmazur1962,prigogine1967}, \AL{40}; Onsager reciprocal relations 는 본
부등식의 근거가 아니라 transport $L_{\alpha\beta}$ 대칭으로 \S\ref{sec:ch4_transport}$\cdot$\AL{47} 에 재배치):
\begin{equation}
\boxed{\;T\dot S_{\irr}=\sum_\alpha J_\alpha X_\alpha\ \ge\ 0\;.}
\label{eq:ch4_fluxforce}
\end{equation}
\textbf{2법칙 의미(무비약).} $\dot S_\irr$ 는 비가역 과정이 생성하는 entropy 율이며, 고립계에서 entropy 가
감소할 수 없다는 2법칙이 곧 $\dot S_\irr\ge0$ 이다. 엄밀히 2법칙은 \emph{합} $\sum_\alpha J_\alpha X_\alpha\ge0$
을 보장하며(cross-coupling 이 있으면 개별 항 $J_\alpha X_\alpha$ 가 음수일 수 있다 — Onsager off-diagonal
$L_{\alpha\beta}$), 비음은 합 차원에서 정의된다. 본 장 계면 transition 항은 \emph{단일 공액 쌍}
($J_n\leftrightarrow\mathcal A_n$, $I_j\leftrightarrow\eta_j$)이라 cross-coupling 이 없으므로 \emph{항별}로도 $\ge0$
임을 \S\ref{sec:ch4_tr} 의 부호 보조정리(두 양수 차$\times$로그비)에서 별도로 증명한다. 좌변에 $T$ 가 곱해진
것은 entropy production $\dot S_\irr$[W/K]을 \emph{열률}[W]로 환산했기 때문이다($T\dot S_\irr$ 가 dissipated
power). \textbf{차원 검산.} $J_\alpha X_\alpha$ 의 곱이 [W]이려면 force 와 flux 가 공액(예: molar flux
$[\mathrm{mol/s}]\times$ molar affinity $[\mathrm{J/mol}]=[\mathrm{J/s}]=\mathrm W$)이어야 한다 — 이 공액성이
flux$\times$force 의 정의 자체에 들어 있다.

\textbf{전기화학 반응의 specialization.} 음극 반응을 molar reaction flux $J_n$ [mol/s]와 generalized reaction
affinity $\mathcal A_n$ [J/mol]로 쓰면 식~\eqref{eq:ch4_fluxforce} 의 한 쌍이
\begin{equation}
\boxed{\;T\dot S_{\irr,n}=J_n\,\mathcal A_n\ \ge\ 0\;.}
\label{eq:ch4_JA}
\end{equation}
\textbf{$J_n\mathcal A_n$ 이 $I_n\eta_n$ 보다 기본형인 이유.} $I_n=FJ_n$, $\mathcal A_n=F\eta_n$ convention 이
일관되면 $J_n\mathcal A_n=I_n\eta_n$ 형과 연결되나, ``$I_n\eta_n$''의 부호는 전류 방향$\cdot$전극 전위
convention 에 종속된다(Chapter 2 가 $q_\irr=|I|\eta$ 로 부호를 양정치화한 이유와 같다). 따라서 본 장은
\emph{positive-dissipation} 형 식~\eqref{eq:ch4_JA} 를 더 기본으로 두고, magnitude 표현
$\dot{\mathcal Q}_{\irr,n}^{\mathrm{mag}}=|I_n||\eta_n|$ 은 \emph{부호가 정렬됐을 때의 크기} 표현이지 entropy
production 의 일반 정의가 아님을 명시한다(\AL{42}). 이로써 ``$I\eta$ 단독''(부호 미정렬)을 비가역 열로 쓰는
것을 차단한다(RB\_CHARTER step1: $q_\irr=|I|\eta\ge0$ 형).

% ====================================================================
\subsection{Transition-level entropy production (미시) 과 거시식의 정합}\label{sec:ch4_tr}
% ====================================================================
이 절이 본 장의 핵심이다 — \emph{Chapter 3 의 forward/backward rate $r_j^\pm$ 가 어떻게 거시 비가역 발열
$n_j^{\eff}I_j\eta_j$ 로 환원되는가}.

\textbf{forward/backward unidirectional flux.} Chapter 3 transition kinetics(식 (L2))의 net flux 를 forward 와
backward 의 \emph{unidirectional} 성분으로 분해한다(\AL{41}):
\begin{equation}
\mathcal J_j^+=r_j^+\,(1-\xi_j),\qquad \mathcal J_j^-=r_j^-\,\xi_j,\qquad
\frac{\dd\xi_j}{\dd t}=\mathcal J_j^+-\mathcal J_j^-.
\label{eq:ch4_Jpm}
\end{equation}
$\mathcal J_j^+$ 는 미점유$\to$점유(방전 방향) 단위시간 진행, $\mathcal J_j^-$ 는 그 역방향이다(Chapter 3
식 $\dd\xi_j/\dd t=r_j^+(1-\xi_j)-r_j^-\xi_j$ 와 동일물; 본 장은 두 성분을 \emph{따로} 보존한다). \textbf{차원
검산.} $\mathcal J_j^\pm=r_j^\pm[\mathrm{s^{-1}}]\times$무차원 fraction $=\mathrm{s^{-1}}$, net 도 $\mathrm{s^{-1}}$
— 정합. 두 성분 모두 비음($r_j^\pm\ge0$, $0\le\xi_j\le1$).

\subsubsection{transition entropy production 의 network 형 (무비약)}\label{sec:ch4_trep}
\textbf{물리.} 비평형 열역학의 network/master-equation 이론에서, 두 상태 사이를 오가는 한 transition 의 entropy
production 은 \emph{net flux}와 \emph{forward/backward flux 비의 로그}의 곱이다(\cite{schnakenberg1976,prigogine1967},
\AL{41},\AL{43}). 이 ``$(\text{net flux})\times\ln(\text{flux 비})$'' 구조가 master equation 의 국소 detailed
balance 에서 따라오며, transition $j$ 의 진행이 동반하는 비가역 발열률은(몰 함량 $Q_{j,\tot}/F$ 와 $RT$ 로
extensive 화)
\begin{equation}
\boxed{\;\dot{\mathcal Q}_{\irr,j}^{\tr}=\frac{Q_{j,\tot}}{F}\,RT\,\big(\mathcal J_j^+-\mathcal J_j^-\big)\,
\ln\!\frac{\mathcal J_j^+}{\mathcal J_j^-}\ \ge\ 0\;.}
\label{eq:ch4_tr_entropy}
\end{equation}
\textbf{$\ge0$ 의 무비약 증명(부호 보조정리; 2법칙).} $\mathcal J_j^+,\mathcal J_j^->0$ 인 두 양수에 대해
$(\mathcal J_j^+-\mathcal J_j^-)\ln(\mathcal J_j^+/\mathcal J_j^-)\ge0$ 임을 직접 보인다. (i) $\mathcal J_j^+>\mathcal J_j^-$
이면 차 $>0$ 이고 비 $>1$ 이라 $\ln>0$ — 곱 $>0$. (ii) $\mathcal J_j^+<\mathcal J_j^-$ 이면 차 $<0$ 이고 비
$<1$ 이라 $\ln<0$ — 음$\times$음 $=$ 곱 $>0$. (iii) $\mathcal J_j^+=\mathcal J_j^-$(평형, net flux $0$)이면
차 $=0$ 이라 곱 $=0$. (경계 $\mathcal J_j^-\!\to0^+$: 차 $\to\mathcal J_j^+>0$, $\ln\to+\infty$ 이므로 곱
$\to+\infty\ge0$ — domain guard 식~\eqref{eq:ch4_tr_entropy} boundbox 로 정칙화). 세 경우 모두 $\ge0$ 이며, 이것이 식~\eqref{eq:ch4_tr_entropy} 의 2법칙 정합이다(``두
양수의 (차)$\times$(로그비)는 항상 음이 아니다''). \textbf{차원 검산.} $(Q_{j,\tot}/F)[\mathrm{mol}]\times
RT[\mathrm{J/mol}]\times(\mathcal J_j^+-\mathcal J_j^-)[\mathrm{s^{-1}}]\times\ln(\cdot)[\text{무차원}]
=\mathrm{mol\cdot(J/mol)\cdot s^{-1}}=\mathrm{J/s}=\mathrm W$ — 발열률의 올바른 단위. ($RT$ 가 $\mathrm{J/mol}$,
$Q_{j,\tot}/F$ 가 $\mathrm{mol}$ 이라 두 몰 인자가 곱해져 extensive [W]가 됨 — Chapter 1 $\rho_G$[mol/J]$\cdot$
Chapter 2 $\sigma_j$[W/mol] 의 ``intensive$\leftrightarrow$extensive 몰 인자 정합'' 교훈 계승.)

\textbf{extensive 화의 다리(무비약; $k_B\!\to\!R$, $\tfrac12$ 부재 사유).} Schnakenberg 의 단일 directed
transition(2-state site)에 대한 per-site entropy production 은 $\dot s_j^{\mathrm{site}}=k_B(\mathcal J_j^+-\mathcal J_j^-)
\ln(\mathcal J_j^+/\mathcal J_j^-)$ [W/K]다 — transition $j$ 를 \emph{단일 directed edge} 로 한 번만 세므로 양방향
edge 이중합의 $\tfrac12$ 인자는 \emph{없다}(이중합 $\tfrac12\sum_{\text{edges}}$ 와 단일-쌍 $\sum_j$ 는 동일물의
다른 셈법; 본 장은 후자). 여기에 (a) $T$ 를 곱해 per-site 발열률 [W], (b) site 수 $N_{\mathrm{site},j}=N_A(Q_{j,\tot}/F)$
(몰 함량 $\times$ Avogadro)를 곱해 독립 site 합(mean-field 전제: site 간 \emph{무상관} — entropy production 이
site 수에 extensive; site-site 상관이 있으면 보정항 발생), (c) $N_A k_B=R$ 로 묶으면 $N_A(Q_{j,\tot}/F)\,k_B T=(Q_{j,\tot}/F)RT$
가 되어 식~\eqref{eq:ch4_tr_entropy} 의 prefactor 가 무유도 점프 없이 따라온다(per-site $k_B$ $\to$ per-mole $R$
치환이 site 수 몰 인자와 정확히 짝).

\begin{verifybox}
\textbf{미시 entropy production $=$ 거시 $n_j^{\eff}I_j\eta_j$ (Ch2$\cdot$Ch3 정합, 무비약; CHARTER step1$\cdot$2$\cdot$3).}
transition chemical affinity 를 network thermo 정의 $\mathcal A_j^{\chem}=RT\ln(\mathcal J_j^+/\mathcal J_j^-)$
(\AL{43})로 두고, 그것을 Chapter 3 의 미시 정합으로 전개한다.

\emph{(1단계) flux 비를 rate 비$\times$점유 odds 로 분해.} 식~\eqref{eq:ch4_Jpm} 에서
$\dfrac{\mathcal J_j^+}{\mathcal J_j^-}=\dfrac{r_j^+(1-\xi_j)}{r_j^-\xi_j}=\dfrac{r_j^+}{r_j^-}\cdot\dfrac{1-\xi_j}{\xi_j}$.
따라서 $\ln(\mathcal J_j^+/\mathcal J_j^-)=\ln(r_j^+/r_j^-)+\ln[(1-\xi_j)/\xi_j]=\ln(r_j^+/r_j^-)-\ln[\xi_j/(1-\xi_j)]$.

\emph{(2단계) Chapter 3 detailed balance 대입(★ 두 전제 명시).} Chapter 3 의 detailed balance ln 형
$\ln(r_j^+/r_j^-)=(V_n-U_j)/w_j$ 를 넣는다. \textbf{이 ln 형은 두 전제에서만 닫힌형으로 성립한다}(Chapter 3
verifybox 한정과 \emph{동일 입도}): (i) ratio $r_j^+/r_j^-$ 가 $\xi_j$\emph{-무관}(즉 $C_j\equiv0$; Chapter 3 의
대칭 split, \AL{34})이라야 \emph{평형} odds 관계가 임의 \emph{비평형} $\xi_j$ 로 외삽되어 닫힌형이 되고, (ii)
\emph{내부 전위} $V_n$ 기준이다(detailed balance 는 평형 전하보존식의 해 $V_n$ 에서 성립). $C_j\not\equiv0$($r^+/r^-$
의 $\xi_j$-의존)이면 우변에 $RT\,C_j(\xi_j)$ 항이 추가되어 아래 닫힘이 깨지므로, 그 경우 미시=거시 환원은 BOUNDED 다
(\AL{44}; 본 장 기본 전개는 $C_j\equiv0$ 대칭 split). 전제 (i),(ii) 하에서
\begin{equation}
\mathcal A_j^{\chem}=RT\,\ln\!\frac{\mathcal J_j^+}{\mathcal J_j^-}
=RT\Big[\frac{V_n-U_j}{w_j}-\ln\frac{\xi_j}{1-\xi_j}\Big]
=\frac{RT}{w_j}\Big[(V_n-U_j)-w_j\ln\frac{\xi_j}{1-\xi_j}\Big].
\label{eq:ch4_affinity_raw}
\end{equation}

\emph{(3단계) 대괄호 $=$ overpotential $\eta_j$ (★ 기본형 $V_{n,\drive}=V_n$ 전제 명시; CHARTER step2).}
대괄호 안은 $(V_n-U_j)-w_j\ln[\xi_j/(1-\xi_j)]$ 다. 여기서 \S\ref{sec:ch4_eta_single} 의 $V_{\eq,j}(\xi_j)=U_j
+w_j\ln[\xi_j/(1-\xi_j)]$ 를 빼는 형으로 묶으면 $(V_n-U_j)-w_j\ln[\xi_j/(1-\xi_j)]=V_n-V_{\eq,j}(\xi_j)$. 본 장
규약(식~\eqref{eq:ch4_eta_def})의 $\eta_j=V_{n,\drive}-V_{\eq,j}(\xi_j)$ 와 일치시키려면 \textbf{기본형
$V_{n,\drive}=V_n$}($\eta_{n,\drive}=0$, RB\_CHARTER step2)을 \emph{여기서 명시적으로 전제}한다 — 그래야
detailed balance 가 쓴 $V_n$ 과 dissipation force 가 쓴 $V_{n,\drive}$ 가 같은 무대에서 닫힌다. 이 전제 하에서
$(V_n-U_j)-w_j\ln[\xi_j/(1-\xi_j)]=V_{n,\drive}-V_{\eq,j}(\xi_j)=\eta_j$. (강구동 $V_{n,\drive}\ne V_n$ 에서는
detailed balance 의 $V_n$ 과 force 의 $V_{n,\drive}$ 가 갈리므로 이 닫힘이 일반적으로 성립하지 않는다 —
BOUNDED, \S\ref{sec:ch4_general}.)

\emph{(4단계) $n_j^{\eff}$ 흡수 + 부호 인자 $s_{\phi,j}$ 복원(★ CHARTER step1$\cdot$3).} $RT/w_j$ 에 Chapter 3 의
$n_j^{\eff}=RT/(Fw_j)$(식 (L4))를 적용하면 $RT/w_j=n_j^{\eff}F$ 다. \emph{위 1$\sim$3단계 대수는 $s_{\phi,j}$ 를
포함하지 않으므로}, boxed 결과의 일반 $s_{\phi,j}$ 인자는 Chapter 3 명시형과의 \emph{표기 일관}을 위한 것이며,
충전 branch($s_{\phi,j}=-1$)에서 logistic 지수 부호로부터의 \emph{유도}는 Chapter 5 로 이월한다(본 장 유도는 방전
$s_{\phi,j}=+1$ 한정). 이 전제 하에서 식~\eqref{eq:ch4_affinity_raw} 는
\begin{equation}
\boxed{\;\mathcal A_j^{\chem}=s_{\phi,j}\,n_j^{\eff}F\,\eta_j\;,\qquad
\eta_j=V_{n,\drive}-V_{\eq,j}(\xi_j)\;.}
\label{eq:ch4_affinity_eta}
\end{equation}
($\eta_j$ 는 keybox 식~\eqref{eq:ch4_eta_def} 와 \emph{동일}하게 $s_{\phi,j}$ 를 \emph{포함하지 않으며}, 부호
인자 $s_{\phi,j}$ 는 $\mathcal A_j^{\chem}$ 의 scale 인자로만($\eta_j$ \emph{밖}) 붙는다; 방전 $s_{\phi,j}=+1$
기본이며 충전 branch 부호 \emph{유도}는 Chapter 5 로 이월한다.)
이는 Chapter 3 의 명시형 $\mathcal A_j=s_{\phi,j}n_j^{\eff}F(V_{n,\drive}-U_j)$ 와 \emph{같은 $s_{\phi,j}$,
같은 $n_j^{\eff}F$ scale} 을 비가역 \emph{heat-force} $\eta_j$ 에 적용한 형이다(rate-affinity 는 중심 $U_j$
기준, heat-force 는 국소 평형 $V_{\eq,j}$ 기준 — \S\ref{sec:ch4_eta_single} boundbox 의 구분 보존).

\emph{(5단계) entropy production 을 발열로(미시=거시).} 식~\eqref{eq:ch4_tr_entropy} 를 $\dot n_j$ 와
$\mathcal A_j^{\chem}$ 로 다시 쓰면 $\dot{\mathcal Q}_{\irr,j}^{\tr}=(Q_{j,\tot}/F)(\mathcal J_j^+-\mathcal J_j^-)
\cdot RT\ln(\mathcal J_j^+/\mathcal J_j^-)=\dot n_j\,\mathcal A_j^{\chem}$ 다($\dot n_j=(Q_{j,\tot}/F)(\dd\xi_j/\dd t)
=(Q_{j,\tot}/F)(\mathcal J_j^+-\mathcal J_j^-)$). 여기에 식~\eqref{eq:ch4_affinity_eta} 를 넣으면 $F$ 가 상쇄되어
\begin{equation}
\boxed{\;\dot{\mathcal Q}_{\irr,j}^{\tr}=\dot n_j\,\mathcal A_j^{\chem}
=\frac{Q_{j,\tot}}{F}\frac{\dd\xi_j}{\dd t}\cdot n_j^{\eff}F\,\eta_j
=n_j^{\eff}\,I_j\,\eta_j\;.}
\label{eq:ch4_micro_macro}
\end{equation}
즉 \textbf{미시 transition entropy production 이 거시 flux$\times$overpotential 형으로 정확히 환원}되며, ideal
width $w_j=RT/F$($n_j^{\eff}=1$)에서 $\dot{\mathcal Q}_{\irr,j}^{\tr}=I_j\eta_j$ 다. \textbf{부호(2법칙
재확인).} $\xi_j<\xi_{\eq,j}\Rightarrow$ (정의상) $V_{\eq,j}(\xi_j)<V_{n,\drive}\Rightarrow\eta_j>0$, 동시에
$\dd\xi_j/\dd t=k_j^{\mathrm{phys}}(\xi_{\eq,j}-\xi_j)>0\Rightarrow I_j>0$ — 둘이 같은 부호라 $n_j^{\eff}I_j\eta_j>0$
($n_j^{\eff}>0$). 반대로 $\xi_j>\xi_{\eq,j}$ 이면 둘 다 음이라 곱은 여전히 양; 평형서 $\eta_j\to0,\ I_j\to0$
동시 소멸. \textbf{이로써 Ch2(거시 flux$\times$overpotential)$\cdot$Ch3(미시 forward/backward rate, $n_j^{\eff}$)
$\cdot$Ch4(entropy production)가 한 표현 $n_j^{\eff}I_j\eta_j$ 로 닫힌다}(P3-6 미시=거시 정합 통과).
\textbf{위계(CHARTER step3).} \emph{Chapter 2 의 $\sum_j J_{n,j}F\eta_j=\sum_j I_j\eta_j$ 는 이 결과의
$n_j^{\eff}=1$(ideal-width) 특수해}이고, 일반 dissipation 은 $\sum_j n_j^{\eff}I_j\eta_j$ 다(\S\ref{sec:ch4_general})
— 두 장은 모순이 아니라 일반/특수 위계다(Chapter 3 의 $n_j^{\eff}=RT/(Fw_j)$ refinement 와 동일 맥락).
\end{verifybox}

\textbf{유도가 쓴 전제의 명시(★ 이중정의 차단 재확인).} 위 verifybox 는 (i) detailed balance(Ch3, \emph{내부}
$V_n$ 기준 ln 비), (ii) $n_j^{\eff}=RT/(Fw_j)$(Ch3), (iii) \textbf{기본형 $V_{n,\drive}=V_n$}(RB\_CHARTER
step2, 3단계서 명시 전제)을 사용했다. (iii)이 없으면 detailed balance 가 쓴 $V_n$ 과 dissipation force 가 쓴
$V_{n,\drive}$ 가 서로 달라 $\mathcal A_j^{\chem}=n_j^{\eff}F\eta_j$ 가 닫히지 않는다 — 이 전제를 사용 시점에
박는 것이 Phase A CRITICAL($\eta$ 이중정의)의 해소다. 또 (iv) Chapter 3 의 keystone 한정(Level A $=$ Level B
의 detailed-balance 극한; $\xi_{\sstat,j}\to\xi_{\eq,j}$ 가 $C_j$ 의 $V_n$$\cdot$$\xi_j$-무의존 partition 에서만
성립)을 상속한다. 강구동 branch 비대칭(non-detailed-balance, Chapter 5)에서는 $\mathcal A_j^{\chem}$ 가 일반화된다.

\begin{boundbox}
\textbf{log singularity 는 floor 아닌 domain guard (\AL{49}).} $\mathcal J_j^\pm\to0$ 에서 $\ln(\mathcal J_j^+/\mathcal J_j^-)$
발산은 coarse-graining resolution 에 해당하는 작은 positive floor 로 처리한다 — 이는 heat fitting parameter 가
아니라 \emph{numerical domain guard} 다(GS-4 의 cap 금지와 구분: 물리 모델의 임계값이 아니라 수치 정의역 보호).
\textbf{게다가} 식~\eqref{eq:ch4_micro_macro} 형 $n_j^{\eff}I_j\eta_j$ 를 쓰면 $\mathcal J_j^\pm$ 가 평형 부근에서
작아도 $I_j=Q_{j,\tot}(\dd\xi_j/\dd t)\to0$ 으로 \emph{매끄럽게} 0 이 되어(그리고 $\eta_j\to0$) 발산 자체가
나타나지 않는다 — 즉 거시형은 미시형의 log 발산을 자동으로 정칙화한다. 따라서 실제 발열 계산은 거시형
$n_j^{\eff}I_j\eta_j$ 를 쓰고, 미시형 식~\eqref{eq:ch4_tr_entropy} 는 \emph{출처(entropy production)와 부호
($\ge0$) 의 근거}로 둔다(BOUNDED).
\end{boundbox}

\textbf{채택 조건(후보식의 검증).} 식~\eqref{eq:ch4_tr_entropy}--\eqref{eq:ch4_micro_macro} 를 실제 발열항으로
\emph{확정}하려면 (1) $r_j^\pm$ 가 실제 coarse-grained transition rate, (2) $\xi_j$ 가 progress coordinate,
(3) detailed balance $\xi_{\eq,j}$ 정합, (4) calorimetry/rest-relaxation 독립 검증이 필요하다. 따라서 본 식은
해석 가능한 물리식이되, 독립 검증 없이 \emph{모든} relaxation heat 를 이 항으로 단정하지 않는다(Chapter 1$\cdot$2
의 forward-only$\cdot$식별성 원칙 계승).

% ====================================================================
\subsection{비가역열 일반형 $\sum_j n_j^{\eff}I_j\eta_j\ge0$ (CHARTER step3 일반/특수 위계)}\label{sec:ch4_general}
% ====================================================================
\textbf{물리.} 식~\eqref{eq:ch4_micro_macro} 의 per-transition 결과를 모든 동시 진행 transition 에 대해 합하면
음극의 \emph{총} 계면 반응 비가역 발열률이다(\AL{45}):
\begin{equation}
\boxed{\;\dot{\mathcal Q}_{\irr,n}=\sum_{j=1}^{N_p} n_j^{\eff}\,I_j\,\eta_j\ \ge\ 0\;,\qquad
n_j^{\eff}=\frac{RT}{F\,w_j(T)}\;.}
\label{eq:ch4_qirr_general}
\end{equation}
단 항별 $\ge0$ 은 detailed-balance 기본형($V_{n,\drive}=V_n$, $n_j^{\eff}$ 일반 허용)에서 성립한다(\S\ref{sec:ch4_tr}
verifybox 항별 부호 논증); 강구동 branch 비대칭(non-detailed-balance)에서는 $\mathcal A_j^{\chem}$ 가 일반화되어
항별 부호가 일반적으로 보장되지 않으므로 BOUNDED 로 Chapter 5 에서 다룬다.
\textbf{$\ge0$ 의 무비약 확인(항별 + 합).} \S\ref{sec:ch4_tr} verifybox 에서 각 항 $n_j^{\eff}I_j\eta_j$ 는
$\xi_j\lessgtr\xi_{\eq,j}$ 양쪽에서 $I_j$ 와 $\eta_j$ 가 같은 부호라 항별로 $\ge0$ 이다(detailed-balance
기본형 전제). 따라서 비음항들의 합 $\sum_j n_j^{\eff}I_j\eta_j\ge0$ 이 2법칙으로 보장된다. \textbf{차원 검산.}
$n_j^{\eff}$(무차원)$\times I_j[\mathrm A]\times\eta_j[\mathrm V]=\mathrm{A\cdot V}=\mathrm W$ — 발열률 단위.
\textbf{부호 양정치 형(CHARTER step1).} 위계의 \emph{환원점은 항별 합} $\sum_j I_j\eta_j$($n_j^{\eff}=1$
ideal-width)로 고정한다. \emph{단일 등가 과전압} 형 $|I|\eta$($\eta\equiv\sum_j(J_{n,j}F/|I|)\eta_j$)는 이와
별개의 표기로, \emph{단일 우세 transition} $\cdot$ ideal-width 극한에서만 이 합이 하나의 전류$\cdot$과전압 크기
곱으로 축약된 형이다(Chapter 2 의 $q_\irr=|I|\eta\ge0$ 와 정합). 어느 표기든 ``$I_j\eta_j$ 단독''(부호 미정렬)이
아니라 부호가 이미 정렬된 양정치 합임을 유지한다.

\begin{keybox}
\textbf{일반형$\leftrightarrow$Ch2 특수형 위계 (CHARTER step3).} 식~\eqref{eq:ch4_qirr_general} 의 \emph{일반형}
$\sum_j n_j^{\eff}I_j\eta_j$ 와 Chapter 2 의 \emph{특수형} $\sum_j I_j\eta_j$ 의 관계는 \emph{모순이 아니라
일반/특수 위계}다:
\begin{equation}
\underbrace{\dot{\mathcal Q}_{\irr,n}=\sum_j n_j^{\eff}I_j\eta_j}_{\text{Ch4 일반형}}
\ \xrightarrow{\ n_j^{\eff}=RT/(Fw_j)\to1\ (w_j=RT/F,\ \text{ideal})\ }\
\underbrace{\sum_j I_j\eta_j}_{\text{Ch2 특수형}}.
\label{eq:ch4_hierarchy}
\end{equation}
Chapter 2 는 ideal-width($n_j^{\eff}=1$)를 기본으로 써서 $\sum_j I_j\eta_j$ 를 얻었고(그 boundbox 에서 일반형을
본 장으로 forward-ref), 본 장은 effective width($w_j\ne RT/F$, 따라서 $n_j^{\eff}\ne1$)를 허용하는 일반형을
닫는다. 그 일반형이 미시 transition entropy production(식~\eqref{eq:ch4_tr_entropy})에서 유도되므로, $n_j^{\eff}$
가중은 임의 피팅 인자가 아니라 \emph{detailed-balance 정합이 강제한} 물리 인자다(Chapter 3 \S$n_j^{\eff}$
정의장과 동일 출처).
\end{keybox}

\textbf{과전압의 물리 성분 분해(Ch2 계승).} 총 과전압 $\eta$(또는 transition별 $\eta_j$)는 charge-transfer
($\eta_\ct$, Chapter 3 의 effective barrier 활성화), ohmic($\eta_\ohm$, Chapter 1 의 $R_n$ 분극),
transport/concentration($\eta_\conc$) 기여의 합으로 분해되며 각각 $\ge0$ 이다. 단 이 성분들을 \emph{발열 저항}
으로 쓸 때는 Chapter 3 의 $R_n\ne R_\ct\ne R_{n,\eff}$ 구분을 유지한다(\S\ref{sec:ch4_transport}).

% ====================================================================
\subsection{가역 열 (1): 평형 전위 온도계수 (Bernardi)}\label{sec:ch4_revOCV}
% ====================================================================
\textbf{물리.} 가역 열은 평형 entropy 가 준다 — 비가역 dissipation 과 \emph{분리}되는 열역학 기원 항이다.
Chapter 2 가 전하 보존식에서 유도한 음극 평형 전위 온도계수를 그대로 받는다($Q_\cell,Q_{j,\tot}$ 기준 용량
고정, \AL{20})\textemdash 외부 lookup 이 아니라 \emph{전하 보존식의 평형 특수해 미분}이다(P3-2):
\begin{equation}
\left(\frac{\partial V_{n,\OCV}}{\partial T}\right)_q
=-\,\frac{\dfrac{\partial Q_\bg}{\partial T}+\displaystyle\sum_j Q_{j,\tot}\left(\dfrac{\partial\xi_{\eq,j}}{\partial T}\right)_{V_n}}
{\dfrac{\partial Q_\bg}{\partial V_n}+\displaystyle\sum_j Q_{j,\tot}\dfrac{\partial\xi_{\eq,j}}{\partial V_n}}.
\label{eq:ch4_ocv_dvdt}
\end{equation}
방전 기준 음극 reduced reversible heat 는(Chapter 2 (L6) 의 Bernardi balance)
\begin{equation}
\boxed{\;\dot{\mathcal Q}_{\rev,n}^{\OCV}=s_{\rev,n}^{\conv}\,|I|\,T\left(\frac{\partial V_{n,\OCV}}{\partial T}\right)_q\;.}
\label{eq:ch4_rev_ocv}
\end{equation}
\textbf{차원 검산.} $|I|\,T\,(\partial V/\partial T)=\mathrm A\cdot\mathrm K\cdot(\mathrm{V/K})=\mathrm{A\,V}=\mathrm W$
— 가역 발열률 단위. $s_{\rev,n}^{\conv}\in\{+1,-1\}$ 는 음극 전위$\cdot$전류 방향$\cdot$heat-generation-positive
규약을 명시한 뒤 고정하는 bookkeeping factor 다(피팅 X; RB\_CHARTER step1 의 $q_\rev$ 부호 convention).
\textbf{부호가변성.} $q_{\rev,j}=I_j\Pi_j$ ($\Pi_j=T\partial U_j/\partial T$, 식 (L6))는 $\partial U_j/\partial T$
와 전류 방향에 따라 흡열/발열 모두 가능하다 — 비가역 열($\ge0$)과 달리 가역 열은 부호가변이며, 그래서 둘을 한
경험항으로 섞지 않는다(Chapter 2 keystone 의 ``가역열$=$entropy(방향) / 비가역열$=$소산(속도)'' 분업).

\textbf{금지되는 연결(P3-1).}
\begin{equation}
\dot{\mathcal Q}_{\rev,n}\ \not\propto\ |I|\,T\,\frac{\dd V_{n,\app}}{\dd T}.
\label{eq:ch4_forbid}
\end{equation}
$V_{n,\app}=V_n+s_I|I|R_n$ 에는 분극$\cdot$kinetic lag$\cdot$transport loss$\cdot$hysteresis-like gap$\cdot R_n(T)$
가 섞이므로 그 경로 미분은 reversible entropy coefficient 가 아니다 — 이를 가역열로 쓰면 비가역 소산을 가역
entropy 로 오계상한다(Chapter 2 의 금지 연결 계승; 식~\eqref{eq:ch4_forbid}).

\begin{verifybox}
\textbf{Bernardi 가역열 $=$ 미시 reaction entropy 합 (Ch2 정합 재확인, 무비약).} Chapter 2 (식 (L6))의
$\Delta S_{\rxn,j}=s_{\phi,j}F\,\partial U_j/\partial T$(식 (L6) 동치형, 부호 인자 $s_{\phi,j}$ 는 이 환산에서만
도입)를 per-transition 가역열에 넣는다. transition $j$ 의 reaction rate 는 표 정의대로 $\dot n_j=I_j/F$
[mol/s](Faraday 법칙; 부호 인자 없음)이고, 그 entropic 열률은 $T\Delta S_{\rxn,j}\dot n_j$ 이므로
\begin{equation}
q_{\rev,j}=T\,\Delta S_{\rxn,j}\,\frac{I_j}{F}
=s_{\phi,j}\,I_j\,T\,\frac{\partial U_j}{\partial T}=s_{\phi,j}\,I_j\,\Pi_j
\ \xrightarrow{\ s_{\phi,j}=+1\ (\text{방전})\ }\ I_j\,\Pi_j,
\label{eq:ch4_qrev_micro}
\end{equation}
(부호 인자 $s_{\phi,j}$ 는 $\dot n_j$ 가 아니라 $\Delta S_{\rxn,j}$ 환산에서 나오며 \emph{잔존}한다 — 방전
$s_{\phi,j}=+1$ 에서만 $I_j\Pi_j$ 와 일치하고, 충전 부호는 Chapter 5). 즉 \textbf{미시 reaction entropy
$\Delta S_{\rxn,j}$ 의 합} $\dot{\mathcal Q}_{\rev}=\sum_j s_{\phi,j}I_j\Pi_j$(방전 $\sum_j I_j\Pi_j=\sum_j I_j
T\partial U_j/\partial T$)가 거시 Bernardi 형 $s_{\rev,n}^{\conv}|I|T(\partial V_{n,\OCV}/\partial T)_q$ 와 같은
\emph{reaction-entropy 정보}로 정합한다 — 단 이 정합은 \S\ref{sec:ch4_revxi} 준평형 극한 한정이며, 본 verifybox
는 reaction-entropy 정보의 \emph{동일성}을 보일 뿐 cell-level 수치 등식을 닫지 않는다(닫힘 조건은
\S\ref{sec:ch4_revxi} 식~\eqref{eq:ch4_consistency} 와 그 width$\cdot D_q$ BOUNDED). Chapter 2 의 double-counting
차단 정합식과 동일 제약이다. \textbf{비가역과의 대비.} 본 절 가역열은 \emph{reaction} entropy
$\Delta S_{\rxn,j}$(평형, $\partial U/\partial T$)에서, \S\ref{sec:ch4_tr} 비가역열은 transition entropy
production(non-equilibrium, $\eta_j$)에서 나온다 — 둘은 별개 entropy 이며(Chapter 2 \AL{24}: reaction entropy
$\ne$ activation entropy), 한 뿌리($U_j(T)$ 와 $V_{\eq,j}(\xi_j)$)에서 갈라진 가역/비가역 두 갈래다.
\end{verifybox}

% ====================================================================
\subsection{가역 열 (2): transition entropy basis 와 double counting 금지}\label{sec:ch4_revxi}
% ====================================================================
Chapter 1$\cdot$2 의 핵심 내부 상태변수는 $\xi_j$ 다. 따라서 상변이 transition 의 entropy 변화는 \emph{진행률
시간율} $\dd\xi_j/\dd t$ 를 통해서도 표현된다 — 식~\eqref{eq:ch4_rev_ocv} 의 OCV 계수 방식과 \emph{같은 entropy
정보}를 더 미시적으로 쓴 것이다(\AL{20}). $\Delta S_j(\equiv\Delta S_{\rxn,j})$ = $j$번째 effective transition
의 방전 방향 mol electron(또는 Li-equiv) 1 mol 당 entropy change 로 두면
\begin{equation}
\dot{\mathcal Q}_{\rev,\xi}=\sum_{j=1}^{N_p}s_{\rev,\xi,j}^{\conv}\,T\,\Delta S_j\,\frac{Q_{j,\tot}}{F}\,\frac{\dd\xi_j}{\dd t}.
\label{eq:ch4_rev_xi}
\end{equation}
\textbf{차원 검산.} $T[\mathrm K]\cdot\Delta S_j[\mathrm{J/(mol\,K)}]\cdot(Q_{j,\tot}/F)[\mathrm{mol}]\cdot
(\dd\xi_j/\dd t)[\mathrm{s^{-1}}]=\mathrm{J/s}=\mathrm W$ — 가역 발열률. 직접 입력은 $\dd\xi_j/\dd t$(휴지 구간서
정전류 좌표 변환 금지, \S\ref{sec:ch4_rest}).

\textbf{double counting 금지(택일 + 정합 제약; \AL{25}).} 식~\eqref{eq:ch4_rev_ocv} 와 식~\eqref{eq:ch4_rev_xi}
는 \emph{같은 평형 entropy 정보}를 서로 다른 수준에서 표현하므로, 둘을 독립 발열항으로 동시에 자유롭게 더하면
double counting 이다(평형 $V_{n,\OCV}$ 온도계수에 이미 transition/background entropy 가 반영). 따라서 택일한다:
\begin{enumerate}[label=\Alph*),leftmargin=2.0em]
\item \textbf{OCV 중심}: 전체 가역 열을 식~\eqref{eq:ch4_rev_ocv} 로 대표; $\Delta S_j$ 는 분해 해석/파라미터
  interpretation 용으로만 쓴다.
\item \textbf{Transition entropy 중심}: $\Delta S_j,h_\bg$ 로 분해하되, OCV $\partial V/\partial T$ 는 독립항이
  아니라 $\Delta S_j,h_\bg$ 를 묶는 \emph{equilibrium consistency condition}(Chapter 2 정합식)으로 둔다:
\end{enumerate}
\begin{equation}
\boxed{\;s_{\rev,n}^{\conv}|I|\,T\left(\frac{\partial V_{n,\OCV}}{\partial T}\right)_q
=s_\bg^{\conv}\,T\,h_\bg\,\dot Q_\bg^{\prog}
+\sum_j s_{\rev,\xi,j}^{\conv}\,T\,\Delta S_j\,\frac{Q_{j,\tot}}{F}\,\frac{\dd\xi_{\eq,j}}{\dd t}\;.}
\label{eq:ch4_consistency}
\end{equation}
\begin{boundbox}
\textbf{이 정합은 reaction-entropy 성분에 한정한다(strict 등호 아님; \AL{25}).} 식~\eqref{eq:ch4_consistency} 의
정합은 \emph{reaction-entropy 성분($\propto\partial U_j/\partial T$)에 한정}해 성립한다. 좌변
$(\partial V_{n,\OCV}/\partial T)_q$ 는 식~\eqref{eq:ch4_ocv_dvdt} 의 $(\partial\xi_{\eq,j}/\partial T)_{V_n}$ 경유로
추가로 logistic width 온도계수 항 $-(V_n-U_j)(\partial w_j/\partial T)/w_j^2$ 와 분모 $D_q=\partial Q_\bg/\partial V_n
+\sum_j Q_{j,\tot}\partial\xi_{\eq,j}/\partial V_n$ 정규화를 포함하므로, 두 방식은 \emph{부분적으로만} 같은
정보다 — width$\cdot D_q$ 기여는 OCV 방식 고유다. 따라서 double-counting 차단은 reaction-entropy 성분에
부과되며, strict 등호는 $\partial w_j/\partial T=0$(또는 등온$\cdot$준정적 $+$ 단일 우세 transition) 전제 하에서만
닫힌다. 아래 준평형 caveat 는 $\dd\xi_j/\dd t$ 와 $\dd\xi_{\eq,j}/\dd t$ 의 차이를 닫을 뿐, 완전 평형에서도
잔존하는 width 온도계수 항($\partial w_j/\partial T\ne0$ 기여)을 막지 못한다.
\end{boundbox}
\textbf{부호 규약 정합 주의.} 좌$\cdot$우변의 $s^{\conv}$($s_{\rev,n}^{\conv},s_\bg^{\conv},s_{\rev,\xi,j}^{\conv}$)
는 모두 \emph{동일한 heat-positive 규약}에서 유도돼야 본 제약이 부호 모순 없이 닫힌다(Chapter 2 의 정합 제약과
동일 — 서로 다른 bookkeeping 규약으로 독립 도입하면 같은 reaction entropy 가 좌우에서 반대 부호로 들어와 제약이
깨진다). \textbf{준평형 전제.} 좌변 OCV 계수는 \emph{평형}($\xi_j=\xi_{\eq,j}$) 미분이고 우변 가역열 항
(식~\eqref{eq:ch4_rev_xi})은 \emph{실제} 진행률 $\dd\xi_j/\dd t$ 를 쓰므로, 두 표현이 ``같은 entropy 정보''로
정합하는 것은 \emph{준평형 극한}($\xi_j\approx\xi_{\eq,j}$, 따라서 $\dd\xi_j/\dd t\approx\dd\xi_{\eq,j}/\dd t$)에
한정된다 — 가역열의 정의역 자체가 준평형이므로 이 한정은 모순이 아니라 적용범위 명시다.

% ====================================================================
\subsection{Background heat 와 coordinate shift}\label{sec:ch4_bg}
% ====================================================================
$Q_\bg$ 는 배경 누적 용량이지 entropy 자체가 아니다. background reversible heat 를 쓰려면 별도 background
\emph{reversible-heat potential coefficient} 가 필요하다(\AL{26}). 여기서 $V_{\bg,\eq}$ 는 background 용량
$Q_\bg(V_n,T)$ 를 \emph{평형 $V_n$ 축으로 본} 가역분 전위좌표($\partial Q_\bg/\partial V_n>0$ 로 일대일), $h_\bg$
는 그 온도계수다:
\begin{equation}
h_\bg(V_n,T)\equiv\left(\frac{\partial V_{\bg,\eq}}{\partial T}\right)_{\mathrm{bg}},\qquad
\dot{\mathcal Q}_{\rev,\bg}=s_\bg^{\conv}\,T\,h_\bg(V_n,T)\,\dot Q_\bg^{\prog}.
\label{eq:ch4_qrev_bg}
\end{equation}
\textbf{차원 검산.} $\dot{\mathcal Q}_{\rev,\bg}=s_\bg^{\conv}\,T[\mathrm K]\cdot h_\bg[\mathrm{V/K}]\cdot
\dot Q_\bg^{\prog}[\mathrm{C/s}=\mathrm A]=\mathrm{K\cdot(V/K)\cdot A}=\mathrm{V\,A}=\mathrm W$ — 발열률 단위($h_\bg$
가 entropy $[\mathrm{J/(mol\,K)}]$ 가 아니라 \emph{per-charge 전위계수} $[\mathrm{V/K}]$ 라 transition 항의 $/F$
가 흡수돼 별도 $/F$ 없이 [W]가 닫힌다 — 이 점이 명칭을 ``background entropy coefficient'' 가 아니라
``background reversible-heat potential coefficient'' 로 두는 이유다).
\textbf{total derivative 와 progress 의 구분(무비약, 중요).} 전하 보존식 (L1) 을 시간 미분하면 $Q_\bg$ 가
$(V_n,T)$ 의 함수이므로 전미분 연쇄법칙으로 $\dot Q_\bg^{\mathrm{tot}}=(\partial Q_\bg/\partial V_n)\dd V_n/\dd t
+(\partial Q_\bg/\partial T)\dd T/\dd t$ 다. 이를 progress 와 coordinate shift 로 나눈다:
\begin{equation}
\dot Q_\bg^{\mathrm{tot}}=\dot Q_\bg^{\prog}+\dot Q_\bg^{\coord},\qquad
\dot Q_\bg^{\prog}=\frac{\partial Q_\bg}{\partial V_n}\frac{\dd V_n}{\dd t},\qquad
\dot Q_\bg^{\coord}=\frac{\partial Q_\bg}{\partial T}\frac{\dd T}{\dd t}.
\label{eq:ch4_qbg_split}
\end{equation}
여기서 \emph{$\dot Q_\bg^{\prog}$(전위 구동 progress)만} heat-generating 반응 진행으로 식~\eqref{eq:ch4_qrev_bg}
에 들어가고, $\dot Q_\bg^{\coord}$(온도 변화에 따른 capacity-coordinate shift)는 그대로 reaction rate 로 넣지
않는다 — 이 구분을 빠뜨리면 단순 온도 drift 가 가역 발열로 오계상된다(Chapter 2 의 bgsplit 계승). 등온 정전류
구간서는 실용적으로 $\dot Q_\bg^{\prog}\approx|I|-\sum_j Q_{j,\tot}(\dd\xi_j/\dd t)$ 로 쓸 수 있으나(전류 보존,
Chapter 3 식), \emph{비등온서는 $\dot Q_\bg^{\coord}$ 항 때문에 이 근사를 그대로 쓰지 않는다}.
\begin{equation}
\dot Q_\bg^{\prog}\ \approx\ |I|-\sum_j Q_{j,\tot}\frac{\dd\xi_j}{\dd t}\qquad(\text{등온 정전류 한정}).
\label{eq:ch4_qbg_iso}
\end{equation}

% ====================================================================
\subsection{Transport 및 물리적 병목 발열}\label{sec:ch4_transport}
% ====================================================================
\textbf{물리.} 강구동 제한을 arbitrary cap 이 아니라 \emph{물리적 병목}으로 두는 Chapter 1$\cdot$3 의 원칙을
발열에도 유지한다. 선형 near-equilibrium 에서 transport dissipation 은 $I^2R$ 형으로 근사되나, 일반적으로는
flux$\times$force 가 더 기본이다(\AL{47}). 선형 근사:
\begin{equation}
\dot{\mathcal Q}_{\transp,n}\approx I^2 R_{\transp,n}\qquad(\text{선형 near-equilibrium},\ R_{\transp,n}>0\Rightarrow\ge0),
\label{eq:ch4_transport_i2r}
\end{equation}
(여기서 $I$ 는 transport 채널을 통과하는 \emph{외부} 전류; 채널 flux 환산이 필요하면 $J_\ell$ 형
식~\eqref{eq:ch4_transport_ff} 을 쓴다. $R_{\transp,n}>0$ 이라 $I^2R_{\transp,n}\ge0$.)
고율/nonlinear transport:
\begin{equation}
\boxed{\;\dot{\mathcal Q}_{\transp,n}=\sum_\ell J_\ell\,X_\ell\ \ge\ 0\;,}
\label{eq:ch4_transport_ff}
\end{equation}
$J_\ell$=물질 flux, $X_\ell$=그 공액 force(chemical-potential / electrolyte-potential / concentration
gradient). 이는 식~\eqref{eq:ch4_fluxforce} 의 transport 채널 specialization 이며 2법칙으로 $\ge0$ 이다(porous
insertion electrode 의 thermal/transport 발열 모델은 \cite{thomasnewman2003}).
\textbf{비중복(double-count 금지).} 본 절 $\dot{\mathcal Q}_{\transp,n}$ 은 계면 반응 affinity 에 흡수되지
\emph{않는} 전해질/고체상 transport 채널(bulk$\cdot$분리막 ohmic)에 한정하며, \S\ref{sec:ch4_general} 계면
$\eta_\conc/\eta_\ohm$($\sum_j n_j^{\eff}I_j\eta_j$ 에 이미 포함)과 \emph{비중복}이다 — 두 항을 함께 쓸 때 ohmic
dissipation 을 이중 계상하지 않는다.
\textbf{차원 검산.} $I^2R=\mathrm{A^2\cdot\Omega}=\mathrm{A^2\cdot V/A}=\mathrm{A\,V}=\mathrm W$; flux$\times$force
$J_\ell X_\ell$ 도 공액쌍이라 [W] — 정합.
\begin{boundbox}
\textbf{세 저항의 구분 유지(\AL{47}; Ch1$\cdot$3 계승).} Chapter 1 의 $R_n$ 은 apparent voltage-axis
coefficient($V_{n,\app}-V_n=s_I|I|R_n$)이므로 $I^2R$ heat resistance 로 \emph{직접} 쓰지 않는다:
\begin{equation}
R_n\ne R_{n,\eff},\qquad R_n\ne R_{\ct,n},\qquad R_{\ct,n}\ne R_{n,\eff}.
\label{eq:ch4_Rsep}
\end{equation}
(i) $R_n$ $=$ apparent voltage-axis shift, (ii) $R_{\ct,n}=RT/(FA_\eff j_{0,n})$ $=$ charge-transfer
small-signal(Chapter 3), (iii) $R_{n,\eff}$ $=$ heat-generation effective($\dot{\mathcal Q}_\irr=I^2R_{n,\eff}$
형, 선형 응답 한정). $R_{n,\eff}$ 사용은 independent physical/calorimetric constraint 가 있을 때만 허용한다
(BOUNDED; 동시 자유 피팅 금지 — Chapter 3 의 식별성 chain).
\end{boundbox}

% ====================================================================
\subsection{휴지 구간 relaxation heat (Ch2 후보의 확정 분해)}\label{sec:ch4_rest}
% ====================================================================
\textbf{물리.} 휴지 구간에서는 $I=0,\ \dd q/\dd t=0$ 이지만 내부 진행률은 (L2) 로 계속 완화된다:
$\dd\xi_j/\dd t=r_j^+(1-\xi_j)-r_j^-\xi_j\ne0$. 외부 전류에 의한 ohmic heat 는 사라지나(전류 $0$), 내부
자유에너지 감소와 entropy production 은 남는다. Chapter 2 가 \emph{후보}로만 남긴 휴지 relaxation heat 를
본 장에서 가역/비가역으로 \emph{분해}한다(\AL{48}):
\begin{equation}
\dot{\mathcal Q}_{\xi,\rlx}=\dot{\mathcal Q}_{\rev,\xi}^{\rlx}+\dot{\mathcal Q}_{\irr,\xi}^{\rlx}.
\label{eq:ch4_rest_split}
\end{equation}
\textbf{비가역 성분은 식~\eqref{eq:ch4_micro_macro} 가 그대로 적용된다.} 휴지 중에도 \emph{내부} transition
current $I_j=Q_{j,\tot}(\dd\xi_j/\dd t)\ne0$ 이므로(외부 $|I|=0$ 이라도 transition 들이 서로 다른 속도로
완화하며 내부 전류를 주고받음), 비가역 후보는
\begin{equation}
\dot{\mathcal Q}_{\irr,\xi}^{\rlx}=\sum_j n_j^{\eff}\,I_j\,\eta_j\ \ge\ 0\qquad(\text{외부 }|I|=0,\ I_j=Q_{j,\tot}\dot\xi_j\ne0),
\label{eq:ch4_rest_irr}
\end{equation}
이고(식~\eqref{eq:ch4_qirr_general} 와 동형, 단 $\sum_j I_j$ 가 외부 전류가 아니라 \emph{내부} 전류이며, 총 내부
전류 보존 $\sum_j I_j+I_\bg^{\prog}=0$ 은 \emph{별개} 진술이다). \textbf{이 $\ge0$ 의 근거는 전류 보존이 아니다}:
휴지서도 모든 transition 이 \emph{공통 무대} $V_n$ 을 공유하므로 각 항에서 $I_j\,(\propto\xi_{\eq,j}-\xi_j)$ 와
$\eta_j\,(=V_n-V_{\eq,j}(\xi_j))$ 가 같은 부호이고(위 미시=거시 verifybox 의 항별 부호 논증을 휴지구간에 그대로
재적용), 따라서 항별 $n_j^{\eff}I_j\eta_j\ge0$ 에서 합 $\ge0$ 이 따라온다. 가역 후보는 transition entropy
basis(식~\eqref{eq:ch4_rev_xi}, 단 휴지서 $\dd\xi_j/\dd t$ 직접 사용)이며 명시적으로
\begin{equation}
\boxed{\;\dot{\mathcal Q}_{\rev,\xi}^{\rlx}=\sum_{j=1}^{N_p}s_{\rev,\xi,j}^{\conv}\,T\,\Delta S_j\,\frac{Q_{j,\tot}}{F}\,\frac{\dd\xi_j}{\dd t}\qquad(\text{외부 }|I|=0,\ \dd\xi_j/\dd t\ \text{직접})\;.}
\label{eq:ch4_rest_rev}
\end{equation}
\textbf{좌표 변환 금지.} 휴지/가변전류 구간에서는 $|I|/Q_\cell$ 이 시간 상수가
아니므로 정전류 좌표 변환 $\dd\xi_j/\dd t=(|I|/Q_\cell)\dd\xi_j/\dd q$ 를 쓰지 않고 $\dd\xi_j/\dd t$ 를 직접 쓴다.
\begin{flagbox}
\textbf{가역/비가역 분리 비율은 독립 검증 필요(\AL{48}).} 식~\eqref{eq:ch4_rest_split} 의 비가역 성분
(식~\eqref{eq:ch4_rest_irr})은 entropy production 으로 $\ge0$ 임이 보장되나, 가역 성분(transition entropy
exchange)과의 \emph{분리 비율}은 forward/backward rate 의 detailed-balance 정합 정도에 달려 있다. detailed
balance 가 정확히 성립하면 식~\eqref{eq:ch4_micro_macro} 가 비가역 성분을 닫지만, branch 비대칭(Chapter 5)에서는
$\mathcal A_j^{\chem}$ 가 일반화되어 분리 비율이 달라진다. 따라서 독립 calorimetry/rest-temperature response
없이 가역/비가역을 임의 분리하지 않는다(Chapter 2 가 후보로 남긴 이유의 확정 분해이되, \emph{정량} 비율은
검증 전제).
\end{flagbox}

% ====================================================================
\subsection{전기--열 coupling 계산 순서와 온도 되먹임}\label{sec:ch4_order}
% ====================================================================
온도는 Chapter 1--3 의 다음 항들에 되먹임된다: $Q_\bg(V_n,T)$, $\xi_{\eq,j}(V_n,T)$, $U_j(T)$, $w_j(T)$,
$r_j^\pm(\cdot,T)$, $k_j^{\mathrm{phys}}(\cdot,T)$, $R_n(q,T,|I|)$, $\rho_j(g;T)$, $V_{n,\OCV}(q,T)$. 따라서
비등온 발열 계산에서 $V_n,\xi_j,T$ 를 독립 후처리하지 않고 같은 시간 적분 루프에서 갱신한다(P3-3 순환 의존성).
계산 순서:
\begin{enumerate}[label=\arabic*),leftmargin=2.0em]
\item 외부 전류로 $q(t)$ 갱신: $\dd q/\dd t=|I|/Q_\cell$.
\item 현재 $q,T,\{\xi_j\}$ 에서 전하 보존식 (L1) 으로 $V_n$ root-find (P3-3: implicit algebraic — 수치 solver $=$
  Ch1 부록 B 로 닫힘; 논리 공백$\cdot$물리 충돌 아님).
\item $V_n,T$ 에서 $\xi_{\eq,j}$, $\mathcal A_j$, $r_j^\pm$ 또는 $k_j^{\mathrm{phys}}$, 그리고
  $\eta_j=V_{n,\drive}-V_{\eq,j}(\xi_j)$(기본형 $V_{n,\drive}=V_n$) 계산.
\item $\dd\xi_j/\dd t$, $I_j$, $\dot n_j$ 계산. 가역열 option B$\cdot$consistency(식~\eqref{eq:ch4_consistency})
  입력을 위해 $\dd\xi_{\eq,j}/\dd t=(\partial\xi_{\eq,j}/\partial V_n)(\dd V_n/\dd t)+(\partial\xi_{\eq,j}/\partial T)(\dd T/\dd t)$
  (chain rule, $T\cdot V_n$ 의존)도 함께 산출.
\item 비가역 열: transition entropy production / 일반형 $\sum_j n_j^{\eff}I_j\eta_j$(식~\eqref{eq:ch4_qirr_general})
  / transport flux$\times$force(식~\eqref{eq:ch4_transport_ff}).
\item 가역 열: OCV coefficient(식~\eqref{eq:ch4_rev_ocv}) 또는 transition entropy basis(식~\eqref{eq:ch4_rev_xi})
  중 택일(정합 제약 식~\eqref{eq:ch4_consistency}).
\item background heat 사용 시 $\dot Q_\bg^{\prog}$ 와 $\dot Q_\bg^{\coord}$ 분리(식~\eqref{eq:ch4_qbg_split}).
\item 내부 열원 $\dot{\mathcal Q}_{\src,n}$(식~\eqref{eq:ch4_decomp}) 계산.
\item 열수지식(식~\eqref{eq:ch4_balance})으로 $T(t)$ 갱신.
\item 갱신 $T(t)$ 는 위 되먹임 항들에 반영(Chapter 2 순환).
\end{enumerate}
(수치 적분기$\cdot$root-find solver 구현과 수렴 판정은 본 이론 장 범위 밖이며 Ch1 부록 B 로 분리한다; GS-4 의
이론식/피팅 실무 분리.)
\begin{boundbox}
\textbf{되먹임이 ICA tail 에 미치는 영향(Ch2 계승).} 비가역 발열 $\to T\uparrow\to k_j\uparrow$
($L=|I|/(Q_\cell k_j)\downarrow$, Chapter 1 spectrum shift) $\to$ 꼬리 단축; 동시에 ideal $w_j=RT/F\uparrow$ 로
평형 peak 폭 증가($n_j^{\eff}=RT/(Fw_j)$ 변화). 두 효과가 경쟁하므로 비등온 ICA tail 을 단일 부호로 단정하지
않는다(Chapter 1$\cdot$2 의 온도 tail 3계층 분해와 정합). 본 장 비가역 발열 $\sum_j n_j^{\eff}I_j\eta_j$ 가
이 되먹임의 열원 항이다.
\end{boundbox}

% ====================================================================
\subsection{식별성 및 필요한 실험 제약}\label{sec:ch4_ident}
% ====================================================================
물리 우선 모델에서도 식별성은 사라지지 않는다. 해결책은 임의 clipping 이 아니라 독립 실험과 물리적 prior 다.
다음 항들은 강하게 상관된다:
\begin{equation}
\eta_n\leftrightarrow\mathcal A_n\leftrightarrow r_j^\pm\leftrightarrow k_{j,\lim}\leftrightarrow R_{n,\eff}\leftrightarrow\Delta S_j\leftrightarrow h_\bg\leftrightarrow\rho_j(g;T).
\label{eq:ch4_ident}
\end{equation}
(여기서 $k_{j,\lim}$ 은 Chapter 3 의 강구동 mobility floor, 즉 rate-limited plateau 에서의 $k_j^{\mathrm{phys}}$
상한이다.)
\begin{longtable}{p{0.36\textwidth} >{\raggedright\arraybackslash}p{0.52\textwidth}}
\toprule 실험/자료 & 제한 가능한 항 \\ \midrule \endhead
저율 OCV 온도 series & $(\partial V_{n,\OCV}/\partial T)_q$, $\Delta S_j$, $h_\bg$ \\
pulse/current interrupt & immediate polarization $=$ ohmic $+R_\ct$ 합(분리 불가; 일부 $R_n$) \\
EIS/small-signal & $R_\ct$ 단독 분리,\brk transport time constant,\brk film (EIS anchor $\to$ pulse 잔차 순서) \\
calorimetry & $\dot{\mathcal Q}_{\irr}$, $\dot{\mathcal Q}_{\rev}$,\brk heat scale \\
rest relaxation 온도 응답 & $\dot{\mathcal Q}_{\xi,\rlx}$ 후보(가역/비가역 분리) 검증 \\
rate series & rate limitation, $r_j^\pm$, $k_{j,\lim}$, $\rho_j(g;T)$ \\
\bottomrule
\end{longtable}
독립 자료 없이 동시 자유 결정 금지: (1) OCV coefficient vs transition entropy basis(double counting),
(2) $R_n,R_\ct,R_{n,\eff}$, (3) $r_j^\pm,k_{j,\lim},\rho_j(g;T)$, (4) $\Delta S_j,h_\bg$ 와 rest relaxation
heat, (5) transport heat vs hysteresis-like voltage loss.

\textbf{calorimetry(또는 rest-temperature 응답) 부재 시 미식별 클래스(FLAGGED).} 발열 절대 스케일과 가역/비가역
분리를 닫으려면 열량/온도 응답 자료가 필요하다. 이 자료가 없으면 다음 세 항은 \emph{미식별}이다: (a) 발열 절대
스케일(순수 전기 측정은 $\eta_j\cdot I_j$ 의 \emph{상대} 곱만 줄 뿐 절대 [W]를 닫지 못함), (b) 가역/비가역 분리
비율(\S\ref{sec:ch4_rest} flagbox), (c) $R_{n,\eff}$. 이 경우 발열식은 \emph{상대 분해}로만 쓰며, 절대 스케일과
분리 비율은 FLAGGED 로 표기하고 established 로 쓰지 않는다(GS-3).
\begin{practicebox}
초기 탐색서 heat residual / effective heat gain / capped heat correction / apparent $I^2R$ 가 피팅을
안정화할 수 있으나 \emph{물리 모델 기본식이 아니다}. 논문용 해석에서는 transport / charge-transfer / entropy
coefficient / relaxation entropy production / hysteresis loss 로 분해하거나 단순 empirical residual 로 명시하고
주 결론에서 제외한다(GS-4; solver/numerical 구현은 Ch1 부록 B).
\end{practicebox}

% ====================================================================
\subsection{Chapter 4 의 가정 원장 (AL-40$\sim$49)}\label{sec:ch4_al}
% ====================================================================
본 장에서 사용한 가정을 통합 ledger(RB\_AL\_MASTER) Ch4 그룹(AL-40$\sim$49)으로 정리한다. AL-40,41 은 기등재분
(Foundation), AL-42$\sim$49 는 본 장 S2 에서 채운 신규 정의다(Phase A 적발 dangling AL-H 계열의 실제 정의 부여
포함: H1$\to$AL-40/42, H2$\to$AL-41/43). tier = GROUNDED/BOUNDED/FLAGGED.
{\footnotesize
\setlength{\tabcolsep}{3pt}
\begin{longtable}{>{\raggedright\arraybackslash}p{0.045\textwidth} >{\raggedright\arraybackslash}p{0.335\textwidth} >{\raggedright\arraybackslash}p{0.115\textwidth} >{\raggedright\arraybackslash}p{0.225\textwidth} >{\raggedright\arraybackslash}p{0.155\textwidth}}
\toprule AL\# & 가정 진술 & tier & 근거 cite (DOI/ISBN) & 사용 식 \\ \midrule \endhead
AL-40 & entropy production $T\dot S_\irr=\sum_\alpha J_\alpha X_\alpha\ge0$ (flux--force,\brk 2법칙; 합 비음 — Onsager reciprocal 은 본 부등식 근거 아님,\brk AL-47 transport) & GROUNDED & degrootmazur\brk 1962(book),\brk prigogine1967(book) & \eqref{eq:ch4_fluxforce} \\
AL-41 & transition-level entropy production $=$ network/master-equation 형 $(\mathcal J^+{-}\mathcal J^-)\ln(\mathcal J^+/\mathcal J^-)\ge0$ & GROUNDED & schnakenberg\brk 1976 (DOI 10.1103/\brk RevModPhys.48.571) & \eqref{eq:ch4_Jpm},\allowbreak\eqref{eq:ch4_tr_entropy} \\
AL-42 & 전기화학 entropy production $=J_n\mathcal A_n\ge0$; $I_n\eta_n$ 은 부호 convention 종속 magnitude($J_n\mathcal A_n$ 이 기본형);\brk rate-affinity $\mathcal A_j$(중심 $U_j$) $\ne$ heat-force $\eta_j$(국소 평형) & GROUNDED & degrootmazur\brk 1962(book),\brk prigogine1967(book),\brk bazant2013 (DOI 10.1021/\brk ar300145c) & \eqref{eq:ch4_eta_def},\allowbreak\eqref{eq:ch4_JA} \\
AL-43 & transition chemical affinity $\mathcal A_j^{\chem}=RT\ln(\mathcal J_j^+/\mathcal J_j^-)$; 부호 보조정리 $(\mathcal J^+{-}\mathcal J^-)\ln(\mathcal J^+/\mathcal J^-)\ge0$ (두 양수 차×로그비) & GROUNDED & schnakenberg\brk 1976,\brk prigogine1967(book) (DOI 10.1103/\brk RevModPhys.48.571) & \eqref{eq:ch4_tr_entropy},\allowbreak\eqref{eq:ch4_affinity_raw} \\
AL-44 & ★ 미시$=$거시: DB(AL-30)$+n_j^\eff{=}RT/(Fw_j)$(AL-34) $\Rightarrow\mathcal A_j^{\chem}{=}n_j^\eff F\eta_j$, $\dot{\mathcal Q}_{\irr,j}^{\tr}{=}n_j^\eff I_j\eta_j$ (★ $\eta_j{=}V_{n,\drive}{-}V_{\eq,j}$, 기본형 $V_{n,\drive}{=}V_n$ 전제 — CHARTER step2 이중정의 해소; $s_{\phi,j}$ 복원 step1) & GROUNDED$^\dagger$ ($\dagger$ 닫힘은 $C_j{\equiv}0$ 대칭 split $+$ 기본형 $V_{n,\drive}{=}V_n$ 조건부; 강구동 BOUNDED) & schnakenberg\brk 1976,\brk bazant2013,\brk degrootmazur\brk 1962(book) (DOI 10.1103/\brk RevModPhys.48.571;\brk 10.1021/\brk ar300145c) & \eqref{eq:ch4_affinity_eta},\allowbreak\eqref{eq:ch4_micro_macro} \\
AL-45 & ★ 비가역열 일반형 $\dot{\mathcal Q}_\irr{=}\sum_j n_j^\eff I_j\eta_j\ge0$ (2법칙);\brk Ch2 $\sum_j I_j\eta_j$ 는 $n^\eff{=}1$ ideal-width 특수해 (CHARTER step3 일반/특수 위계) & GROUNDED & degrootmazur\brk 1962(book),\brk schnakenberg\brk 1976 (DOI 10.1103/\brk RevModPhys.48.571) & \eqref{eq:ch4_qirr_general},\allowbreak\eqref{eq:ch4_hierarchy} \\
AL-46 & Ch2 3-분해 → Ch4 4-분해 정련: $\dot{\mathcal Q}_{\irr}^{\mathrm{Ch2}}\supset\dot{\mathcal Q}_{\irr}^{\mathrm{Ch4}}+\dot{\mathcal Q}_{\transp}$ (계면반응 EP 와 transport dissipation 분리; 이중계상 금지) & BOUNDED & bernardi1985,\brk rao1997 (DOI 10.1149/\brk 1.2113792;\brk 10.1149/\brk 1.1837884) & \eqref{eq:ch4_decomp},\allowbreak\eqref{eq:ch4_ch2refine} \\
AL-47 & transport/병목 발열 $=$ flux--force $\sum_\ell J_\ell X_\ell$ (선형 near-eq 서 $I^2R_\transp$);\brk transport coupling $L_{\alpha\beta}$ 대칭 $=$ Onsager reciprocal; $R_n\ne R_{n,\eff}\ne R_\ct$ 유지($R_{n,\eff}$ 는 독립 calorimetric 제약 시만) & BOUNDED & degrootmazur\brk 1962(book),\brk onsager1931,\brk thomasnewman\brk 2003,\brk newman,\brk bard\brk faulkner (DOI 10.1103/\brk PhysRev.37.405;\brk 10.1149/\brk 1.1531194;\brk ISBN 978-0-\brk 471-\brk 47756-3;\brk 978-0-\brk 471-\brk 04372-0) & \eqref{eq:ch4_transport_ff},\allowbreak\eqref{eq:ch4_Rsep} \\
AL-48 & 휴지 relaxation heat $\dot{\mathcal Q}_{\xi,\rlx}{=}\dot{\mathcal Q}_{\rev,\xi}^\rlx{+}\dot{\mathcal Q}_{\irr,\xi}^\rlx$; 외부 $I{=}0$ 라도 내부 $I_j{\ne}0$ 의 dissipation $n_j^\eff I_j\eta_j$ 존재; 가역/비가역 정량분리는 독립 calorimetry 필요 (Ch2 AL-29 후보 확정) & BOUNDED & schnakenberg\brk 1976,\brk degrootmazur\brk 1962(book) (DOI 10.1103/\brk RevModPhys.48.571) & \eqref{eq:ch4_rest_split},\allowbreak\eqref{eq:ch4_rest_irr} \\
AL-49 & log singularity($\mathcal J^\pm{\to}0$)는 coarse-graining resolution positive floor (numerical domain guard,\brk heat fitting 아님); $n_j^\eff I_j\eta_j$ 형은 평형서 $I_j{\to}0$ 로 매끄럽게 0 (발산 없음) & BOUNDED & schnakenberg\brk 1976 (DOI 10.1103/\brk RevModPhys.48.571) & \eqref{eq:ch4_tr_entropy} boundbox \\
\bottomrule
\end{longtable}
}
\textbf{AGP(Assumption Grounding Policy).} 모든 본문 가정식은 AL-\# 를 인용한다. 본 장에 FLAGGED \emph{전용}
항목은 없으며(AL-48 의 가역/비가역 정량 분리 비율은 BOUNDED+검증 전제), entropy production$\cdot$flux--force$\cdot$
network thermo 는 GROUNDED, 정련/transport/식별성/domain-guard 는 BOUNDED(유효범위 병기)다. 임의
clip/cap/softplus/step 정의식 0(GS-4; log 발산은 식~\eqref{eq:ch4_tr_entropy} domain guard + 거시형 정칙화).
solver/numerical 구현 0(Ch1 부록 B 분리). \textbf{DOI 정확 귀속(★ Ch3 오귀속 교훈 계승):} degrootmazur1962$\cdot$
prigogine1967 $=$ 교과서(DOI 없음), onsager1931 $=$ 10.1103/PhysRev.37.405, schnakenberg1976 $=$
10.1103/RevModPhys.48.571 — Onsager DOI 를 degrootmazur 책에 붙이는 류의 오귀속 0건.

% ====================================================================
\subsection{Chapter 4 자체 검수표}\label{sec:ch4_selfcheck}
% ====================================================================
{\small
\begin{longtable}{>{\raggedright\arraybackslash}p{0.74\textwidth} >{\raggedright\arraybackslash}p{0.16\textwidth}}
\toprule 검수 항목 & 결과 \\ \midrule \endhead
Chapter 1--3(RB) 를 수정 없이 적층; 전하 보존식이 발열로 대체되지 않음(P3-6) & 통과(\S\ref{sec:ch4_inherit}) \\
발열 항을 fitting residual 아닌 물리적 heat source term 으로 정의 & 통과(\S\ref{sec:ch4_decomp}) \\
★ $\eta_j$ 기준 전위 단일화 $\eta_j{=}V_{n,\drive}{-}V_{\eq,j}$, 기본형 $V_{n,\drive}{=}V_n$ (CHARTER step2 이중정의 해소) & 통과(\S\ref{sec:ch4_eta_single},\brk keybox) \\
비가역 열을 flux--force / entropy production 에서 출발($J_n\mathcal A_n\ge0$ 기본형) & 통과(식~\eqref{eq:ch4_JA}) \\
$I\eta$ 를 부호 convention 종속 magnitude 와 구분(``$I\eta$ 단독'' 회피; $q_\irr{=}|I|\eta$ 양정치) & 통과(\S\ref{sec:ch4_irr},\ref{sec:ch4_general}) \\
transition EP 식에 $Q_{j,\tot}/F$ 포함, $(\mathcal J^+{-}\mathcal J^-)\ln(\cdot)\ge0$ 무비약 2법칙 증명 & 통과(식~\eqref{eq:ch4_tr_entropy}) \\
★ 미시 transition EP $=$ 거시 $n_j^{\eff}I_j\eta_j$ 정합(무비약 5단계; $s_{\phi,j}$ 복원 step1) & 통과(verifybox,식~\eqref{eq:ch4_micro_macro}) \\
★ 미시$=$거시 대입서 기본형 $V_{n,\drive}{=}V_n$ 전제 명시(이중정의 차단 재확인) & 통과(verifybox 3단계, 전제 명시 단락) \\
★ 비가역열 일반형 $\sum_j n_j^{\eff}I_j\eta_j\ge0$ + Ch2 특수형($n^\eff{=}1$) 위계(CHARTER step3) & 통과(식~\eqref{eq:ch4_qirr_general},\eqref{eq:ch4_hierarchy}) \\
ideal width 서 $\to I_j\eta_j$(Ch2), 평형서 $I_j{\to}0$ 로 log 발산 회피(domain guard) & 통과(boundbox,\AL{49}) \\
가역 열을 OCV coefficient 또는 transition entropy basis 에서 출발 & 통과(\S\ref{sec:ch4_revOCV},\ref{sec:ch4_revxi}) \\
★ Bernardi 가역열 $=$ 미시 reaction entropy 합 정합 재확인(reaction $\ne$ activation entropy) & 통과(verifybox \S\ref{sec:ch4_revOCV}) \\
$V_{n,\app}$ 온도 미분을 가역 열로 직접 쓰지 않음(P3-1) & 통과(식~\eqref{eq:ch4_forbid}) \\
OCV 방식과 entropy basis double counting 정합 제약(동일 heat-positive 규약) & 통과(식~\eqref{eq:ch4_consistency}) \\
$Q_\bg$ progress 와 coordinate shift 구분; 등온 근사 한정 & 통과(식~\eqref{eq:ch4_qbg_split},\eqref{eq:ch4_qbg_iso}) \\
$R_n,R_\ct,R_{n,\eff}$ 구분;\brk transport heat 를 flux--force/제한된 $I^2R$ 로만 & 통과(식~\eqref{eq:ch4_Rsep}) \\
휴지 relaxation heat 가역/비가역 확정 분해; 외부 $I{=}0$ 라도 내부 $I_j{\ne}0$ (Ch2 후보 확정) & 통과(\S\ref{sec:ch4_rest}) \\
전기--열 coupling 계산 순서 + 온도 되먹임 명시(P3-3 순환) & 통과(\S\ref{sec:ch4_order}) \\
모든 식 차원$\cdot$부호$\cdot$2법칙$\cdot$극한 검산 명시 & 통과(\S\ref{sec:ch4_tr},\ref{sec:ch4_general} 등) \\
``자명/clearly/쉽게/obviously'' 본문 0건(G-noleap) & 통과 \\
임의 heat correction/capped heat 를 기본식에 넣지 않음(실무 팁 격하;\brk GS-4) & 통과 \\
모든 가정 AL-40$\sim$49 인용;\brk BOUNDED 유효범위 병기;\brk DOI 정확 귀속(오귀속 0) & 통과(\S\ref{sec:ch4_al}) \\
Ch5/Ch1 부록 B 이월 항목 분리(충전 branch 부호 → Ch5) & 통과(\S\ref{sec:ch4_transmit}) \\
\bottomrule
\end{longtable}
}

% ====================================================================
\subsection{Chapter 5 로 전달되는 항목 (수치$\cdot$피팅은 Ch1 부록 B)}\label{sec:ch4_transmit}
% ====================================================================
\begin{enumerate}[label=\arabic*),leftmargin=2.0em]
\item \textbf{Chapter 5}: 충전 branch $\dot{\mathcal Q}_{\rev}$ 부호 convention($s_{\phi,j}^b$ branch 부호
  \emph{유도}); branch-dependent entropy coefficient; hysteresis heat vs polarization heat 분리;
  charge/discharge asymmetric affinity/rate($C_j$ 의 detailed-balance 이탈 $\Rightarrow\xi_{\sstat,j}\ne\xi_{\eq,j}$);
  full-cell voltage gap heat 해석. 본 장 일반형 $\sum_j n_j^{\eff}I_j\eta_j$ 의 부호 양정치는 detailed-balance
  기본형 전제이며, branch 비대칭에서 $\mathcal A_j^{\chem}$ 의 일반화가 Chapter 5 과제다.
\item \textbf{Ch1 부록 B}: $\dd\xi_j/\dd t$ 기반 heat-rate 계산; stiff kinetics$+$thermal ODE 결합 수치 해법;
  직접 $g$-분포 적분서 heat-rate 동시 계산; calorimetry 부재 시 발열 항을 피팅하지 않고 forward prediction
  으로만 두는 제한(Chapter 1$\cdot$2 의 forward-only 원칙).
\end{enumerate}

% ====================================================================
\subsection{Chapter 4 의 결론}\label{sec:ch4_concl}
% ====================================================================
첫째, 비가역 열은 임의 heat correction 이 아니라 flux$\times$force entropy production $T\dot S_\irr=\sum_\alpha
J_\alpha X_\alpha\ge0$(식~\eqref{eq:ch4_fluxforce}), 전기화학적으로 $T\dot S_{\irr,n}=J_n\mathcal A_n\ge0$
(식~\eqref{eq:ch4_JA})에서 출발한다. ``$I\eta$ 단독''(부호 미정렬)은 entropy production 의 일반 정의가 아니다.

둘째, \textbf{미시 transition entropy production $\dot{\mathcal Q}_{\irr,j}^{\tr}=(Q_{j,\tot}/F)RT(\mathcal J_j^+-\mathcal J_j^-)
\ln(\mathcal J_j^+/\mathcal J_j^-)\ge0$ 은 거시 $n_j^{\eff}I_j\eta_j$ 로 정확히 환원}된다(식~\eqref{eq:ch4_micro_macro}).
이 환원은 (i) Chapter 3 detailed balance, (ii) $n_j^{\eff}=RT/(Fw_j)$, (iii) \textbf{기본형 $V_{n,\drive}=V_n$}
(★ $\eta$ 기준 전위 단일화, CHARTER step2)을 전제로 닫히며, $s_{\phi,j}$ 부호 인자가 명시된다(★ CHARTER step1).
이로써 Ch2(거시 flux$\times$overpotential)$\cdot$Ch3(미시 forward/backward rate, $n_j^{\eff}$)$\cdot$Ch4
(entropy production)가 한 표현으로 닫히고, 평형서 $I_j\to0$ 로 발산 없이 0 이 된다.

셋째, \textbf{비가역 발열 일반형은 $\dot{\mathcal Q}_\irr=\sum_j n_j^{\eff}I_j\eta_j\ge0$}(식~\eqref{eq:ch4_qirr_general})
이며, \emph{Chapter 2 의 $\sum_j I_j\eta_j$ 는 그 ideal-width($n_j^{\eff}=1$) 특수해}다(★ CHARTER step3 일반/
특수 위계, 모순 아님). $n_j^{\eff}$ 가중은 임의 피팅 인자가 아니라 detailed-balance 정합이 강제한 물리 인자다.

넷째, 가역 열은 평형 전위 온도계수(Bernardi, 식~\eqref{eq:ch4_rev_ocv}) 또는 transition entropy basis
(식~\eqref{eq:ch4_rev_xi})에서 출발하며, 두 방식을 독립항으로 동시에 자유롭게 더하면 double counting 이다
(정합 제약 식~\eqref{eq:ch4_consistency}). Bernardi 가역열이 미시 reaction entropy $\Delta S_{\rxn,j}$ 의 합과
정합함을 재확인했고(verifybox \S\ref{sec:ch4_revOCV}), reaction entropy 는 activation entropy 와 별개다.

다섯째, Chapter 1 의 $R_n$ 은 apparent voltage-axis coefficient 이므로 heat-generation resistance 가 아니다
($R_n\ne R_\ct\ne R_{n,\eff}$, 식~\eqref{eq:ch4_Rsep}). transport/병목 발열은 flux$\times$force(식~\eqref{eq:ch4_transport_ff})
또는 독립 제약된 $I^2R_{n,\eff}$ 로만 쓰며, 강구동 rate limitation 을 soft cap 으로 절단하지 않는다(GS-4).

여섯째, 휴지 구간에서도 외부 $|I|=0$ 이지만 내부 transition flux $I_j\ne0$ 의 dissipation $\sum_j n_j^{\eff}I_j\eta_j\ge0$
이 남는다(식~\eqref{eq:ch4_rest_irr}; Chapter 2 가 후보로 남긴 relaxation heat 의 확정 분해). 본 장의
$\dot{\mathcal Q}_{\irr},\dot{\mathcal Q}_{\rev},\dot{\mathcal Q}_{\rlx},\dot{\mathcal Q}_{\transp}$ 는 Chapter 5
의 charge/discharge branch 및 hysteresis heat 해석으로 전달된다. 이로써 Chapter 1 의 ``열역학$=$방향(가역/
entropy), 동역학$=$속도(비가역/소산)'' 분업이 entropy-production 수준에서 \emph{정량}되되 앞 장과 충돌하지 않는다.

\subsection*{Chapter 5 로의 전달}
본 장이 확정한 발열 일반형($\dot{\mathcal Q}_\irr=\sum_j n_j^{\eff}I_j\eta_j\ge0$, transition entropy
production, Bernardi 가역열)과 상태변수 위에서: Chapter 5 는 $C_j$ 의 detailed-balance 이탈
($\xi_{\sstat,j}\ne\xi_{\eq,j}$)로 충방전 히스테리시스의 branch 의존 target 과 부호($s_{\phi,j}^b$)를 유도하고,
hysteresis heat 와 polarization heat 를 분리한다. Ch1 부록 B 은 본 장 발열률의 수치 적분과 calorimetry 부재 시
forward-prediction 제한을 다룬다.



% ====================================================================
% ===== Chapter 5 body (원본 graphite_ica_ch5_rebuilt.tex, byte-verbatim; heading 1단계 demote) =====
% ====================================================================
\clearpage
\section{Chapter 5: 충방전 히스테리시스 (충전 부호 유도 $\cdot$ $\Delta V_\obs\ne\Delta V_\hys$)}\label{chap:ch5}

% ====================================================================
\subsection*{서 — 인과 사슬에서 ``히스테리시스''의 자리}
\addcontentsline{toc}{subsection}{서 — 히스테리시스의 자리}
% ====================================================================
Chapter 1 의 keystone 은 ``\textbf{열역학이 방향(평형 target $\xi_{\eq}$)을, 동역학이 속도(mobility $k_j$)를}''
전담한다는 분업이었다(Level A). Chapter 3 은 forward/backward rate 가 detailed balance 를 만족하면 그 분업이
미시적으로 정확함을 보였다($\xi_{\sstat,j}\to\xi_{\eq,j}$). Chapter 4 는 그 위에서 발열 일반형
$\dot{\mathcal Q}_\irr=\sum_j n_j^{\eff}I_j\eta_j\ge0$ 을 닫되, \emph{방전}($s_{\phi,j}=+1$)만 다루고 충전
branch 의 부호 인자 $s_{\phi,j}^b$ 유도는 \textbf{본 장으로 명시적으로 이월}했다(Ch.4 verifybox 4단계).
\textbf{Chapter 5 의 히스테리시스는 바로 Level A 분업이 \emph{부분적으로 무너지는} 영역이다}: 계가 metastable
branch 자유에너지 위에 갇히면 \emph{target 자체가 이력(branch memory)을 획득}해 $\xi_{\tar,j}^b\ne\xi_{\eq,j}$
가 되고, ``방향''이 더 이상 순수 평형 열역학만으로 정해지지 않는다. 즉 히스테리시스는 새 피팅 곡선이 아니라,
\emph{detailed balance 의 이탈}(Ch.3 의 $C_j$ 가 branch 의존)로 인과 사슬 안에서 발생한다.

본 장은 이 확장을 \textbf{여섯 단계}로 푼다. 각 단계는 학부 수준(전기화학 평형, logistic, 미적분 부호·극한,
2법칙 $\dot S_\irr\ge0$)으로 따라갈 수 있게 중간을 생략하지 않는다.
\begin{enumerate}[label=\textbf{(\arabic*)},leftmargin=2.2em]
\item \textbf{충방전을 한 좌표로 묶는다.} signed current $I_s$ 와 signed state coordinate $q_\stt$ 로 방전/충전/
  휴지를 한 시간 적분에 넣되(\S\ref{sec:ch5_signed}), Chapter 1--4 의 $|I|$(방전 magnitude)와 $I_s$(부호 있음)
  를 반드시 구분한다.
\item \textbf{★ 충전 branch 의 부호 인자를 \emph{유도}한다.} 방전 $s_{\phi,j}=+1$ 에 대해 충전
  $s_{\phi,j}^b=-1$ 을 \emph{logistic 지수의 부호 정합}으로부터 유도한다(\S\ref{sec:ch5_chargesign}). 충전서
  $\mathcal A_j$, $\eta_j$, $q_\irr$ 의 부호가 어떻게 뒤집히고, 그럼에도 $n_j^{\eff}I_j\eta_j^b\ge0$(2법칙)이
  \emph{부호 상쇄}로 여전히 성립함을 \emph{유도}로 보인다(주장 금지; Phase A HIGH 해소, CHARTER step1).
\item \textbf{target 이 이력을 얻는다.} true equilibrium $\xi_{\eq,j}$ 와 metastable branch target
  $\xi_{\tar,j}^b$ 를 구분하고(\S\ref{sec:ch5_branch}), branch target 이 Chapter 3 의 nonequilibrium
  stationary $\xi_{\sstat,j}$ 임을 보인다 — 히스테리시스는 mobility($k_j$)가 아니라 \emph{target 이동}이다
  (dreyer2010 many-particle, sethna1993 disorder pinning).
\item \textbf{관측 gap 은 히스테리시스가 아니다.} $\Delta V_\obs\ne\Delta V_\hys$(\S\ref{sec:ch5_apparent}):
  관측 loop 폭에는 polarization($R_n$, Ch.1)$\cdot$kinetic lag$\cdot$transport$\cdot$aging 이 섞이며, true
  thermodynamic hysteresis(평형 분지 차)는 분리 식별된 경우에만 본문 모델에 둔다.
\item \textbf{loop 면적 = 히스테리시스 소산열.} $q_\hys=\oint V\,\dd Q$ 의 분리된 성분을 entropy production
  과 정합시킨다(\S\ref{sec:ch5_hysheat}, Ch.4 의 $\ge0$ 형 계승). loop area 전체를 true hysteresis heat 로
  단정하지 않는다.
\item \textbf{branch 발열$\cdot$휴지$\cdot$full-cell 을 정렬한다.} Chapter 4 의 열원 분해를 branch 로 확장하고
  (\S\ref{sec:ch5_branchheat}--\ref{sec:ch5_rest}), full-cell 관측이 음극 단독과 다름을 명시한다
  (\S\ref{sec:ch5_fullcell}).
\end{enumerate}

\noindent\textbf{한 문장 요지.} \emph{히스테리시스는 mobility 의 변화가 아니라 target 의 이력 의존(Level A
분업의 부분 붕괴)이며, 충전 branch 는 부호 인자 $s_{\phi,j}^b=-1$ 을 logistic 지수에서 \emph{유도}함으로써
방전과 한 구조에 닫히고, 그 부호 뒤집힘 속에서도 비가역 발열 $n_j^{\eff}I_j\eta_j^b\ge0$ 은 부호 상쇄로
2법칙을 보존한다.} 본 장은 Chapter 1--4 기본식을 수정하지 않고 내부 state 로 방전 기준 $\xi_j$ 를 유지하며,
수치 해법$\cdot$validation pipeline 은 Ch1 부록 B 로 분리한다.

\begin{groundingbox}
\textbf{지배 표준(GS, Ch1--4 계승) 및 히스테리시스 근거 원칙.} \textbf{(GS-1)} 모든 식은 물리적 의미를 prose 로
먼저 말한 뒤 수식화하고 중간 단계를 생략하지 않는다(학부 무비약). \textbf{(GS-2)} 결론은 출판 수준 엄밀성.
\textbf{(GS-3)} 모든 가정은 통합 Assumption Ledger(\AL{\#}: 논문/교과서 근거 + 검증 DOI)를 인용하며, 근거가
없으면 FLAGGED 로 표기하고 established 로 쓰지 않는다. \textbf{(GS-4)} 임의 clip/cap/softplus/step-function
정의식과 근거 없는 수치 임계값을 물리 모델로 쓰지 않는다. \textbf{히스테리시스 근거 원칙.} 충전 표기 보조
변수는 허용하나 독립 state 를 임의로 늘리지 않는다; voltage gap 을 곧바로 hysteresis 로 등치하지 않는다
($\Delta V_\obs\ne\Delta V_\hys$); branch target 은 metastable free-energy/nucleation/domain memory 로 해석
가능할 때만 본문 모델; 근거 없는 empirical branch offset / arbitrary hysteresis curve 는 실무 팁으로 격하;
hysteresis heat 는 polarization/transport/kinetic lag 와 분리된 path-dependent loss 가 식별된 경우에만 후보.
\textbf{충전 부호 인자 $s_{\phi,j}^b$ 는 주장하지 않고 logistic 지수에서 유도한다}(CHARTER step1, Phase A
HIGH 해소).
\end{groundingbox}

% ====================================================================
\subsection{Chapter 1--4 에서 전달받는 고정 기준식}\label{sec:ch5_inherit}
% ====================================================================
본 장은 Chapter 1$\cdot$2$\cdot$3$\cdot$4(RB 재구성본)의 다음을 \textbf{수정 없이} 입력으로 받는다(프로젝트
검수 P3-6). 각 식 번호는 앞 장의 boxed 결과를 가리키며, 본 장은 그 위에 충방전 branch 계층을 \emph{적층}한다
(앞 장 본문 기본식을 충전으로 \emph{고쳐 쓰지} 않는다).
\begin{linkbox}
\textbf{(L1) 전하 보존식(중심 상태식).} $Q_\cell\,q=Q_\bg(V_n,T)+\sum_{j=1}^{N_p}Q_{j,\tot}\xi_j$(방전 기준
축약형). $V_n$ 은 그 해(외부 입력 아님). 충전 포함 시 signed coordinate 로 일반화한다(\S\ref{sec:ch5_signed}).\\
\textbf{(L2) forward/backward kinetics(Ch3).} $\dfrac{\dd\xi_j}{\dd t}=r_j^+(1-\xi_j)-r_j^-\xi_j$; detailed
balance $r_j^+/r_j^-=\xi_{\eq,j}/(1-\xi_{\eq,j})$, ln 형 $\ln(r_j^+/r_j^-)=(V_n-U_j)/w_j$; nonequilibrium
stationary $\xi_{\sstat,j}=r_j^+/(r_j^++r_j^-)$; rate ratio 일반형 $\ln(r_j^+/r_j^-)=C_j(\xi_j,T)+\mathcal A_j/RT$.
\textbf{branch 비대칭은 $\xi_{\sstat,j}\ne\xi_{\eq,j}$ 인 영역}이다(\S\ref{sec:ch5_branch}).\\
\textbf{(L3) 구동력$\cdot$과전압$\cdot$dissipation(Ch3$\cdot$Ch4).} 중심 affinity $\mathcal A_j=s_{\phi,j}n_j^{\eff}
F(V_{n,\drive}-U_j)$($s_{\phi,j}=+1$ 방전); heat-force(국소 평형 이탈) $\eta_j=s_{\phi,j}[V_{n,\drive}-V_{\eq,j}(\xi_j)]$,
$V_{\eq,j}(\xi_j)=U_j+w_j\ln[\xi_j/(1-\xi_j)]$; 비가역 열 $\dot{\mathcal Q}_{\irr,j}=n_j^{\eff}I_j\eta_j\ge0$
(미시$=$거시, Ch.4); $n_j^{\eff}=RT/(Fw_j)$.\\
\textbf{(L4) keystone 한정(Ch1$\cdot$Ch3).} Level A $=$ Level B 의 detailed-balance \emph{극한}이며 일치는
\emph{정상 target} $\xi_{\sstat,j}\to\xi_{\eq,j}$ \emph{한정}이다(무조건 ``$A=B$'' 아님). 그 조건은 (i) $C_j$
의 $V_n$$\cdot$$\xi_j$-무의존(대칭 split), (ii) $V_{n,\drive}=V_n$. \emph{이 둘 중 하나가 깨지면 branch}
($\xi_{\sstat,j}\ne\xi_{\eq,j}$)이며 — \textbf{바로 그 깨짐을 본 장이 다룬다}.\\
\textbf{(L5) 열원 분해(Ch4).} $\dot{\mathcal Q}_{\src,n}=\dot{\mathcal Q}_{\irr,n}+\dot{\mathcal Q}_{\rev,n}
+\dot{\mathcal Q}_{\rlx,n}+\dot{\mathcal Q}_{\transp,n}$; 비가역 $\ge0$(2법칙), 가역 부호가변; $R_n\ne R_\ct\ne R_{n,\eff}$.
본 장서 branch 항 $\dot{\mathcal Q}_{\hys,n}$ 을 추가한다(\S\ref{sec:ch5_branchheat}).\\
\textbf{(L6) 네 전위 구분.} 내부 $V_n$(보존식 해), apparent $V_{n,\app}=V_n+s_I|I|R_n$($s_I=\mathrm{sgn}(I_s)$,
$+1$ 방전$\cdot$$-1$ 충전; 분극 포함, 평형 전위 아님), 구동 $V_{n,\drive}$(기본 $=V_n$), 평형
$V_{n,\OCV}$(평형 보존식 해).
\end{linkbox}
임의 empirical branch offset / arbitrary hysteresis curve 를 기본 물리식으로 쓰지 않는다(GS-4); 충전 부호는
유도(\S\ref{sec:ch5_chargesign}), branch target 은 metastable free-energy 근거 시만 본문 모델.

\textbf{신규 기호(Ch1--4 위에 추가).} 본 장에서만 새로 도입하는 기호를 모은다(앞 장 기호는 (L1)--(L6) 과 동일물).
\begin{longtable}{p{0.16\textwidth} p{0.13\textwidth} >{\raggedright\arraybackslash}p{0.63\textwidth}}
\toprule 기호 & 단위 & 의미 \\ \midrule \endhead
$I_s$ & A & signed external current ($>0$ 방전, $<0$ 충전, $=0$ 휴지); $|I_s|$ 가 Ch1--4 의 $|I|$ \\
$s_I$ & --- & 전류 방향 부호 $=\mathrm{sgn}(I_s)$ ($+1$ 방전, $-1$ 충전);\brk apparent 전위 식~\eqref{eq:ch5_vapp_branch} 의 부호 인자 \\
$b(t)$ & --- & branch index $\in\{\dis,\chg,\rst\}$ (방전/충전/휴지) \\
$s_{\phi,j}^b$ & --- & branch 부호 인자 ($s_{\phi,j}^\dis=+1$, $s_{\phi,j}^\chg=-1$; \S\ref{sec:ch5_chargesign} 에서 \emph{유도}) \\
$q_\stt$ & --- & signed internal state coordinate $\dd q_\stt/\dd t=I_s/Q_\cell$ (방전 증가, 충전 감소) \\
$\zeta_j$ & --- & 충전 직관용 보조 $=1-\xi_j$ (\textbf{독립 state 아님}; 종속 변수) \\
$\xi_{\tar,j}^b$ & --- & branch target progress (metastable);\brk branch-local Ch3 $\xi_{\sstat,j}$ 형식(본 장 \S\ref{sec:ch5_branch} 구성) \\
$U_j^b,\ w_j^b$ & V & branch center potential $\cdot$ branch 전이폭 (metastable free-energy 기원) \\
$G_j^b$ & J/mol & transition $j$ 의 branch(metastable) free-energy \\
$h_j(t)$ & --- & hysteresis state $\in[-1,1]$ ($+1$ 방전 memory, $-1$ 충전 memory, $0$ relaxed) \\
$h_{\tar,j}^b$ & --- & 현재 branch 의 memory target $=s_{\phi,j}^b$ ($b=\dis\to+1$, $b=\chg\to-1$; 새 자유 파라미터 아님) \\
$\Delta U_j^\hys$ & V & 방전--충전 branch center 차(signed);\brk magnitude 강제 시 부호 인자 $s_{h,j}\in\{+1,-1\}$ 로 $\Delta U_j^\hys=s_{h,j}|\Delta U_j^\hys|$ \\
$\lambda_j^b$ & 1/s & hysteresis state $h_j$ relaxation rate (domain rearrangement/depinning time scale);\brk rest 구간 특수화 $\lambda_j^{b=\rst}$ (식~\eqref{eq:ch5_h_rest_relax}) \\
$k_j^\rst$ & 1/s & rest 구간 $\xi_j$ 완화율(식~\eqref{eq:ch5_rest_eq});\brk Ch1 mobility $k_j$ 와 별개의 휴지 relaxation 상수 \\
$n_j^{\eff,b}$ & --- & branch 유효 전자수 $\equiv RT/(Fw_j^b)$ (Ch1--4 $n_j^\eff{=}RT/(Fw_j)$ 의 branch 폭 $w_j^b$ 판) \\
$\eta_{j,\prog}^b$ & V & branch-local 진행 과전압 $=s_{\phi,j}^b[V_{n,\drive}-V_{\tar,j}^b(\xi_j)]$ \\
$\Delta V_\obs,\Delta V_\hys$ & V & 관측 voltage gap $\cdot$ true thermodynamic hysteresis gap \\
$\Delta V_\pol,\Delta V_\kin,\Delta V_\transp$ & V & 분극$\cdot$kinetic lag$\cdot$transport 기여 gap \\
$\Delta V_{\mathrm{aging}}$ & V & aging/side reaction 누적 변화 gap (비가역 누적, $|I_s|\to0$ 외삽으로 미분리) \\
$E_{\mathrm{loop}}$ & J & 전압--용량 loop 면적 $\oint V\,\dd Q$ (소산 + 분극 + lag 혼합) \\
$\dot{\mathcal Q}_{\hys,n}^b$ & W & branch path-dependent hysteresis 발열률 \\
$V_\cell^b,\ V_p^b$ & V & branch full-cell 전압 $\cdot$ 양극 전위 \\
$\eta_\loss^b$ & V & branch-dependent signed loss/overpotential aggregate \\
$s^{b,\conv}$ & --- & heat-positive bookkeeping factor $\in\{+1,-1\}$ (피팅 X; 단일 규약, \emph{branch 무의존} — branch 부호 반전은 $\dd\xi_j/\dd t$ 가 전담; 식~\eqref{eq:ch5_branch_rev},\allowbreak\eqref{eq:ch5_loop_area}) \\
\bottomrule
\end{longtable}
$F=96485\ \mathrm{C/mol}$, $R=8.314\ \mathrm{J/(mol\,K)}$, $\mathrm{J/C}=\mathrm V$. 열률 dot 표기는 시간율 [W].

% ====================================================================
\subsection{Signed current $\cdot$ branch index $\cdot$ signed state coordinate}\label{sec:ch5_signed}
% ====================================================================
\textbf{물리.} Chapter 1--4 는 방전 한 방향만 다뤘으므로 전류 크기 $|I|$ 만으로 충분했다. 충방전을 한 구조에
넣으려면 \emph{부호 있는} 전류가 필요하다. signed current 를 도입한다:
\begin{equation}
I_s>0:\ \text{방전(discharge)},\quad I_s<0:\ \text{충전(charge)},\quad I_s=0:\ \text{휴지(rest)};\qquad |I_s|=\text{크기}.
\label{eq:ch5_signed_current}
\end{equation}
Chapter 1--4 의 $|I|$ 는 \emph{방전 기준 magnitude} 였으므로, 앞 장 식에 $I_s$ 를 그대로 넣지 않고 $|I_s|$ 로
대응시킨다($|I|\equiv|I_s|$). branch index $b(t)\in\{\dis,\chg,\rst\}$ 를 부호로 정한다:
$I_s>0\Rightarrow b=\dis$, $I_s<0\Rightarrow b=\chg$, $I_s=0\Rightarrow b=\rst$. 내부 상태 좌표는 \emph{부호
있는} signed coordinate 다:
\begin{equation}
\boxed{\;\frac{\dd q_\stt}{\dd t}=\frac{I_s}{Q_\cell}\;}\qquad(\text{방전서 증가},\ \text{충전서 감소}).
\label{eq:ch5_qstate}
\end{equation}
\textbf{차원 검산.} $I_s[\mathrm A=\mathrm{C/s}]/Q_\cell[\mathrm C]=\mathrm{s^{-1}}$, 좌변 $\dd q_\stt/\dd t$ 도
$\mathrm{s^{-1}}$($q_\stt$ 무차원) — 정합. \textbf{부호 검산.} $I_s>0$(방전)서 $\dd q_\stt/\dd t>0$,
$I_s<0$(충전)서 $\dd q_\stt/\dd t<0$ — 방전이 진행 좌표를 키우고 충전이 되돌린다는 직관과 일치. $Q_\cell$ 은
기준 용량이며 irreversible capacity loss(LLI/LAM)/electrode balancing error 를 이 식에 임의 흡수하지 않는다
(필요 시 별도 capacity-balancing 변수; \AL{52}). 전하 보존식은 적분 상수 모호성을 피하려 \emph{기준점 포함
상대형}이 안전하다:
\begin{equation}
Q_\cell(q_\stt-q_{\stt,\mathrm{ref}})=\big[Q_\bg(V_n,T)-Q_\bg(V_{n,\mathrm{ref}},T_{\mathrm{ref}})\big]
+\sum_j Q_{j,\tot}(\xi_j-\xi_{j,\mathrm{ref}}).
\label{eq:ch5_relbalance}
\end{equation}
\textbf{검산(방전 극한 환원).} 방전 단독($I_s=|I_s|>0$, 기준점을 방전 시작으로)에서, \emph{이상 쿨롱효율}
($\eta_C=1$)$\cdot$$\dd Q_\bg=0$ 극한이면 $q_\stt-q_{\stt,\mathrm{ref}}\to q$ 가 되어 식~\eqref{eq:ch5_relbalance}
가 Chapter 1 의 전하 보존식 (L1) 로 환원된다(L1 환원; 앞 장과 충돌 없음). \emph{일반 경우}에는 $q_\stt$(외부
전류 적분)와 $q$(보존식 해)가 부반응(LLI/LAM/background)만큼 어긋나며, 이 차이를 식~\eqref{eq:ch5_qstate} 에
임의 흡수하지 않는다(위 $Q_\cell$ 비흡수 원칙과 일관, \AL{52}). 실험 plotting 용량은 내부 좌표와 \emph{별개}로 둔다: $Q_\dis=\int_\dis|I_s|\,\dd t$,
$Q_\chg=\int_\chg|I_s|\,\dd t$ (ICA/DVA overlay 용; 내부 상태방정식은 $q_\stt$ 를 따른다).

\subsubsection{내부 state $\xi_j$ 유지와 보조 변수 $\zeta_j$}\label{sec:ch5_xikeep}
\textbf{물리.} 기본 state 는 \emph{방전 기준} $\xi_j$ 하나다($\xi_j=0$ 미진행, $\xi_j=1$ 완료). 방전서 대체로
$\dd\xi_j/\dd t>0$, 충전서 $\dd\xi_j/\dd t<0$ 이다(같은 $\xi_j$ 가 양방향으로 움직인다). 충전 방향을 직관적으로
표기할 때 $\zeta_j\equiv1-\xi_j$ 를 쓸 수 있으나, 이는 \textbf{독립 state 가 아니라 종속 보조 변수}다
($\dd\zeta_j/\dd t=-\dd\xi_j/\dd t$). branch 전환 시 $\xi_j$ 를 재초기화하지 않으며(\S\ref{sec:ch5_rest}),
수치적분$\cdot$전하보존$\cdot$발열의 기본 입력은 항상 $\xi_j,\dd\xi_j/\dd t$ 다(\AL{52}; 독립 state 임의
증식 금지 — GS 근거 원칙).

% ====================================================================
\subsection{★ 충전 branch 부호 인자 $s_{\phi,j}^b$ 의 유도 (CHARTER step1, Ch4 이월 수령)}\label{sec:ch5_chargesign}
% ====================================================================
이 절이 본 장의 \textbf{핵심}이며, Chapter 4 가 명시적으로 본 장에 이월한 책임(verifybox 4단계: ``충전 branch
$s_{\phi,j}=-1$ 에서 logistic 지수 부호로부터의 \emph{유도}는 Chapter 5 로 이월'')을 수령한다. \textbf{목표}:
방전 $s_{\phi,j}^\dis=+1$ 에 대해 충전 $s_{\phi,j}^\chg=-1$ 을 \emph{logistic 지수의 부호 정합}으로부터
\emph{유도}하고(주장 금지), 충전서 $\mathcal A_j$, $\eta_j$, $q_\irr$ 의 부호가 어떻게 뒤집히는지, 그럼에도
$n_j^{\eff}I_j\eta_j^b\ge0$(2법칙)이 \emph{부호 상쇄}로 여전히 성립함을 보인다.

\subsubsection{logistic 지수의 부호 정합에서 $s_{\phi,j}^b$ 가 나온다}\label{sec:ch5_sign_derive}
\textbf{출발(Ch1 평형식, 무비약).} Chapter 1 의 전기화학 평형 조건(식 (iii), \S 평형 진행률)은 화학퍼텐셜이
전극 전위와 균형을 이루는 것이었다: $\mu(\xi_{\eq,j})=s_{\phi,j}F(V_n-U_j^0)$. ideal 극한에서 이는
$RT\ln[\xi_{\eq,j}/(1-\xi_{\eq,j})]=s_{\phi,j}F(V_n-U_j)$ 였다(Ch1 식~eq:eqbal\_ideal). 여기서 부호 인자
$s_{\phi,j}$ 는 \emph{선택한 진행 방향의 전이 구동력 부호}를 정하는 양으로, Chapter 1 은 방전을 기본으로
$s_{\phi,j}=+1$ 로 두었다. 본 절은 \emph{충전 branch 에서 이 부호가 무엇이어야 하는지}를 동일 평형식의 부호
정합으로 \emph{역으로 묻는다}.

\textbf{물리적 요구(branch 진행 방향).} ``$\mathcal A_j>0\Leftrightarrow$ 선택한 branch 의 진행이 촉진''이라는
규약(Ch3 식~eq:ch3\_Aj 의 $s_{\phi,j}$ 선택 규칙: ``$\mathcal A_j>0$ 이면 그 방향 진행 촉진'')을 충전에 그대로
적용한다. 방전 branch 의 진행은 $\xi_j$ \emph{증가}($\dd\xi_j/\dd t>0$, 미점유$\to$점유), 충전 branch 의 진행은
$\xi_j$ \emph{감소}($\dd\xi_j/\dd t<0$, 점유$\to$미점유)다(\S\ref{sec:ch5_xikeep}). 따라서 \emph{같은 부호 규약}이
방전과 충전에서 서로 \emph{반대 방향}의 $\xi_j$ 변화를 촉진해야 한다.

\textbf{유도(부호 정합, 한 줄씩).} 평형식 $RT\ln[\xi_{\eq,j}/(1-\xi_{\eq,j})]=s_{\phi,j}^b F(V_n-U_j^b)$ 를
$\xi_{\eq,j}$ 에 대해 풀면(Chapter 1 의 (iv) 단계와 동일 대수) logistic
\begin{equation}
\xi_{\eq,j}^b(V_n,T)=\frac{1}{1+\exp\!\big[-s_{\phi,j}^b\,(V_n-U_j^b)/w_j^b\big]}.
\label{eq:ch5_logistic_branch}
\end{equation}
이 곡선의 $V_n$-기울기는 Chapter 1 식~eq:dxidV 와 같은 logistic 항등식으로
\begin{equation}
\frac{\partial\xi_{\eq,j}^b}{\partial V_n}=\frac{s_{\phi,j}^b\,\xi_{\eq,j}^b(1-\xi_{\eq,j}^b)}{w_j^b}.
\label{eq:ch5_dxidV_branch}
\end{equation}
$\xi_{\eq,j}^b(1-\xi_{\eq,j}^b)>0$, $w_j^b>0$ 이므로 \emph{이 기울기의 부호는 정확히 $s_{\phi,j}^b$ 의 부호}다.
이제 두 branch 의 물리적 진행 방향을 부호로 강제한다.
\begin{enumerate}[label=(\roman*),leftmargin=2.2em]
\item \textbf{방전 branch.} 방전은 $V_n$ 이 평형 $V_{\eq,j}$ 보다 \emph{위}로($V_n>U_j$) 구동될수록 진행(점유
  증가)이 촉진되며 $\xi_j$ 가 증가한다. Chapter 1 의 규약대로 $V_n>U_j$ 에서 $\xi_{\eq,j}>1/2$(점유 우세)가 되려면
  식~\eqref{eq:ch5_logistic_branch} 지수 안 $-s_{\phi,j}^b(V_n-U_j)/w_j<0$, 즉 $s_{\phi,j}^b(V_n-U_j)>0$ 이어야
  하고 $V_n>U_j$ 에서 이는 $s_{\phi,j}^\dis=+1$ 을 강제한다. (Chapter 1$\sim$4 가 쓴 값과 일치 — 환원 검산.)
\item \textbf{충전 branch.} 충전은 그 \emph{역과정}이다: 전위가 \emph{반대} 방향으로 구동될 때(리튬화/충전
  방향) 선택한 충전 transition 의 진행($\xi_j$ \emph{감소})이 촉진돼야 한다. ``$\mathcal A_j>0\Leftrightarrow$
  충전 방향 진행 촉진'' 규약을 만족하려면, 같은 전위 편차 $(V_n-U_j^b)$ 에 대해 \emph{방전과 반대 부호}의
  구동력이 충전을 밀어야 한다. 즉 충전 branch 의 진행 affinity 가 방전과 부호가 반대가 되려면 지수 안 부호
  인자가 뒤집혀야 하고, 이는 곧
\end{enumerate}
\begin{equation}
\boxed{\;s_{\phi,j}^\dis=+1,\qquad s_{\phi,j}^\chg=-1\;.}
\label{eq:ch5_schg}
\end{equation}
\textbf{이 부호는 \emph{유도}된 것이지 가정된 것이 아니다}: 식~\eqref{eq:ch5_dxidV_branch} 의 기울기 부호가
branch 진행 방향(방전 $\dd\xi_j/\dd t>0$ vs 충전 $\dd\xi_j/\dd t<0$)을 결정하고, ``$\mathcal A_j>0\Leftrightarrow$
진행 촉진''이라는 \emph{단 하나의 규약}을 두 branch 에 일관 적용하면 충전에서 $s_{\phi,j}^\chg=-1$ 이
\emph{필연}으로 따라온다. 방전 부호를 따로 가정하지 않아도, 같은 평형식의 부호 정합이 방전($+1$)과 충전($-1$)을
\emph{동시에} 고정한다(\AL{53}).

\begin{boundbox}
\textbf{$\xi_{\eq,j}^b$ vs metastable $\xi_{\tar,j}^b$ 의 구분(여기서는 부호만).} 식~\eqref{eq:ch5_logistic_branch}
는 \emph{branch 부호 인자의 유도}를 위해 logistic 형을 쓴 것이며, 충전/방전 평형 곡선이 \emph{실제로 다른지}
($U_j^\dis\ne U_j^\chg$, true thermodynamic hysteresis)는 \S\ref{sec:ch5_branch} 의 별개 물리(metastable
branch)다. 본 절의 결론(식~\eqref{eq:ch5_schg})은 \emph{부호 정합}이며, $U_j^b,w_j^b$ 의 branch 차이를
가정하지 않는다(그것은 \S\ref{sec:ch5_branch} 의 metastability 가 정한다).
\end{boundbox}

\subsubsection{충전서 $\mathcal A_j$, $\eta_j$, $q_\irr$ 의 부호 뒤집힘 (무비약)}\label{sec:ch5_sign_flip}
\textbf{구동력 $\mathcal A_j$.} Chapter 3 명시형 $\mathcal A_j^b=s_{\phi,j}^b n_j^{\eff}F(V_{n,\drive}-U_j^b)$
에 식~\eqref{eq:ch5_schg} 를 넣으면, \emph{같은} 전위 편차 $(V_{n,\drive}-U_j^b)$ 에 대해
\begin{equation}
\mathcal A_j^\dis=+\,n_j^{\eff}F(V_{n,\drive}-U_j),\qquad
\mathcal A_j^\chg=-\,n_j^{\eff}F(V_{n,\drive}-U_j^\chg),
\label{eq:ch5_A_branch}
\end{equation}
즉 부호 인자가 전체 affinity 의 부호를 뒤집는다(\emph{표기}: 부호 유도 단계서 방전 center 는 $U_j^\dis\equiv U_j$
로 적고, center 의 branch 분기 $U_j^\dis\ne U_j^\chg$ 자체는 가정하지 않으며 \S\ref{sec:ch5_branch} metastability 로
이연한다 — 여기 $U_j^\chg$ 는 충전 branch center 의 \emph{자리}만 표시). \textbf{heat-force $\eta_j$.} 마찬가지로 Chapter 4 의
heat-force $\eta_j^b=s_{\phi,j}^b[V_{n,\drive}-V_{\eq,j}^b(\xi_j)]$ 도 충전서 부호가 뒤집힌다:
\begin{equation}
\eta_j^\dis=+\,[V_{n,\drive}-V_{\eq,j}(\xi_j)],\qquad
\eta_j^\chg=-\,[V_{n,\drive}-V_{\eq,j}^\chg(\xi_j)].
\label{eq:ch5_eta_branch}
\end{equation}
\textbf{전류 $I_j$ 의 부호.} transition current $I_j=Q_{j,\tot}\,\dd\xi_j/\dd t$ 는 $Q_{j,\tot}>0$ 이므로
$\dd\xi_j/\dd t$ 의 부호를 따른다: 방전서 $\dd\xi_j/\dd t>0\Rightarrow I_j>0$, 충전서 $\dd\xi_j/\dd t<0\Rightarrow
I_j<0$. \textbf{비가역 발열 $q_\irr$.} $q_\irr=|I|\eta\ge0$ 형(부호 양정치화)은 Chapter 2 에서 도입되어
Chapter 4 의 미시$=$거시 결과로 닫힌 것이며(RB\_CHARTER step1 이 본 장에 그 충전 branch 근거를 이월), 본 장은
이를 branch 부호 정합으로 \emph{유도}한다 — 단순 ``$I_j\eta_j$''를 그대로 쓰면 충전서 부호가 어떻게 되는지를 먼저
따져야 한다(아래 verifybox).

\begin{verifybox}
\textbf{★ 충전서도 $n_j^{\eff}I_j\eta_j^b\ge0$ 이 \emph{부호 상쇄}로 성립함 (CHARTER step1; 주장 아닌 유도).}
충전 branch($s_{\phi,j}^\chg=-1$)에서 비가역 발열 항의 부호를 \emph{대수적으로} 따라간다. 충전이 진행 중이면
계가 충전 target 을 향해 가므로, 현재 $\xi_j$ 가 그 충전 평형 $\xi_{\eq,j}^\chg$ 보다 \emph{크다}
($\xi_j>\xi_{\eq,j}^\chg$, 아직 덜 빠진 상태). 두 인자의 부호를 각각 본다.

\emph{(1) $I_j$ 의 부호.} 완화식 $\dd\xi_j/\dd t=k_j^{\mathrm{phys}}(\xi_{\eq,j}^\chg-\xi_j)$ 에서
$\xi_j>\xi_{\eq,j}^\chg\Rightarrow\xi_{\eq,j}^\chg-\xi_j<0\Rightarrow\dd\xi_j/\dd t<0\Rightarrow I_j=Q_{j,\tot}
\,\dd\xi_j/\dd t<0$. (충전서 전류가 음 — 식~\eqref{eq:ch5_signed_current} 의 $I_s<0$ 과 정합.)

\emph{(2) $\eta_j^\chg$ 의 부호.} 충전 branch 의 국소 평형 전위는 branch logistic 의 \emph{역함수}이므로 부호 인자가
로그 항에 실린 $V_{\eq,j}^\chg(\xi_j)=U_j^\chg-w_j^\chg\ln[\xi_j/(1-\xi_j)]$ 다($s_{\phi}^\chg=-1$ → $\xi_j$ 에
\emph{감소}형; \S\ref{sec:ch5_kinetics} 의 $V_{\tar}^b$ 정의와 동일). 따라서 $\xi_j>\xi_{\eq,j}^\chg$ 이면 감소형에
의해 $V_{\eq,j}^\chg(\xi_j)<V_{\eq,j}^\chg(\xi_{\eq,j}^\chg)=V_{n,\drive}$, 즉 $[V_{n,\drive}-V_{\eq,j}^\chg(\xi_j)]>0$.
여기에 충전 부호 인자를 곱하면 $\eta_j^\chg=s_{\phi,j}^\chg[V_{n,\drive}-V_{\eq,j}^\chg(\xi_j)]=(-1)\cdot[\,>0\,]<0$.

\emph{(3) 곱의 부호(상쇄).} 따라서 $I_j<0$ \emph{이고} $\eta_j^\chg<0$ — \textbf{두 인자가 \emph{같은}(둘 다
음) 부호}라 그 곱은 양이다:
\begin{equation}
\boxed{\;n_j^{\eff}\,I_j\,\eta_j^\chg=\underbrace{n_j^{\eff}}_{>0}\,\underbrace{(<0)}_{I_j}\,\underbrace{(<0)}_{\eta_j^\chg}\;>\;0\;.}
\label{eq:ch5_qirr_charge}
\end{equation}
\textbf{핵심(부호 상쇄).} 충전이 \emph{두 인자의 부호를 동시에} 뒤집으므로($I_j$ 는 $\dd\xi_j/\dd t<0$ 에서,
$\eta_j^\chg$ 는 $s_{\phi,j}^\chg=-1$ \emph{와} $V_{\eq}^\chg$ 감소형에서) 그 곱은 \emph{여전히 양}이다 — 음$\times$음$=$양.
\emph{더 robust 한 근거(부호 우연이 아님)}: 두 인자는 독립이 아니라 둘 다 같은 편차 $(\xi_{\eq,j}^b-\xi_j)$ 에 종속한다
— $I_j\propto(\xi_{\eq,j}^b-\xi_j)$ 이고, $\eta_j^b$ 도 $V_{\eq}^b(\xi_j)$ 의 단조성을 통해 같은 편차와 동부호이므로
($s_\phi^b$ 가 $V_{\eq}^b$ 의 단조 방향과 $\eta$ 의 부호 인자에 \emph{같이} 들어가 짝이 맞는다), $I_j\eta_j^b\ge0$ 은
flux 와 그 공액력(conjugate force)의 conjugacy 의 귀결이다. 이것이
``$s_{\phi,j}^b$ 의 부호 상쇄로 $n_j^{\eff}I_j\eta_j^b\ge0$ 이 성립''의 \emph{유도}다. 방전($+1$,
식~\eqref{eq:ch5_eta_branch} 첫 식)에서는 $\xi_j<\xi_{\eq,j}\Rightarrow I_j>0$ \emph{이고} $\eta_j^\dis>0$
이라 양$\times$양$=$양(Ch.4 verifybox 의 항별 부호 논증과 동일). \textbf{즉 양 branch 모두 $I_j$ 와 $\eta_j^b$
가 같은 부호}이며, 평형서 $\eta_j^b\to0,\ I_j\to0$ 동시 소멸로 매끄럽게 0 이 된다(발산 없음). 이로써
충전에서도 2법칙 $\sum_j n_j^{\eff}I_j\eta_j^b\ge0$ 이 \emph{유도}로 보존된다(Chapter 4 일반형의 충전 확장;
\AL{54}). 단순 ``$I_j\eta_j$''가 아니라 \emph{부호 정렬이 유도로 보장된} 양정치 합이다(RB\_CHARTER step1:
$q_\irr=|I|\eta\ge0$ 형의 충전 branch 근거).
\end{verifybox}

\textbf{가역 열의 부호가변성(대비).} 비가역 발열이 양 branch 에서 $\ge0$ 인 것과 달리, 가역 열
$q_{\rev,j}^b=I_j\Pi_j$($\Pi_j=T\partial U_j^b/\partial T$, Ch.4 식 (L6))는 $I_j$ 의 부호(방전 $+$, 충전 $-$)와
$\partial U_j^b/\partial T$ 의 부호에 따라 \emph{흡열/발열 모두} 가능하다. 충전과 방전에서 $I_j$ 부호가
뒤집히므로 \emph{같은} entropy 계수 $\partial U_j^b/\partial T$ 라도 가역 열의 부호가 branch 마다 반대로 나온다
— 이것이 충방전 발열 비대칭의 한 물리 원천이며, 비가역(항상 $\ge0$)과 혼동하면 안 된다(\S\ref{sec:ch5_branchheat};
\AL{54}).

\begin{boundbox}
\textbf{$s_{\phi,j}^b$ 유도의 유효범위(detailed-balance 기본형).} 위 verifybox 의 부호 상쇄는 \emph{detailed
balance 기본형}($C_j$ 의 $V_n$$\cdot$$\xi_j$-무의존, $V_{n,\drive}=V_n$; Ch.3 keystone 한정, L4)에서 닫힌다.
강구동 branch 비대칭($C_j$ 의 detailed-balance 이탈, \S\ref{sec:ch5_branch})에서는 $\xi_{\tar,j}^b\ne\xi_{\eq,j}$
이므로 ``현재 평형'' 기준을 $\xi_{\tar,j}^b$ 의 국소 평형 $V_{\tar,j}^b(\xi_j)$ 로 바꿔야 하고, 그 경우 부호
논증은 $\eta_{j,\prog}^b=s_{\phi,j}^b[V_{n,\drive}-V_{\tar,j}^b(\xi_j)]$ 로 일반화된다(\S\ref{sec:ch5_kinetics};
같은 부호 상쇄 구조 유지, BOUNDED). 즉 부호 양정치는 ``target 기준 과전압''에 대해 성립하며, true equilibrium
과 branch target 의 차(히스테리시스 loss)는 \emph{별도} 에너지다(\S\ref{sec:ch5_hysheat} 이중계상 금지).
\end{boundbox}

% ====================================================================
\subsection{True equilibrium vs metastable branch target (Ch3 $\xi_{\sstat}$ 와의 연결)}\label{sec:ch5_branch}
% ====================================================================
\textbf{물리.} Chapter 1/3 의 $\xi_{\eq,j}(V_n,T)$ 는 \emph{true thermodynamic equilibrium} target 이다.
충방전 경로의 branch-dependent curve 는 다음 중 \emph{하나 이상}의 결과일 수 있다: (1) metastable phase path,
(2) nucleation barrier / phase-boundary pinning, (3) particle/domain memory, (4) finite-rate kinetic lag,
(5) concentration polarization, (6) ohmic / charge-transfer polarization, (7) aging / side reaction. 따라서
branch target 도입 시 \emph{단순 fitting curve 인지 metastable free-energy branch 인지}를 먼저 구분한다(이
구분 없이 충방전 곡선 차를 모두 ``열역학 히스테리시스''로 부르면 polarization 과 혼동된다 —
\S\ref{sec:ch5_apparent}).

\textbf{many-particle 기원(dreyer2010).} 삽입전지 히스테리시스의 열역학적 기원은 \emph{개별 입자}의 곡선이
아니라 \emph{다입자 집합}의 거동에서 나온다: 각 입자가 1차 상전이(비단조 화학퍼텐셜)를 가지면, 집합은 충전 시와
방전 시 \emph{서로 다른 입자 부분집합}이 순차적으로 변태하여 충/방전이 다른 평형 분지를 좇는다
(\cite{dreyer2010}, \AL{50}). 즉 $\xi_{\eq}^{\chg}\ne\xi_{\eq}^{\dis}$ 는 단일 입자의 비가역성이 아니라
\emph{다입자/mosaic 열역학}의 귀결이며, loop area 가 곧 한 입자의 소산열인 것은 아니다(loop area $\ne$ true
hys; \S\ref{sec:ch5_hysheat}).

\textbf{disorder pinning 기원(sethna1993).} 무질서가 도입한 1차 상전이는 \emph{barrier hierarchy}(여러
크기의 nucleation/depinning 장벽)를 만들어, 외부 구동이 한 장벽을 넘을 때마다 avalanche 형 변태가 일어나고
경로가 \emph{이력}을 획득한다(\cite{sethna1993}, \AL{51}). 이 disorder-driven first-order 전이의 hysteresis
는 흑연 staging 의 도메인 불균일$\cdot$결정립계$\cdot$phase-boundary pinning 과 대응하며, branch target 이
``방전 memory''와 ``충전 memory'' 사이에서 갈리는 물리 근거다.

\textbf{branch potential(물리 해석 가능 시).} branch 자유에너지 $G_j^b$ 로 branch potential 을 정의한다
(\cite{dreyer2010,sethna1993}, \AL{55}):
\begin{equation}
U_j^b(T,h_j)\ \sim\ \frac{1}{F}\frac{\partial G_j^b}{\partial n_j},
\label{eq:ch5_branch_potential}
\end{equation}
($\partial G_j^b/\partial n_j$[J/mol]$/F$[C/mol]$=$[J/C]$=$[V], 차원 정합; $h_j$ 의존은 \S\ref{sec:ch5_hstate}).
그때 branch target 은 식~\eqref{eq:ch5_logistic_branch} 형의 branch 버전이다:
\begin{equation}
\xi_{\tar,j}^b(V_n,T,h_j)=\frac{1}{1+\exp\!\big[-s_{\phi,j}^b\,(V_n-U_j^b(T,h_j))/w_j^b(T)\big]}.
\label{eq:ch5_branch_target}
\end{equation}
\textbf{부호 정합.} 식~\eqref{eq:ch5_branch_target} 의 지수에는 \S\ref{sec:ch5_chargesign} 에서 유도한
$s_{\phi,j}^b$ 가 들어가, branch target 의 $V_n$-단조 방향이 branch 진행 방향과 정합한다(방전 증가, 충전 감소).

\begin{keybox}
\textbf{★ branch target $=$ Chapter 3 의 $\xi_{\sstat}$ (detailed balance 이탈의 표현).} (\emph{앞 장 식 번호
규약}: Chapter 3 의 식~eq:ch3\_xiss 처럼 다른 장을 가리키는 번호는 본서로 결합 컴파일할 때 해소되며, 본 장 단독
컴파일에선 평문으로 둔다 — 본 장 내부 식만 \texttt{\textbackslash eqref} 를 쓴다. 이하 1회 적용.) branch 안에서 branch-local detailed
balance $r_j^{+,b}/r_j^{-,b}=\xi_{\tar,j}^b/(1-\xi_{\tar,j}^b)$ 를 쓰면 Chapter 3 의 nonequilibrium stationary
target 정의(식~eq:ch3\_xiss)가 그대로 $\xi_{\sstat,j}^b=r_j^{+,b}/(r_j^{+,b}+r_j^{-,b})=\xi_{\tar,j}^b$ 를 준다.
이 대수는 \emph{정의 일관성 확인(항등)}이다 — 정의 식~\eqref{eq:ch5_branch_db}($r^{+,b}/r^{-,b}=\xi_\tar^b/(1-\xi_\tar^b)$)
를 $\xi_{\sstat}=r^+/(r^++r^-)=[1+r^-/r^+]^{-1}$ 에 대입하면 $[1+(1-\xi_\tar)/\xi_\tar]^{-1}=\xi_\tar$ 로 돌아오는
\emph{무모순 점검}이지, $\xi_{\sstat}^b=\xi_\tar^b$ 를 독립 유도하는 것이 아니다(순환 회피). 즉 \emph{Chapter 3 의
비평형 정상 target 이 곧 branch target} 이며, \textbf{$\xi_{\tar,j}^b\ne\xi_{\eq,j}$ 는 Chapter 3 의 $C_j$ 가
true-equilibrium 값에서 branch 만큼 \emph{이탈}}했음을 뜻한다(Level A 분업에서 ``방향''이 metastable 자유에너지로
인해 이력을 획득). \textbf{히스테리시스는 이 \emph{target 이동}이지, mobility($k_j$)의 변화가 아니다.}

\textbf{실질 keystone 붕괴 논증(L4 두 조건과 1:1).} ``$\xi_{\sstat,j}^b\ne\xi_{\eq,j}$''가 생기는 것은 정의
재대입이 아니라 \emph{다음 위배}들 때문이다: (a) \textbf{$C_j^b$ 의 detailed-balance 이탈}(midpoint 에서 $C_j^b\ne0$
또는 $C_j^b$ 가 $h_j$-의존) — keystone 조건 (i)``$C_j$ 의 대칭 split'' 위배; (b) \textbf{$r_j^{+,b}/r_j^{-,b}$ 의
$\xi_j$-의존}(닫힌형이 아닌 음함수 정상점, path memory) — 역시 조건 (i) 의 $\xi_j$-무의존 위배; (c)
\textbf{$V_{n,\drive}\ne V_n$}(구동 과전압이 평형 보존식 해와 어긋남, polarization 기원) — keystone 조건 (ii)
``$V_{n,\drive}=V_n$'' 위배이며, 이 (c) 는 평형 분지 이동이 아니라 \emph{구동 과전압}이라 \S\ref{sec:ch5_apparent}
에서 polarization 으로 분리한다. (a)/(b) 가 조건 (i) 에, (c) 가 조건 (ii) 에 대응 — L4 의 두 조건과 1:1 이다.
Chapter 3 keystone 한정(L4)이 ``branch 비대칭 $\xi_{\sstat,j}\ne\xi_{\eq,j}$ 은 $C_j$ 가 detailed balance 를
이탈하거나 $r_j^+/r_j^-$ 가 $\xi_j$-의존일 때 발생하며 Chapter 5 로 넘긴다''고 명시한 그 영역이 바로 여기다.
\end{keybox}
\begin{boundbox}
\textbf{근거 요구(\AL{55}).} $\xi_{\tar,j}^b$ 를 쓰는 것 자체는 물리 모델이 아니다. $U_j^b,w_j^b,h_j$ 가
metastable branch / domain memory / nucleation barrier / phase-boundary history 와 연결돼야 한다
(dreyer2010 many-particle, sethna1993 disorder pinning). 근거 없으면 $\xi_{\tar,j}^b$ 는 empirical branch
fit 이며 본문 기본식이 아니라 실무 팁/비교 모델로 내린다(GS-4; \S\ref{sec:ch5_ident} practicebox).
\end{boundbox}

\subsubsection{Ch3$\to$Ch5 전달: detailed balance$\cdot$keystone 이 깨지는 자리}\label{sec:ch5_dbbreak}
\textbf{Chapter 3 가 본 장으로 넘긴 것(무비약 재확인).} Chapter 3 의 keystone 검증(verifybox)은 ``Level A $=$
Level B 의 detailed-balance 극한''이 \emph{두 조건}에서만 닫힘을 보였다: (i) $C_j$ 의 $V_n$$\cdot$$\xi_j$-무의존
(대칭 split), (ii) $V_{n,\drive}=V_n$. 그 한정의 마지막 문장이 ``branch 비대칭 $\xi_{\sstat,j}\ne\xi_{\eq,j}$
은 $\cdots$ Chapter 5 로 넘긴다''였다. 본 절이 그 \emph{넘겨받은} 부분을 명시한다. Chapter 3 의 rate ratio
일반형(식~eq:ch3\_ratio\_B)
\begin{equation}
\ln\frac{r_j^{+,b}}{r_j^{-,b}}=C_j^b(\xi_j,T)+\frac{\mathcal A_j^b}{RT}
\label{eq:ch5_ratio_branch}
\end{equation}
에서, true equilibrium 정합은 ``$C_j$ 가 $\xi_j$-무관 \emph{이고} midpoint 에서 0''을 요구했다(Ch.3
식~eq:ch3\_neff). \textbf{branch 가 생기는 정확한 조건}은 이 둘 중 하나의 위배다:
\begin{enumerate}[label=(\alph*),leftmargin=2.2em]
\item \textbf{$C_j^b$ 의 branch 이탈.} $C_j^b\ne0$ at midpoint, 또는 $C_j^b$ 가 branch(메모리 $h_j$)에 의존하면
  detailed-balance 우변이 true-equilibrium odds 와 갈라져 $\xi_{\tar,j}^b\ne\xi_{\eq,j}$ 다. 이는 metastable
  free-energy(다른 $U_j^b,w_j^b$)나 activity 비 $a_j^{+,b}/a_j^{-,b}$ 의 branch 의존에서 온다.
\item \textbf{$r_j^{+,b}/r_j^{-,b}$ 의 $\xi_j$-의존.} ratio 가 $\xi_j$ 에 의존하면 $\xi_{\sstat,j}^b$ 가
  닫힌형이 아니라 \emph{음함수}(implicit fixed point)가 되어, 같은 $V_n$ 에서도 경로(이력)에 따라 다른 정상점에
  머문다 — path memory 의 미시 표현이다(\S\ref{sec:ch5_hstate} 의 $h_j$).
\end{enumerate}
\textbf{무엇이 \emph{안} 깨지는가(중요).} Chapter 3 의 \emph{forward/backward 구조 자체}($\dd\xi_j/\dd t=
r_j^{+,b}(1-\xi_j)-r_j^{-,b}\xi_j$, mobility $=r^++r^-$, $\beta_j$ 가 ratio 의 $\mathcal A$-기울기서 상쇄)는
branch 안에서도 그대로 성립한다. 깨지는 것은 ``ratio $=$ \emph{true}-equilibrium odds''라는 \emph{정합 제약}
뿐이며, 그 자리에 ``ratio $=$ \emph{branch}-target odds''(식~\eqref{eq:ch5_branch_db})가 들어선다. 즉 Level A
의 ``mobility$=$속도'' 부분은 보존되고, ``열역학$=$방향(true equilibrium)'' 부분만 ``방향$=$branch target(이력
의존)''으로 \emph{부분} 붕괴한다 — 이것이 ``Level A 분업의 부분 붕괴''의 정확한 의미다(\AL{55}).

% ====================================================================
\subsection{Hysteresis state variable 과 metastability}\label{sec:ch5_hstate}
% ====================================================================
\textbf{물리.} branch 가 ``이력''을 가지려면 그 이력을 담는 내부 변수가 필요하다. transition 별 hysteresis state
$h_j(t)\in[-1,1]$ 를 도입한다: $h_j=+1$ 방전 branch memory, $h_j=-1$ 충전 branch memory, $h_j=0$
relaxed/중간(\AL{56}). branch center 는 이 메모리에 \emph{선형}으로 의존하는 signed hysteresis parameter 로
둔다:
\begin{equation}
U_j^\hys(T,h_j)=U_j^0(T)+\frac{\Delta U_j^\hys(T)}{2}\,h_j,\qquad \Delta U_j^\hys\in\mathbb R.
\label{eq:ch5_hys_center}
\end{equation}
\textbf{부호 검산.} $h_j=+1$(방전 memory)서 $U_j^\hys=U_j^0+\Delta U_j^\hys/2$, $h_j=-1$(충전 memory)서
$U_j^\hys=U_j^0-\Delta U_j^\hys/2$ — 두 branch center 의 차가 정확히 $\Delta U_j^\hys$ 다. (\emph{표기}: 이
$U_j^\hys(T,h_j)$ 가 다른 절의 branch center $U_j^b$ 와 동일물이다 — $b$ 가 정한 memory $h_j=s_{\phi,j}^b$ 를
넣은 값이 곧 $U_j^b$, 즉 $U_j^\hys\equiv U_j^b$.)
$\Delta U_j^\hys>0$ 이면 방전 branch center 가 충전보다 높은 convention(반대 관측이면 $\Delta U_j^\hys<0$ 로
피팅; 양 magnitude 를 강제하려면 별도 부호 인자 $s_{h,j}$ 를 명시하고 $\Delta U_j^\hys=s_{h,j}|\Delta U_j^\hys|$).
hysteresis state 의 동역학:
\begin{equation}
\frac{\dd h_j}{\dd t}=\lambda_j^b\big[h_{\tar,j}^b-h_j\big],
\label{eq:ch5_h_dyn}
\end{equation}
$h_{\tar,j}^b$ 는 현재 branch 의 memory target 으로 $h_{\tar,j}^b=s_{\phi,j}^b$ 이다($b=\dis$ 면 $+1$, $b=\chg$ 면
$-1$ 로 끌림; \S\ref{sec:ch5_chargesign} 의 부호 인자와 동일물이라 \emph{새 자유 파라미터가 아니다}). $\lambda_j^b$
는 \emph{fitting speed 가 아니라} domain rearrangement / phase-boundary depinning / nucleation /
microstructural memory 의 relaxation time scale 을 대표한다(\AL{56}; rest relaxation 자료 없으면 식별성 약함,
\S\ref{sec:ch5_ident}). \textbf{차원 검산.} $\lambda_j^b[\mathrm{s^{-1}}]\times$무차원 $[h]=\mathrm{s^{-1}}$,
좌변도 $\mathrm{s^{-1}}$ — 정합. \textbf{$[-1,1]$ 유지.} $h_{\tar,j}^b\in\{-1,+1\}$ 로 끌리는 1차 완화이므로
초기 $h_j\in[-1,1]$ 이면 해가 구간을 벗어나지 않는다(완화식의 단조 수렴; clip 으로 강제하지 않음 — GS-4).

% ====================================================================
\subsection{Branch-dependent forward/backward kinetics}\label{sec:ch5_kinetics}
% ====================================================================
\textbf{물리.} Chapter 3 의 forward/backward 구조를 branch 위에서 \emph{형태 그대로} 유지한다(깨지는 것은
정합 제약뿐, \S\ref{sec:ch5_dbbreak}):
\begin{equation}
\frac{\dd\xi_j}{\dd t}=r_j^{+,b}(1-\xi_j)-r_j^{-,b}\xi_j.
\label{eq:ch5_branch_fb}
\end{equation}
branch free-energy landscape 가 정의되면(식~\eqref{eq:ch5_branch_potential}) branch 안 \emph{local} detailed
balance 는 true equilibrium 이 아니라 branch target odds 를 만족한다:
\begin{equation}
\frac{r_j^{+,b}}{r_j^{-,b}}=\frac{\xi_{\tar,j}^b}{1-\xi_{\tar,j}^b}
\label{eq:ch5_branch_db}
\end{equation}
(global equilibrium 에 대한 detailed balance 가 아니라 \emph{metastable branch surface} 위의 local balance).
전위 구동력이 명시되면 branch 진행 affinity 는 \S\ref{sec:ch5_chargesign} 에서 유도한 부호 인자로(branch 유효
전자수 $n_j^{\eff,b}\equiv RT/(Fw_j^b)$, Ch1--4 $n_j^\eff{=}RT/(Fw_j)$ 의 branch 폭 $w_j^b$ 판)
\begin{equation}
\mathcal A_{j,\prog}^b=s_{\phi,j}^b\,n_j^{\eff,b}F\,[V_{n,\drive}-V_{\tar,j}^b(\xi_j)]=n_j^{\eff,b}F\,\eta_{j,\prog}^b,
\qquad \eta_{j,\prog}^b\equiv s_{\phi,j}^b\,[V_{n,\drive}-V_{\tar,j}^b(\xi_j)],
\label{eq:ch5_branch_affinity}
\end{equation}
$V_{\tar,j}^b(\xi_j)=U_j^b+s_{\phi,j}^b\,w_j^b\ln[\xi_j/(1-\xi_j)]$ 는 \emph{branch-local} 평형 이탈 기준
전위(\S\ref{sec:ch5_branch} 의 metastable 평형). \emph{부호 인자 $s_{\phi,j}^b$ 가 로그 항에 실리는 근거}: 이는
branch logistic $\xi_{\tar,j}^b=[1+e^{-s_{\phi,j}^b(V_{n,\drive}-U_j^b)/w_j^b}]^{-1}$ 의 \emph{역함수}이고, 역으로
풀면 $V_{\tar,j}^b-U_j^b=(w_j^b/s_{\phi,j}^b)\ln[\xi_j/(1-\xi_j)]=s_{\phi,j}^b w_j^b\ln[\xi_j/(1-\xi_j)]$
($1/s_{\phi}^b=s_\phi^b$, $\pm1$)다. 따라서 방전($s_\phi^\dis{=}+1$)에서는 $\xi_j$ 에 \emph{증가}형(Ch.4 형으로 환원),
충전($s_\phi^\chg{=}-1$)에서는 \emph{감소}형이다 — 이 감소형이 아래 부호 상쇄의 정의적 근거다. $\eta_{j,\prog}^b>0\Rightarrow b$-branch 방향 진행 촉진(방전서 $\dd\xi_j/\dd t>0$, 충전서
$\dd\xi_j/\dd t<0$). \textbf{affinity$\leftrightarrow$rate 부호 사슬.} 식~\eqref{eq:ch5_branch_fb} 의 두 율을
식~\eqref{eq:ch5_branch_db}($r^{+,b}/r^{-,b}=\xi_\tar^b/(1-\xi_\tar^b)$)로 정리하면 $\dd\xi_j/\dd t=
(r_j^{+,b}+r_j^{-,b})(\xi_{\tar,j}^b-\xi_j)$ (mobility $\times$ target 편차)다. 충전($\xi_j>\xi_{\tar,j}^\chg$,
즉 $\eta_{j,\prog}^\chg=s_\phi^\chg[V_{n,\drive}-V_{\tar,j}^\chg]>0$ 으로 충전 방향 구동)에서는 $\xi_{\tar,j}^\chg
-\xi_j<0\Rightarrow\dd\xi_j/\dd t<0$ — 진행 과전압의 양($\eta_{j,\prog}^\chg>0$)이 $\xi_j$ 의 \emph{감소}를 밀어
부호 사슬이 닫힌다. \textbf{이 $\eta_{j,\prog}^b$ 는 Chapter 3$\cdot$4 의 true-equilibrium 기반 $\eta_j$ 와
\emph{자동으로 같지 않다}} — branch target $V_{\tar,j}^b$ 가 true 평형 $V_{\eq,j}$ 에서 이탈했기 때문이다
(\S\ref{sec:ch5_branch} keybox). \textbf{부호 검산(verifybox 일반화).} \S\ref{sec:ch5_chargesign} verifybox 의
부호 상쇄가 ``true 평형''을 ``branch target''으로 바꾼 형으로 그대로 성립한다: 충전서 $I_j<0$ \emph{이고}
$\eta_{j,\prog}^\chg<0$ 이라 곱 $\ge0$.
\begin{boundbox}
\textbf{branch-local dissipation vs metastable free-energy loss(★ 이중계상 금지).} Chapter 4 의 미시$=$거시
결과는 branch 안에서 $\dot{\mathcal Q}_{\irr,j}^b=n_j^{\eff,b}I_j\,\eta_{j,\prog}^b\ge0$ 로 일반화된다($\eta$ 를
branch-local 이탈로). 그러나 이는 \emph{branch surface 위의 kinetic dissipation} 일 뿐이며, branch 가 결국
true equilibrium 으로 풀릴 때 방출되는 \emph{metastable 자유에너지 차이}(히스테리시스 loss,
\S\ref{sec:ch5_hysheat})와 \emph{구분}된다. 둘을 합치면 같은 에너지를 이중 계상한다. 즉 ``현재 branch 위에서
target 을 좇는 소산''($\eta_{j,\prog}^b$)과 ``branch 자체가 true 평형 대비 갖는 자유에너지 잉여''($V_{\tar,j}^b
-V_{\eq,j}$)는 서로 다른 에너지 항이다(\AL{57}).
\end{boundbox}

% ====================================================================
\subsection{Apparent 전위와 polarization 의 branch 확장 ($\Delta V_\obs\ne\Delta V_\hys$)}\label{sec:ch5_apparent}
% ====================================================================
\textbf{물리.} 관측되는 충방전 전압 차(loop 폭)를 곧바로 히스테리시스로 부르면 polarization 과 혼동된다. 본
절이 그 \emph{분리}를 못박는다. branch 별 apparent 음극 전위는 Chapter 1 의 apparent 전위(식~eq:Vapp)를 signed
current 로 확장한 것이다:
\begin{equation}
V_{n,\app}^b=V_n+\Delta V_{n,\pol}^b,\qquad
\Delta V_{n,\pol}^b\approx I_s\,R_n^b(q_\stt,T,|I_s|)\ \ \text{(small-signal/secant)}.
\label{eq:ch5_vapp_branch}
\end{equation}
\textbf{이는 polarization 표현이지 hysteresis offset 이 아니다.} $R_n^b>0$ 이어도 \emph{$I_s$ 의 부호}에 따라
전압 이동 \emph{방향}이 바뀐다: 방전($I_s>0$)서 $\Delta V_{n,\pol}^b>0$, 충전($I_s<0$)서 $\Delta V_{n,\pol}^b<0$.
즉 polarization 은 $|I_s|\to0$ 에서 사라지지만(전류 차단 시 0), \emph{true hysteresis 는 $|I_s|\to0$ 에서도
남는} 평형 분지 차다 — 이 점이 둘을 구분하는 \emph{조작적} 기준이다(\S\ref{sec:ch5_ident} 의 저율/GITT 비교).
\textbf{한정.} $|I_s|\to0$ 외삽은 $\Delta V_\pol$ 만 제거하며, kinetic lag($\Delta V_\kin$, 유한 완화시간 잔류)와
aging($\Delta V_{\mathrm{aging}}$, 비가역 누적)은 외삽으로 분리되지 않는다 — 따라서 외삽 단독으로 $\Delta V_\hys$
를 분리하려면 lag/aging 이 무시 가능한 GITT 완화 종점을 함께 써야 한다(\S\ref{sec:ch5_ident}).
관측 voltage gap 의 분해. \textbf{각 성분은 loop \emph{폭}(branch 차)이다}: $\Delta V_X\equiv X^\dis-X^\chg$ 로
정의하여 식~\eqref{eq:ch5_vapp_branch} 의 signed single-branch 표현과 입도를 통일한다(분해식 좌변 $\Delta V_\obs$
가 loop 폭이므로 우변 각 항도 폭이어야 한다):
\begin{equation}
\Delta V_\obs=\Delta V_\hys+\Delta V_\pol+\Delta V_\kin+\Delta V_\transp+\Delta V_{\mathrm{aging}},
\label{eq:ch5_gap_decomp}
\end{equation}
특히 polarization 성분은 두 branch 의 \emph{부호 상쇄}로 더해진다: $\Delta V_\pol\equiv\Delta V_{n,\pol}^\dis
-\Delta V_{n,\pol}^\chg\approx(+|I_s|R_{n,\pol})-(-|I_s|R_{n,\pol})=2|I_s|R_{n,\pol}>0$ — single-branch 부호
($\pm$)가 \emph{차}를 취하면 양의 폭으로 합쳐진다. 여기서 \textbf{$R_n^b$ 는 단독 ohmic 저항이 아니라 합성 분극
저항} $R_{n,\pol}^b\equiv R_n^b+R_\ct^b$(ohmic $+$ charge-transfer secant)로 읽는다 — Chapter 3 의 $R_n\ne R_\ct$
구분과 충돌하지 않게 명시한다(\AL{58}).
\begin{longtable}{p{0.24\textwidth} >{\raggedright\arraybackslash}p{0.64\textwidth}}
\toprule 성분 & 물리적 의미 \\ \midrule \endhead
$\Delta V_\hys$ & \textbf{true thermodynamic hysteresis}: 충/방전 평형 분지 차 $U_j^\dis-U_j^\chg$ (metastable free-energy; $|I_s|\to0$ 잔존) \\
$\Delta V_\pol$ & polarization($R_n^b$,\brk ohmic/charge-transfer); $|I_s|\to0$ 소멸, $I_s$ 부호 의존 \\
$\Delta V_\kin$ & kinetic lag($\xi_j$ 가 target 못 따라감;\brk Ch.1 $L$ 꼬리) \\
$\Delta V_\transp$ & transport/concentration limitation (Ch.4 transport heat 와 짝) \\
$\Delta V_{\mathrm{aging}}$ & aging/side reaction 누적 변화 \\
\bottomrule
\end{longtable}
따라서 본 장의 기본 원칙은
\begin{equation}
\boxed{\;\Delta V_\obs\ne\Delta V_\hys\;.}
\label{eq:ch5_obs_not_hys}
\end{equation}
\textbf{Ch1$\cdot$Ch3 정합.} $\Delta V_\pol$ 은 Chapter 1 의 $R_n$(apparent voltage-axis shift)$\cdot$Chapter 3
의 $R_\ct$(charge-transfer)에서 오고, $\Delta V_\hys$ 만이 평형 분지 차다 — Chapter 3 의 $R_n\ne R_\ct\ne R_{n,\eff}$
구분이 여기서 ``polarization gap $\ne$ hysteresis gap''으로 이어진다(\AL{58}). gap 전체를 한 항으로 피팅하면
polarization 과 true hysteresis 가 co-linear 라 식별 불가능하다(\S\ref{sec:ch5_ident}).

% ====================================================================
\subsection{Branch-dependent ICA/DVA 좌표}\label{sec:ch5_icadva}
% ====================================================================
\textbf{물리.} 충방전 ICA/DVA 를 한 그림에 겹칠 때 좌표를 명시하지 않으면 부호와 peak 방향이 혼동된다. signed
capacity 와 branch positive capacity 를 구분한다:
\begin{equation}
Q_s(t)=Q_{s,\mathrm{ref}}+\int_{t_{\mathrm{ref}}}^t I_s(t')\,\dd t',\qquad
Q_b=\begin{cases}Q_\dis & (b=\dis)\\ Q_\chg & (b=\chg)\end{cases}
\label{eq:ch5_signed_cap}
\end{equation}
($Q_s$ 는 signed, $Q_b$ 는 각 branch 의 양 용량). branch 별 apparent DVA 는 $(\dd V_{n,\app}^b/\dd Q_b)_b$,
signed thermodynamic 분석은 $\dd V_n/\dd Q_s$ 를 \emph{별도로} 쓴다. \textbf{ICA($\dd Q/\dd V$) branch 좌표.}
DVA 의 역수로, signed thermodynamic ICA 와 branch-positive apparent ICA 를 병기한다:
\begin{equation}
\boxed{\;\Big(\frac{\dd Q_b}{\dd V_{n,\app}^b}\Big)_b=\Big(\frac{\dd V_{n,\app}^b}{\dd Q_b}\Big)_b^{-1}
\quad(\text{branch-positive}),\qquad
\frac{\dd Q_s}{\dd V_n}=\Big(\frac{\dd V_n}{\dd Q_s}\Big)^{-1}\quad(\text{signed thermodynamic}).\;}
\label{eq:ch5_branch_ica}
\end{equation}
\textbf{충전 branch 부호 반전(유도).} 충전서 $\dd q_\stt/\dd t<0$(식~\eqref{eq:ch5_qstate}, $I_s<0$)라 signed
좌표 $Q_s$ 가 \emph{감소}하므로, 같은 물리 peak 라도 $\dd Q_s/\dd V_n$ 의 부호가 방전 대비 뒤집혀 \emph{음의
lobe} 로 나타난다($\dd Q_s<0$ 이 분자 부호를 반전). branch-positive $Q_b$ 는 항상 증가(각 branch 양 용량)라
$(\dd Q_b/\dd V_{n,\app}^b)_b$ 는 부호 반전 없이 양 branch 를 같은 향으로 그린다. \textbf{좌표를 명시하지 않은
charge/discharge ICA overlay 는 부호와 peak 방향 혼동을 만든다}(\AL{58}). \textbf{Ch1 정합.} 방전 단독서
$Q_s=Q_\ext$ 이므로 Chapter 1 의 ICA/DVA($Q_\ext$ 기준, 식~eq:fiteq — 앞 장 식 번호 규약대로 결합 컴파일 시
해소, 본 장 단독서 평문)와 정확히 일치한다(앞 장과 충돌 없음).

% ====================================================================
\subsection{Full-cell 관측 전압과 음극 peak}\label{sec:ch5_fullcell}
% ====================================================================
\textbf{물리.} 실험에서 직접 읽는 것은 음극 단독 전위가 아니라 full-cell 전압이다. full-cell 전압은 양극/음극
전위와 signed loss 의 합이다:
\begin{equation}
V_\cell^b=V_p^b-V_n^b-\eta_\loss^b,
\label{eq:ch5_fullcell}
\end{equation}
$\eta_\loss^b$=branch-dependent signed loss/overpotential aggregate. \textbf{exact$\leftrightarrow$apparent 부호
폐합($\eta_\loss^b$ 의 정의).} apparent 식 $V_{\cell,\app}^b=V_{p,\app}^b-V_{n,\app}^b$ 는 손실항을 \emph{제거}한
표현인데, exact 식~\eqref{eq:ch5_fullcell} 의 $-\eta_\loss^b$ 와 폐합하려면 $\eta_\loss^b$ 가 양극$\cdot$음극
분극 \emph{차}로 흡수돼야 한다. 따라서
\begin{equation}
\eta_\loss^b\equiv(\Delta V_{p,\pol}^b-\Delta V_{n,\pol}^b)+\eta_{\loss,\mathrm{res}}^b,
\label{eq:ch5_etaloss_def}
\end{equation}
여기서 \emph{양극 분극} $\Delta V_{p,\pol}^b=I_s R_p^b$(음극 식~\eqref{eq:ch5_vapp_branch} 와 \emph{동일 secant
규약}, $V_{p,\app}^b=V_p^b+\Delta V_{p,\pol}^b$), $\eta_{\loss,\mathrm{res}}^b$ 는 $V_{p,\app}^b,V_{n,\app}^b$ 에
\emph{흡수되지 않은} 잔여 transport$\cdot$lag 손실이다. (검산: $V_{\cell,\app}^b=(V_p^b+\Delta V_{p,\pol}^b)
-(V_n^b+\Delta V_{n,\pol}^b)=V_p^b-V_n^b-(\Delta V_{n,\pol}^b-\Delta V_{p,\pol}^b)$ 이며 exact 식~\eqref{eq:ch5_fullcell}
와 $\eta_{\loss,\mathrm{res}}^b=0$ 일 때 일치한다.) \textbf{두 표현 동시 사용 시 polarization 중복 계상 금지} —
$V_{p,\app}^b,V_{n,\app}^b$ 에 이미 포함된 분극(식~\eqref{eq:ch5_vapp_branch})을 다시 $\eta_\loss^b$ 의 $\Delta V_\pol$
부분에 \emph{이중}으로 넣지 않는다(잔여항 $\eta_{\loss,\mathrm{res}}^b$ 만 별도, \AL{58}). signed capacity 기준
full-cell DVA(\emph{양극 좌표 환산 주의}: $Q_s$ 는 음극 signed 좌표이며 양극은 stoichiometric coupling $\dd Q_p=
-\dd Q_s$ 로 환산되어 $\dd V_p^b/\dd Q_s=(\dd V_p^b/\dd Q_p)(\dd Q_p/\dd Q_s)=-\dd V_p^b/\dd Q_p$):
\begin{equation}
\frac{\dd V_\cell^b}{\dd Q_s}=\frac{\dd V_p^b}{\dd Q_s}-\frac{\dd V_n^b}{\dd Q_s}-\frac{\dd\eta_\loss^b}{\dd Q_s}
=-\frac{\dd V_p^b}{\dd Q_p}-\frac{\dd V_n^b}{\dd Q_s}-\frac{\dd\eta_\loss^b}{\dd Q_s}.
\label{eq:ch5_fullcell_dva}
\end{equation}
따라서 full-cell graphite peak 는 양극$\cdot$음극$\cdot$loss 가 섞인 관측 신호다(여기서 transport 손실은
$\eta_\loss^b$ 의 잔여항 $\eta_{\loss,\mathrm{res}}^b$ 에 포함된 \emph{부분집합}이지 독립 항이 아니다). reference
electrode / half-cell reconstruction / electrode balancing 없이 full-cell peak 를 음극 peak 로 \emph{직접
등치하지 않는다}(Chapter 1 의 ``OCV 를 읽는 흐름으로 회귀 X''와 같은 입도의 경계).

% ====================================================================
\subsection{Hysteresis energy 와 heat (loop area = 소산열)}\label{sec:ch5_hysheat}
% ====================================================================
\textbf{물리.} 한 충방전 cycle 의 전압--용량 loop 면적은 그 cycle 에서 소산된 에너지다:
\begin{equation}
E_{\mathrm{loop}}=\oint V\,\dd Q.
\label{eq:ch5_loop_area}
\end{equation}
\textbf{차원 검산.} $V[\mathrm V]\times Q[\mathrm C]=\mathrm{V\cdot C}=\mathrm J$ — 에너지. \textbf{부호$\cdot$적분방향
$\cdot$$\ge0$.} 가역 준정적 cycle 에서는 $\oint V\,\dd Q=0$(2법칙 등호, 경로 무관)이고 비가역 cycle 에서 $\oint
V\,\dd Q>0$ 이다. 양정치를 보장하려면 적분 방향을 $q_\stt$ 의 1주기(방전 증가 $\to$ 충전 감소)로 \emph{고정}하고
부호를 \S\ref{sec:ch5_branchheat} 의 $s^{b,\conv}$ bookkeeping 과 정렬한다; 방향$\cdot$부호가 미정렬이면
$E_{\mathrm{loop}}=|\oint V\,\dd Q|$ 로 절댓값을 취한다. \textbf{가역 극한 검산.} $|I_s|\to0$ 으로 polarization
$\to0$ 이고 평형 분지 차도 0(분지 비대칭 없음)이면 loop 가 닫혀 $E_{\mathrm{loop}}\to0$ — 가역 극한에서 소산이
사라진다는 2법칙 한계와 일치. 그러나 이 면적은
true hysteresis + polarization loss + kinetic lag + transport loss + side reaction 을 \emph{모두} 포함한다
(식~\eqref{eq:ch5_gap_decomp} 의 모든 성분이 loop 폭에 기여). \textbf{따라서 $E_{\mathrm{loop}}$ 전체를 true
hysteresis heat 로 단정하지 않는다}(loop area $\ne$ true hys; dreyer2010 의 many-particle 결론과 정합 —
\AL{50}). 물리적으로 \emph{분리된} hysteresis loss voltage 가 식별된 경우에만 후보 발열항으로 둔다:
\begin{equation}
\dot{\mathcal Q}_{\hys,n}^b=|I_s|\,\Delta V_{\hys\text{-}\loss,n}^b,\qquad \Delta V_{\hys\text{-}\loss,n}^b\ge0,
\label{eq:ch5_hys_heat}
\end{equation}
더 일반적으로는 Chapter 4 의 flux--force 형으로 $\dot{\mathcal Q}_{\hys,n}^b=J_\hys^b\,\mathcal A_\hys^b\ge0$
이다(\AL{59}; 2법칙). (\emph{J$\leftrightarrow$W 연결}: 식~\eqref{eq:ch5_hys_heat} 의 $\dot{\mathcal Q}_{\hys,n}^b$
는 \emph{발열률}[$|I_s|[\mathrm A]\times\Delta V[\mathrm V]=\mathrm W$]이고, 이를 1주기 적분한 $\oint\dot{\mathcal Q}_{\hys,n}^b\,\dd t$[J]
가 식~\eqref{eq:ch5_loop_area} 의 $E_{\mathrm{loop}}$ 중 \emph{분리된 hysteresis 성분}에 해당한다 — 율[W] 과
loop 에너지[J] 의 차원이 시간 적분으로 이어진다.) \textbf{단순 signed $I_s\,\Delta V_\hys^b$ 는 부호 convention 이 정렬된 경우에만} 쓴다
— Chapter 4 가 ``$I\eta$ 단독''(부호 미정렬)을 비가역 열로 쓰는 것을 차단한 것과 같은 원칙이며, 본 장
\S\ref{sec:ch5_chargesign} 의 부호 상쇄 유도로 $|I_s|\,\Delta V_{\hys\text{-}\loss}^b\ge0$ 형이 정당화된다.
\textbf{entropy production 정합(Ch4).} hysteresis loss 가 metastable 자유에너지의 비가역 방출이면, 그것은
Chapter 4 의 entropy production $T\dot S_\irr=\sum_\alpha J_\alpha X_\alpha\ge0$ 의 한 채널(path-dependent
free-energy relaxation)이며 2법칙으로 $\ge0$ 이다 — 즉 loop area 의 \emph{분리된 hysteresis 성분}은 Chapter 4
의 비가역 발열 framework 안에서 닫힌다(새 피팅항이 아님).
\begin{boundbox}
\textbf{★ 중복 계상 금지(\AL{57},\AL{59}).} $\dot{\mathcal Q}_{\hys,n}^b$ 를 별도 항으로 쓰려면 같은 voltage-gap
성분을 $\eta_\loss^b$, $R_n^b$, $R_{n,\eff}^b$, transport heat, 그리고 \S\ref{sec:ch5_kinetics} 의 branch-local
kinetic dissipation $n_j^{\eff}I_j\eta_{j,\prog}^b$ 에 \emph{동시에} 넣지 않는다. 특히 branch-local kinetic
dissipation(target 을 좇는 소산)과 hysteresis loss(branch 가 true 평형으로 풀릴 때의 자유에너지 방출)는
\emph{다른 에너지}이므로(\S\ref{sec:ch5_kinetics} boundbox) 둘을 합산하면 같은 에너지를 두 번 센다. 분리되지
않은 loop area 전체는 true hysteresis heat 가 아니다.
\end{boundbox}

% ====================================================================
\subsection{Branch-dependent heat 구조}\label{sec:ch5_branchheat}
% ====================================================================
\textbf{물리.} Chapter 4 의 열원 분해(식 (L5))를 branch 로 확장하고 hysteresis 항을 추가한다(\AL{59}):
\begin{equation}
\dot{\mathcal Q}_{\src,n}^b=\dot{\mathcal Q}_{\irr,n}^b+\dot{\mathcal Q}_{\rev,n}^b+\dot{\mathcal Q}_{\rlx,n}^b
+\dot{\mathcal Q}_{\transp,n}^b+\dot{\mathcal Q}_{\hys,n}^b.
\label{eq:ch5_branch_heat}
\end{equation}
\begin{longtable}{p{0.31\textwidth} >{\raggedright\arraybackslash}p{0.57\textwidth}}
\toprule 항 & 물리적 근거 (부호) \\ \midrule \endhead
$\dot{\mathcal Q}_{\irr,n}^b$ & $J^b\mathcal A^b\ge0$,\brk overpotential-driven entropy production $=\sum_j n_j^{\eff}I_j\eta_{j,\prog}^b\ge0$ (\S\ref{sec:ch5_chargesign} 부호 상쇄로 양 branch $\ge0$) \\
$\dot{\mathcal Q}_{\rev,n}^b$ & branch별 entropy coefficient 또는 transition entropy (부호가변; 충전서 $I_j<0$, 즉 $\dd\xi_j/\dd t<0$ 로 방전과 부호 반전 — $I_j=Q_{j,\tot}\,\dd\xi_j/\dd t$) \\
$\dot{\mathcal Q}_{\rlx,n}^b$ & metastable branch relaxation 또는 transition-network entropy production \\
$\dot{\mathcal Q}_{\transp,n}^b$ & diffusion/electrolyte/finite-current dissipative loss ($\ge0$) \\
$\dot{\mathcal Q}_{\hys,n}^b$ & polarization 과 분리된 path-dependent hysteresis loss ($\ge0$; \S\ref{sec:ch5_hysheat}) \\
\bottomrule
\end{longtable}
\textbf{가역 열(transition entropy basis, 충전 부호 명시).} Chapter 4 의 transition entropy basis(식~eq:ch4\_rev\_xi)
를 branch 로 쓰면
\begin{equation}
\dot{\mathcal Q}_{\rev,\xi}^b=\sum_{j=1}^{N_p}s_{\rev,\xi,j}^{b,\conv}\,T\,\Delta S_j^b\,\frac{Q_{j,\tot}}{F}\,\frac{\dd\xi_j}{\dd t}.
\label{eq:ch5_branch_rev}
\end{equation}
\textbf{차원 검산.} $T[\mathrm K]\Delta S_j^b[\mathrm{J/(mol\,K)}](Q_{j,\tot}/F)[\mathrm{mol}](\dd\xi_j/\dd t)[\mathrm{s^{-1}}]
=\mathrm{J/s}=\mathrm W$ — 가역 발열률. \textbf{충전 부호 명시.} 충전서 $\dd\xi_j/\dd t<0$ 이 될 수 있으므로
($\xi_j$ 감소), 같은 $\Delta S_j^b$ 라도 가역 열의 부호가 branch 마다 반대로 나온다 — 이것이
\S\ref{sec:ch5_chargesign} 의 ``가역 열 부호가변성''(verifybox 직후 단락)의 정량 표현이다. 충전서 $\dd\xi_j/\dd t<0$
은 $I_j=Q_{j,\tot}\,\dd\xi_j/\dd t<0$ 과 동치이고, 이는 다시 $q_\stt$ 감소($I_s<0$, 식~\eqref{eq:ch5_qstate})와
\S\ref{sec:ch5_chargesign} 의 ``$I_j<0$'' 표현과 \emph{동일 기제}다(셋이 한 부호 사슬). \textbf{부호 책임 일원화.}
$s_{\rev,\xi,j}^{b,\conv}$ 는 heat-positive 규약을 명시한 뒤 고정하는 bookkeeping factor 인데, 이 인자는
\emph{branch 에 무의존}(단일 heat-positive 규약)이며 — branch 에 따른 부호 반전은 \emph{오직} $\dd\xi_j/\dd t$ 가
담당한다(부호를 $s^{b,\conv}$ 와 $\dd\xi/\dd t$ 양쪽에 동시에 싣는 \emph{이중계상}을 피한다). 즉
$s_{\rev,\xi,j}^{b,\conv}\equiv s_{\rev,\xi,j}^{\conv}$(피팅 X; Chapter 4 의 $s^\conv$ 와 동일 규약, 서로 다른
규약으로 독립 도입 금지). \textbf{Hysteresis heat 와 reversible heat 가 branch heat asymmetry 를
\emph{동시에} 설명하지 않도록 독립 제약이 필요}하다(\S\ref{sec:ch5_ident}; calorimetry).

% ====================================================================
\subsection{Rest 구간과 branch switching}\label{sec:ch5_rest}
% ====================================================================
\textbf{물리.} 휴지 구간에서는 $I_s=0,\ \dd q_\stt/\dd t=0$ 이지만 내부 $\xi_j,h_j,V_n$ 은 계속 relax 한다
(Chapter 4 의 휴지 relaxation heat 와 짝). $\xi_j$ 가 무엇을 향해 풀리는지에 두 후보가 있다(완화율 $k_j^\rst$
[1/s] 은 rest 구간 $\xi_j$ 완화율이며, Ch1 mobility $k_j$ 와는 별개의 휴지 relaxation 상수다). true equilibrium
으로 relax(메모리 소멸):
\begin{equation}
\frac{\dd\xi_j}{\dd t}=k_j^\rst\big[\xi_{\eq,j}(V_n,T)-\xi_j\big],
\label{eq:ch5_rest_eq}
\end{equation}
metastable memory 유지(branch target 으로 relax):
\begin{equation}
\frac{\dd\xi_j}{\dd t}=k_j^\rst\big[\xi_{\tar,j}^\rst(V_n,T,h_j)-\xi_j\big].
\label{eq:ch5_rest_meta}
\end{equation}
무엇이 맞는지는 rest relaxation voltage + thermal response 로 제한한다(둘은 \emph{같은 데이터로} 판별 가능 —
true 평형으로 풀리면 OCV 가 두 branch 에서 한 값으로 수렴, memory 유지면 갈린 채 남음). hysteresis state 도
두 후보다:
\begin{align}
\text{memory-retaining:}\quad &\frac{\dd h_j}{\dd t}=0,\label{eq:ch5_h_rest_mem}\\
\text{relaxing:}\quad &\frac{\dd h_j}{\dd t}=-\lambda_j^{b=\rst}\,h_j\quad(h_j\to0\ \text{단조 감쇠};\ \text{식~\eqref{eq:ch5_h_dyn} 의 }b{=}\rst,\ h_{\tar,j}^\rst{=}0\text{ 특수화}).\label{eq:ch5_h_rest_relax}
\end{align}
rest relaxation 데이터가 없으면 두 모델은 식별 곤란하다(\AL{56}; 임의 택일 금지). \textbf{두 후보 짝의 결합.}
$\xi_j$-target 후보(eq vs meta, 식~\eqref{eq:ch5_rest_eq}--\eqref{eq:ch5_rest_meta})와 $h_j$ 후보(memory-retaining
vs relaxing, 식~\eqref{eq:ch5_h_rest_mem}--\eqref{eq:ch5_h_rest_relax})는 독립 두 짝이라 교차조합($2\times2$)이
가능하나, 물리적으로 ``$\xi_j$ true-eq relax''와 ``$h_j\to0$''(메모리 소멸)이 \emph{같은 점근}(완전 relaxed)으로,
``$\xi_j$ branch-target relax''와 ``$\dd h_j/\dd t=0$''(메모리 유지)이 \emph{같은 점근}으로 짝지어진다 — 두 짝을
독립으로 피팅하지 않고 한 물리 가설(소멸/유지)로 묶어 식별한다. \textbf{branch switching 연속성(중요).} branch
전환 시 \emph{state} 변수 $\xi_j,h_j,V_n$ 은 \emph{연속}이며 재초기화하지 않는다 — 물리 상태는 전류 부호가
바뀐다고 점프하지 않는다. 반면 \emph{구동} 측 부호 인자 $s_{\phi,j}^b$(및 $h_{\tar,j}^b$)는 branch 전환서
\emph{불연속}으로 뒤집힌다($+1\leftrightarrow-1$) — 즉 불연속은 구동 규약에만 있고 state 적분 변수에는 없다. 전하
보존식(식~\eqref{eq:ch5_relbalance})은 매 시간점 다시 풀어 $V_n$ 을 얻으므로, branch 전환은 $I_s$ 부호 변화로만
들어가고 state 변수의 불연속을 만들지 않는다(\AL{52}).

% ====================================================================
\subsection{식별성 및 필요한 실험 제약}\label{sec:ch5_ident}
% ====================================================================
\textbf{물리.} branch 항들은 강하게 상관되어, 같은 충방전 곡선만으로 동시 결정할 수 없다. 해결책은 임의
clipping 이 아니라 독립 실험과 물리적 prior 다(Chapter 1--4 의 forward-only$\cdot$식별성 원칙 계승).
\textbf{식별 scope.} 식별 대상 파라미터는 $\{U_j^b,\,w_j^b,\,k_j^b,\,\lambda_j^b,\,\Delta S_j^b,\,R_n^b\}$ 이다 —
충전 부호 인자 $s_{\phi,j}^b$ 는 \AL{53} 에서 \emph{유도}된 상수($\pm1$)라 식별 대상이 \emph{아니고}, loop 면적
$E_{\mathrm{loop}}$ 는 피팅 파라미터가 아니라 \emph{관측량}이다(식별 입력).
\begin{longtable}{p{0.34\textwidth} >{\raggedright\arraybackslash}p{0.54\textwidth}}
\toprule 상관되는 항 & 주의점 \\ \midrule \endhead
$\Delta V_\hys$ 와 $\Delta V_\pol$ & voltage gap 을 모두 hysteresis 로 해석 금지($\Delta V_\obs\ne\Delta V_\hys$; $|I_s|\to0$ 외삽으로 분리) \\
$U_j^b$ 와 $R_n^b$ & branch offset(평형 분지)과 polarization shift 가 peak 위치를 모두 변화 \\
$k_j^b,\ h_j,\ \lambda_j^b$ & kinetic lag 와 memory relaxation 이 tail/peak shift 를 모두 설명 가능 — 실험축 분리: $k_j^b\leftarrow$ rate series(율 의존), $\lambda_j^b\leftarrow$ rest relaxation 시간상수, $h_j\leftarrow$ rest 후 재기동(메모리 잔존 여부) \\
$\Delta S_j^b$ 와 hysteresis heat & branch heat asymmetry 를 가역 열과 hysteresis loss 가 중복 설명 위험 \\
full-cell $V_p$ 와 음극 $V_n$ & full-cell 서 양극/음극 contribution 중첩(reference electrode 필요) \\
\bottomrule
\end{longtable}
\textbf{필요 실험.} 저율 charge/discharge 비교(GITT/저율 OCV — true hysteresis vs polarization 분리), rest
relaxation(memory 소멸 vs 유지 판별), current interrupt/pulse(immediate polarization), rate series(kinetic
lag vs hysteresis), temperature series($\Delta S_j^b$, $U_j^b(T)$), reference electrode/half-cell
reconstruction(full-cell 분해), calorimetry/temperature response(가역/비가역/hysteresis heat 분리).
\textbf{$U_j^b$ 의 polarization$\cdot$lag-독립 고정(양방향 GITT).} branch center $U_j^b$ 를 polarization$\cdot$kinetic
lag 와 독립으로 고정하려면, 각 SOC 격자에서 충$\cdot$방전 \emph{양방향} GITT 완화 종점($I=0$ OCV)을 측정해
$U_j^\dis-U_j^\chg=\Delta V_\hys$ 를 직접 추출한다(완화 종점이라 kinetic lag 가 배제됨). \emph{저율 연속곡선의
$|I_s|\to0$ 외삽 단독}은 polarization 만 제거하여 hysteresis 와 kinetic lag 를 분리하지 못하므로, 양방향 완화
종점 측정과 병용해야 한다.
\begin{practicebox}
초기 비교에서 충전/방전을 별도 empirical branch curve 로 피팅할 수 있으나 \emph{물리 모델 기본식이 아니다}.
그 curve 가 metastable free-energy branch / domain memory / nucleation barrier / polarization / transport /
aging 중 무엇인지 분리되지 않으면 논문용 물리 결론에 쓰지 않고, 단순 empirical residual 로 명시한 뒤 주 결론에서
제외한다(GS-4; solver/numerical 구현과 validation pipeline 은 Ch1 부록 B).
\end{practicebox}

% ====================================================================
\subsection{Chapter 5 의 가정 원장 (AL-50$\sim$59)}\label{sec:ch5_al}
% ====================================================================
본 장에서 사용한 가정을 통합 ledger(RB\_AL\_MASTER) Ch5 그룹(AL-50$\sim$59)으로 정리한다. AL-50,51 은 기등재분
(Foundation), AL-52$\sim$59 는 본 장 S2 에서 채운 신규 정의다(consolidated\_ch5 의 dangling \AL{Y1} 의 실제
정의 부여 포함: Y1$\to$AL-50/51/55 — 원천은 \texttt{Claude/docs} 의 consolidated\_ch5 에 있던 미정의 \AL{Y1}
이며, 본 장에서 AL-50/51/55 로 재지정한 추적은 master(RB\_AL\_MASTER)의 Correction 주석에 박았다). tier $=$
GROUNDED/BOUNDED/FLAGGED.
{\footnotesize
\setlength{\tabcolsep}{3pt}
\begin{longtable}{>{\raggedright\arraybackslash}p{0.045\textwidth} >{\raggedright\arraybackslash}p{0.335\textwidth} >{\raggedright\arraybackslash}p{0.115\textwidth} >{\raggedright\arraybackslash}p{0.225\textwidth} >{\raggedright\arraybackslash}p{0.155\textwidth}}
\toprule AL\# & 가정 진술 & tier & 근거 cite (DOI/ISBN) & 사용 식 \\ \midrule \endhead
AL-50 & 삽입전지 히스테리시스 $=$ metastable/다입자(many-particle,\brk mosaic) 열역학 기원;\brk loop area $\ne$ true hys & GROUNDED & dreyer2010 (DOI 10.1038/\brk nmat2730) & \eqref{eq:ch5_branch_potential},\allowbreak\eqref{eq:ch5_loop_area},\allowbreak\eqref{eq:ch5_hys_heat} \\
AL-51 & disorder-driven first-order 전이의 hysteresis/hierarchy(nucleation/depinning barrier 위계,\brk avalanche,\brk path memory) & GROUNDED & sethna1993 (DOI 10.1103/\brk PhysRevLett.70.3347) & \eqref{eq:ch5_branch_potential},\allowbreak\eqref{eq:ch5_hys_center},\allowbreak\eqref{eq:ch5_h_dyn} \\
AL-52 & signed current $I_s$ / signed coordinate $q_\stt$($\dd q_\stt/\dd t{=}I_s/Q_\cell$); 내부 state $\xi_j$ 유지($\zeta_j{=}1{-}\xi_j$ 종속);\brk LLI/LAM 임의 흡수 금지;\brk branch switching 연속(재초기화 X) & BOUNDED & doyle1993,\brk newman (DOI 10.1149/\brk 1.2221597;\brk ISBN 978-0-\brk 471-\brk 47756-3) & \eqref{eq:ch5_signed_current},\allowbreak\eqref{eq:ch5_qstate},\allowbreak\eqref{eq:ch5_relbalance} \\
AL-53 & ★ 충전 branch 부호 $s_{\phi,j}^\chg{=}-1$ 을 logistic 지수 부호 정합으로 \emph{유도}(``$\mathcal A_j{>}0\Leftrightarrow$ 진행 촉진'' 규약을 양 branch 일관 적용; 방전 $+1$ 과 동시 고정) & GROUNDED & bard\brk faulkner,\brk newman,\brk bazant2013 (ISBN 978-0-\brk 471-\brk 04372-0;\brk 978-0-\brk 471-\brk 47756-3;\brk DOI 10.1021/\brk ar300145c) & \eqref{eq:ch5_logistic_branch},\allowbreak\eqref{eq:ch5_dxidV_branch},\allowbreak\eqref{eq:ch5_schg} \\
AL-54 & ★ 충전서도 $n_j^\eff I_j\eta_j^b\ge0$(2법칙) — 부호 상쇄(충전 $I_j{<}0$ \emph{이고} $\eta_j^\chg{<}0$, 음$\times$음$=$양)로 유도; 가역 열 $q_\rev^b{=}I_j\Pi_j$ 는 $I_j$ 부호 반전으로 branch 부호가변 & GROUNDED & degrootmazur\brk 1962(book),\brk onsager1931,\brk schnakenberg\brk 1976 (DOI 10.1103/\brk PhysRev.37.405;\brk 10.1103/\brk RevModPhys.48.571) & \eqref{eq:ch5_A_branch},\allowbreak\eqref{eq:ch5_eta_branch},\allowbreak\eqref{eq:ch5_qirr_charge} \\
AL-55 & ★ branch target $\xi_{\tar,j}^b{=}$ Ch3 $\xi_{\sstat,j}^b$(branch-local DB;\brk true-eq 에서 $C_j$ 이탈);\brk Level A ``방향'' 부분 붕괴(target 이력 의존,\brk mobility 불변);\brk metastable free-energy 근거 시만 본문 모델 & GROUNDED/\brk BOUNDED & dreyer2010,\brk sethna1993,\brk bazant2013 (DOI 10.1038/\brk nmat2730;\brk 10.1103/\brk PhysRevLett.70.3347;\brk 10.1021/\brk ar300145c) & \eqref{eq:ch5_branch_target},\allowbreak\eqref{eq:ch5_branch_db},\allowbreak\eqref{eq:ch5_ratio_branch} \\
AL-56 & hysteresis state $h_j\in[-1,1]$ 1차 완화($\dd h_j/\dd t{=}\lambda_j^b[h_\tar^b{-}h_j]$); $\lambda_j^b$ $=$ domain rearrangement/depinning time scale(피팅속도 아님);\brk rest 시 eq vs meta 두 후보(독립 검증 필요) & BOUNDED & sethna1993,\brk dreyer2010 (DOI 10.1103/\brk PhysRevLett.70.3347;\brk 10.1038/\brk nmat2730) & \eqref{eq:ch5_hys_center},\allowbreak\eqref{eq:ch5_h_dyn},\allowbreak\eqref{eq:ch5_rest_eq}--\eqref{eq:ch5_h_rest_relax} \\
AL-57 & ★ branch-local kinetic dissipation($n_j^\eff I_j\eta_\prog^b$,\brk target 좇는 소산) $\ne$ metastable 자유에너지 loss($V_\tar^b{-}V_\eq$,\brk branch 가 true 평형 풀릴 때); 합산 시 이중계상 & BOUNDED & dreyer2010,\brk degrootmazur\brk 1962(book),\brk bazant2013 (DOI 10.1038/\brk nmat2730;\brk 10.1021/\brk ar300145c) & \eqref{eq:ch5_branch_affinity} boundbox, \eqref{eq:ch5_hys_heat} boundbox \\
AL-58 & ★ $\Delta V_\obs\ne\Delta V_\hys$(loop 폭 분해: hys/pol/kin/transp/aging);\brk polarization 은 $|I_s|{\to}0$ 소멸$\cdot$$I_s$ 부호 의존,\brk true hys 는 $|I_s|{\to}0$ 잔존;\brk full-cell $\ne$ 음극 단독;\brk ICA/DVA 좌표 명시 & BOUNDED & dreyer2010,\brk bard\brk faulkner,\brk newman (DOI 10.1038/\brk nmat2730;\brk ISBN 978-0-\brk 471-\brk 04372-0;\brk 978-0-\brk 471-\brk 47756-3) & \eqref{eq:ch5_vapp_branch},\allowbreak\eqref{eq:ch5_gap_decomp},\allowbreak\eqref{eq:ch5_obs_not_hys},\allowbreak\eqref{eq:ch5_fullcell} \\
AL-59 & hysteresis heat $=$ loop area 의 \emph{분리된} 성분만($J_\hys^b\mathcal A_\hys^b\ge0$,\brk flux--force; $|I_s|\Delta V_{\hys\text{-}\loss}^b$ 는 부호정렬 시);\brk branch heat 분해;\brk loop area 전체 $\ne$ true hys heat; 중복계상 금지 & GROUNDED/\brk BOUNDED & dreyer2010,\brk degrootmazur\brk 1962(book),\brk onsager1931 (DOI 10.1038/\brk nmat2730;\brk 10.1103/\brk PhysRev.37.405) & \eqref{eq:ch5_hys_heat},\allowbreak\eqref{eq:ch5_branch_heat},\allowbreak\eqref{eq:ch5_branch_rev} \\
\bottomrule
\end{longtable}
}
\textbf{AGP(Assumption Grounding Policy).} 모든 본문 가정식은 AL-\# 를 인용한다. 본 장에 FLAGGED \emph{전용}
항목은 없으며(branch target 의 metastable 해석 가능성$\cdot$가역/비가역/hysteresis heat 분리 비율은 BOUNDED $+$
독립 검증 전제), 히스테리시스의 metastable/many-particle$\cdot$disorder 기원과 충전 부호 유도$\cdot$2법칙 보존은
GROUNDED, signed-coordinate$\cdot$polarization 분리$\cdot$중복계상 금지$\cdot$식별성은 BOUNDED(유효범위 병기)다.
임의 empirical branch offset / arbitrary hysteresis curve / clip/cap/step 정의식 0(GS-4; empirical branch fit
은 practicebox 로 격하). solver/numerical 구현 0(Ch1 부록 B 분리). \textbf{DOI 정확 귀속(★ Ch3$\cdot$Ch4 오귀속 교훈
계승):} dreyer2010 $=$ 10.1038/nmat2730, sethna1993 $=$ 10.1103/PhysRevLett.70.3347, degrootmazur1962 $=$
교과서(DOI 없음), onsager1931 $=$ 10.1103/PhysRev.37.405, schnakenberg1976 $=$ 10.1103/RevModPhys.48.571 —
Onsager DOI 를 degrootmazur 책에 붙이는 류의 오귀속 0건.

% ====================================================================
\subsection{Chapter 5 자체 검수표}\label{sec:ch5_selfcheck}
% ====================================================================
{\small
\begin{longtable}{>{\raggedright\arraybackslash}p{0.74\textwidth} >{\raggedright\arraybackslash}p{0.16\textwidth}}
\toprule 검수 항목 & 결과 \\ \midrule \endhead
Chapter 1--4(RB) 를 수정 없이 적층; 전하 보존식이 충전으로 고쳐 쓰이지 않음(P3-6) & 통과(\S\ref{sec:ch5_inherit}) \\
signed current convention 을 먼저 정의; $I_s$ vs $|I_s|$ 구분; 방전 환원 검산 & 통과(식~\eqref{eq:ch5_signed_current},\allowbreak\eqref{eq:ch5_relbalance}) \\
$q_\stt$ 와 reference form 명시;\brk LLI/LAM 임의 흡수 금지 & 통과(식~\eqref{eq:ch5_qstate},\allowbreak\eqref{eq:ch5_relbalance}) \\
방전 기준 $\xi_j$ continuity 유지; $\zeta_j$ 는 종속 보조 변수(독립 state 증식 X) & 통과(\S\ref{sec:ch5_xikeep}) \\
★ 충전 branch 부호 $s_{\phi,j}^\chg{=}-1$ 을 logistic 지수 부호 정합으로 \emph{유도}(주장 X;\brk CHARTER step1) & 통과(식~\eqref{eq:ch5_dxidV_branch},\allowbreak\eqref{eq:ch5_schg}) \\
★ 충전서 $\mathcal A_j,\eta_j,q_\irr$ 부호 뒤집힘 + $n_j^\eff I_j\eta_j^b\ge0$ 부호 상쇄로 유도(2법칙) & 통과(verifybox,식~\eqref{eq:ch5_qirr_charge}) \\
가역 열 부호가변성($I_j$ 부호 반전으로 branch 발열 비대칭) 명시 & 통과(\S\ref{sec:ch5_chargesign},\ref{sec:ch5_branchheat}) \\
true equilibrium target 과 metastable branch target 구분(many-particle/disorder 기원) & 통과(\S\ref{sec:ch5_branch}) \\
★ branch target $=$ Ch3 $\xi_{\sstat}$(detailed balance 이탈); 히스테리시스 $=$ target 이동(mobility 아님) & 통과(keybox) \\
★ Ch3$\to$Ch5 전달: DB/keystone 이 깨지는 정확한 조건($C_j$ 이탈/$\xi_j$-의존); 안 깨지는 것(fb 구조) 구분 & 통과(\S\ref{sec:ch5_dbbreak}) \\
branch target 을 arbitrary fitting curve 로 두지 않음(metastable 근거 요구) & 통과(boundbox) \\
hysteresis state $h_j$ 에 물리적 memory 의미 + 부호$\cdot$차원$\cdot$구간 검산 & 통과(\S\ref{sec:ch5_hstate}) \\
branch forward/backward kinetics 를 Ch3 구조와 연결;\brk branch-local DB; 부호 검산 & 통과(식~\eqref{eq:ch5_branch_db},\allowbreak\eqref{eq:ch5_branch_affinity}) \\
★ branch-local kinetic dissipation vs metastable 자유에너지 loss 구분(이중계상 금지) & 통과(\S\ref{sec:ch5_kinetics} boundbox) \\
★ $\Delta V_\obs\ne\Delta V_\hys$ 명시; $R_n^b$ 는 polarization($I_s$ 부호 의존, $|I_s|{\to}0$ 소멸) & 통과(식~\eqref{eq:ch5_obs_not_hys}) \\
ICA/DVA 좌표 명시(signed vs branch positive); 방전 단독서 $Q_s{=}Q_\ext$ 환원 & 통과(\S\ref{sec:ch5_icadva}) \\
full-cell apparent 전위와 internal-loss double counting 금지; 음극 단독과 구분 & 통과(\S\ref{sec:ch5_fullcell}) \\
★ loop area $=$ 소산열이되 전체 $\ne$ true hys heat;\brk flux--force 후보;\brk entropy production 정합(Ch4) & 통과(\S\ref{sec:ch5_hysheat}) \\
branch-dependent heat 를 물리 근거 항으로 제한; 충전 $\dd\xi/\dd t<0$ sign 명시 & 통과(\S\ref{sec:ch5_branchheat}) \\
rest relaxation 의 eq vs meta 두 후보 구분;\brk branch switching $\xi_j,h_j,V_n$ continuity & 통과(\S\ref{sec:ch5_rest}) \\
모든 식 차원$\cdot$부호$\cdot$극한$\cdot$2법칙 검산 명시 & 통과(\S\ref{sec:ch5_signed},\ref{sec:ch5_chargesign},\ref{sec:ch5_hysheat} 등) \\
``자명/clearly/쉽게/obviously'' 본문 0건(G-noleap) & 통과 \\
모든 가정 AL-50$\sim$59 인용;\brk BOUNDED 유효범위 병기;\brk DOI 정확 귀속(오귀속 0) & 통과(\S\ref{sec:ch5_al}) \\
empirical branch fit 을 실무 팁으로 격하;\brk solver 0 & 통과(\S\ref{sec:ch5_ident}) \\
Ch1 부록 B 이월 항목 분리 & 통과(\S\ref{sec:ch5_transmit}) \\
\bottomrule
\end{longtable}
}

% ====================================================================
\subsection{실데이터 피팅(Ch1 부록 B)으로 전달되는 항목}\label{sec:ch5_transmit}
% ====================================================================
\begin{enumerate}[label=\arabic*),leftmargin=2.0em]
\item \textbf{시간 적분.} signed current $I_s$ 와 $q_\stt$ 기반 시간 적분; branch switching 시 $\xi_j,h_j,V_n$
  continuity(재초기화 X); branch 별 root finding 으로 매 시간점 $V_n$ 결정.
\item \textbf{분리 검증.} hysteresis vs polarization 구분 검증 절차($|I_s|\to0$ 외삽, GITT, rest relaxation);
  가역/비가역/hysteresis heat 의 calorimetry 분리; full-cell 조합 시 양극/음극 coordinate matching.
\item \textbf{empirical 격하.} empirical branch fit 을 \emph{validation candidate} 로만 취급(metastable
  free-energy/domain memory 분리 안 되면 주 결론서 제외; Chapter 1--4 의 forward-only 원칙).
\end{enumerate}

% ====================================================================
\subsection{Chapter 5 의 결론}\label{sec:ch5_concl}
% ====================================================================
첫째, Chapter 1--4 의 방전 기준 모델은 signed current $I_s$ 와 signed state coordinate $q_\stt$ 도입으로 충전
branch 까지 확장된다(식~\eqref{eq:ch5_signed_current},\allowbreak\eqref{eq:ch5_qstate}). 내부 state 는 방전
기준 $\xi_j$ 를 유지하고($\zeta_j=1-\xi_j$ 는 종속 보조 변수), branch switching 시 $\xi_j,h_j,V_n$ 은 연속이며
방전 단독서 Chapter 1 식으로 정확히 환원된다(앞 장과 충돌 없음).

둘째, \textbf{충전 branch 의 부호 인자 $s_{\phi,j}^\chg=-1$ 은 \emph{유도}된다}(식~\eqref{eq:ch5_schg}). logistic
기울기의 부호가 정확히 $s_{\phi,j}^b$ 이고(식~\eqref{eq:ch5_dxidV_branch}), ``$\mathcal A_j>0\Leftrightarrow$
진행 촉진''이라는 단 하나의 규약을 방전($\xi_j$ 증가)과 충전($\xi_j$ 감소)에 일관 적용하면 방전 $+1$ 과 충전
$-1$ 이 동시에 고정된다(주장이 아니라 부호 정합 — CHARTER step1, Ch.4 이월 수령). 충전서 $\mathcal A_j,\eta_j,
q_\irr$ 의 부호가 뒤집히지만, \textbf{$I_j$ 와 $\eta_j^b$ 가 \emph{동시에} 부호를 뒤집어(음$\times$음$=$양)
비가역 발열 $n_j^{\eff}I_j\eta_j^b\ge0$ 이 부호 상쇄로 2법칙을 보존}한다(식~\eqref{eq:ch5_qirr_charge}, verifybox).
가역 열은 $I_j$ 부호 반전으로 branch 부호가변이며, 이것이 충방전 발열 비대칭의 한 원천이다.

셋째, true equilibrium target $\xi_{\eq,j}$ 와 metastable branch target $\xi_{\tar,j}^b$ 를 구분한다.
\textbf{branch target 은 Chapter 3 의 nonequilibrium stationary $\xi_{\sstat,j}$ 이며}(식~\eqref{eq:ch5_branch_db}),
$\xi_{\tar,j}^b\ne\xi_{\eq,j}$ 는 detailed balance 의 branch 이탈이다 — 히스테리시스는 mobility 변화가 아니라
\emph{target 의 이력 의존}(Level A 분업의 부분 붕괴)이다. 그 기원은 many-particle/mosaic 열역학(dreyer2010)과
disorder-driven first-order 전이의 barrier hierarchy(sethna1993)이며, Chapter 3 의 keystone 이 깨지는 정확한
자리($C_j$ 의 detailed-balance 이탈 / $r_j^+/r_j^-$ 의 $\xi_j$-의존)를 명시했다(\S\ref{sec:ch5_dbbreak}).

넷째, 관측 voltage gap 은 hysteresis 가 아니다(\textbf{$\Delta V_\obs\ne\Delta V_\hys$}, 식~\eqref{eq:ch5_obs_not_hys}).
gap 에는 polarization($R_n^b$; $|I_s|\to0$ 소멸, $I_s$ 부호 의존), kinetic lag, transport, aging 이 섞이며,
true thermodynamic hysteresis(평형 분지 차)만이 $|I_s|\to0$ 에서 잔존한다 — 이 조작적 차이가 둘을 분리한다.

다섯째, loop 면적 $E_{\mathrm{loop}}=\oint V\,\dd Q$ 는 소산 에너지이되 그 전체가 true hysteresis heat 는
\emph{아니다}(dreyer2010 정합). 분리된 hysteresis loss 만이 flux--force $J_\hys^b\mathcal A_\hys^b\ge0$
(Chapter 4 의 entropy production 채널)으로 후보 발열항이 되며, branch-local kinetic dissipation 과 합산하면
같은 에너지를 이중 계상한다(식~\eqref{eq:ch5_hys_heat} boundbox). full-cell ICA/DVA 는 양극/음극/branch signed
loss/transport/polarization 이 섞인 관측 신호라 음극 단독과 직접 같지 않다.

여섯째, Chapter 5 는 물리적 path dependence 와 metastability 를 중심으로 충방전 확장 구조를 닫되, Chapter 1--4
기본식을 수정하지 않는다. 충전 부호 유도$\cdot$2법칙 보존$\cdot$target 이력 의존$\cdot$loop-area 분리가 한 인과
사슬 안에서 정합하며, Ch1 부록 B 에서 이 구조를 실제 수치 해법$\cdot$validation pipeline$\cdot$GITT/저율/고율
데이터 비교로 연결한다.

\subsection*{실데이터 피팅(Ch1 부록 B)으로의 전달}
본 장이 확정한 충방전 구조(signed $I_s\cdot q_\stt$, 충전 부호 $s_{\phi,j}^b$, branch target $\xi_{\tar,j}^b
=\xi_{\sstat,j}^b$, hysteresis state $h_j$, $\Delta V_\obs\ne\Delta V_\hys$, loop-area 분리 발열)와 상태변수
위에서: Ch1 부록 B 은 stiff kinetics $+$ branch switching 의 시간 적분, 매 시간점 root-find 로 $V_n$ 결정,
hysteresis vs polarization 분리 검증 pipeline, calorimetry 부재 시 발열$\cdot$branch curve 를 피팅하지 않고
forward prediction 으로만 두는 제한을 다룬다(Chapter 1--4 의 forward-only 원칙).



% ====================================================================
% ===== 통합 참고문헌 (33 키 union; Ch1-first dedupe; refs 문건과 동일) =====
% ====================================================================
\clearpage
\section{통합 참고문헌 (Chapter 1$\sim$5 union, 33 키)}\label{sec:full_bib}
아래 33 키는 Chapter 1$\sim$5 의 \texttt{\textbackslash bibitem} 합집합(중복 제거; Ch1 우선)이며, 별도 통합 references 문건(\texttt{graphite\_ica\_refs\_rebuilt.tex})의 bibliography 와 동일하다. 통합 Assumption Ledger(AL-1$\sim$69)는 그 문건에 수록.
\begin{thebibliography}{99}
\bibitem{doyle1993} M. Doyle, T. F. Fuller, J. Newman, ``Modeling of Galvanostatic Charge and Discharge
  of the Lithium/Polymer/Insertion Cell,'' \emph{J. Electrochem. Soc.} \textbf{140}, 1526 (1993).
  DOI: 10.1149/1.2221597.
\bibitem{newman} J. Newman, K. E. Thomas-Alyea, \emph{Electrochemical Systems}, 3rd ed., Wiley (2004).
  ISBN 978-0-471-47756-3.
\bibitem{mckinnon1983} W. R. McKinnon, R. R. Haering, ``Physical Mechanisms of Intercalation,'' in
  \emph{Modern Aspects of Electrochemistry} No.~15, Plenum, New York (1983), p.~235.
\bibitem{bazant2013} M. Z. Bazant, ``Theory of Chemical Kinetics and Charge Transfer Based on
  Nonequilibrium Thermodynamics,'' \emph{Acc. Chem. Res.} \textbf{46}, 1144 (2013).
  DOI: 10.1021/ar300145c.
\bibitem{eyring1935} H. Eyring, ``The Activated Complex in Chemical Reactions,'' \emph{J. Chem. Phys.}
  \textbf{3}, 107 (1935). DOI: 10.1063/1.1749604.
\bibitem{evanspolanyi1938} M. G. Evans, M. Polanyi, ``Inertia and Driving Force of Chemical Reactions,''
  \emph{Trans. Faraday Soc.} \textbf{34}, 11 (1938). DOI: 10.1039/TF9383400011.
\bibitem{bardfaulkner} A. J. Bard, L. R. Faulkner, \emph{Electrochemical Methods}, 2nd ed., Wiley (2001).
  ISBN 978-0-471-04372-0.
\bibitem{marcus1956} R. A. Marcus, ``On the Theory of Oxidation-Reduction Reactions Involving Electron
  Transfer. I,'' \emph{J. Chem. Phys.} \textbf{24}, 966 (1956). DOI: 10.1063/1.1742723.
\bibitem{onsager1931} L. Onsager, ``Reciprocal Relations in Irreversible Processes. I,'' \emph{Phys. Rev.}
  \textbf{37}, 405 (1931). DOI: 10.1103/PhysRev.37.405.
\bibitem{degrootmazur1962} S. R. de Groot, P. Mazur, \emph{Non-Equilibrium Thermodynamics},
  North-Holland (1962); Dover (1984).
\bibitem{dahn1991} J. R. Dahn, ``Phase Diagram of Li$_x$C$_6$,'' \emph{Phys. Rev. B} \textbf{44}, 9170
  (1991). DOI: 10.1103/PhysRevB.44.9170.
\bibitem{ohzuku1993} T. Ohzuku, Y. Iwakoshi, K. Sawai, ``Formation of Lithium-Graphite Intercalation
  Compounds...,'' \emph{J. Electrochem. Soc.} \textbf{140}, 2490 (1993). DOI: 10.1149/1.2220849.
\bibitem{funabiki1999ea} A. Funabiki, M. Inaba, T. Abe, Z. Ogumi, ``Nucleation and Phase-Boundary
  Movement upon Stage Transformation in Lithium-Graphite Intercalation Compounds,''
  \emph{Electrochim. Acta} \textbf{45}, 865 (1999). DOI: 10.1016/S0013-4686(99)00290-X.
\bibitem{funabiki1999jes} A. Funabiki, M. Inaba, T. Abe, Z. Ogumi, ``Stage Transformation of
  Lithium-Graphite Intercalation Compounds Caused by Electrochemical Lithium Intercalation,''
  \emph{J. Electrochem. Soc.} \textbf{146}, 2443 (1999). DOI: 10.1149/1.1391953.
\bibitem{baessler1993} H. B\"assler, ``Charge Transport in Disordered Organic Photoconductors,''
  \emph{Phys. Status Solidi B} \textbf{175}, 15 (1993). DOI: 10.1002/pssb.2221750102.
\bibitem{svare2000} I. Svare, S. W. Martin, F. Borsa, ``Stretched Exponentials with $T$-Dependent
  Exponents from Fixed Distributions of Energy Barriers for Relaxation Times in Fast-Ion Conductors,''
  \emph{Phys. Rev. B} \textbf{61}, 228 (2000). DOI: 10.1103/PhysRevB.61.228.
\bibitem{johnston2006} D. C. Johnston, ``Stretched Exponential Relaxation Arising from a Continuous Sum
  of Exponential Decays,'' \emph{Phys. Rev. B} \textbf{74}, 184430 (2006). DOI: 10.1103/PhysRevB.74.184430.
\bibitem{lindsey1980} C. P. Lindsey, G. D. Patterson, ``Detailed Comparison of the Williams-Watts and
  Cole-Davidson Functions,'' \emph{J. Chem. Phys.} \textbf{73}, 3348 (1980). DOI: 10.1063/1.440530.
\bibitem{jcp2017} (사용자 PhD) K. Lee, S. Lee, C. H. Choi, S. Lee, ``Effects of External Electric Field and Anisotropic Long-Range Reactivity on Charge Separation Probability,''
  \emph{J. Chem. Phys.} \textbf{147}, 144111 (2017). DOI: 10.1063/1.5000882.
\bibitem{lee2011} (Refs.~6) S. Lee, C. Y. Son, J. Sung, S.-H. Chong, ``Communication: Propagator for Diffusive Dynamics of an Interacting Molecular Pair,''
  \emph{J. Chem. Phys.} \textbf{134}, 121102 (2011). DOI: 10.1063/1.3565476.
\bibitem{son2013} (Refs.~7) C. Y. Son, J. Kim, J.-H. Kim, J. S. Kim, S. Lee, ``An Accurate Expression for the Rates of Diffusion-Influenced Bimolecular Reactions with Long-Range Reactivity,''
  \emph{J. Chem. Phys.} \textbf{138}, 164123 (2013). DOI: 10.1063/1.4802584.
\bibitem{dubarry2022} M. Dubarry, D. Anseán, ``Best practices for incremental capacity analysis,''
  \emph{Front. Energy Res.} \textbf{10}, 1023555 (2022). DOI: 10.3389/fenrg.2022.1023555.
\bibitem{fly2020} A. Fly, R. Chen, ``Rate dependency of incremental capacity analysis (dQ/dV)...,''
  \emph{J. Energy Storage} \textbf{29}, 101329 (2020). DOI: 10.1016/j.est.2020.101329.
\bibitem{brenan1996} K. E. Brenan, S. L. Campbell, L. R. Petzold, \emph{Numerical Solution of Initial-Value
  Problems in Differential-Algebraic Equations}, SIAM Classics in Applied Mathematics 14 (1996).
  DOI: 10.1137/\brk 1.9781611971224.
\bibitem{hindmarsh2005} A. C. Hindmarsh, P. N. Brown, K. E. Grant, S. L. Lee, R. Serban, D. E. Shumaker,
  C. S. Woodward, ``SUNDIALS: Suite of Nonlinear and Differential/Algebraic Equation Solvers,''
  \emph{ACM Trans. Math. Softw.} \textbf{31}, 363 (2005). DOI: 10.1145/\brk 1089014.1089020.
\bibitem{bernardi1985} D. Bernardi, E. Pawlikowski, J. Newman, ``A General Energy Balance for Battery
  Systems,'' \emph{J. Electrochem. Soc.} \textbf{132}, 5 (1985). DOI: 10.1149/\brk 1.2113792.
\bibitem{thomasnewman2003} K. E. Thomas, J. Newman, ``Thermal Modeling of Porous Insertion Electrodes,''
  \emph{J. Electrochem. Soc.} \textbf{150}, A176 (2003). DOI: 10.1149/\brk 1.1531194.
\bibitem{rao1997} L. Rao, J. Newman, ``Heat-Generation Rate and General Energy Balance for Insertion
  Battery Systems,'' \emph{J. Electrochem. Soc.} \textbf{144}, 2697 (1997). DOI: 10.1149/1.1837884.
\bibitem{reynier2004} Y. Reynier, R. Yazami, B. Fultz, ``Thermodynamics of Lithium Intercalation into
  Graphites and Disordered Carbons,'' \emph{J. Electrochem. Soc.} \textbf{151}, A422 (2004).
  DOI: 10.1149/1.1646152.
\bibitem{prigogine1967} I. Prigogine, \emph{Introduction to Thermodynamics of Irreversible Processes},
  3rd ed., Interscience (1967). [entropy production; 교과서, DOI 없음]
\bibitem{schnakenberg1976} J. Schnakenberg, ``Network Theory of Microscopic and Macroscopic Behavior of
  Master Equation Systems,'' \emph{Rev. Mod. Phys.} \textbf{48}, 571 (1976). DOI: 10.1103/\brk RevModPhys.48.571.
\bibitem{dreyer2010} W. Dreyer, J. Jamnik, C. Guhlke, R. Huth, J. Mo\v{s}kon, M. Gaber\v{s}\v{c}ek,
  ``The Thermodynamic Origin of Hysteresis in Insertion Batteries,'' \emph{Nature Materials} \textbf{9},
  448 (2010). DOI: 10.1038/\brk nmat2730.
\bibitem{sethna1993} J. P. Sethna, K. Dahmen, S. Kartha, J. A. Krumhansl, B. W. Roberts, J. D. Shore,
  ``Hysteresis and Hierarchies: Dynamics of Disorder-Driven First-Order Phase Transformations,''
  \emph{Phys. Rev. Lett.} \textbf{70}, 3347 (1993). DOI: 10.1103/\brk PhysRevLett.70.3347.
\end{thebibliography}

% ====================================================================
\section*{통합 노트 (Ch1~5 병합 메타; 구 Ch6 해체)}
\addcontentsline{toc}{section}{통합 노트}
% ====================================================================
\begin{itemize}[leftmargin=1.6em]
\item \textbf{원본 5개 챕터 수정 0 (body byte-verbatim).} 각 챕터의 \texttt{\textbackslash maketitle}
  이후$\sim$\texttt{\textbackslash begin\{thebibliography\}} 직전 본문을 그대로 복사(요약$\cdot$축약 0).
  label 충돌 0, cross-chapter 평문 상호참조 현행 유지(수정 0).
\item \textbf{Ch6 해체.} 별도 Chapter 6 없음 — 구 Chapter 6(실데이터 피팅 워크플로)은 Ch1
  \emph{부록 B}(\texttt{sec:ch6\_*} 라벨)에 흡수되어 본 통합본에 Ch1 의 일부로 포함된다.
  본문의 ``구 Chapter 6''$\cdot$``부록 B'' 언급은 그 흡수 설명이다.
\item \textbf{heading 1단계 demote.} 챕터 = \texttt{\textbackslash section{Chapter N}}, 챕터 내부
  \texttt{\textbackslash section}$\to$\texttt{\textbackslash subsection},
  \texttt{\textbackslash subsection}$\to$\texttt{\textbackslash subsubsection}
  (article 유지, \texttt{\textbackslash part}/\texttt{\textbackslash chapter} 미사용).
\item \textbf{theorem 환경 caption 통일.} \texttt{verifybox}$\cdot$\texttt{linkbox} 는 union 1회
  정의이므로 대표 caption 으로 통일했다 — box 내부 \emph{물리 내용$\cdot$식은 변경 0}, heading label 문구만 통일.
\item \textbf{통합 bibliography.} 33 키 union(Ch1 우선 dedupe; refs 문건과 동일 서지$\cdot$DOI).
\end{itemize}

\end{document}
